\documentclass[11pt,a4paper]{article}
\usepackage[margin=25mm,headheight=14pt]{geometry}
\usepackage[T1]{fontenc}
\usepackage{lmodern,microtype}
\usepackage{amsmath,amssymb,amsthm,mathtools}
\usepackage{booktabs,longtable,array,tabularx}
\usepackage{enumitem,fancyhdr,tocloft}
\usepackage[hidelinks,bookmarksnumbered]{hyperref}
\hypersetup{pdftitle={Subset Sums in a Fourfold Sumset},pdfsubject={Complete mathematical descent to the upper bound 86},pdfauthor={}}
\setlength{\parindent}{1.5em}
\setlength{\parskip}{0.2em}
\setlength{\emergencystretch}{2em}
\allowdisplaybreaks[3]
\setlist{itemsep=0.2em,topsep=0.4em}
\setcounter{tocdepth}{1}
\setlength{\cftsecnumwidth}{2.2em}
\setlength{\cftbeforesecskip}{0.7em}
\setlength{\cftbeforepartskip}{1.4em}
\pagestyle{fancy}
\fancyhf{}
\fancyhead[L]{\small Subset sums in a fourfold sumset}
\fancyhead[R]{\small Upper bound 86}
\fancyfoot[C]{\thepage}
\renewcommand{\headrulewidth}{0pt}
\newtheorem{theorem}{Theorem}[section]
\newtheorem{lemma}[theorem]{Lemma}
\newtheorem{proposition}[theorem]{Proposition}
\newtheorem{corollary}[theorem]{Corollary}
\theoremstyle{remark}
\newtheorem{remark}[theorem]{Remark}
\newtheorem{question}[theorem]{Question}
\newcommand{\FS}{\operatorname{FS}}
\newcommand{\card}[1]{\lvert #1\rvert}
\newcommand{\Zint}[2]{[#1,#2]_{\mathbb Z}}
\newcommand{\rhoD}{\varrho}
\newcommand{\Z}{\mathbb Z}
\newcommand{\R}{\mathbb R}
\newcommand{\tri}[1]{T_{#1}}
\newcommand{\hull}{\operatorname{hull}}
\begin{document}
\begin{titlepage}
\centering
\vspace*{1.9cm}
{\LARGE\bfseries Subset Sums in a Fourfold Sumset\par}
\vspace{0.7cm}
{\Large Proof Supplement for the Upper Bound 86\par}
\vspace{0.6cm}
{\large 5 October 2026\par}
\vspace{1.8cm}
\begin{minipage}{0.86\textwidth}
\setlength{\parskip}{0.6em}\setlength{\parindent}{0pt}


This supplement proves that every finite real set $A$ whose subset
sums lie in its fourfold sumset has at most $86$ elements. It supplies the
complete descent used as an input to the earlier report establishing
$71\le N\le85$; the later exclusion of size $86$ and the lower
construction are outside its scope.

The complete finite calculations have been replayed using exact
integer or rational arithmetic. Their inputs, coverage obligations and
verification record are specified in Section~\ref{completion:verification}.

Write $\FS(A)=\{\sum_{x\in X}x:X\subseteq A\}$, with the empty sum
included, and $4A=\{a_1+a_2+a_3+a_4:a_i\in A\}$, with repetition allowed.
The notation $N$ denotes the supremum of $|A|$ over feasible finite real
sets. The preliminary bound proved here shows that this supremum is finite.
We write $T_r=r(r+1)/2$ for triangular numbers and use AP to mean
an arithmetic progression.

The argument separates the proof that actual sets satisfy each finite
system from the exact nonnegative row combinations that contradict those
systems. The electronic certificate data are an integral part of this
computer-assisted proof.
\end{minipage}
\vfill
{\small Mathematical exposition and complete dependency specification\par}
\end{titlepage}
\tableofcontents
\clearpage
\section*{How the proof is organized}
\addcontentsline{toc}{section}{How the proof is organized}
The proof begins with a universal finite cap and then excludes every
remaining cardinality in descending order. It never assumes that deleting
arbitrary elements preserves feasibility.
\begin{enumerate}
\item Part~\ref{part:foundations} proves the real-to-integer reduction,
covering and peeling estimates, and the preliminary bound $N\le115$.
\item Part~\ref{part:upper96} proves all exclusions from size $115$ through
size $97$, including the common-lattice, defect and hole lemmas used later.
\item Part~\ref{part:upper86} proves the successive exclusions of sizes
$96,95,\ldots,87$, with each actual-set reduction and every selected
finite alternative specified.
\item The appendices give the fixed-slot certificate and the complete
smaller coefficient pools, together with exact specifications of the
larger electronic pools and branch trees.
\end{enumerate}
The only externally cited mathematical inputs are the published statements
explicitly reproduced with their hypotheses in the proof. References to
prior versions of the research are not used as mathematical premises.



\section{The logical role of the finite certificates}\label{completion:certificate-principle}
The finite calculations in this proof are exact arithmetic arguments, not
numerical evidence for an inequality. We state explicitly how they are
used. None of the following principles asserts that a feasible point of
a necessary model is realized by a set.

\begin{lemma}[Exact nonnegative row combination]\label{completion:farkas}
Let $M\in\mathbb Q^{p\times q}$ and $b\in\mathbb Q^p$.
If $w\in\mathbb Q_{\ge0}^p$ satisfies
\[
 w^{\mathsf T}M\ge0\quad\hbox{coordinatewise},\qquad
 w^{\mathsf T}b<0,
\]
then there is no $x\in\mathbb R_{\ge0}^q$ satisfying $Mx\le b$.
\end{lemma}
\begin{proof}
Such an $x$ would give
$0\le(w^{\mathsf T}M)x\le w^{\mathsf T}b<0$.
Multiplying a rational $w$ by a common positive denominator permits all
weights to be integers. Zero weights may be omitted.
\end{proof}

\begin{lemma}[Exhaustive necessary models]\label{completion:coverage}
Fix a cardinality $n$. Suppose every feasible integer $n$-element set
containing zero determines a member of a specified finite set of roots.
For each root, suppose the written construction provides a finite tree
of necessary systems with the following properties.
\begin{enumerate}
\item Every coordinate assigned to an actual set has the asserted sign
and satisfies every row at the root.
\item At every internal node, every actual set satisfying its hypotheses
satisfies the hypotheses of at least one child. All these children are
retained. Additional rows are proved under that child's hypotheses.
\item Every leaf has a row combination as in
Lemma~\ref{completion:farkas} for its actual necessary system.
\end{enumerate}
Then no such integer set exists.
\end{lemma}
\begin{proof}
An alleged set selects a root and then, by the second property, a child
at each internal node. Finiteness yields a leaf. Its actual coordinates
contradict Lemma~\ref{completion:farkas}.
\end{proof}

Overlapping alternatives are harmless: disjointness of the mathematical
children is unnecessary. Their union must cover the parent. If an
integer-valued expression $f$ is split at an integer $t$, the alternatives
$f\le t$ and $f\ge t+1$ are exhaustive. This split is invalid for a merely
real auxiliary variable; for a real expression one must instead use
alternatives whose union covers all real values. The integrality of each
split expression is therefore part of the applicability proof.

The programs regenerate the labeled coefficient rows from the current
root and branch, read the supplied multipliers, and perform the indicated
integer or rational operations. A successful multiplier check does not
prove that its rows apply to an actual set; that is the first and second
property of Lemma~\ref{completion:coverage}, established in the written
chapters. Conversely, a written classification without a contradiction
at every leaf does not exclude its roots.

For the compact endpoints, assignment words select an ordinary multiplier
or mark a root for a supplemental tree. Completeness requires equality
between the actual marked root set and the union of the supplemental
parent sets, with no missing parent. The checker also verifies the complete
required child sets and the contradiction at every leaf. Counts alone,
an archived success field, a digest, or a previous run's terminal message
cannot substitute for these mathematical checks.


\clearpage\part{Global reductions and a preliminary bound}\label{part:foundations}

\section{Definitions and the cardinality reduction}\label{base:sec:global-reduction}
For a finite real set $A$, write
\[
 \FS(A)=\left\{\sum_{x\in X}x:X\subseteq A\right\},\qquad
 4A=\{a_1+a_2+a_3+a_4:a_i\in A\}.
\]
The empty subset contributes zero, and repetitions are allowed in $4A$.
Call $A$ feasible if $\FS(A)\subseteq4A$. We will first prove the finite
bound $|A|\le115$. The subsequent endpoint arguments exclude every
cardinality from $115$ down through $87$. The reduction below explains why
this proves the bound for all finite real sets.

\subsection{Adjoining zero}

\begin{lemma}[Adjoining zero]\label{base:lem:add-zero}
If $\FS(A)\subseteq4A$ and $0\notin A$, then $A\cup\{0\}$ also satisfies the
inclusion.
\end{lemma}

\begin{proof}
Adjoining zero creates no new subset-sum value, while the fourfold sumset can
only grow:
\[
 \FS(A\cup\{0\})=\FS(A)\subseteq4A\subseteq4(A\cup\{0\}).
\]
\end{proof}

It is therefore enough, for a uniform upper bound, to consider sets containing
zero.

\subsection{Integralization}

\begin{theorem}[Integralization]\label{base:thm:integerization}
For every $n\geq1$, a real $n$-element solution exists if and only if an
integral $n$-element solution exists.  Zero can be preserved.
\end{theorem}

\begin{proof}
Only the forward direction needs proof.  Write
$A=\{a_1<\cdots<a_n\}$.  For every $X\subseteq A$, choose one representation
\[
 \sum_{a_i\in X}a_i=a_{j_1}+a_{j_2}+a_{j_3}+a_{j_4},
\]
with repetition of indices allowed.  These finitely many identities form a
homogeneous linear system $Mx=0$ with integer coefficients, and
$a=(a_1,\ldots,a_n)$ is a solution.  The kernel is a rational linear
subspace, so its rational points are dense.  The relatively open region
$x_1<\cdots<x_n$ containing $a$ therefore contains a rational solution.
Clearing denominators gives distinct integers and preserves all selected
representations.  If $a_i=0$ is to be preserved, append $x_i=0$ before taking
the rational point.
\end{proof}


Rationalization is an existence argument. It preserves distinctness and a
prescribed zero; it does not imply that the original real set lies on a
one-dimensional rational lattice. Clearing denominators means multiplying
by a positive common denominator, so all strict order inequalities survive.

\begin{lemma}[Descending exact-cardinality reduction]\label{base:lem:descending-cardinality}
Let $m\ge1$. Suppose every feasible real set has at most $m$ elements and
there is no feasible integer set of exactly $m$ elements containing zero.
Then every feasible real set has at most $m-1$ elements.
\end{lemma}
\begin{proof}
A counterexample would have exactly $m$ elements. If it omitted zero,
Lemma~\ref{base:lem:add-zero} would produce a feasible set of size $m+1$,
contrary to the assumed bound. It therefore contains zero, and
Theorem~\ref{base:thm:integerization} produces an integer counterexample with
the same size and with zero preserved. This contradicts the exact-size
exclusion.
\end{proof}

Equivalently, a global finite cap $M$ together with exclusions of every
zero-containing integer size in $\{87,\ldots,M\}$ proves $|A|\le86$:
adjoin zero first, if necessary, and rationalize the resulting set.
Its cardinality is between $87$ and $M$ for any putative counterexample.
Neither formulation deletes arbitrary elements from a feasible set.
In particular, later subsets $D,D_1,E$ need not themselves be feasible;
all uses of feasibility are applied to subset sums of the original $A$.

\section{The four-extreme-element covering count}\label{base:sec:cover}

For a finite set $F$, put
\[
 R(F)=\{x+y:x,y\in F,\ x\ne y\},
 \qquad \Delta(F)=\card{2F\setminus R(F)}.
\]
Every value in $2F\setminus R(F)$ is a diagonal value $2x$, and
$x\mapsto2x$ is injective.  Hence
\begin{equation}
 \Delta(F)\leq\card F.
 \label{base:eq:delta-trivial}
\end{equation}

From a feasible set $A$ containing zero, take the four largest positive
elements and the four smallest negative elements, taking every element on a
side that has fewer than four.  Call their union $B$, and put
\[
 D=A\setminus B,\qquad d=\card D.
\]
Then $0\in D$ and $\card B\leq8$.

\begin{lemma}[Four-extreme-element covering]\label{base:lem:cover-count}
If $\FS(A)\subseteq4A$, then
\begin{equation}
 \boxed{\card{\FS(D)}\leq2\card{2D}+18d+85.}
 \label{base:eq:sharp-cover}
\end{equation}
\end{lemma}

\begin{proof}
We first work on the positive side.  Let
$p_1>p_2>p_3>p_4>0$ be the four largest positive elements, and put
$T_p=p_1+p_2+p_3+p_4$.  If $x\in\FS(D)$ and $x>0$, then
$T_p+x\in4A$.  Write a representation in decreasing order:
\[
 T_p+x=z_1+z_2+z_3+z_4,
 \qquad z_1\geq z_2\geq z_3\geq z_4.
\]
Classifying by the first failure of $z_i\leq p_i$ gives only
\[
\begin{array}{c|l}
2&2p_1+u+v,\\
3&p_1+2p_2+u\quad\text{or}\quad3p_2+u,\\
4&p_1+p_2+2p_3,\ p_1+3p_3,\ 2p_2+2p_3,\ p_2+3p_3,\ 4p_3.
\end{array}
\]
Subtracting $T_p$ and substituting $A=D\sqcup B$ gives the table
\[
\begin{array}{c|ccc}
\toprule
\text{source}&\text{translates of }2D&\text{translates of }D&\text{fixed slots}\\
\midrule
(2p_1-T_p)+A+A&1&8&36\\
(p_1+2p_2-T_p)+(A\setminus\{p_1\})&0&1&7\\
(3p_2-T_p)+(A\setminus\{p_1\})&0&1&7\\
\text{five fixed shapes}&0&0&5\\
\midrule
\text{total}&1&10&55\\
\bottomrule
\end{array}
\]
The negative side gives the same table.  Without deletions, the two sides and
zero yield
\begin{equation}
 \card{\FS(D)}\leq2\card{2D}+20d+111.
 \label{base:eq:raw-cover}
\end{equation}

If one side has fewer than four elements, all elements of that sign lie in
$B$, and $D$ has no subset sum of that sign.  The single remaining table is
already stronger than \eqref{base:eq:sharp-cover}.  We may therefore assume that
both sides have four extreme elements.  Write the absolute values on the
negative side as $q_1>q_2>q_3>q_4>0$.

Among the twenty translates of $D$, Appendix~\ref{base:app:slots} gives two
distinct pairs of centers $u\leq v$.  For any such pair,
\begin{equation}
 \card{(u+D)\cap(0,\infty)}+
 \card{(v+D)\cap(-\infty,0)}\leq d,
 \label{base:eq:halfline}
\end{equation}
because for one $z\in D$ the conditions $z>-u$ and $z<-v$ cannot both hold.
The two pairs reduce the linear contribution from $20d$ to $18d$.

There are $110$ fixed slots in the two tables.  Appendix~\ref{base:app:slots}
places them in one common universe and gives $26$ mutually disjoint groups.
The total score in every group is coordinatewise nonpositive in eight
strictly positive gap variables, so at least one slot in each group lies on
the wrong open half-line.  At most $110-26=84$ fixed slots survive.  Adding
zero gives the constant $85$, proving \eqref{base:eq:sharp-cover}.
\end{proof}

\section{Two four-element peeling steps}\label{base:sec:peeling}

\begin{lemma}[Four-element peeling]\label{base:lem:peeling}
Let $F$ be a finite real set containing zero and at least four nonzero
elements.  Four nonzero elements can be deleted to obtain $F'$ such that
\[
 \card{\FS(F)}\geq\card{\FS(F')}+\card{R(F)}-1.
\]
\end{lemma}

\begin{proof}
Choose the two largest positive and the two smallest negative elements,
filling with as-yet unselected nonzero elements if one side has fewer than
two.  Call the selected set $H$, put $F'=F\setminus H$, and define
\[
 \alpha=\sum_{h\in H,\,h<0}h,
 \qquad \beta=\sum_{h\in H,\,h>0}h,
 \qquad t=\beta-\alpha>0.
\]
The extreme choice gives $R(F)\subseteq[\alpha,\beta]$, while
\[
 \alpha+\FS(F')\subseteq\FS(F),
 \qquad \beta+\FS(F')\subseteq\FS(F).
\]
In every residue class modulo $t$ occupied by $\FS(F')$, the union of these
two translates contains at least one more point than $\FS(F')$.  In the
length-$t$ interval $[\alpha,\beta]$, each residue class occurs at most once,
except that the endpoints have the same residue.  Thus $R(F)$ occupies at
least $\card{R(F)}-1$ residue classes.  Add one point of $R(F)$ from every
class not already supplied by the two translates.  The stated inequality
follows.
\end{proof}

If $d\leq8$, then $\card A\leq d+8\leq16$, so the desired upper bound is
already trivial.  Hence assume $d\geq9$.  The set $D$ then has at least eight
nonzero elements, and after the first deletion the remaining set still has at
least four.  We may therefore apply the lemma twice:
\[
 D\supset D_1\supset E,
 \qquad \card{D_1}=d-4,
 \qquad \card E=d-8,
 \qquad 0\in E.
\]
Deleting four elements affects at most
\[
 4(d-4)+\binom42=4d-10
\]
unordered pairs of distinct elements.  Hence
\[
 \card{R(D_1)}\geq\card{R(D)}-(4d-10).
\]
If $s=\card{\FS(E)}$, the two peeling inequalities give
\begin{equation}
 \card{\FS(D)}\geq s+2\card{R(D)}-4d+8.
 \label{base:eq:peel-lower}
\end{equation}
Since $\card{2D}=\card{R(D)}+\Delta(D)$, comparison with
\eqref{base:eq:sharp-cover} eliminates $R(D)$ and gives
\begin{equation}
 \boxed{s\leq22d+2\Delta(D)+77.}
 \label{base:eq:core-upper}
\end{equation}

\section{The rigorous upper bound \texorpdfstring{$N\leq115$}{N <= 115}}
\label{base:sec:upper}

Write
\[
 E=P\cup\{0\}\cup(-Q),
\]
where $P,Q\subset\mathbb R_{>0}$, and put
\[
 a=\card P,\qquad b=\card Q,
 \qquad a+b=d-9.
\]

For an $r$-element positive real set $X$, define its subset-sum excess and its
doubled-value defect by
\[
 \mu(X)=\card{\FS(X)}-\left(1+\frac{r(r+1)}2\right),
 \qquad
 \delta_0(X)=\card{\{x\in X:2x\notin R(X\cup\{0\})\}}.
\]
The elementary increasing-chain argument gives $\mu(X)\geq0$.

We use the following published stability theorem in precisely its stated
range: if $r\geq4$, $0\leq M\leq r-4$, and
\[
 \card{\FS(X)}\leq1+\frac{r(r+1)}2+M,
\]
then $X=cB$ for some $c>0$ and an $r$-element set
$B\subset\mathbb Z_{>0}$ satisfying
\[
 \sum_{b\in B}b\leq\frac{r(r+1)}2+M.
\]
This is Theorem~1.1 of Carpenter--Defant--Kravitz~\cite{CDK2026}.

We also use Lemma~3.6 and Theorem~1.5 of
Balandraud--Girard--Griffiths--Hamidoune~\cite{BGGH2013}.  For a finite
subset of a finite abelian group, these give
\begin{equation}
 U\cap(-U)=\varnothing\quad\Longrightarrow\quad
 \card{\FS(U)}\geq2\card U,
 \label{base:eq:antisymmetric-bound}
\end{equation}
and, if $S=-S$, $\card S=2p\geq4$, and $\FS(S)$ is aperiodic, then
\begin{equation}
 \card{\FS(S)}\geq p(p-1)+5.
 \label{base:eq:symmetric-bound}
\end{equation}
Here a finite set is \emph{aperiodic} if only translation by zero fixes it.

\begin{lemma}[Aperiodic subset sums]\label{base:lem:aperiodic-linear}
Let $G$ be an abelian group and let
$R\subseteq G\setminus\{0\}$ have $m\geq2$ elements.  If $\FS(R)$ is
aperiodic, then
\begin{equation}
 \boxed{\card{\FS(R)}\geq2m-2.}
 \label{base:eq:aperiodic-linear}
\end{equation}
\end{lemma}

\begin{proof}
First suppose that $G$ is finite.  No member of $R$ has order two: otherwise
$\FS$ of that singleton, and hence $\FS(R)$, has a nonzero period.  Decompose
\[
 R=S\sqcup U,\qquad S=-S,\quad \card S=2p,\quad
 U\cap(-U)=\varnothing,\quad \card U=u.
\]
The subset-sum set of every subset of $R$ is aperiodic, since any of its
periods would also be a period after addition of the complementary
subset-sum set.  Kneser's addition theorem, in the form
$|C+D|\ge|C+H|+|D+H|-|H|$ with $H$ the stabilizer of $C+D$,
therefore has trivial stabilizer here and
gives
\[
 \card{\FS(R)}\geq\card{\FS(S)}+\card{\FS(U)}-1.
\]
If $p=0$, \eqref{base:eq:antisymmetric-bound} gives the stronger bound $2m$.
If $p=1$, then $\card{\FS(S)}=3$, and the right side is at least
$2u+2=2m-2$ (with the empty-$U$ case even larger).  If $p\geq2$ and
$u\geq1$, equations \eqref{base:eq:antisymmetric-bound} and
\eqref{base:eq:symmetric-bound} give
\[
 p(p-1)+5+2u-1
 =2m-2+(p-2)(p-3)\geq2m-2.
\]
If $p\geq2$ and $u=0$, the difference between the bound
$p(p-1)+5$ and $2m-2=4p-2$ is $p^2-5p+7>0$.

For a general $G$, work in the finitely generated subgroup generated by
$R$ and put $K=\FS(R)$. A finitely generated abelian group is a product
of a finite group and a free abelian group; reduction of the free coordinates
modulo a sufficiently large integer gives a finite quotient whose map is
injective on the finite set $K+K-K$.  It also preserves aperiodicity: a nonzero period
of the image of $K$ can be written as the image of $k_1-k_0$; injectivity on
$K+K-K$ then lifts every translated point back to $K$, producing a nonzero
period of $K$.  The finite-group result in this quotient proves
\eqref{base:eq:aperiodic-linear}.
\end{proof}

Order a positive set as $X_i=\{x_1<\cdots<x_i\}$ and put
\[
 d_i=\card{(x_i+\FS(X_{i-1}))\setminus\FS(X_{i-1})},
 \qquad y_i=d_i-i.
\]
Indeed, modulo $x_i$ the numbers $0,x_1,\ldots,x_{i-1}$ lie
in distinct residue classes. Each occupied class of $\FS(X_{i-1})$
has a largest point, whose translate by $x_i$ is new. Thus $d_i\ge i$,
so $y_i\ge0$, and telescoping gives
\begin{equation}
 \mu(X_i)=\sum_{j=1}^i y_j.
 \label{base:eq:increment-sum}
\end{equation}
We shall also use the exact equality case in Lemma~2.1 of
Carpenter--Defant--Kravitz~\cite{CDK2026}: for $i\geq4$,
\begin{equation}
 y_i=0\quad\Longleftrightarrow\quad
 X_i=t\{1,2,\ldots,i\}\quad\text{for some }t>0.
 \label{base:eq:zero-increment}
\end{equation}

\begin{lemma}[Unit-increment inverse lemma]\label{base:lem:unit-increment}
If $i\geq6$ and $y_i=1$, then there are $t>0$ and $q\in\{1,\ldots,i\}$
such that
\begin{equation}
 X_i=t\bigl(\{1,2,\ldots,i+1\}\setminus\{q\}\bigr).
 \label{base:eq:unit-increment}
\end{equation}
\end{lemma}

\begin{proof}
Put $x=x_i$, reduce $x_1,\ldots,x_{i-1}$ modulo $x$ in the circle group
$\mathbb R/x\mathbb Z$, and call the resulting $(i-1)$-element set $R$.
Let $\rho=\card{\FS(R)}$.  Distinct occupied residues give distinct maximal
$x$-chains in $\FS(X_{i-1})$, while
$0,x_1,\ldots,x_{i-1}$ occupy distinct residues.  Hence
\[
 i\leq\rho\leq d_i=i+1.
\]
If $\FS(R)$ were aperiodic, Lemma~\ref{base:lem:aperiodic-linear} would give
$\rho\geq2i-4>i+1$, a contradiction.

Let $H$ be the nontrivial stabilizer of $\FS(R)$, write
$h=\card H$ and $\rho=hq_0$, and put $v=\card{R\cap H}$.  The stabilizer $H$ is finite, since it acts faithfully by translations
on the finite nonempty set $\FS(R)$. Iterated Kneser applied to the
$(i-1)$ sets $\{0,r\}$ gives
\[
 \rho\geq h(i-v),\qquad v\leq h-1,\qquad h+q_0\geq i+1.
\]
If $q_0\geq2$, then
\[
 \rho=hq_0\geq2(h+q_0)-4\geq2i-2>i+1,
\]
again impossible.  Thus $q_0=1$, so $\FS(R)=H$ and $R\subseteq H$.
The circle group has a unique subgroup of order $h$; with $t=x/h$ and
representatives in $(0,x)$,
\[
 X_i=t\bigl(\{1,\ldots,h\}\setminus D\bigr),\qquad
 \card D=h-i,\qquad i\leq h\leq i+1,
\]
where $D\subseteq\{1,\ldots,h-1\}$.  The case $h=i$ is the homogeneous
progression and has $y_i=0$.  Therefore $h=i+1$ and $D=\{q\}$, proving
\eqref{base:eq:unit-increment}.
\end{proof}

\begin{lemma}[Uniform defect versus excess]\label{base:lem:defect-excess}
For every finite positive real set $X$,
\begin{equation}
 \boxed{2\delta_0(X)-\mu(X)\leq10.}
 \label{base:eq:defect-excess}
\end{equation}
\end{lemma}

\begin{proof}
Use the prefixes above and write $\delta_i=\delta_0(X_i)$ and
$\Phi_i=2\delta_i-\mu(X_i)$.  Adding $x_i$ can remove old defects and can
create only the one new defect $x_i$, so
\begin{equation}
 \Phi_i\leq\Phi_{i-1}+2-y_i.
 \label{base:eq:defect-recurrence}
\end{equation}
For $i\leq5$, the trivial bounds $\delta_i\leq i$ and $\mu(X_i)\geq0$
give $\Phi_i\leq10$.

Proceed by induction for $i\geq6$.  If $y_i\geq2$, equation
\eqref{base:eq:defect-recurrence} preserves the bound.  If $y_i=0$, the published
increment-equality case \eqref{base:eq:zero-increment} makes $X_i$ a homogeneous
progression; then $\delta_i=1$, $\mu(X_i)=0$, and $\Phi_i=2$.

It remains to consider $y_i=1$.  By Lemma~\ref{base:lem:unit-increment}, after
scaling we have $X_i=[1,i+1]_{\mathbb Z}\setminus\{q\}$.  Every present
$k\in[2,i-1]_{\mathbb Z}$ has two disjoint representations
\[
 2k=(k-1)+(k+1)=(k-2)+(k+2),
\]
where $0$ is allowed; deleting one $q$ cannot destroy both.  The largest
element $i+1$ is defective.  The only possible additional defect is $1$
when $q=2$, or $i$ when $q=i-1$.  Thus $\delta_i\leq2$, and
$\delta_i=2$ is possible only for these two values of $q$.

For $q=2$ and $q=i-1$, the integer totals are respectively
\[
 T_i+i-1\quad\hbox{and}\quad T_i+2,
 \qquad T_i=\frac{i(i+1)}2.
\]
Both exceed $T_i+1$.  If $\mu(X_i)\leq1$, the published stability theorem
with $M=\mu(X_i)\leq1\leq i-4$ would give an integer normalization of total
at most $T_i+1$.  Since both displayed sets contain $1$, every integer
normalization has total at least the displayed total, a contradiction.
Hence $\mu(X_i)\geq2$ whenever $\delta_i=2$.  In all cases
$\Phi_i\leq2$, completing the induction and the proof.
\end{proof}

Put $\mu_P=\mu(P)$, $\mu_Q=\mu(Q)$, and $\mu=\mu_P+\mu_Q$.  For finite
real sets $U,V$, the standard increasing chain gives
$\card{U+V}\geq\card U+\card V-1$.  Hence
\begin{equation}
 s=\card{\FS(P)-\FS(Q)}\geq\Lambda+\mu,
 \qquad
 \Lambda:=1+\frac{a(a+1)}2+\frac{b(b+1)}2.
 \label{base:eq:core-lower}
\end{equation}

The set $D\setminus E$ has eight elements.  If $x\in P$ contributes to
$\Delta(D)$, then $2x\notin R(P\cup\{0\})$; the analogous statement holds
on the negative side.  Counting the eight deleted elements and zero separately
therefore gives
\begin{equation}
 \Delta(D)\leq\delta_0(P)+\delta_0(Q)+9.
 \label{base:eq:delta-split}
\end{equation}
By Lemma~\ref{base:lem:defect-excess},
\begin{align}
 2\Delta(D)-\mu
 &\leq(2\delta_0(P)-\mu_P)+(2\delta_0(Q)-\mu_Q)+18\notag\\
 &\leq10+10+18=38.
 \label{base:eq:joint-rebate}
\end{align}
Combining \eqref{base:eq:core-upper}, \eqref{base:eq:core-lower}, and
\eqref{base:eq:joint-rebate} yields
\begin{equation}
 \Lambda\leq22d+115.
 \label{base:eq:lambda-upper}
\end{equation}

For fixed $a+b=d-9$, $\Lambda$ is minimized when $a$ and $b$ are as balanced
as possible, so
\[
 \Lambda\geq\frac{(d-8)^2+3}{4}.
\]
Consequently
\[
 \frac{(d-8)^2+3}{4}\leq22d+115,
 \qquad\text{or equivalently}\qquad
 d^2-104d-393\leq0.
\]
The quadratic is $39>0$ at $d=108$ and is strictly increasing for
$d\geq108$.  Hence $d\leq107$.  Since $\card B\leq8$,
\[
 \card A=d+\card B\leq107+8=115.
\]
We have proved this bound for zero-containing real sets. Applying
Lemma~\ref{base:lem:add-zero} to a set that omits zero yields the same bound
(and in fact the bound $114$ for such a set). Thus no cardinality greater
than $115$ can occur.

\begin{theorem}[Preliminary global bound]\label{base:thm:upper115}
Every finite real feasible set has at most $115$ elements.
\end{theorem}

The following sections sharpen the structural estimates before excluding
$108,\ldots,115$ in one argument.

\begin{lemma}[Small-increment inverse lemma]\label{base:lem:small-increment}
Let $i\geq6$ and $0\leq y_i\leq i-5$.  There are an integer $h$ and a set
$D_0\subseteq[1,h-1]_{\mathbb Z}$ such that, for some $t>0$,
\[
 i\leq h\leq i+y_i,\qquad \card{D_0}=h-i,\qquad
 X_i=t\bigl([1,h]_{\mathbb Z}\setminus D_0\bigr).
\]
\end{lemma}

\begin{proof}
Repeat the circle-group argument in Lemma~\ref{base:lem:unit-increment}, now with
$d_i=i+y_i$.  If the projected subset-sum set were aperiodic, its cardinality
$\rho$ would satisfy
\[
 \rho\geq2i-4>i+y_i.
\]
For a nontrivial stabilizer, retain the notation $h,q_0,v$ from that proof.
If $q_0\geq2$, then
\[
 \rho=hq_0\geq2(h+q_0)-4\geq2i-2>i+y_i.
\]
Thus $q_0=1$, the projected subset sums form the subgroup of order
$h=\rho$, and $i\leq h\leq i+y_i$.  Taking the representatives in $(0,x_i)$
gives the asserted interval with $h-i$ deletions.
\end{proof}



\clearpage\part{The descent from 115 to 96}\label{part:upper96}

\section{Scope and the descending implication}
The preceding chapter proves integralization, the covering and peeling
lemmas, the positive-prefix increment estimates, and $N\leq115$.
We now exclude every integer cardinality from 115 down through 97.
At an endpoint $n$, adjoining zero is used only after $n+1$ has been
excluded. All inverse theorems are applied after integralization.

The finite verifications below have two logically separate parts.
First, the arguments assign every actual feasible set to an explicitly
listed linear system and prove that each branch covers its parent.
Second, nonnegative integer combinations contradict every resulting
system. The complete smaller multiplier pools are printed, and the
larger pools and trees are identified exactly in Appendix~\ref{u96:app:upper96}.
A numerical optimization status is never a premise of this proof.

\section{Two stronger defect estimates}\label{u96:sec:new-defect}
For a positive set $X$, write
\[
\Phi(X)=2\delta_0(X)-\mu(X),\qquad
\rhoD(X)=2-\Phi(X)=\mu(X)+2-2\delta_0(X).
\]
The uniform estimate already proved is $\Phi\leq10$.

\begin{lemma}[Deleted intervals]\label{u96:lem:deleted}
Let $X=[1,r+q]_{\mathbb Z}\setminus D$, where $|D|=q$, $r+q\notin D$,
and $r\geq\max(5,4q+1)$. Then $\Phi(X)\leq2$.
For $q=1,2$ and $r\geq7$, one also has $\delta_0(X)\leq2q$.
\end{lemma}
\begin{proof}
The case $q=0$ is immediate. Put $h=r+q$ and
$C=\{h-d:d\in D\}$. The integer total excess is
\[
e=\sum X-T_r=\sum_{c\in C}c-T_{q-1}.
\]
There are $\min(k,h-k)$ disjoint representations of $2k$ by distinct
members of $[0,h]_{\mathbb Z}$ when $k\ne h/2$; for $k=h/2$, use $0+h$.
Therefore a defective present $k$ is either in $[1,q]$ or in $[r,h]$.
There are at most $q+1$ high defects. For $q\geq2$ there are at most $q-1$
low defects: having all of $1,\ldots,q$ present and defective would make
2 both absent (from defectivity of 1) and present. For $q=1$, the direct
one-deletion argument above gives at most two defects. These observations
prove $\delta_0(X)\leq2q$ and the last assertion.

Since $r>q$, no integer greater than one can divide all $r$ elements
of this interval with $q$ deletions; hence $X$ has gcd one.
If $\mu(X)\leq r-4$, the stability theorem in Section~\ref{base:sec:upper}
implies $\mu(X)=e$: any integer normalization of a primitive integer set
is an integer dilation, and the support lies in $[0,\sum X]$.
Suppose first that there are no low defects. Apart from $h$, let the $m$
high defects be $h-t$. Their presence and defectivity imply
$t\notin C$ and $2t\in C$. Replacing these distinct depths $2t$ by their
absent halves $t$ leaves $q$ distinct positive integers. Hence
\[
e\geq q+\sum t\geq q+m\geq2m=2\delta_0(X)-2.
\]
If there are $\ell\geq1$ low defects, their forced holes $2k$ have distinct
depths at least $r-q$, spaced by at least two. Thus
\[
e\geq\ell(r-q)+\ell(\ell-1)+T_{q-\ell}-T_{q-1}.
\]
Subtracting $2q+2\ell$ gives
\[
\ell(r-2q-2)+\tfrac32\ell(\ell-1)-q\geq0
\]
under $r\geq4q+1$. Since $\delta_0(X)\leq q+1+\ell$, this again implies
$e\geq2\delta_0(X)-2$.
Finally, if $\mu(X)>r-4$, then
\[
\Phi(X)\leq4q-(r-3)\leq2.
\]
All cases give the claim.
\end{proof}

\begin{theorem}[Threshold 14]\label{u96:thm:phi14}
For every positive real set $X$ of size $r\geq14$,
\[
\boxed{\Phi(X)\leq2,\qquad\rhoD(X)\geq0.}
\]
\end{theorem}
\begin{proof}
For $9\leq i\leq14$, put $q_i=\lfloor(i-1)/4\rfloor$.
If $y_i\leq q_i$, the small-increment inverse lemma expresses $X_i$, up to
scaling, as an initial integer interval with at most $q_i$ deletions.
Lemma~\ref{u96:lem:deleted} gives $\Phi(X_i)\leq2$.
If $y_i\geq q_i+1$, the defect recurrence gives
$\Phi(X_i)\leq\Phi(X_{i-1})+1-q_i$.
Starting from the uniform bound 10 at size 8 gives the following successive
upper bounds:
\[
\begin{array}{c|rrrrrrr}
i&8&9&10&11&12&13&14\\\hline
\Phi(X_i)\text{ upper bound}&10&9&8&7&6&4&2.
\end{array}
\]
After size 14, increments $y_i\geq2$ preserve the bound. For $y_i=0$ or 1,
the equality and unit-increment arguments in Section~\ref{base:sec:upper}
already give $\Phi(X_i)\leq2$.
\end{proof}

\begin{theorem}[Penalty for a non-progression support]\label{u96:thm:nonap}
If $X$ is a positive real set of size $r\geq6$ and $\FS(X)$ is not an
arithmetic progression, then
\[
\boxed{\rhoD(X)\geq r-14.}
\]
\end{theorem}
\begin{proof}
At $r=6$ the uniform bound gives $\rhoD\geq-8=r-14$.
Induct on $r\geq7$. If $y_r\leq2$, the small-increment inverse lemma gives
a primitive deleted interval with at most two deletions. It has at most
four defects. Its non-progression support implies $\mu\geq r-3$ by the
stability theorem, since $\mu\leq r-4$ would force a full interval in the
primitive normalization. Consequently $\rhoD\geq r-9\geq r-14$.

Now let $y_r\geq3$. If the preceding support is not an arithmetic
progression, the induction hypothesis and
$\rhoD(X_r)\geq\rhoD(X_{r-1})+y_r-2$ give the result.
If the preceding support is $t[0,L]_{\mathbb Z}$, its translate by $x_r$
must be disjoint from it; an overlap would produce an arithmetic progression
again. Hence
\[
y_r=|\FS(X_{r-1})|-r\geq\frac{(r-1)(r-2)}2.
\]
Using $\rhoD(X_{r-1})\geq-8$ gives
$\rhoD(X_r)\geq (r-1)(r-2)/2-10\geq r-14$.
\end{proof}

\section{A sharper effective hole threshold}\label{u96:sec:holes}
For a finite integer set $S$, its hull-hole count is
$\max S-\min S+1-|S|$. For a set on a specified lattice $u\mathbb Z$,
divide by $u$ before counting holes.

We use Lev's central-interval theorem~\cite{Lev1998}: if
$B\subseteq[1,L]_{\mathbb Z}$, $|B|=s$, and $L\leq3s/2-2$, then
\[
[2L-2s+1,\ \sum B-(2L-2s+1)]_{\mathbb Z}\subseteq\FS(B).
\]
Write $\mathcal S(L,H)=[0,L]_{\mathbb Z}\setminus(H\cup(L-H))$.
The endpoint bridge proved in Lemma~\ref{u96:lem:bridge-proof} states that an
integer prefix of size $m$ with support $\mathcal S(L,H)$ preserves the same
endpoint holes upon adjoining its next element with excess at most $E$, if
$2h<L$ and
\[
m\geq M(h,q,E):=
\max\left\{h,1+\left\lfloor
\frac{1+\sqrt{9+8(E+2q+2h)}}2\right\rfloor\right\},
\quad h=\max H,\quad q=|H|.
\]
Here $h=q=0$ is allowed. We need this explicit bridge, not an unspecified
``sufficiently large'' threshold from a classification theorem.

\begin{lemma}[Small increments with counted holes]\label{u96:lem:counted-holes}
If $i\geq2E+5$ and $0\leq y_i=e\leq E$, then
\[
X_i=t([1,i+e]_{\mathbb Z}\setminus D),\qquad |D|=e,
\]
and, after division by $t$,
\[
\FS(t^{-1}X_i)=\mathcal S(L,H),\qquad
H\subseteq[1,2e]_{\mathbb Z},\qquad |H|\leq e.
\]
\end{lemma}
\begin{proof}
For $e=0$, the increment equality case gives the assertion directly.
Assume $e>0$. The small-increment inverse lemma first gives the same form with $\delta$
deletions, $0\leq\delta\leq e$, and largest coefficient $i+\delta$.
Remove that coefficient. The remaining $i-1$ integers lie in
$[1,i+\delta-1]$ and satisfy Lev's density condition, because
$i\geq2\delta+5$. Their support therefore has endpoint state
$H\subseteq[1,2\delta]$. Its total $L_0$ satisfies
\[
L_0\geq T_{i-1}>i+\delta+4\delta.
\]
Adding $i+\delta$ overlaps the two central intervals and leaves exactly
the same endpoint holes. Its support increment is $i+\delta$; comparison
with $i+e$ gives $\delta=e$.
Finally every lower hole in $[1,2e]$ is absent from the coefficient set
itself. Since $2e<i+e$ and there are exactly $e$ deleted coefficients,
$H\subseteq D$ and $|H|\leq e$. This last count halves the former bound
$|H|\leq2e$.
\end{proof}

\begin{theorem}[Effective hole bound]\label{u96:thm:holes103}
Let $X$ be a positive real set of size $r\geq35$ with
$\FS(X)\subseteq u\mathbb Z$. If its hull-hole count $h$ satisfies
\[
14<h<|\FS(X)|-1,
\]
then
\[
\boxed{\mu(X)\geq r+103.}
\]
\end{theorem}
\begin{proof}
Suppose $\mu(X)\leq r+102$. Divide by $u$ to obtain an integer set $B$.
It is primitive: a gcd of at least two would leave at least one missing
lattice point between each two successive support values, contradicting
$h<|\FS(X)|-1$. The support is not a full integer interval.

We first obtain a tail of positive increments with a uniform budget.
If no prefix of size at least three has an arithmetic-progression support,
then $y_j\geq1$ for $j\geq4$, and
$\sum_{j=4}^r(y_j-1)\leq105$.
Otherwise let $k\geq3$ be the last size with such a support. The next
translate is disjoint, so
\[
y_{k+1}\geq T_k-k,
\qquad
\mu(B)\geq r-1+\frac{k(k-3)}2.
\]
Thus $k(k-3)/2\leq103$, whence $k\leq15$. All later increments are
positive, and
$\sum_{j=k+1}^r(y_j-1)\leq102+k\leq117$.
In either case there is $r_0\leq16$ such that
\begin{equation}
y_j\geq1\ (j\geq r_0),\qquad
\sum_{j=r_0}^r(y_j-1)\leq117.\label{u96:eq:tail117}
\end{equation}

The window $19\leq j\leq r$ has at least 17 indices. If every increment
there were at least eight, it would contribute at least $7\cdot17=119$
to \eqref{u96:eq:tail117}. Therefore some $i_0\geq19$ has $1\leq y_{i_0}=e\leq7$.
Lemma~\ref{u96:lem:counted-holes} gives
\[
B_{i_0}=t([1,i_0+e]_{\mathbb Z}\setminus D),\quad
H\subseteq[1,14],\quad |H|\leq7.
\]
The coefficient set has gcd one since $i_0>e$, so $t$ is a positive integer.
If $t>1$, primitivity of $B$ gives a first later coefficient outside
$t\mathbb Z$. At that step the two supports are disjoint, and
\[
y_j-1\geq\frac{i_0(i_0-1)}2-1\geq170>117,
\]
contradicting \eqref{u96:eq:tail117}. Hence $t=1$.

Every subsequent increment is at most 118, and
\[
M(14,7,118)=19,\qquad T_{19}=190>28.
\]
The endpoint bridge now preserves $H$ through the whole remaining tail.
The final number of holes is $2|H|\leq14$, contradicting $h>14$.
\end{proof}

\section{A common argument for all eight endpoints}\label{u96:sec:endpoints}
\begin{lemma}[Unpaired absolute values]\label{u96:lem:unpaired}
Let $0\in F$, let $e=\card F$, and write the distinct nonzero absolute values
as $0<b_1<\cdots<b_r$.  Let $m_i\in\{1,2\}$ count how many of
$\{b_i,-b_i\}$ lie in $F$.  If
\[
 p=\card{\{i:m_i=2\}},\qquad q=\card{\{i:m_i=1\}},
\]
and
\[
 J=\card{\{(i,j):i<j,\ m_i=1,\ m_j=2\}},
\]
then
\begin{equation}
 \card{\FS(F)}\geq\frac{e^2+3+q^2}{4}+J.
 \label{u96:eq:unpaired}
\end{equation}
\end{lemma}

\begin{proof}
Translate $\FS(F)$ by the sum of the absolute values of the negative members
of $F$.  The resulting set is
\[
 \left\{\sum_{i=1}^r k_i b_i:0\leq k_i\leq m_i\right\}.
\]
Before $b_i$ is introduced, the residues represented modulo $b_i$ by
$0,b_1,\ldots,b_{i-1}$ are distinct.  Extending each corresponding maximal
$b_i$-chain by the $m_i$ new coefficient levels creates at least $i m_i$
new sums.  Therefore
\[
 \card{\FS(F)}\geq1+\sum_{i=1}^r i m_i.
\]
If the paired positions are $s_1<\cdots<s_p$, then
\[
 \sum_{i=1}^r i m_i
 =\frac{r(r+1)}2+\sum_{j=1}^p s_j
 =\frac{r(r+1)}2+\frac{p(p+1)}2+J.
\]
Using $r=p+q$ and $e=2p+q+1$ gives \eqref{u96:eq:unpaired}.
\end{proof}

\begin{lemma}[Four-extreme restricted-sum defect]
\label{u96:lem:four-extreme-defect}
Let
\[
 F=\{x_1<\cdots<x_m\}\subseteq\mathbb Z,\qquad m\geq2,
\]
and let
\[
 a_1<a_2<x_1<\cdots<x_m<b_2<b_1,
 \qquad D=F\cup\{a_1,a_2,b_2,b_1\}.
\]
Let $\eta_F,\eta_D$ be the hole counts of $R(F),R(D)$ in their endpoint
integer hulls, and put
\[
 g_H=4m+6-\bigl(\card{R(D)}-\card{R(F)}\bigr).
\]
Then
\begin{equation}
 \boxed{\eta_D+\eta_F+g_H\geq2m-8.}
 \label{u96:eq:four-extreme-defect}
\end{equation}
\end{lemma}

\begin{proof}
Put $L=x_1+x_2$, $U=x_{m-1}+x_m$, and let
\[
 h_-=L-(a_1+a_2),\qquad
 h_+=(b_1+b_2)-U,\qquad h=h_-+h_+.
\]
Let $W_+$ be the union of the portions of $b_1+F$ and $b_2+F$ strictly
above $U$.  Write
\[
 y=b_2,\quad z=b_1,\quad u=y-x_m,\quad v=z-y,
 \quad e=x_m-x_{m-1}.
\]
All three last quantities are positive and $h_+=2u+v+e$.  If $p,q$ count
the terms above $U$ in $y+F,z+F$, respectively, then adjacency of
$x_{m-1},x_m$ gives $p\leq u+1$, while $q\leq m$.  Also
$W_+\subseteq[U+1,z+x_m]_{\mathbb Z}$.  Combining the cardinality and
interval bounds gives
\[
 2\card{W_+}\leq h_++m+1.
\]
The symmetric lower-end set $W_-$ satisfies
$2\card{W_-}\leq h_-+m+1$.

Let $N_{\rm out}$ count the new restricted sums outside $[L,U]$.  Besides
$W_+$ and $W_-$ only the six pairs among the four new elements can occur, so
\[
 2N_{\rm out}\leq h+2m+14.
\]
If $N=\card{R(D)}-\card{R(F)}$, at most $\eta_F$ further new values can lie
inside the old hull; hence $N\leq N_{\rm out}+\eta_F$.  Since
\[
 \eta_D=h+\eta_F-N,
 \qquad g_H=4m+6-N,
\]
we obtain
\[
 \eta_D+\eta_F+g_H
 =h+2\eta_F+4m+6-2N
 \geq h+4m+6-2N_{\rm out}\geq2m-8.
\]
\end{proof}

\begin{lemma}[Sharper four-extreme defect]\label{u96:lem:four-sharp}
In the notation of Lemma~\ref{u96:lem:four-extreme-defect},
\[
 \boxed{\eta_D+\eta_F+g_H\geq2m-4.}
\]
Consequently, whenever the common-lattice and overlapping-translate
reductions have been established, the lower bounds for $T$ improve to
(the slack quantities are defined in Proposition~\ref{u96:prop:universal-lattice})
\begin{equation}
 T\geq\kappa+G-3z\geq G-2z,\qquad
 T\geq\lceil G/3\rceil,\qquad G=2d-12.
 \label{u96:eq:sharp-G}
\end{equation}
\end{lemma}
\begin{proof}
Keep $L=x_1+x_2$, $U=x_{m-1}+x_m$, $h_-,h_+$ and $W_-,W_+$ from that
lemma. Define the integer intervals
\[
 I_+=(U,b_1+x_m]\cap\mathbb Z,\qquad
 I_-=[a_1+x_1,L)\cap\mathbb Z,
\]
and their unused capacities $H_\pm=|I_\pm|-|W_\pm|\geq0$.
For the upper end, put $u=b_2-x_m$, $v=b_1-b_2$, and
$e=x_m-x_{m-1}$. The two bounds used previously are
\[
 |W_+|\leq u+1+m,\qquad
 |W_+|=u+v+e-H_+.
\]
Adding gives $2|W_+|\leq h_++m+1-H_+$; reflection gives the lower-end
version. This retains the capacity discarded in the earlier proof.

Every mixed new-new sum $a_i+b_j$ outside $[L,U]$ lies in $I_-\cup I_+$.
Let $q$ count the distinct such values not already in $W_-\cup W_+$.
Then
\[
 q\leq4,\qquad q\leq H_-+H_+.
\]
The other two new-new sums, $a_1+a_2$ and $b_1+b_2$, lie beyond these
intervals. Thus the number of new restricted sums outside the old hull is
$N_{\rm out}=|W_-|+|W_+|+2+q$, and
\[
 2N_{\rm out}
 \leq h_-+h_++2m+6-(H_-+H_+)+2q
 \leq h_-+h_++2m+10.
\]
As before, $N=|R(D)|-|R(F)|\leq N_{\rm out}+\eta_F$. Substitution into
\[
 \eta_D+\eta_F+g_H=h_-+h_++2\eta_F+4m+6-2N
\]
proves $2m-4$. Apply this with $F=D_1$, $m=d-4$, and use
$\eta_2\leq z+\ell_2$, $\eta_1\leq2z+\ell_1$ to obtain
\eqref{u96:eq:sharp-G}. Here $z\leq\kappa\leq T$, so $3T\geq G$.
\end{proof}

\begin{lemma}[Universal core-span estimate]\label{u96:lem:span96}
If $X\subset\mathbb Z_{>0}$ has $r\geq15$ elements, then
\[
 \max X\leq r+\rhoD(X)+h(X)+2.
\]
Neither primitivity nor full subset-sum support is required.
\end{lemma}
\begin{proof}
Let $X'$ be the first $r-1$ elements, write $\max X=r+q$, and let $v$
be the number of old diagonal defects destroyed by adjoining this maximum.
The new maximum is defective, so
$\delta_0(X)=\delta_0(X')+1-v$. Also
\[
 \mu(X)-\mu(X')=q-h(X)+h(X').
\]
It follows exactly that
\[
 \rhoD(X)=\rhoD(X')+q-h(X)+h(X')-2+2v.
\]
Here $\rhoD(X')\geq0$ because $r-1\geq14$, and $h(X'),v\geq0$.
Rearranging proves the claim.
\end{proof}

\begin{proposition}[One common-lattice reduction for all endpoints]\label{u96:prop:universal-lattice}
Let $97\leq n\leq115$, and suppose a feasible $n$-element set contains
zero. After integralization and an existence-preserving dilation there
are sets
\[
 A=B\sqcup D,\qquad D\supset D_1\supset E=P\cup\{0\}\cup(-Q)
\]
with $|B|=8$, $d=|D|=n-8$, $|D_1|=d-4$, $e=|E|=n-16$,
and $a=|P|,b=|Q|\geq17$. The set $D$ is integral on one common lattice.
Put $s=|\FS(E)|$, $\Lambda=1+T_a+T_b$, and define
\[
\begin{aligned}
 \kappa&=s-\Lambda-\mu(P)-\mu(Q),&
 \rho_P&=\rhoD(P),\quad\rho_Q=\rhoD(Q),\\
 \vartheta&=\delta_0(P)+\delta_0(Q)+8-\Delta(D),&
 c&=2|2D|+18d+85-|\FS(D)|,\\
 \ell_1&=|\FS(D)|-|\FS(D_1)|-|R(D)|+1,&
 \ell_2&=|\FS(D_1)|-s-|R(D_1)|+1,\\
 g&=4d-10-(|R(D)|-|R(D_1)|),&
 T&=\kappa+\ell_1+\ell_2+g.
\end{aligned}
\]
All these quantities are nonnegative, and the exact budget is
\begin{equation}
 \rho_P+\rho_Q+2\vartheta+c+T=22d+97-\Lambda=:\Omega.
 \label{u96:eq:budget-new}
\end{equation}
Let $h_P,h_Q,z,w$ be the hull-hole counts of $\FS(P),\FS(Q),
\FS(E),\FS(D)$, respectively, on this lattice. Then
\begin{gather}
 h_P,h_Q\leq\kappa,\qquad z=h_P+h_Q-\kappa,\qquad
 0\leq z\leq\min(h_P,h_Q)\leq\kappa,\label{u96:eq:holes-core-new}\\
 0\leq w\leq4z,\qquad
 T\geq\kappa+G-3z\geq G-2z,\qquad
 T\geq\lceil G/3\rceil,\qquad T\geq z,\quad G=2d-12.
 \label{u96:eq:universal-holes}
\end{gather}
If $z>0$, both one-sided supports are non-progressions and
\[
 \rho_P+\rho_Q\geq a+b-28=n-45,\qquad
 \mu(P)+\mu(Q)\geq a+b-6=n-23.
\]
\end{proposition}
\begin{proof}
The notation $B,D,D_1,E$ is obtained by the covering and two peeling
steps. If $|B|\leq7$, one sign is absent from $D$; the nonzero part of
$E$ is then one-sided and has at least $n-16\geq81$ elements. Hence
\[
 s\geq1+T_{81}=3322>24\cdot115+77\geq24d+77,
\]
contradicting the preliminary core bound. Thus $|B|=8$.
The unpaired-value lemma and $s\leq24e+269$ give
\[
 q^2\leq96e+1073-e^2\leq2288\qquad(81\leq e\leq99).
\]
The number of opposite pairs is $(e-1-q)/2\geq17$; the endpoint $e=81$
uses the fact that $q$ is even. Thus $a,b\geq17$ and $E$ contains an
opposite pair. Zero is consequently not a diagonal exception in $D$,
so $\Delta(D)\leq\delta_0(P)+\delta_0(Q)+8$.
The covering, peeling and restricted-pair estimates prove nonnegativity
of the displayed slacks; $\rho_P,\rho_Q\geq0$ follows from
Theorem~\ref{u96:thm:phi14}. Substituting their definitions and cancelling
$|2D|=|R(D)|+\Delta(D)$ gives \eqref{u96:eq:budget-new} exactly.

\emph{The inverse theorem on the core.}
We use the following sufficient form of Stanchescu's two-set theorem
\cite[Corollary~2]{Stanchescu1996}. If finite integer sets $U,V$, each
of size at least two, satisfy
\[
 |U+V|<|U|+|V|+\min(|U|,|V|)-3,
\]
then their hull cardinalities on a common difference lattice are at most
$|U+V|-|V|+1$ and $|U+V|-|U|+1$, respectively.
Apply it to $U=\FS(P)$ and $V=-\FS(Q)$ after integralization.
Write $m=a+b=n-17$, $(a,b)=(\lfloor m/2\rfloor-k,
\lceil m/2\rceil+k)$ up to reflection, and
\[
 I=k^2+(m\bmod2)k,\qquad
 \Lambda_0=1+T_{\lfloor m/2\rfloor}+T_{\lceil m/2\rceil},
 \qquad\Omega_0=22d+97-\Lambda_0.
\]
Then $\Lambda=\Lambda_0+I$, $\Omega=\Omega_0-I$.
Nonnegativity gives $I\leq\Omega_0$. The table below includes every
possible integer $k$, from zero to its stated maximum. Its inverse
margin is the minimum of $T_{\lfloor m/2\rfloor-k}-2-(\Omega_0-I)$
over precisely that range. It is strictly positive, so
$\kappa\leq\Omega<T_{\min(a,b)}-2\leq\min(|U|,|V|)-3$.
This is a slightly stronger sufficient condition than required by the
inverse theorem. Both $U,V$ contain zero; after division by the common
step they, and hence $P,Q$, are integral. The two hull bounds give
$h_P,h_Q\leq\kappa$. Comparing the hull of $U+V$ with the two individual
hulls gives $z=h_P+h_Q-\kappa$, proving \eqref{u96:eq:holes-core-new}.

\emph{Forcing both peeled layers onto that lattice.}
Only $P,Q,E$ are asserted integral at this stage. Set
$\rho=\rho_P+\rho_Q$, $\mu=\mu(P)+\mu(Q)$ and
$R_p=\max P,R_q=\max Q$. Lemma~\ref{u96:lem:span96}, applied to the
integer core sides, gives
\[
 R_p+R_q\leq m+\rho+h_P+h_Q+4=m+\rho+\kappa+z+4.
\]
Since $R(E)$ is integral, $|R(E)|\leq2(R_p+R_q)+1$.
Regardless of whether the four elements of $D_1\setminus E$ are on
the lattice, their restricted pairs add at most $4e+6$ values. Thus,
using $s=\Lambda+\mu+\kappa$ and
$C=2m+4e+14-\Lambda$, we obtain
\begin{align*}
 |\FS(D_1)|-2s
 &\leq C+2\rho-\mu+\kappa+2z+\ell_2\\
 &\leq C+3\Omega-\rho-\mu.
\end{align*}
The last inequality uses $z\leq\kappa$ and
$\rho+\kappa+\ell_2\leq\Omega$ once, keeping the same quantities.
If $z=0$, the first line is at most $C+2\Omega$, because
$\rho\leq\Omega$ and $\mu\geq0$.
If $z>0$, each one-sided support has positive hull-hole count smaller
than its cardinality minus one. An integer arithmetic progression of
step at least two has at least that many holes, so neither support is
an arithmetic progression. The non-progression theorem gives
$\rho\geq m-28$. The stability theorem in the foundation chapter gives
$\mu(X)\geq r-3$ for a non-progression support of an $r$-element set:
otherwise $\mu(X)\leq r-4$ places its support in an interval of
exactly $1+T_r+\mu(X)$ integers, forcing it to fill that interval.
Hence $\mu\geq m-6$.

The two relevant balanced upper bounds are
$C_0+2\Omega_0$ and $C_0+3\Omega_0-(2m-34)$, where
$C_0=2m+4e+14-\Lambda_0$. Imbalance decreases them by $3I$ and $4I$.
Both are strictly negative in every row below. Therefore
$|\FS(D_1)|<2s$. Any element of $D_1\setminus E$ off the core lattice
would put two disjoint $s$-point translates of $\FS(E)$ inside
$\FS(D_1)$, a contradiction. This proves integrality of $D_1$.

For the other layer,
\[
 |\FS(D)|-2|\FS(D_1)|\leq4d-10+\ell_1-\ell_2-s
 \leq4d-10+\Omega-s<0.
\]
The last strict inequality follows from the last column of the table,
which gives $\Lambda_0-(4d-10+\Omega_0)>0$.
The same disjoint-translate argument puts all of $D$ on the lattice.

\begin{center}\small
\begin{tabular}{rrrrrrrr}
\toprule
$n$&$\Lambda_0$&$\Omega_0$&$k_{\max}$&Inv.&$z=0$&$z>0$&Layer 1\\\midrule
115&2451&0&0&1223&$-1845$&$-2007$&2033\\
114&2402&27&4&981&$-1748$&$-1881$&1961\\
113&2353&54&7&854&$-1651$&$-1755$&1889\\
112&2305&80&8&770&$-1557$&$-1633$&1819\\
111&2257&106&10&695&$-1463$&$-1511$&1749\\
110&2210&131&10&643&$-1372$&$-1393$&1681\\
109&2163&156&12&581&$-1281$&$-1275$&1613\\
108&2117&180&12&535&$-1193$&$-1161$&1547\\
107&2071&204&14&486&$-1105$&$-1047$&1481\\
106&2026&227&14&446&$-1020$&$-937$&1417\\
105&1981&250&15&408&$-935$&$-827$&1353\\
104&1937&272&16&371&$-853$&$-721$&1291\\
103&1893&294&17&335&$-771$&$-615$&1229\\
102&1850&315&17&299&$-692$&$-513$&1169\\
101&1807&336&18&264&$-613$&$-411$&1109\\
100&1765&356&18&230&$-537$&$-313$&1051\\
99&1723&376&19&196&$-461$&$-215$&993\\
98&1682&395&19&163&$-388$&$-121$&937\\
97&1641&414&20&131&$-315$&$-27$&881\\\bottomrule
\end{tabular}
\end{center}
All entries are evaluations of the displayed quadratic formulas; no
set search or optimization enters this table.

\emph{Hole propagation and the common lower budget.}
At the second peeling step write $\alpha_2,\beta_2$ for the sums of the
two removed negative and positive elements. Then $R(D_1)$ has endpoints
$\alpha_2,\beta_2$. The translates
$\alpha_2+\FS(E)$ and $\beta_2+\FS(E)$ are inside $\FS(D_1)$.
The strict $2s$ inequality proves they overlap. Their hulls cover the
new subset-sum hull and leave at most $2z$ holes.
If $u=\beta_2-\alpha_2$, their union has at least $s+u-z$ points.
Comparing with
$|\FS(D_1)|=s+|R(D_1)|-1+\ell_2$
gives $\eta_2\leq z+\ell_2$, where $\eta_2$ is the hull-hole count
of $R(D_1)$. The first peeling step similarly gives at most $4z$
subset-sum holes and $\eta_1\leq2z+\ell_1$ restricted-sum holes.
Thus
\begin{equation}
 \eta_2\leq z+\ell_2,\qquad\eta_1\leq2z+\ell_1.
 \label{u96:eq:prop-new}
\end{equation}
Lemma~\ref{u96:lem:four-sharp}, applied with $F=D_1$, $|F|=d-4$, gives
$\eta_1+\eta_2+g\geq2d-12=G$. Hence
$T\geq\kappa+G-3z\geq G-2z$.
Also $z\leq\kappa\leq T$, so $3T\geq G$ and $T\geq z$.
This proves all the claimed conclusions.
\end{proof}


\begin{theorem}[Excluding sizes 108 through 115]\label{u96:thm:upper107}
No feasible set containing zero has cardinality $108,109,\ldots,115$.
Consequently every feasible set has at most 107 elements.
\end{theorem}
\begin{proof}
Use Proposition~\ref{u96:prop:universal-lattice} and deliberately weaken its
sharp constant to $G=2d-16$. This preserves every inequality
$T+\rho_P+\rho_Q\leq\Omega$, $T\geq G-2z$ and
$T\geq\lceil G/3\rceil$ used below. Put $\rho=\rho_P+\rho_Q$.
\emph{The endpoint arithmetic.}
The minimum $\Lambda_0$ occurs when $a,b$ are as balanced as possible.
Set $\Omega_0=22d+97-\Lambda_0$, so $\Omega\leq\Omega_0$.
\begin{center}
\begin{tabular}{rrrrrrr}
\toprule
$n$&$d$&$e$&$\Lambda_0$&$\Omega_0$&$G$&$\lceil G/3\rceil$\\
\midrule
115&107&99&2451&0&198&66\\
114&106&98&2402&27&196&66\\
113&105&97&2353&54&194&65\\
112&104&96&2305&80&192&64\\
111&103&95&2257&106&190&64\\
110&102&94&2210&131&188&63\\
109&101&93&2163&156&186&62\\
108&100&92&2117&180&184&62\\
\bottomrule
\end{tabular}
\end{center}
The first three rows are already impossible since $T\leq\Omega_0$.
In all remaining rows $G>\Omega_0$, so \eqref{u96:eq:universal-holes} forces $z>0$.
Both $h_P,h_Q$ are then positive. They are at most 180 and are smaller than
$|U|-1,|V|-1$. Thus neither $U$ nor $V$ can itself be an arithmetic
progression: a progression of step at least two on this lattice has at
least $|U|-1$, respectively $|V|-1$, holes.
Theorem~\ref{u96:thm:nonap} therefore gives
\begin{equation}
\rho_P+\rho_Q\geq a+b-28=n-45.\label{u96:eq:nonap-budget}
\end{equation}
For $n=112$, this leaves $T\leq80-67=13<64$;
for $n=111$, it leaves $T\leq106-66=40<64$.

Only $108,109,110$ remain. Combining \eqref{u96:eq:universal-holes} with the exact
imbalance in $\Lambda$ gives the following bounds:
\begin{center}
\small
\begin{tabular}{rrrrrr}
\toprule
$n$&$\min(a,b)$&$T$ upper bound&$z$ lower bound&
$\rho_P+\rho_Q$ lower bound&$T+\rho_P+\rho_Q$\\
\midrule
110&$\geq39$&66&61&117&$\geq180>131$\\
109&$\geq37$&92&47&118&$\geq180>156$\\
108&$\geq35$&117&34&119&$\geq181>180$\\
\bottomrule
\end{tabular}
\end{center}
For completeness, write $(a,b)=(46-k,47+k)$ at 110,
$(46-k,46+k)$ at 109, and $(45-k,46+k)$ at 108, allowing reflection.
The respective excesses $\Lambda-\Lambda_0$ are $k(k+1),k^2,k(k+1)$.
The requirement $\Omega\geq\lceil G/3\rceil$ bounds $k$ by $7,9,10$,
which proves the size column. The next two columns follow from
$T\leq\Omega_0-(n-45)$ and $2z\geq G-T$.

In every row $h_P,h_Q\geq z>14$. Theorem~\ref{u96:thm:holes103} applies to
each side and gives $\mu(P)\geq a+103$ and $\mu(Q)\geq b+103$.
Using $\delta_0(P)\leq a$ and $\delta_0(Q)\leq b$ gives
\[
\rho_P+\rho_Q\geq(105-a)+(105-b)=227-n.
\]
This proves the strong lower-bound column. The last column contradicts
\eqref{u96:eq:budget-new} in each case, including a strict one-unit contradiction
at 108. All eight endpoints are excluded.


The preliminary bound $N\leq115$ and these eight exclusions give the
claim for sets containing zero. A feasible set omitting zero extends by
zero to a larger feasible set, so it too has at most 107 elements.
\end{proof}


\section{An elementary endpoint bridge and a parameterized hole test}\label{u96:sec:audit}
\begin{lemma}[Elementary endpoint bridge]\label{u96:lem:bridge-proof}
The endpoint bridge stated in Section~\ref{u96:sec:holes} follows from its explicit threshold $M(h,q,E)$.
\end{lemma}
\begin{proof}
Let $S=\mathcal S(L,H)$ be the support of $m$ positive integers. Then
$L\geq T_m$, $|S|=L+1-2q$, and $x>\max B\geq m\geq h$.
If $x\geq L-2h$, the part of $x+S$ above $L$ has at least
$\min(x,L+1)-2q\geq L-2h-2q$ points. But the definition of $M$ gives
\[
 T_m-2h-2q>m+1+E,
\]
contradicting the permitted increment. Thus $h<x<L-2h$.
The integer intervals $[h+1,L-h-1]$ and $[x+h+1,x+L-h-1]$
overlap or abut, because the integral $x<L-2h$. They cover the
whole central interval $[h+1,L+x-h-1]$. Since $x>h$, the lower holes stay $H$
and the upper holes are their reflections about $(L+x)/2$.
Consequently $S\cup(x+S)=\mathcal S(L+x,H)$, its cardinality increases by
$x$, and $x=m+1+e_x$ follows. This also supplies a short self-contained proof
of the bridge used in the earlier hole estimate.
\end{proof}

\begin{lemma}[Parameterized effective hole test]\label{u96:lem:hole-parameter}
Let $C\geq0$ and $b\geq1$ be integers. Define
\[
 K=\max\{k\geq3:k(k-3)/2\leq C+1\},\quad
 D=C+\max(3,K),\quad E=D+1.
\]
Suppose integers $J,r$ satisfy
\begin{gather*}
 J\geq\max\{K+1,2b+5,M(2b,b,E)\},\\
 (r-J+1)b>D,\qquad T_{J-1}-1>D.
\end{gather*}
If a positive $r$-element set has, on a specified lattice, a hull-hole count
$h$ satisfying $2b<h<|\FS|-1$, then $\mu\geq r+C+1$.
\end{lemma}
\begin{proof}
Assume instead $\mu\leq r+C$, and divide by the lattice step. The strict
upper bound on $h$ makes the integer set primitive, and its support is not
an arithmetic progression. The tail argument of
Theorem~\ref{u96:thm:holes103}, with $102$ replaced by $C$, gives a starting
index at most $K+1$ (or 4 if no prefix of size at least three has progression
support) after which
\[
 y_j\geq1,\qquad \sum(y_j-1)\leq D.
\]
Here a last progression prefix of size $k$ costs at least
$r-1+k(k-3)/2$, so $k\leq K$; without such a prefix the budget is $C+3$.
The window $J,\ldots,r$ contains some increment $1\leq y_i\leq b$.
Lemma~\ref{u96:lem:counted-holes} gives a deleted interval, with endpoint holes
in $[1,2b]$ and at most $b$ lower holes. Its scale is an integer. A first
later coefficient off this scale would cost
$y_j-1\geq T_{i-1}-1\geq T_{J-1}-1>D$; primitivity therefore makes the
scale one. Every later increment is at most $E$. The endpoint bridge
preserves the holes, since $i\geq J$ and $L\geq T_J>4b$.
The final hole count is at most $2b$, a contradiction.
\end{proof}

The following are exact applications, with no asymptotic size convention:
\begin{center}
\begin{tabular}{rrrrrrrl}
\toprule
$r\geq$&$C$&$b$&$K$&$D$&$J$&$M(2b,b,D+1)$&Conclusion when $2b<h<|\FS|-1$\\
\midrule
35&102&7&15&117&19&19&$\mu\geq r+103$\\
36&111&8&16&127&21&20&$\mu\geq r+112$\\
37&118&8&17&135&21&20&$\mu\geq r+119$\\
\bottomrule
\end{tabular}
\end{center}
For example, the last two window costs are $16\cdot8=128>127$ and
$17\cdot8=136>135$, while $T_{20}-1=209$ exceeds both budgets.

\section{The complete-core reduction at size 107}
\label{u96:sec:107-structure}
The following reduction supplies the complete core used to exclude size 107.
\begin{theorem}[The complete-core reduction at size 107]\label{u96:thm:107-core}
A hypothetical feasible 107-element set has an existence-preserving
normalization in which $D,D_1,E$ are integral, their subset-sum supports
are full intervals, and
\[
 (a,b)=(45-k,45+k)\quad\text{up to reflection},\qquad
 0\leq k\leq4,\qquad \rho_P+\rho_Q+\kappa+c\leq22-k^2.
\]
\end{theorem}
\begin{proof}
Size 108 is already excluded, so a feasible 107-element set contains zero.
Use Proposition~\ref{u96:prop:universal-lattice}, retaining the weaker
$G=182$ for compatibility with the following exact arithmetic. Then
$\Lambda=2071+k^2$, $\Omega=204-k^2$, and $T\geq61$.

If $z>0$, non-progression gives $\rho\geq62$, so $k\leq9$,
$a,b\geq36$, $T\leq142$ and $z\geq20$.
The size-36 hole row gives $\rho\geq138$; then $k^2\leq5$ and
$a,b\geq43$. The size-37 row gives $\rho\geq152$, contradicting
$61+152>204$. Thus $z=0$.

The stronger inequality $T\geq\kappa+182$ now gives
$\rho+\kappa+c\leq22-k^2$, hence $k\leq4$ and $a,b\geq41$.
If $h_P$ or $h_Q$ were positive, the corresponding support would be a
non-progression and would have defect at least $41-14=27>22$.
Thus $h_P=h_Q=\kappa=0$. The overlap and hole propagation in
Proposition~\ref{u96:prop:universal-lattice} give full interval supports
for $E,D_1,D$.
\end{proof}




\section{A further general improvement: \texorpdfstring{$N\leq104$}{N at most 104}}\label{u96:sec:upper104}
This section closes the three cases $107,106,105$. The reduction at size
107 in Theorem~\ref{u96:thm:107-core} is used as an intermediate implication;
that size is no longer an unresolved candidate. The new argument keeps the
\emph{sum} of the losses in the covering translates, and uses the total
length of the subset-sum hull. It does not require an exhaustive search of
integer sets or any new external inverse theorem.

\subsection{Lattice and hole bookkeeping at sizes 105 and 106}
\begin{lemma}\label{u96:lem:lower-lattice}
Suppose a feasible integer set of size $n\in\{105,106\}$ contains zero.
The notation and exact budget identity of Theorem~\ref{u96:thm:upper107}
continue to apply, with $d=n-8$, $e=n-16$, and $a+b=n-17$.
After a common lattice normalization, $D$ is integral. If $w$ is the
hull-hole count of $\FS(D)$, then
\begin{gather}
 0\leq w\leq4z,\qquad
 T\geq\kappa+G-3z\geq G-2z,\qquad
 T\geq\lceil G/3\rceil,\quad G=2d-16,\label{u96:eq:lower-holes}\\
 \rho_P+\rho_Q+c+T\leq\Omega=22d+97-\Lambda.\label{u96:eq:lower-budget}
\end{gather}
If $z>0$, then $\rho_P+\rho_Q\geq n-45$.
\end{lemma}
\begin{proof}
Apply Proposition~\ref{u96:prop:universal-lattice} at the stated size(s).
The displayed parameters are obtained from $d=n-8$, $e=n-16$,
$a+b=n-17$, $\Lambda=1+T_a+T_b$ and $\Omega=22d+97-\Lambda$.
Where the statement uses $G=2d-16$ rather than $2d-12$, it is a weaker
consequence of the same proved inequalities. If needed, zero is present
because the next larger size has already been excluded; the proposition
itself assumes zero and does not use that endpoint exclusion.
\end{proof}

Three additional applications of Lemma~\ref{u96:lem:hole-parameter} will be
used below. They are displayed with all threshold data; the earlier rows
are included to make the successive deductions checkable.
\begin{center}\small
\begin{tabular}{rrrrrrrr}
\toprule
$r\geq$&$C$&$b$&$K$&$D$&$J$&$M(2b,b,D+1)$&$(r-J+1)b$\\
\midrule
33&89&7&15&104&19&18&105\\
35&102&7&15&117&19&19&119\\
36&111&8&16&127&21&20&128\\
37&118&8&17&135&21&20&136\\
38&126&8&17&143&21&21&144\\
40&143&9&18&161&23&22&162\\
\bottomrule
\end{tabular}
\end{center}
Each last entry exceeds $D$, and $T_{J-1}-1>D$. In each row the conclusion
is $\mu(X)\geq r+C+1$ if $2b<h<|\FS(X)|-1$. In particular, applying one
row to both sides gives
\begin{equation}
 \rho_P+\rho_Q\geq2C+6-(a+b),\label{u96:eq:two-hole-penalty}
\end{equation}
using only $\delta_0(X)\leq|X|$.

\subsection{A certificate that uses all translate losses}
The following lemma is the main new estimate. Its two inequalities are
explicit nonnegative combinations of covering constraints; numerical
linear programming is not a premise of the proof.

\begin{lemma}[Outer translate certificate]\label{u96:lem:outer-certificate}
In the covering setup for a feasible $A$, use the four outer elements on each side,
$p_1>p_2>p_3>p_4>0$ and $-q_1<-q_2<-q_3<-q_4<0$.
Suppose $D,E$ are integral, $E=P\cup\{0\}\cup(-Q)$, and $D\setminus E$
has four positive and four negative elements, all beyond the corresponding
extreme of $E$. Write
\[
\begin{gathered}
 R_p=\max P,\quad R_q=\max Q,\quad
 M_p=\max D,\quad M_q=-\min D,\\
 S_p=\sum_{x\in D,\ x>0}x,\qquad
 S_q=\sum_{x\in D,\ x<0}|x|,\qquad
 q_p=R_p-a,\quad q_q=R_q-b,\\
 F=S_p+S_q+c-2(M_p+M_q)-2(R_p+R_q).
\end{gathered}
\]
Assume $c<d$. Then there are exactly eighteen contributing $D$-translate
slots, counted with their original multiplicities, after the two pairings.
Their individual losses sum to at most
$c$. All eight outer elements are integers. Reflect the whole set so that
the first paired slot that contributes is the positive one.

If all eighteen centers lie in their respective intervals $[0,S_p]$ or
$[0,S_q]$, one of the following bounds holds, with $\xi=0$:
\begin{align}
 F&\geq F_A(a,b)-\xi,&
 F_A(a,b)&=\tfrac12T_a+T_b+6a+8b+85,\label{u96:eq:FA}\\
 F&\geq F_B(a,b)-\xi,&
 F_B(a,b)&=\tfrac23T_a+T_b+6a+\tfrac{22}3b+\tfrac{236}3.
 \label{u96:eq:FB}
\end{align}
The same alternatives hold with $\xi=\max(q_p,q_q)$ if at most one center
is outside its respective interval and $c<d-4$.
In all cases, if $w$ is the hull-hole count of $\FS(D)$, then
\begin{equation}
 F\leq18d+126+w.\label{u96:eq:F-upper}
\end{equation}
\end{lemma}
\begin{proof}
For a positive slot with center $v$, count points in
$(v+D)\cap[1,S_p]$; on the negative side use $(v-D)\cap[1,S_q]$.
These hull counts dominate the actual contributions to $\FS(D)$.

\emph{The two pairings and integrality.}
The two centers in the first pair, expressed on positive half-lines, are
$c_p-q_1$ and $c_q-p_1$. Their sum is
$-(p_2+p_3+p_4+q_2+q_3+q_4)<-(M_p+M_q)$, so at most one of the two
slots even meets its open half-line. The sixteen unpaired slots must all
contribute, since an empty one would cost $d>c$.

If $\beta_p+\beta_q\leq0$, the second pair is $\beta_p,\beta_q$.
The unpaired centers $c_p-q_2$ and $c_q-p_2$ contribute, so their sum is
greater than $-(M_p+M_q)$. Consequently
\[
 \beta_p+\beta_q
 <M_p+M_q-2(p_3+p_4+q_3+q_4)<-(M_p+M_q).
\]
If $\beta_p+\beta_q>0$, the second paired centers instead have sum
$-(\beta_p+\beta_q)-2(p_3+p_4+q_3+q_4)<-(M_p+M_q)$.
Thus at most one member of each pair contributes. Neither pair can be
empty, again because that would cost $d$. The bound is now the sum of
eighteen actual slot counts, $2|2D|$, 84 fixed scores, and zero.
If $\lambda_v$ denotes $d$ minus the appropriate hull count, this proves
\begin{equation}
 \sum_v\lambda_v\leq c.\label{u96:eq:loss-sum}
\end{equation}
Every active center is integral, since it has an intersection with the
integral set $\FS(D)$. In particular $b_i=c_p+p_i$ is integral for all four
$i$. The identities $c_p=-b_1+b_2+b_3+b_4$ and $p_i=b_i-c_p$ make all
$p_i$ integral; the reflected identities do the same for $q_i$.

\emph{Quantile inequalities.}
Split $\lambda_v=l_v+u_v$ into losses below 1 and above the appropriate
$S$. For a positive-side center $0\leq v\leq S_p$, distinctness and
integrality give
\begin{equation}
 v\geq b+5-l_v,\qquad
 v\leq S_p-a-4+u_v.\label{u96:eq:quantile}
\end{equation}
For example, at most $v-1$ negative magnitudes can be smaller than the
positive integer $v$; the case $v=0$ includes zero among the lower losses.
The negative-side versions interchange $a,b$ and $p,q$.

If $v<0$ and $c<d-4$, the slot has more than four points. Therefore it
uses a core point, and $-v<R_p$. Among the integers $1,\ldots,-v$, at
least $-v-q_p$ belong to $P$. Counting them as lower losses gives
\[
 v\geq b+5-l_v-q_p.
\]
If $v>S_p$, the same argument on $Q$ gives
$v\leq S_p-a-4+u_v+q_q$. These are the only corrections needed. On the
negative side the roles of $q_p,q_q$ reverse. If there is at most one
exceptional slot, any combination with loss coefficients in $[0,1]$ costs
at most one additional $\xi=\max(q_p,q_q)$.

\emph{The two possible configurations.}
Reflect so that $A_0=c_p-q_1$ is the active first paired center.
In the branch $\beta_p+\beta_q\leq0$, put $B_0=c_q-p_2$; both centers
meet their half-lines, and
\[
 \beta_p=-(A_0+B_0)-2(p_3+p_4)-q_2-q_3-q_4<-M_p.
\]
Thus the active second center is $\beta_q$. In the other branch,
$\beta_p$ is unpaired and contributes, and
$p_2-q_1=A_0+\beta_p+2(p_3+p_4)>2M_p$.
It follows that $c_q-p_2<-M_q$, so the active second paired center is
$c_p-q_2$. This classification did not assume that centers are positive.

In the first configuration use the lower inequalities for the five slots
\[
 c_p-q_1,\quad c_q-p_2,\quad c_q-p_3,\quad\alpha_p,\quad\beta_q,
\]
each with coefficient one, and the upper inequality for $b_1=c_p+p_1$
with coefficient $1/2$. Subtracting the lower centers from $b_1/2$ gives
\[
 \tfrac12p_2+\tfrac52p_3+\tfrac32p_4+3q_3+3q_4.
\]
Equations~\eqref{u96:eq:loss-sum}--\eqref{u96:eq:quantile}, with the possible
correction $\xi$, therefore give
\begin{equation}
 \tfrac12 S_p+c
 \geq\tfrac92M_p+6M_q+\tfrac72a+2b+44-\xi.
 \label{u96:eq:cert-A-mid}
\end{equation}
Here $p_2\geq M_p+3$, $p_3\geq M_p+2$, $p_4\geq M_p+1$, and their
negative-side analogues were used.

In the second configuration use lower centers
$c_p-q_1,c_q-p_3,c_q-p_4,\beta_q$ with coefficient one, and $\beta_p$
with coefficient $1/3$, together with $b_1$'s upper inequality with
coefficient $1/3$. The corresponding difference of centers is
$2p_3+2p_4+3q_3+3q_4$. Hence
\begin{equation}
 \tfrac13 S_p+c
 \geq4M_p+6M_q+\tfrac{10}3a+\tfrac43b+38-\xi.
 \label{u96:eq:cert-B-mid}
\end{equation}
All used losses have coefficient at most one and refer to distinct slots.

Finally,
\begin{gather*}
 S_p\geq T_a+3R_p+6+M_p,\quad
 S_q\geq T_b+3R_q+6+M_q,\\
 M_p\geq R_p+4,\quad M_q\geq R_q+4,\quad R_p\geq a,\quad R_q\geq b.
\end{gather*}
Substitution in \eqref{u96:eq:cert-A-mid} or \eqref{u96:eq:cert-B-mid} gives
\eqref{u96:eq:FA} or \eqref{u96:eq:FB}, respectively.

For the upper bound, every sum involving $E$, and every mixed-sign pair
from $D\setminus E$, lies in
$[-M_q-R_q,M_p+R_p]$. Only the ten positive-positive and ten
negative-negative outlier pairs can lie outside. Thus
\[
 |2D|\leq M_p+M_q+R_p+R_q+21.
\]
Since $|\FS(D)|=S_p+S_q+1-w$, the definition of $c$ gives
$S_p+S_q+c=2|2D|+18d+84+w$. This proves \eqref{u96:eq:F-upper}.
\end{proof}

\begin{lemma}[A small combined excess budget]\label{u96:lem:rho23}
For a positive set of size $r\geq39$ with $\rhoD(X)\leq23$,
\[
 \mu(X)\leq\rhoD(X)+14+2\lfloor\rhoD(X)/3\rfloor.
\]
Consequently, if $z=0$, $a,b\geq39$, and
$\rho_P+\rho_Q+\kappa\leq23$, then $q_p+q_q\leq65$.
\end{lemma}
\begin{proof}
Let $j$ be the last index in $[18,r]$ with $y_j\leq4$. If it did not
exist, nonnegativity of $\rhoD(X_{17})$ and the increment recurrence would
give $\rhoD(X_r)\geq3(r-17)>23$. The later increments give
$r-j\leq\lfloor\rhoD(X_r)/3\rfloor\leq7$, so $j\geq32$.
The small-increment inverse lemma and the deleted-interval defect bound
give $\delta_0(X_j)\leq8$. Each later coefficient adds at most one defect,
proving the first formula. On the common lattice,
$q_p\leq\mu(P)+h_P$ and $q_q\leq\mu(Q)+h_Q$.
Since $h_P+h_Q=\kappa$ when $z=0$, summing gives
$q_p+q_q\leq23+28+2\lfloor23/3\rfloor=65$.
\end{proof}

\subsection{Excluding 107, 106, and 105}
\begin{theorem}[Improved general upper bound]\label{u96:thm:upper104}
Every feasible real set has at most 104 elements. In particular,
\[
 \boxed{N\leq104.}
\]
\end{theorem}
\begin{proof}
Theorem~\ref{u96:thm:upper107} gives $N\leq107$. We exclude the next
three sizes in descending order. At each size a set omitting zero would
extend to the already excluded next size, so it suffices to treat sets
containing zero. Integralization is applied before any lattice inverse
theorem; the reductions remain existence-preserving.

\emph{Size 107.} Theorem~\ref{u96:thm:107-core} gives $z=w=0$,
$(a,b)=(45-k,45+k)$ up to reflection, $0\leq k\leq4$, and
$c\leq22-k^2$. Every active center is in its hull interval: a center below
zero or above its $S$ loses at least $\min(a+5,b+5)>c$ points.
Lemma~\ref{u96:lem:outer-certificate} applies with $\xi=0$.
The minimum of $F_A$ over both orientations is $4361/2$, and that of
$F_B$ is 2351; both exceed $18\cdot99+126=1908$.
Thus 107 is impossible.

\emph{Size 106, many core holes.} Here
\[
 (a,b)=(44-k,45+k),\quad
 \Omega=227-k(k+1),\quad G=180,\quad T\geq60,
\]
up to reflection. If $z\geq15$, the non-progression penalty gives
$\rho_P+\rho_Q\geq61$, hence $k\leq9$ and $a,b\geq35$.
The size-35 hole row gives $\rho_P+\rho_Q\geq121$.
Now $k(k+1)\leq46$, so $a,b\geq38$, and
$T\leq106$ forces $z\geq37$. The size-38 row gives
$\rho_P+\rho_Q\geq169$. This contradicts $60+169>227$.

\emph{Size 106, at most fourteen core holes.} If $z=0$, then
$c\leq47-k(k+1)$ and $k\leq6$.
If $1\leq z\leq14$, then
\[
 c\leq47-k(k+1)+2z-61\leq14-k(k+1),
\]
so $k\leq3$. In either case $38\leq a,b\leq51$, $w\leq56$, and
$c<\min(a+5,b+5)$. All centers are in their hull intervals.
The minimum certificate bounds are $F_A\geq2099$ and
$F_B\geq6865/3$, while
$F\leq18\cdot98+126+56=1946$. This excludes 106.

\emph{Size 105, many core holes.} Now
\[
 (a,b)=(44-k,44+k),\quad
 \Omega=250-k^2,\quad G=178,\quad T\geq60.
\]
Suppose $z\geq15$. The initial penalty is
$\rho_P+\rho_Q\geq60$, giving $k\leq11$ and $a,b\geq33$.
Applying the hole rows successively gives the following exact deductions.
The indicated minimum size is available \emph{before} its row is applied.
\begin{center}\small
\begin{tabular}{rrrrl}
\toprule
Minimum side size&$C$&New $\rho_P+\rho_Q\geq$&New $k\leq$&Further information\\
\midrule
33&89&96&9&Next side size at least 35\\
35&102&122&8&$z\geq25$, next size at least 36\\
36&111&140&7&$z\geq34$, next size at least 37\\
37&118&154&6&$z\geq41$, next size at least 38\\
38&126&170&4&$z\geq49$, next size at least 40\\
40&143&204&--&$60+204>250$\\
\bottomrule
\end{tabular}
\end{center}
The first two rows need $z>14$, the next three $z>16$, and the last
$z>18$; the displayed deductions meet these strict conditions. All
one-sided hole counts are below their support cardinalities minus one.
Thus this case is impossible.

\emph{Size 105, one to fourteen core holes.} The budget now gives
\[
 c\leq72-k^2+2z-60\leq40-k^2.
\]
Hence $k\leq6$, $38\leq a,b\leq50$, and every active center lies in its
hull interval. The certificate gives $F_A\geq4135/2$ and
$F_B\geq6745/3$, both larger than
$18\cdot97+126+56=1928$.

\emph{Size 105, no core holes.} Here $w=0$, $k\leq8$, and the stronger
form $T\geq\kappa+178$ yields
\begin{equation}
 \rho_P+\rho_Q+\kappa+c\leq72-k^2.\label{u96:eq:105-zero-budget}
\end{equation}
We have $c\leq72<d-4=93$. A center outside its hull interval loses at
least $49-k$ points. There is at most one such center, since
$2(49-k)>72-k^2$ for every integer $k$.
If there is none, $36\leq a,b\leq52$ and the minimum bounds
$F_A\geq2040$, $F_B\geq6718/3$ already exceed 1872.

If there is one exceptional center, then $c\geq49-k$ and
\[
 \rho_P+\rho_Q+\kappa\leq23+k-k^2\leq23.
\]
Nonnegativity forces $k\leq5$, so $a,b\geq39$.
Lemma~\ref{u96:lem:rho23} gives $\xi=\max(q_p,q_q)\leq65$.
The corrected certificate yields
\[
 F\geq\min\{F_A-65,F_B-65\}
 \geq\min\{4037/2,6571/3\}>1872,
\]
again contradicting \eqref{u96:eq:F-upper}. This excludes 105 and completes
the proof.
\end{proof}

For audit purposes, all the last numerical comparisons are collected here.
The ranges include both orientations after the reflection in the certificate.
Each minimum follows by substituting $b=n-17-a$ into an explicit quadratic;
these quadratics decrease throughout the listed ranges of $a$.
\begin{center}\small
\begin{tabular}{lrrrrr}
\toprule
Case&Range of $a$&$\xi$&$\min(F_A-\xi)$&$\min(F_B-\xi)$&Upper bound for $F$\\
\midrule
107&41--49&0&$4361/2$&2351&1908\\
106, $0\leq z\leq14$&38--51&0&2099&$6865/3$&1946\\
105, $1\leq z\leq14$&38--50&0&$4135/2$&$6745/3$&1928\\
105, $z=0$, no exception&36--52&0&2040&$6718/3$&1872\\
105, $z=0$, one exception&39--49&65&$4037/2$&$6571/3$&1872\\
\bottomrule
\end{tabular}
\end{center}

\subsection{An additional reduction for the next case, 104}
The lattice step itself can also be extended to the remaining endpoint.
This proposition alone does not exclude that endpoint. It supplies the
lattice step used in the exclusion in Section~\ref{u96:sec:upper102}.
\begin{proposition}\label{u96:prop:104-lattice}
If a feasible real set has 104 elements, integralization followed by a
common-core-lattice normalization gives the lattice and hole conclusions
of Lemma~\ref{u96:lem:lower-lattice}, with
\[
 d=96,\quad e=88,\quad (a,b)=(43-k,44+k),\quad
 \Omega=272-k(k+1),\quad G=176,\quad T\geq59,
\]
up to reflection. In particular $D$ is integral in this normalization and
$w\leq4z$. Integrality of the outer eight elements in this normalization
is not asserted here.
\end{proposition}
\begin{proof}
Apply Proposition~\ref{u96:prop:universal-lattice} at the stated size(s).
The displayed parameters are obtained from $d=n-8$, $e=n-16$,
$a+b=n-17$, $\Lambda=1+T_a+T_b$ and $\Omega=22d+97-\Lambda$.
Where the statement uses $G=2d-16$ rather than $2d-12$, it is a weaker
consequence of the same proved inequalities. If needed, zero is present
because the next larger size has already been excluded; the proposition
itself assumes zero and does not use that endpoint exclusion.
\end{proof}
With the older constant, the balanced zero-hole loss budget is
$\Omega-G=96=d$. Section~\ref{u96:sec:upper102} improves the four-extreme
constant, reduces this budget to 92, and handles all out-of-hull centers.
This closes the obstruction left by the preceding version.

\section{Two more exclusions: a general upper bound of 102}\label{u96:sec:upper102}
We exclude sizes 104 and 103. The first exclusion is analytic.
For the second, the only additional ingredient is a finite family of
\emph{explicit integer linear certificates}, printed in
Appendix~\ref{u96:app:one-empty}. Their verification uses exact arithmetic and
does not depend on a numerical optimization solver. All earlier theorems
are retained as intermediate results, with their original weaker constants.

\subsection{Improving the four-extreme estimate by four}
The sharp four-extreme estimate was established in Lemma~\ref{u96:lem:four-sharp}.

\subsection{Absorbing every out-of-hull center}
The earlier outer certificate required all centers, or all but one, to lie
in their target hulls. That restriction is unnecessary if the full span of
$D$ is retained until the end.

\begin{lemma}[Unrestricted active-center certificate]\label{u96:lem:active-sharp}
In the setup and notation of Lemma~\ref{u96:lem:outer-certificate}, suppose all
sixteen unpaired covering units and both paired units are nonempty. Then
all eight outer elements are integral on the common lattice. After the
reflection specified there, either
\[
 F\geq F_A(a,b)\quad\hbox{or}\quad F\geq F_B(a,b),
 \qquad F\leq18d+126+w.
\]
There is no restriction on the number of centers outside their hulls.
In particular, the conclusion applies whenever $c<d$.
\end{lemma}
\begin{proof}
The two-pair classification and integrality proof in
Lemma~\ref{u96:lem:outer-certificate} use only nonemptiness, not $c<d$ itself.
They leave eighteen individual active slots with total loss at most $c$.
Set
\[
 U=M_p-(a+4),\qquad V=M_q-(b+4),\qquad
 q_p=R_p-a,\quad q_q=R_q-b.
\]
All are nonnegative, and $q_p\leq U$, $q_q\leq V$.
For an active positive-side center $v$, let $l_v,u_v$ count the elements
of $v+D$ below 1 and above $S_p$, respectively. Then
\begin{equation}
 \begin{split}
 v&\geq b+5-l_v-U\,\mathbf1_{v<0},\\
 v&\leq S_p-a-4+u_v+V\,\mathbf1_{v>S_p}.
 \end{split}\label{u96:eq:global-quantile}
\end{equation}
For example, if $v<0$, nonemptiness gives $-v<M_p$. Among the integers
$1,\ldots,-v$, at least $-v-U$ occur in the positive part of $D$; together
with all negative elements and zero this gives the first inequality.
If $v>S_p$, use the negative magnitudes smaller than $v-S_p$ for the
second. In-hull centers satisfy the old inequalities. Reflection exchanges
$U,V$ and $a,b$.

In branch A, the five lower constraints used in the old certificate have
at most two $U$ corrections and three $V$ corrections; the upper constraint
has coefficient $1/2$ and costs at most $V/2$. Thus their total correction
is at most $2U+7V/2$. Retaining the span terms when substituting the mass
bounds into \eqref{u96:eq:cert-A-mid} gives
\[
 F\geq F_A+3U+5V-\tfrac12q_p+q_q-(2U+\tfrac72V)
 \geq F_A+\tfrac12U+\tfrac32V\geq F_A.
\]
In branch B, the total correction is at most $4U/3+10V/3$. The analogous
substitution in \eqref{u96:eq:cert-B-mid} gives
\[
 F\geq F_B+\tfrac83U+5V+q_q-(\tfrac43U+\tfrac{10}3V)
 \geq F_B+\tfrac43U+\tfrac53V\geq F_B.
\]
These estimates allow every selected center to be exceptional at once.
The upper bound for $F$ is unchanged.
\end{proof}

\subsection{Extending the common lattice to size 103}


\begin{proposition}[The 103-element lattice reduction]\label{u96:prop:103-lattice}
Suppose sizes at least 104 have been excluded. A feasible 103-element set
can be replaced, preserving existence, by an integer set containing zero.
After common-core normalization, $D$ is integral and
\[
 \begin{gathered}
 d=95,\quad e=87,\quad(a,b)=(43-k,43+k)\text{ up to reflection},\\
 \Lambda=1893+k^2,\quad\Omega=294-k^2,\quad G=178,\quad T\geq60,\\
 w\leq4z,\quad T\geq\kappa+178-3z,\quad
 \rho_P+\rho_Q+c+T\leq294-k^2.
 \end{gathered}
\]
If $z>0$, then $\rho_P+\rho_Q\geq58$.
\end{proposition}
\begin{proof}
Apply Proposition~\ref{u96:prop:universal-lattice} at the stated size(s).
The displayed parameters are obtained from $d=n-8$, $e=n-16$,
$a+b=n-17$, $\Lambda=1+T_a+T_b$ and $\Omega=22d+97-\Lambda$.
Where the statement uses $G=2d-16$ rather than $2d-12$, it is a weaker
consequence of the same proved inequalities. If needed, zero is present
because the next larger size has already been excluded; the proposition
itself assumes zero and does not use that endpoint exclusion.
\end{proof}

\subsection{Many holes: asymmetric use of the effective test}
It is useful to apply different parameter rows to the two sides. A row
$(r_0,C,b_0)$ in the following table means that a side of size $r\geq r_0$
with $2b_0<h<|\FS|-1$ has $\mu\geq r+C+1$ and therefore
$\rhoD\geq C+3-r$. The full threshold data for
Lemma~\ref{u96:lem:hole-parameter} are included.
\begin{center}\small
\begin{tabular}{rrrrrrrr}
\toprule
$r_0$&$C$&$b_0$&$K$&$D$&$J$&$M$&$(r_0-J+1)b_0$\\
\midrule
30&70&7&13&83&19&17&84\\
32&83&7&14&97&19&18&98\\
33&89&7&15&104&19&18&105\\
34&96&8&15&111&21&19&112\\
36&111&8&16&127&21&20&128\\
38&126&8&17&143&21&21&144\\
39&134&9&18&152&23&21&153\\
40&143&9&18&161&23&22&162\\
41&151&9&19&170&23&22&171\\
42&160&9&19&179&23&23&180\\
43&142&7&18&160&21&21&161\\
43&156&8&19&175&22&22&176\\
43&168&9&19&187&23&23&189\\
44&147&7&18&165&21&21&168\\
44&176&9&20&196&23&23&198\\
\bottomrule
\end{tabular}
\end{center}
Every window entry exceeds $D$, and every row also has $T_{J-1}-1>D$.
Applying rows with parameters $C_-,C_+$ to the smaller and larger side
gives
\begin{equation}
 \rho_P+\rho_Q\geq C_-+C_++6-(a+b).
 \label{u96:eq:asymmetric-hole-penalty}
\end{equation}

\begin{lemma}\label{u96:lem:104103-high}
For sizes 104 and 103, once their lattice reductions are available, $z\geq15$
is impossible.
\end{lemma}
\begin{proof}
At size 104, write $(a,b)=(43-k,44+k)$ up to reflection.
The new constants are $G=180$, $T\geq60$, and $\Omega=272-k(k+1)$.
The initial penalty $\rho\geq59$ gives $k\leq11$, so the smaller side is
at least 32 and the larger at least 44. The following successive rows use
\eqref{u96:eq:asymmetric-hole-penalty}:
\begin{center}\small
\begin{tabular}{rrrrrl}
\toprule
Available side sizes&$C_-$&$C_+$&New $\rho\geq$&New $k\leq$&Further consequence\\
\midrule
32,44&83&147&149&7&$z\geq29$, smaller side $\geq36$\\
36,44&111&176&206&2&$z\geq57$, both sides $\geq41$\\
41,41&151&151&221&--&$60+221>272$\\
\bottomrule
\end{tabular}
\end{center}
The successive hole thresholds are $z>14$, $z>18$, $z>18$ and are all met.
For example, after the first row $T\leq123$ and $T\geq180-2z$ give
$z\geq29$. This closes the 104-element many-hole case.

At size 103, $\Omega=294-k^2$, $G=178$, $T\geq60$. The initial penalty
$\rho\geq58$ gives $k\leq13$, a smaller side at least 30, and a larger
side at least 43. The analogous deductions are
\begin{center}\small
\begin{tabular}{rrrrrr}
\toprule
Available side sizes&$C_-$&$C_+$&New $\rho\geq$&New $k\leq$&New $z\geq$\\
\midrule
30,43&70&142&132&10&15\\
33,43&89&142&151&9&18\\
34,43&96&156&172&7&28\\
36,43&111&168&199&5&42\\
38,43&126&168&214&4&49\\
39,43&134&168&222&3&53\\
40,43&143&168&231&1&58\\
42,43&160&168&248&--&Contradiction\\
\bottomrule
\end{tabular}
\end{center}
Here the first two rows require $z>14$, the third $z>16$, and the
remaining rows $z>18$. All conditions hold before their use. The last
contradiction is $60+248>294$.
\end{proof}

\subsection{The one-empty-unit reduction}
\begin{lemma}\label{u96:lem:one-empty}
In the lattice setup of Proposition~\ref{u96:prop:103-lattice}, there is no
feasible 103-element set for which $z=0$ and at least one of the eighteen
covering units is empty.
\end{lemma}
\begin{proof}
The exact budget and \eqref{u96:eq:sharp-G} give
\begin{equation}
 \rho_P+\rho_Q+\kappa+c\leq116-k^2.
 \label{u96:eq:103-zero-budget}
\end{equation}
An empty unit costs $d=95$, so $k\leq4$. Moreover $c\leq116<190=2d$;
there is exactly one empty unit. Put $\delta=c-95$. Then
\begin{equation}
 a+b=86,\qquad39\leq a,b\leq47,\qquad
 0\leq\delta\leq21-k^2\leq21.
 \label{u96:eq:one-empty-parameters}
\end{equation}
Each other unit has at least $95-\delta\geq74$ actual contributions.

\emph{Outer integrality.}
Every nonempty individual slot has an integral center, because both $D$
and $\FS(D)$ are integral. If no one of the eight unpaired centers
$c_p+p_i,c_q+q_i$ is missing, the usual identities give all outer elements
integral. If, for example, $c_p+p_j$ is missing, all four $c_q+q_i$ remain,
so every $q_i$ is integral. The unpaired centers $c_p-q_3,c_p-q_4$ both
remain, giving $c_p\in\mathbb Z$. The other three $p_i$ and then the
identity $c_p=p_1-p_2-p_3-p_4$ give the last one. The other case is its
reflection. Thus all strict order gaps are at least one.

\emph{Pair choices and possible splitting.}
The first pair always has at most one contributing member. When this unit
is present, reflect to make its positive member $c_p-q_1$ active.
In branch B ($\beta_p+\beta_q>0$), the second pair also always has at most
one contributing member. In branch A ($\beta_p+\beta_q\leq0$), the same
is true whenever both unpaired cross-centers $c_p-q_2,c_q-p_2$ remain.
If either of these is the empty unit, we impose no individual constraint
on the second pair; allowing it to split only enlarges the possible set.
If the first pair is empty, both orientations of the second are retained.
Thus no orientation has been silently discarded.

\emph{The relaxed inequalities.}
For every remaining unpaired slot, and every paired slot known to have
one contributing member, the hull loss is at most $\delta\leq21$.
A center outside $[0,S_p]$ or $[0,S_q]$ would lose at least
$\min(a+5,b+5)\geq44$ points. All these centers are therefore in their
hulls and obey the uncorrected quantile inequalities. The sum of all
their lower and upper losses is at most $\delta$; dropping a possibly
split pair preserves this inequality.

Let $a=39+t$, $b=47-t$, where we relax to the continuous range $0\leq t\leq8$.
Use $\delta\leq21$ and the lower tangent
\[
 T_r\geq\tfrac{87}{2}r-\tfrac{1849}{2},
 \qquad T_r-\bigl(\tfrac{87}{2}r-\tfrac{1849}{2}\bigr)
 =\tfrac12(r-43)^2.
\]
Together with the order, span, mass and quantile constraints, these give
the integer coefficient systems printed in Appendix~\ref{u96:app:one-empty}.
There are 33 systems in branch A and 35 in branch B. The appendix supplies
21 explicit nonnegative integer row combinations covering all 68 systems.
In 30 systems they imply a negative upper bound for a nonnegative linear
form. In the remaining 38 they imply
\[
 F=S_p+S_q+95+\delta-2(M_p+M_q)-2(R_p+R_q)\geq1940.
\]
But $z=0$ gives $w=0$ and $F\leq18\cdot95+126=1836$.
This contradiction proves the lemma. The certificate identities, including
their nonnegative residual coefficients, can be checked directly from the
printed tables or by the supplied standard-library verifier.
\end{proof}

\subsection{The new upper bound}
\begin{theorem}[General upper bound 102]\label{u96:thm:upper102}
Every finite real set $A$ satisfying $\FS(A)\subseteq4A$ has at most
102 elements. Hence
\[
 \boxed{N\leq102.}
\]
\end{theorem}
\begin{proof}
Theorem~\ref{u96:thm:upper104} gives $N\leq104$. We exclude 104 first, then
103; at each stage adjoining zero and integralization are legitimate.

For 104, Proposition~\ref{u96:prop:104-lattice} provides the common lattice,
and Lemma~\ref{u96:lem:four-sharp} improves its constants to $G=180$, $T\geq60$.
Lemma~\ref{u96:lem:104103-high} excludes $z\geq15$. If $z=0$, then
$c\leq92-k(k+1)<96=d$ and $k\leq9$. If $1\leq z\leq14$, then
\[
 c\leq92-k(k+1)+2z-59\leq61-k(k+1)<d,
\]
so $k\leq7$. Lemma~\ref{u96:lem:active-sharp} applies in both cases.

For 103, Proposition~\ref{u96:prop:103-lattice} now applies and the many-hole
lemma again excludes $z\geq15$. If $1\leq z\leq14$, then
\[
 c\leq116-k^2+2z-58\leq86-k^2<95=d,
\]
so $k\leq9$ and the active-center certificate applies. If $z=0$, then
$k\leq10$. When all eighteen units are nonempty, the same certificate
applies regardless of whether $c<d$. Any empty-unit case is excluded by
Lemma~\ref{u96:lem:one-empty}.

For completeness, all active-center comparisons are displayed below.
The ranges include both orientations after reflection.
\begin{center}\small
\begin{tabular}{llrrrr}
\toprule
Size&Core holes&Range of $a$&$\min F_A$&$\min F_B$&Upper bound for $F$\\
\midrule
104&$z=0$&34--53&$3971/2$&2195&1854\\
104&$1\leq z\leq14$&36--51&2008&$6596/3$&1910\\
103&$z=0$, all units active&33--53&$3887/2$&$6461/3$&1836\\
103&$1\leq z\leq14$&34--52&1953&$6461/3$&1892\\
\bottomrule
\end{tabular}
\end{center}
Each row contradicts the upper bound. Both remaining sizes are excluded.
\end{proof}

\section{Retaining core mass: a general upper bound of 100}\label{u96:sec:upper100}
The next exclusions are 102 and 101. The principal
change is to retain the exact size and the exact mass of the core, rather
than replacing both by their balanced lower bounds. This simplifies the
lattice step and makes the many-hole case substantially stronger. The
remaining finite certificates allow one empty covering unit and, when
necessary, one out-of-hull center.

\subsection{A shorter, stronger common-lattice argument}
Write $\mu=\mu(P)+\mu(Q)$ and $\rho=\rho_P+\rho_Q$.
The exact identities and bounds already established include
\begin{equation}
 s=\Lambda+\mu+\kappa,\quad
 R_p+R_q\leq a+b+\mu+\kappa+z,\quad
 \mu\leq\rho+2(a+b)-4.
 \label{u96:eq:retain-core}
\end{equation}
The last inequality follows directly from $\delta_0(P)\leq a$ and
$\delta_0(Q)\leq b$.

\begin{proposition}\label{u96:prop:lattice100}
For $n\in\{102,101,100\}$, any feasible integer set containing zero admits
the common-core normalization in which $D$ is integral. In this
normalization all the exact budget and hole conclusions apply, including
\[
 w\leq4z,\quad T\geq\kappa+G-3z\geq G-2z,\quad
 T\geq\lceil G/3\rceil,
 \qquad G=2d-12.
\]
\end{proposition}
\begin{proof}
Apply Proposition~\ref{u96:prop:universal-lattice} at the stated size(s).
The displayed parameters are obtained from $d=n-8$, $e=n-16$,
$a+b=n-17$, $\Lambda=1+T_a+T_b$ and $\Omega=22d+97-\Lambda$.
Where the statement uses $G=2d-16$ rather than $2d-12$, it is a weaker
consequence of the same proved inequalities. If needed, zero is present
because the next larger size has already been excluded; the proposition
itself assumes zero and does not use that endpoint exclusion.
\end{proof}


\subsection{A mass bonus in the active-center certificate}
In the common normalization put
\[
 \varepsilon_p=\mu(P)+h_P=\sum P-T_a,\qquad
 \varepsilon_q=\mu(Q)+h_Q=\sum Q-T_b.
\]
These are exact equalities, not stability approximations.

\begin{lemma}\label{u96:lem:mass-bonus}
When all eighteen covering units are nonempty, the alternatives in
Lemma~\ref{u96:lem:active-sharp} strengthen to
\begin{align*}
 F&\geq F_A(a,b)+\tfrac12\varepsilon_p+\varepsilon_q,\quad\hbox{or}\\
 F&\geq F_B(a,b)+\tfrac23\varepsilon_p+\varepsilon_q.
\end{align*}
Both remain valid with arbitrarily many out-of-hull centers. If $z>0$ in
the reductions of Proposition~\ref{u96:prop:lattice100}, then
\[
 \mu(P)\geq a-3,\quad\mu(Q)\geq b-3,
 \qquad\varepsilon_p\geq a-3+z,\quad
 \varepsilon_q\geq b-3+z.
\]
\end{lemma}
\begin{proof}
The mass estimates in the original certificate can be replaced by
\[
 S_p\geq T_a+\varepsilon_p+3R_p+6+M_p,
 \quad S_q\geq T_b+\varepsilon_q+3R_q+6+M_q.
\]
Their coefficients in the two substitutions are $(1/2,1)$ and $(2/3,1)$.
All the span terms that absorb exceptional-center errors are unchanged.
This proves the two strengthened alternatives.

For a positive $r$-element set, the stability theorem stated in
Section~\ref{base:sec:upper} shows that $\mu\leq r-4$ forces progression
support: take $M=\mu$. In its integer normalization the support has
$1+T_r+\mu$ points, while the sum of all elements is at most $T_r+\mu$;
therefore the support fills that entire integer interval. Consequently a
non-progression support has $\mu\geq r-3$.
When $z>0$, the inverse-theorem bounds ensure
$0<h_P,h_Q<|\FS(P)|-1,|\FS(Q)|-1$ respectively, so both supports are
non-progressions. Finally $h_P,h_Q\geq z$.
\end{proof}

\subsection{Excluding the many-hole branches at 102 and 101}
Only nine applications of the effective-hole lemma are needed, all with
$b_0=7$. Thus every row is available whenever $z\geq15$.
\begin{center}\small
\begin{tabular}{rrrrrrr}
\toprule
$r_0$&$C$&$K$&$D$&$J$&$M$&$7(r_0-J+1)$\\
\midrule
29&63&12&75&19&16&77\\
32&83&14&97&19&18&98\\
34&96&15&111&19&19&112\\
35&102&15&117&19&19&119\\
36&109&16&125&19&19&126\\
38&116&16&132&20&20&133\\
39&122&17&139&20&20&140\\
40&128&17&145&20&20&147\\
42&135&18&153&21&21&154\\
\bottomrule
\end{tabular}
\end{center}
Every last entry exceeds $D$, and $T_{J-1}-1>D$ in each row.
Let $C(r)$ be the largest available $C$ with $r_0\leq r$, or $-4$ if
none is available. The default $-4$ denotes the non-progression bound
$\mu\geq r-3$. Define
\[
 u(r)=r+C(r)+1,\qquad
 v(r)=\max\{r-14,\ u(r)+2-2r\}.
\]
Thus $\mu(P)\geq u(a)$, $\mu(Q)\geq u(b)$ and
$\rho\geq v(a)+v(b)$ in the many-hole branch.

\begin{lemma}\label{u96:lem:high100}
A feasible set of size 102 or 101 cannot have $z\geq15$ in its common
normalization.
\end{lemma}
\begin{proof}
Put $H=\Omega-v(a)-v(b)$ and $t_0=\lceil G/3\rceil$.
The budget gives $T+c\leq H$, so $z\leq T\leq H$ and $c\leq H-t_0$.
If $H<t_0$, this is already impossible. Otherwise the table below verifies
$H-t_0<d$, so all eighteen units are active.

By Lemma~\ref{u96:lem:mass-bonus}, the two differences between the lower bound
for $F$ and its upper bound $18d+126+4z$ are bounded below by
\[
\begin{split}
 F_A(a,b)+\tfrac12u(a)+u(b)-(18d+126)-\tfrac52H,\\
 F_B(a,b)+\tfrac23u(a)+u(b)-(18d+126)-\tfrac73H.
\end{split}
\]
Indeed both expressions decrease as $z$ increases, and $z\leq H$.
Let $\mathcal M$ be the minimum of these two expressions and their two
reflected versions. The following exact table checks every remaining
imbalance. For 102 the initial non-progression budget gives $k\leq13$;
$k=12,13$ yield $H=40,14<t_0=59$. For 101 it gives $k\leq14$;
$k=13,14$ yield $H=47,44<t_0=58$.
\begin{center}\small
\begin{tabular}{rrrr|rrrr}
\toprule
\multicolumn{4}{c|}{$n=102$, $d=94$}&\multicolumn{4}{c}{$n=101$, $d=93$}\\
$k$&$H$&$c\leq$&$\mathcal M$&$k$&$H$&$c\leq$&$\mathcal M$\\
\midrule
0&124&65&$389/2$&0&144&86&$269/2$\\
1&129&70&$305/2$&1&150&92&$179/2$\\
2&125&66&$283/2$&2&147&89&$151/2$\\
3&125&66&116&3&148&90&47\\
4&123&64&97&4&147&89&25\\
5&120&61&81&5&145&87&6\\
6&108&49&96&6&134&76&18\\
7&101&42&93&7&128&70&12\\
8&91&32&100&8&119&61&16\\
9&86&27&89&9&115&57&2\\
10&66&7&130&10&96&38&40\\
11&64&5&$215/2$&11&95&37&$29/2$\\
&&&&12&72&14&$131/2$\\
\bottomrule
\end{tabular}
\end{center}
Every margin is strictly positive, a contradiction. No search over sets
$A$ is involved in this table.
\end{proof}

\subsection{Low holes and one possible exceptional center}
After the preceding lemma, $0\leq z\leq14$. The budgets are
\begin{center}\small
\begin{tabular}{lll}
\toprule
Size&$z=0$&$1\leq z\leq14$\\
\midrule
102&$\rho+\kappa+c\leq139-k(k+1)$&$c\leq110-k(k+1)$\\
101&$\rho+\kappa+c\leq162-k^2$&$c\leq134-k^2$\\
\bottomrule
\end{tabular}
\end{center}
The positive-hole bounds use $\rho\geq n-45$ and $T\geq G-2z$.
In particular $c<2d$: there is at most one empty covering unit.

When no unit is empty, the strengthened active-center certificate gives
the following complete comparisons. For positive $z$ we use only
$\varepsilon_p\geq a-2$, $\varepsilon_q\geq b-2$ and $w\leq56$.
\begin{center}\small
\begin{tabular}{llrrrr}
\toprule
Size&Core holes&Range of $a$&First minimum&Second minimum&Upper $F$\\
\midrule
102&$z=0$&31--54&$3791/2$&$6337/3$&1818\\
102&$1\leq z\leq14$&32--53&1958&$6530/3$&1874\\
101&$z=0$&30--54&$3713/2$&$6215/3$&1800\\
101&$1\leq z\leq14$&31--53&1917&$6406/3$&1856\\
\bottomrule
\end{tabular}
\end{center}
Thus only exactly one empty unit needs further treatment.

\begin{lemma}\label{u96:lem:one-empty100}
The one-empty-unit cases at sizes 102 and 101 are impossible.
\end{lemma}
\begin{proof}
Put $\delta=c-d$. All other units are nonempty and the sum of their losses
is at most $\delta$. The outer-integrality argument in
Lemma~\ref{u96:lem:one-empty} requires only that at most one unit is empty,
so it applies here unchanged. The same is true of the pair choices:
reflect to select the positive first pair when it is present, retain both
possibilities for a nonempty second pair, and drop individual second-pair
constraints in branch A if one of its unpaired cross-centers is missing.

Only the centers retained by those rules are used below. Their hull losses
sum to at most $\delta$. An out-of-hull retained center loses at least
$\min(a,b)+5$ points.

\emph{Size 102.} If $z=0$, then
\[
 k\leq6,\quad a,b\geq36,\quad
 \delta\leq45-k(k+1)<47-k=\min(a,b)+5.
\]
If $z>0$, then $k\leq3$ and $\delta\leq16-k(k+1)$, an even smaller
bound. Thus every retained center is in its hull. Both cases are covered
by the single relaxed regime
\[
 a+b=85,\quad36\leq a,b\leq49,\quad0\leq\delta\leq45,
 \quad F\leq1874.
\]
For this regime the finite certificates in Appendix~\ref{u96:app:upper100}
give a contradiction in every mode, or $F\geq9378/5>1874$.

\emph{Size 101 with $z>0$.} Here
\[
 k\leq6,\quad36\leq a,b\leq48,\qquad
 \delta\leq41-k^2<47-k.
\]
All retained centers again lie in their hulls. The exact mass bonus gives
$\varepsilon_p\geq a-2$, $\varepsilon_q\geq b-2$. In this regime the
certificates give a contradiction or $F\geq1910>1856$.

\emph{Size 101 with $z=0$.} Initially
\[
 k\leq8,\quad34\leq a,b\leq50,\quad\delta\leq69-k^2.
\]
If all retained centers are in their hulls, the certificate lower bound
is $F\geq9103/5>1800$, or the constraints are inconsistent.
Otherwise, since
\[
 2(47-k)-(69-k^2)=(k-1)^2+24>0,
\]
there is at most one exceptional retained center. Its loss is at least
$47-k$, so $k^2-k\leq22$, whence $k\leq5$ and $a,b\geq37$.

We retain the improved range $37\leq a,b\leq47$ in this exceptional
regime. If the exceptional center is on the positive side, the corrected
quantile bounds of \eqref{u96:eq:global-quantile} become
\[
 \begin{array}{ll}
 v<0:&-v-l_v-M_p\leq-d,\qquad l_v\geq b+5,\\
 v>S_p:&v-u_v-S_p-M_q\leq-(d-1),\qquad u_v\geq a+5.
 \end{array}
\]
The reflected formulas cover the negative side. The unmodified opposite
quantile bound is still valid. Appendix~\ref{u96:app:upper100} considers every
retained center and both endpoints; it does not guess the exceptional
location. Every such system is inconsistent or has
$F\geq10957/6>1800$.

All four regimes are exact relaxations of the corresponding configurations.
Their 2512 systems are certified by 42 printed nonnegative integer row
combinations. The exhaustive mode-generation rule and a solver-free
verification procedure are included in the appendix.
\end{proof}

\begin{theorem}[General upper bound 100]\label{u96:thm:upper100}
Every feasible finite real set has at most 100 elements. Thus
\[
 \boxed{N\leq100}.
\]
\end{theorem}
\begin{proof}
Theorem~\ref{u96:thm:upper102} gives $N\leq102$. Exclude size 102 first, then 101.
At each size a set omitting zero could be extended to the already excluded
next size; integralization preserves existence. Proposition~\ref{u96:prop:lattice100}
supplies the common lattice. Lemma~\ref{u96:lem:high100} excludes $z\geq15$;
the displayed active-center comparisons handle $0\leq z\leq14$ when no
unit is empty, and Lemma~\ref{u96:lem:one-empty100} handles the only remaining
possibility. 
\end{proof}

\section{Excluding size 100: a general upper bound of 99}\label{u96:sec:upper99}
The new exclusion retains the integer imbalance in both the loss budget
and the exact triangular mass bounds. It also treats two phenomena that
were absent from the preceding proof: two exceptional retained centers,
and two empty covering units. The latter case requires no assertion that
all eight outer elements are integral.

At size 100, adjoining zero is legitimate because size 101 is already
excluded. Integralization and Proposition~\ref{u96:prop:lattice100} give
\begin{equation}
\begin{gathered}
 d=92,\quad e=84,\quad a+b=83,\quad
 \{a,b\}=\{41-k,42+k\},\quad I=k(k+1),\\
 \Lambda=1765+I,\quad \Omega=356-I,\quad G=172,\\
 \rho+2\vartheta+c+T=\Omega,\quad
 T\geq\kappa+172-3z\geq172-2z,\quad
 T\geq58,\quad T\geq z,\quad w\leq4z.
\end{gathered}\label{u96:eq:99-budget}
\end{equation}
Here $\rho=\rho_P+\rho_Q$, and all slack variables are nonnegative.
If $z>0$, the non-progression penalty is $\rho\geq55$.

\subsection{The many-hole branch}
For each $r\in[26,57]_{\mathbb Z}$ use the two effective-hole rows
$(r,C_7(r),7)$ and $(r,C_9(r),9)$ printed in
Appendix~\ref{u96:app:upper99}. Each satisfies the strict hypotheses of
Lemma~\ref{u96:lem:hole-parameter}. For $z\geq15$ put
\[
 u(r,z)=r+1+\max\bigl(\{C_7(r)\}\cup
       \{C_9(r):z\geq19\}\bigr),\qquad
 v(r,z)=\max\{r-14,u(r,z)+2-2r\}.
\]
These give $\mu(P)\geq u(a,z)$, $\mu(Q)\geq u(b,z)$ and
$\rho\geq v(a,z)+v(b,z)$. The initial budget
$55+58\leq356-I$ implies $k\leq15$, so every required side size is in
the stated range. Define
\[
 C_k(z)=356-I-v(a,z)-v(b,z)-\max(58,z,172-2z).
\]
Then $c\leq C_k(z)$. For all states with $C_k(z)\geq0$ the following table
gives $C_k(z)<92$. Thus all eighteen units are active.

For each orientation $(x,y)=(a,b)$ or $(b,a)$, the two lower bounds for $F$
are
\[
 F_A(x,y)+\tfrac12\bigl(u(x,z)+z\bigr)+u(y,z)+z,
 \quad F_B(x,y)+\tfrac23\bigl(u(x,z)+z\bigr)+u(y,z)+z.
\]
Their minimum, minus $1782+4z$, is denoted by $M_k(z)$, with both
orientations included. Exact evaluation of the integer range
$15\leq z\leq356-I$ gives:
\begin{center}\small
\begin{tabular}{rrrrr}
\toprule
$k$&Smallest possible $z$&Largest possible $z$&Maximum $C_k(z)$&Minimum $M_k(z)$\\
\midrule
0&15&122&64&185\\
1&15&120&62&164\\
2&15&117&59&146\\
3&15&118&60&117\\
4&15&110&52&115\\
5&36&100&42&120\\
6&41&91&33&123\\
7&47&78&20&139\\
8&56&61&3&167\\
9--15&\multicolumn{4}{c}{No state has $C_k(z)\geq0$}\\
\bottomrule
\end{tabular}
\end{center}
The listed endpoints need not assert that every intervening $z$ survives;
all integer values in the original range are checked. Every surviving
margin is positive. Consequently $z\geq15$ is impossible.

\subsection{Low holes and the possible empty units}
Assume henceforth $0\leq z\leq14$. At $z=0$, the stronger first inequality
for $T$ in \eqref{u96:eq:99-budget} gives
\begin{equation}
 \rho+\kappa+2\vartheta+c\leq184-I.
 \label{u96:eq:99-zero}
\end{equation}
For $z>0$, the weaker but useful bound is
\begin{equation}
 c\leq129-I+2z\leq157-I.\label{u96:eq:99-positive}
\end{equation}
Hence there are at most two empty units, and two are possible only when
$z=0$. If no unit is empty, the mass-bonus lemma applies. The minima of
its first alternative are
\[
 \min_{28\leq a\leq55}F_A(a,83-a)=1815>1782
\]
and
\[
 \min_{29\leq a\leq54}
 \bigl(F_A(a,83-a)+(a-2)/2+(81-a)\bigr)
 =3743/2>1838.
\]
The second alternative is larger: indeed
$F_B-F_A=T_a/6-2b/3-19/3>0$ throughout these ranges, and its mass
coefficient is also larger. The upper bounds are $F\leq1782$ at $z=0$
and $F\leq1838$ for $1\leq z\leq14$. This excludes the no-empty case.

Suppose exactly one unit is empty and write $\delta=c-92$.
The integrality and pair-choice arguments in Lemma~\ref{u96:lem:one-empty}
apply unchanged. All retained centers are integral, and the sum of their
hull losses is at most $\delta$. A retained center outside its hull loses
at least
\[
 L=\min(a,b)+5=46-k
\]
points. The following exhaustive distinctions will be used:
\begin{center}\small
\begin{tabular}{llll}
\toprule
Regime&Residual budget&Range of $a$&Exceptional retained centers\\
\midrule
Z&$\delta\leq92-I$&32--51&None\\
P&$\delta\leq65-I$&34--49&None\\
ZE&$\delta\leq92-I$&35--48&Exactly one\\
PE&$\delta\leq65-I$&37--46&Exactly one\\
ZZ&$\delta\leq92-I$&40--43&Exactly two\\
\bottomrule
\end{tabular}
\end{center}
The intervals are derived rather than assumed. Nonnegative residual
budget gives $k\leq9$ in Z and $k\leq7$ in P. One exceptional center
requires $L\leq\delta$, giving $k^2\leq46$ in ZE and $k^2\leq19$ in PE.
Further,
\[
 2L-(65-I)=(k-1/2)^2+107/4>0,
 \qquad 3L-(92-I)=(k-1)^2+45>0.
\]
Thus there are at most one and two exceptions, respectively.

\subsection{Rigidity when two retained centers are exceptional}
In regime ZZ, $2L\leq\delta\leq92-I$ implies $k(k-1)\leq0$.
Therefore $k=0$ or 1, $\delta=2L$, and equality in
\eqref{u96:eq:99-zero} forces
\begin{equation}
 \rho_P=\rho_Q=\kappa=0.\label{u96:eq:99-rigid}
\end{equation}
We need the following elementary equality consequence.

\begin{lemma}[A zero-defect core]\label{u96:lem:zero-core99}
Let $X\subset\mathbb Z_{>0}$ have $r\geq15$ elements, suppose
$\FS(X)=[0,\sum X]_{\mathbb Z}$, and suppose $\rhoD(X)=0$.
Then
\[
 X=[1,r]_{\mathbb Z}\quad\hbox{or}\quad
 X=[1,r+1]_{\mathbb Z}\setminus\{r-1\}.
\]
In particular $\max X\leq r+1$.
\end{lemma}
\begin{proof}
The increment recurrence and $\rhoD(X_{r-1})\geq0$ give $y_r\leq2$.
The small-increment inverse lemma therefore writes $X$, up to scaling,
as $[1,r+q]\setminus D_0$ with $q\leq2$ deletions and the last element
present. Since $1\in\FS(X)$, the smallest element is 1; the scale is 1
in the integral normalization. More explicitly, the deleted interval is
primitive, its scale is an integer because $X$ is integral, and then the
smallest element condition forces that integer to be 1.

Put $h=r+q$ and $C=\{h-d:d\in D_0\}$. The full-support hypothesis gives
\[
 \mu=\sum X-T_r=\sum C-T_{q-1}=2\delta_0(X)-2.
\]
For $q=0$ the assertion is immediate. For $q=1,2$, the deleted-interval
lemma gives $\delta_0(X)\leq2q$, hence $\mu\leq6$.
A low defect $x\leq q$ would force $2x\in D_0$ and a depth at least
$h-2q\geq r-2$ in $C$, implying $\mu\geq r-3>6$.
So there are no low defects. Apart from the maximum $h$, let the $m$
high defects be $h-t$. Then $t\notin C$, while $2t\in C$; these $m$
forced depths are distinct. Replacing each $2t$ by $t$ leaves $q$
distinct positive depths. Hence
\[
 2m=\mu\geq q+\sum t\geq q+T_m,\qquad0\leq m\leq q\leq2.
\]
The only possibilities are $(q,m)=(0,0)$ or $(1,1)$. In the latter,
$\mu=2$ forces $C=\{2\}$, giving the second displayed set.
\end{proof}

In \eqref{u96:eq:99-rigid}, $h_P,h_Q\leq\kappa=0$, so both core supports are
full intervals. The lemma gives $R_p-a,R_q-b\leq1$.
Moreover, each retained center has loss at most $L\leq46$, because the
other exceptional center already costs $L$; its hull count is therefore
at least $92-46=46>4$. The core-point version of the quantile argument
in Lemma~\ref{u96:lem:outer-certificate} is valid for every retained center.
Allowing a correction of 1 at both ends, even for nonexceptional centers,
gives the simultaneous relaxed inequalities
\begin{equation}
 -v-l_v\leq-b-4,\qquad v-u_v-S_p\leq-a-3
 \label{u96:eq:99-two-exception}
\end{equation}
on the positive side, and the reflected versions on the negative side.
This handles all locations of the two exceptions without enumerating them.

\subsection{Two empty units: use zero loss, not outer integrality}
If two units are empty, their combined cost is $184$.
Equation~\eqref{u96:eq:99-zero} forces
\[
 k=0,\quad \{a,b\}=\{41,42\},\quad c=184,
 \quad \rho_P=\rho_Q=\kappa=\vartheta=0.
\]
Every other unit contributes its full capacity of 92. For every retained
individual center the entire translate of $D$ therefore lies in its target
interval. In positive-side notation this gives the stronger inequalities
\begin{equation}
 v\geq M_q+1,\qquad v\leq S_p-M_p,
 \label{u96:eq:99-full-translate}
\end{equation}
and, on the negative side, $v\geq M_p+1$, $v\leq S_q-M_q$.

Both empty units are enumerated among the same eighteen units as before.
The first pair has at most one contributing member and can be oriented
positively when present. In branch B the second pair also has at most one.
In branch A this remains true if both cross-centers are present; if either
cross-center is missing, all individual constraints for the second pair
are dropped. There are 258 modes in branch A and 289 in branch B, hence
547 modes in all, and both values $a=41,42$ are retained.

The outer elements need not be proved integral here. The models relax
their strict order to weak order, $p_i\geq p_{i+1}$,
$q_i\geq q_{i+1}$, $p_4\geq M_p$, $q_4\geq M_q$.
Likewise branch B uses $\beta_p+\beta_q\geq0$.
The inner set $D$ is integral, so $M_p\geq R_p+4$ and its reflection
remain valid. Every retained center satisfies
\eqref{u96:eq:99-full-translate}, independently of outer integrality.
This is the regime denoted by EE below.

\subsection{Exact finite certificates and the conclusion}
Appendix~\ref{u96:app:upper99} specifies every inequality and every mode for
the six regimes. The masses use the exact $T_a,T_b$; the budget uses the
same integer $I=k(k+1)$, with neither expression replaced by a global
minimum. In P and PE, the additional core masses are at least $a-2,b-2$.
In ZE and PE the exceptional center and both possible directions are
enumerated, with the corrected quantile row and its necessary loss bound.

The 52 printed nonnegative integer row combinations give the following
complete certificate results. A bound is listed only for systems not
certified inconsistent; $F$ has constant term 92 in the first five rows
and 184 in the last.
\begin{center}\small
\begin{tabular}{lrrrrr}
\toprule
Regime&Systems&Bound certificates&Inconsistency certificates&Least certified $F$&Upper $F$\\
\midrule
Z&1360&1280&80&$5357/3$&1782\\
P&1088&1024&64&1840&1838\\
ZE&32312&29541&2771&1783&1782\\
PE&23080&21002&2078&1840&1838\\
ZZ&272&256&16&1795&1782\\
EE&1094&1002&92&$5374/3$&1782\\
\bottomrule
\end{tabular}
\end{center}
Every finite certificate is verified coefficient by coefficient using
arbitrary-precision integers. The 59206 models exhaust the low-hole cases
with an empty unit. The numerical optimizer was used only to discover
candidate multipliers; it is absent from the delivered verification path.

\begin{theorem}[General upper bound 99]\label{u96:thm:upper99}
Every finite real set $A$ with $\FS(A)\subseteq4A$ has at most 99 elements.
Consequently
\[
 \boxed{N\leq99}.
\]
\end{theorem}
\begin{proof}
The previous theorem gives $N\leq100$. A hypothetical set of size 100
can be assumed to contain zero, since otherwise adjoining zero gives the
already excluded size 101. The lattice reduction and
\eqref{u96:eq:99-budget} apply. The many-hole table excludes $z\geq15$.
For $0\leq z\leq14$, the active-center bounds exclude no empty unit,
and the six certificate regimes exclude one or two empty units.
There are no other possibilities. 
\end{proof}

\subsection{An additional reduction at the next size, 99}
\begin{proposition}\label{u96:prop:lattice99}
At size 99 the common-lattice and hole conclusions still hold, with
\[
 d=91,\quad e=83,\quad(a,b)=(41-k,41+k),\quad
 \Lambda=1723+k^2,\quad\Omega=376-k^2,\quad G=170.
\]
This proposition does not exclude size 99.
\end{proposition}
\begin{proof}
Apply Proposition~\ref{u96:prop:universal-lattice} at the stated size(s).
The displayed parameters are obtained from $d=n-8$, $e=n-16$,
$a+b=n-17$, $\Lambda=1+T_a+T_b$ and $\Omega=22d+97-\Lambda$.
Where the statement uses $G=2d-16$ rather than $2d-12$, it is a weaker
consequence of the same proved inequalities. If needed, zero is present
because the next larger size has already been excluded; the proposition
itself assumes zero and does not use that endpoint exclusion.
\end{proof}

\section{Retaining residual loss: a general upper bound of 98}\label{u96:sec:upper98}
We exclude size 99. The distinction from the preceding endpoint
is essential: two empty units now leave a residual budget as large as 24,
so the full-translate inequalities used at size 100 cannot be reused.
Instead we prove a span estimate for small-defect cores, bound $2D$ using
several nested interval envelopes, and retain nonemptiness of a potentially
split paired unit in the only mode where dropping it is insufficient.

By Theorem~\ref{u96:thm:upper99}, a hypothetical 99-element set can be assumed
to contain zero. Proposition~\ref{u96:prop:lattice99} gives
\begin{equation}
\begin{gathered}
 d=91,\quad e=83,\quad a+b=82,\quad \{a,b\}=\{41-k,41+k\},\\
 \Lambda=1723+k^2,\quad\Omega=376-k^2,\quad G=170,\\
 \rho+2\vartheta+c+T=\Omega,\qquad
 T\geq\kappa+170-3z\geq170-2z,\\
 T\geq57,\quad T\geq z,\quad w\leq4z,
 \qquad\rho=\rho_P+\rho_Q.
\end{gathered}\label{u96:eq:98-budget}
\end{equation}
All slack quantities are nonnegative. For $z>0$, both side supports are
non-progressions and $\rho\geq54$.

\subsection{A stronger small-defect estimate}
\begin{lemma}\label{u96:lem:six-flexible98}
Let $X$ be a positive real set of size $r\geq15$, and let
$\rho=\rhoD(X)$. If $0\leq\rho<5(r-14)$, then
\[
 \mu(X)\leq\rho+22+2\lfloor\rho/5\rfloor
 \leq\tfrac75\rho+22.
\]
Consequently, a lower bound $\mu(X)\geq u$ implies
\[
 \rho\geq\min\left(5(r-14),
              \left\lceil\frac{5(u-22)}7\right\rceil\right).
\]
\end{lemma}
\begin{proof}
Use the ordered prefixes $X_j$ and the increments $y_j$ from the defect
recurrence. The counting part of Lemma~\ref{u96:lem:deleted}, independently of
its stronger hypothesis for the estimate on $\Phi$, gives
$\delta_0([1,j+q]\setminus H)\leq2q$ for $q\geq2$ and $j>q$.
Indeed every defect belongs to $[1,q]$ or $[j,j+q]$, and there are at most
$q-1$ defects in the former band: if all of $1,\ldots,q$ were present
and defective, defectivity of 1 would require 2 to be absent.
There are at most $q+1$ in the latter band. For $q=1$ the one-deletion
argument gives at most two defects (at the prefix sizes used here);
for $q=0$ there is one.

If every $y_j$ for $15\leq j\leq r$ were at least 7, nonnegativity of
$\rhoD(X_{14})$ and the recurrence would give
$\rho\geq5(r-14)$. Choose the largest $j\geq15$ with $y_j\leq6$.
The small-increment inverse lemma applies, since $6\leq j-5$, and makes
$X_j$, up to scaling, an interval with at most six deletions.
The preceding count gives $\delta_0(X_j)\leq12$.
Each subsequent step adds at most one defect, while all its increments
are at least 7; hence $r-j\leq\lfloor\rho/5\rfloor$ and
$\delta_0(X)\leq12+\lfloor\rho/5\rfloor$.
Now use $\mu=\rho-2+2\delta_0$. The final assertion follows by separating
$\rho\geq5(r-14)$ from its complement.
\end{proof}

For $25\leq r\leq57$, take the effective-hole rows with parameters
$\beta=7,9$ in Appendix~\ref{u96:app:upper98}. Put
\[
\begin{split}
 u(r,z)&=r+1+\max\bigl(\{C_7(r)\}\cup
                         \{C_9(r):z\geq19\}\bigr),\\
 v(r,z)&=\max\left(r-14,\ u(r,z)+2-2r,\
       \min\left(5(r-14),\left\lceil\frac{5(u(r,z)-22)}7\right\rceil\right)\right).
\end{split}
\]
For $z\geq15$ these give $\mu(P)\geq u(a,z)$,
$\mu(Q)\geq u(b,z)$ and $\rho\geq v(a,z)+v(b,z)$.
The initial bound $54+57\leq376-k^2$ gives $k\leq16$, so the side-size
range is valid. Define
\[
 C_k(z)=376-k^2-v(a,z)-v(b,z)-\max(57,z,170-2z).
\]
The minimum of the two mass-bonus alternatives, in both orientations,
minus $1764+4z$ is denoted by $M_k(z)$.
Exact evaluation of all $15\leq z\leq376-k^2$ gives:
\begin{center}\small
\begin{tabular}{rrrrr}
\toprule
$k$&Smallest surviving $z$&Largest surviving $z$&Maximum $C_k(z)$&Minimum $M_k(z)$\\
\midrule
0&15&130&73&151\\
1&15&129&72&$255/2$\\
2&15&127&70&$213/2$\\
3&15&121&64&98\\
4&15&121&64&$143/2$\\
5&30&111&54&$153/2$\\
6&36&99&42&$177/2$\\
7&41&88&31&$197/2$\\
8&48&74&17&117\\
9--16&\multicolumn{4}{c}{No state has $C_k(z)\geq0$}\\
\bottomrule
\end{tabular}
\end{center}
The verifier checks all 4658 states, not just the displayed endpoints.
Since $c\leq73<91$, all eighteen units are active. Each displayed margin
is positive, so Lemma~\ref{u96:lem:mass-bonus} excludes every many-hole case.

\subsection{The low-hole classification}
For $z=0$, equation~\eqref{u96:eq:98-budget} yields
\begin{equation}
 \rho+\kappa+2\vartheta+c\leq206-k^2.
 \label{u96:eq:98-zero}
\end{equation}
For $1\leq z\leq14$ it gives
$c\leq152-k^2+2z\leq180-k^2$.
Thus there are at most two empty units, and two are possible only at
$z=0$. If none is empty, the mass-bonus lemma gives
\[
 \min_{27\leq a\leq55}F_A(a,82-a)=1779>1764,
\]
and, for $z>0$,
\[
 \min_{28\leq a\leq54}
 \bigl(F_A(a,82-a)+(a-2)/2+(80-a)\bigr)
 =3667/2>1820.
\]
The $F_B$ alternative is larger in these ranges. These comparisons allow
arbitrarily many exceptional centers and exclude the no-empty case.

For one empty unit put $\delta=c-91$; for two put $\delta=c-182$.
A retained center outside its target hull loses at least
$L=\min(a,b)+5=46-k$ points. The exhaustive remaining regimes are:
\begin{center}\small
\begin{tabular}{llll}
\toprule
Regime&Budget on $\delta$&Range of $a$&Exceptional retained centers\\
\midrule
Z&$115-k^2$&31--51&None\\
P&$89-k^2$&32--50&None\\
ZE&$115-k^2$&33--49&Exactly one\\
PE&$89-k^2$&34--48&Exactly one\\
ZZ&$115-k^2$&36--46&Exactly two\\
EE&$24-k^2$&37--45&None\\
\bottomrule
\end{tabular}
\end{center}
The ranges follow from $\delta\geq0$, $\delta\geq L$, or
$\delta\geq2L$ as appropriate. In particular
\[
 3L-(115-k^2)=(k-3/2)^2+83/4>0,
 \quad 2L-(89-k^2)=(k-1)^2+2>0.
\]
In EE the entire residual budget is smaller than $L$.
With one empty unit, the outer-integrality and pair-choice proof of
Lemma~\ref{u96:lem:one-empty} applies. With two, only weak outer order is used.

\subsection{Small defect controls the core span}
\begin{lemma}\label{u96:lem:span98}
Suppose $X\subset\mathbb Z_{>0}$ has $r\geq16$ elements,
$\FS(X)=[0,\sum X]_{\mathbb Z}$, and
$0\leq\rhoD(X)\leq r-5$. Then
\[
 \max X\leq r+\rhoD(X)+1.
\]
\end{lemma}
\begin{proof}
Every prefix has full interval support: otherwise its first missing
integer below the next element can never be supplied by later elements.
Equivalently, fullness forces $x_1=1$ and
$x_i\leq1+x_1+\cdots+x_{i-1}$ at every step.
Put $q=x_r-r\geq0$. The last excess increment is $y_r=q$.
If $v$ counts old diagonal defects destroyed by adjoining $x_r$, then
\[
 \rhoD(X)=\rhoD(X_{r-1})+q-2+2v.
\]
Nonnegativity gives $q\leq\rhoD(X)+2$. If equality holds, then
$\rhoD(X_{r-1})=v=0$ and $2\leq q\leq r-3$.
Lemma~\ref{u96:lem:zero-core99} says that the prefix is either $[1,r-1]$ or
$[1,r]\setminus\{r-2\}$.
In the first case its defective maximum $r-1$ acquires the representation
$2(r-1)=x_r+(r-q-2)$, with the second summand in the prefix.
In the second case $r$ and $r-1$ are defective. If $q\ne2$, the identity
$2r=x_r+(r-q)$ destroys the former defect; if $q=2$, the identity
$2(r-1)=x_r+(r-4)$ destroys the latter. All these are sums of distinct
present elements. Thus $v>0$, a contradiction. The equality case is
impossible, proving the stated integer bound.
\end{proof}

In ZZ, two exceptions and \eqref{u96:eq:98-zero} imply
\[
 \rho+\kappa\leq23+2k-k^2,\qquad 0\leq k\leq5.
\]
In EE they imply $\rho+\kappa+\delta\leq24-k^2$, $0\leq k\leq4$.
Each bound on $\rho$ is strictly below $\min(a,b)-14$, so both side
supports are arithmetic progressions by Theorem~\ref{u96:thm:nonap}.
An integer progression of step at least two and cardinality $s$ has at
least $s-1$ hull holes. Since $h_P,h_Q\leq\kappa\leq24$ and
$s-1\geq T_{\min(a,b)}>24$, both steps are one. Thus
$h_P=h_Q=\kappa=0$. The bounds on $\rho$ also meet the hypotheses of
Lemma~\ref{u96:lem:span98}. Adding its two conclusions to the residual budget
gives
\begin{equation}
 R_p+R_q+\delta\leq
 \begin{cases}
 199-k^2,&\text{ZZ},\\
 106-k^2,&\text{EE}.
 \end{cases}\label{u96:eq:98-span-budget}
\end{equation}
In ZZ every retained center loses at most
$\delta-L\leq69+k-k^2\leq69$, so its hull count exceeds four.
The core-point quantile correction is therefore valid at every center.
For positive centers its relaxed form is
\[
 -v-l_v-R_p\leq-87,\qquad
 v-u_v-S_p-R_q\leq-86,
\]
with reflected inequalities for negative centers. These bounds use the
actual core spans, not an unjustified claim that they are at most $a+1,b+1$.

\subsection{Two additional mass and sumset estimates}
In all regimes, distinct positive integers in the core give
$\sum P\geq T_{a-1}+R_p$. The four positive elements of $D\setminus E$
are at least $R_p+1,R_p+2,R_p+3,M_p$. Consequently
\begin{equation}
 S_p\geq T_{a-1}+4R_p+6+M_p,
 \qquad S_q\geq T_{b-1}+4R_q+6+M_q.
 \label{u96:eq:98-strong-mass}
\end{equation}
These inequalities are imposed alongside, not added to, the earlier
mass inequalities.

A useful discrete improvement holds if $M_q-R_q=4$.
Then the four negative magnitudes outside the core are exactly
$R_q+1,\ldots,R_q+4$, and their ten unordered pairs with repetition
have only seven possible sums. All other sums of two elements of $D$
lie in $[-M_q-R_q,M_p+R_p]$, except possibly the ten positive-positive
outlier pairs. Hence
\begin{equation}
 |2D|\leq M_p+M_q+R_p+R_q+18,
 \qquad F\leq1758\quad(z=0).
 \label{u96:eq:98-gap4}
\end{equation}
This improves the usual upper bound $1764$ by six. At the one ZE mode
specified in the appendix we split the integer gap into $M_q-R_q=4$
and $M_q-R_q\geq5$. The first uses \eqref{u96:eq:98-gap4}; the second adds
its gap inequality to the model. This avoids turning the non-strict
bound $F\geq1764$ into a false contradiction.

For EE expose the four inner extremes on each side:
\[
 x_0=R_p<x_1<x_2<x_3<x_4=M_p,
 \quad y_0=R_q<y_1<y_2<y_3<y_4=M_q.
\]
All are integers and successive gaps are at least one. Also
\[
 S_p\geq T_{a-1}+R_p+x_1+x_2+x_3+x_4,
 \quad S_q\geq T_{b-1}+R_q+y_1+y_2+y_3+y_4.
\]
\begin{lemma}[Nested interval envelopes]\label{u96:lem:envelope98}
For $j=0,1,2,3$, put $W_j=x_j+y_j$ and $t_j=2(4-j)$. Then
\begin{equation}
 |2D|\leq2W_j+1+\binom{t_j+1}{2}
 +\sum_{i=j+1}^4
  \bigl(\min(W_j+1,x_i-x_{i-1})+
        \min(W_j+1,y_i-y_{i-1})\bigr).
 \label{u96:eq:98-envelope}
\end{equation}
\end{lemma}
\begin{proof}
Let $D_j$ consist of $E$ and the first $j$ added elements on each side.
Replace $D_j$ by its integer hull $[-y_j,x_j]$.
Its double sum occupies $2W_j+1$ integers. Adjoining the translates
$x_i+[-y_j,x_j]$ in increasing order enlarges their union with this
initial interval by at most $\min(W_j+1,x_i-x_{i-1})$ at step $i$.
For the first step the preceding upper endpoint is $2x_j$; for later
steps it is $x_{i-1}+x_j$. The negative translates give the reflected
bounds. Adding the two one-sided enlargement bounds is still an upper
bound on their combined union, even if some intervals overlap.
The only remaining sums use two of the $t_j$ external elements, with
repetition allowed. There are at most $\binom{t_j+1}{2}$ such sums.
This covers every element of $2D$ and proves the inequality.
\end{proof}
In the model put $K=|2D|$. The minima in \eqref{u96:eq:98-envelope} can be
replaced by nonnegative variables $g$ bounded above by both of their
arguments: an actual set realizes the model by taking $g$ equal to the
minimum. Since $w=0$ in EE, the exact identity defining $c$ gives
\begin{equation}
 S_p+S_q+\delta-2K=18\cdot91+84-182=1540.
 \label{u96:eq:98-mass-cardinality}
\end{equation}
We impose only the corresponding upper inequality, a valid relaxation.

\subsection{The split pair and complete finite certification}
Use the mode labels of Appendix~\ref{u96:app:upper99}. Among the EE modes,
one needs an extra necessary condition: branch A, missing units
$p_0,q_5$. The second pair $\beta_p,\beta_q$ is present, though its
individual constraints were dropped to allow splitting.
Its combined contribution is at least $91-\delta>0$.
Thus at least one of its two individual slots contributes. If the positive
one contributes, its center is integral and its largest translated point
is $\beta_p+M_p\geq1$. If the negative one contributes, likewise
$\beta_q+M_q\geq1$. It is therefore enough to check the two relaxations
\begin{equation}
 \beta_p+M_p\geq1\qquad\text{or}\qquad
 \beta_q+M_q\geq1.
 \label{u96:eq:98-split}
\end{equation}
No assertion about which member contributes, or whether both do, is made.
In particular this step does not require integrality of all eight outer
elements.

Appendix~\ref{u96:app:upper98} gives all rows and all 77 nonnegative integer
row combinations. They certify the following complete case enumeration:
\begin{center}\small
\begin{tabular}{lrrr}
\toprule
Regime&Systems&Strict lower-bound certificates&Inconsistency certificates\\
\midrule
Z&1428&1344&84\\
P&1292&1216&76\\
ZE&39237&36361&2876\\
PE&34620&32056&2564\\
ZZ&748&704&44\\
EE&4932&4129&803\\
\midrule
Total&82257&75810&6447\\
\bottomrule
\end{tabular}
\end{center}
For every system the checker regenerates its rows, verifies a nonnegative
integer combination coefficient by coefficient, and compares the resulting
bound with that system's upper bound using exact fractions. The gap-four
ZE system uses 1758; all other zero-hole systems use 1764, and the positive
regimes use 1820. Thus the displayed system count includes one additional
gap branch and nine additional split-pair branches, and no endpoint or
orientation is inferred by interpolation.

\begin{theorem}\label{u96:thm:upper98}
Every finite real set satisfying $\FS(A)\subseteq4A$ has at most 98
elements. Consequently
\[
 \boxed{N\leq98}.
\]
\end{theorem}
\begin{proof}
The preceding theorem already excludes all sizes greater than 99.
At size 99, adjoining zero is legitimate because size 100 has been
excluded. The common-lattice reduction, the many-hole table, and the
no-empty-unit comparisons leave exactly the six regimes above.
The small-span, mass, quantile, envelope, and split-pair arguments prove
that any remaining actual configuration realizes at least one enumerated
linear system. Every such system has a strict exact certificate of
impossibility. Hence size 99 is impossible. 
\end{proof}

\subsection{A further common-lattice reduction at size 98}
\begin{proposition}\label{u96:prop:lattice98}
At size 98 the common-lattice and hole conclusions hold, with
\[
\begin{gathered}
 d=90,\quad e=82,\quad (a,b)=(40-k,41+k),\quad I=k(k+1),\\
 \Lambda=1682+I,\quad\Omega=395-I,\quad G=168.
\end{gathered}
\]
This is a structural reduction, not an exclusion of size 98.
\end{proposition}
\begin{proof}
Apply Proposition~\ref{u96:prop:universal-lattice} at the stated size(s).
The displayed parameters are obtained from $d=n-8$, $e=n-16$,
$a+b=n-17$, $\Lambda=1+T_a+T_b$ and $\Omega=22d+97-\Lambda$.
Where the statement uses $G=2d-16$ rather than $2d-12$, it is a weaker
consequence of the same proved inequalities. If needed, zero is present
because the next larger size has already been excluded; the proposition
itself assumes zero and does not use that endpoint exclusion.
\end{proof}

\section{Two further exclusions: the general upper bound 96}\label{u96:sec:upper96}
We exclude sizes 98 and 97, in that order. The new finite
proof uses integer inconsistency certificates and, for 19 roots at size
97, exhaustive binary disjunctions. The disjunctions are proved for actual
sets; no claim that they partition every point of the relaxed polyhedron
is needed or made. The certificate files accompanying this report are
part of the finite proof.

\subsection{A span estimate that permits holes}
For a positive integer set $X$, write
$h(X)=\sum X+1-|\FS(X)|$. This is its number of missing integer points
in the full subset-sum hull, on the specified unit lattice.

The universal span estimate of Lemma~\ref{u96:lem:span96} applies here.

Put $m=a+b=n-17$, $d=n-8$, $\rho=\rho_P+\rho_Q$, and
$\mu=\mu(P)+\mu(Q)$. On the common core lattice, the lemma gives
\begin{equation}
 R_p+R_q\leq m+\rho+\kappa+z+4.
 \label{u96:eq:span96}
\end{equation}
After $D$ has been placed on that lattice, write $j$ for the number of
empty covering units and $\delta=c-jd$. Combining the exact budget with
$T\geq\kappa+G-3z$, where $G=2d-12$, gives the useful joint bound
\begin{equation}
 R_p+R_q+\delta\leq m+4+\Omega-G-jd+4z.
 \label{u96:eq:span96-budget}
\end{equation}
The same $z$ occurs in the loss, mass, cardinality, and span constraints.
Replacing it by unrelated worst-case values can lose the contradiction.

\begin{lemma}[Flexible defect bound with an integer inverse]\label{u96:lem:flex96}
Let $X$ be a positive real set of size $r\geq15$, and let
$2\leq\beta\leq10$ be an integer. If
$0\leq\rhoD(X)<(\beta-1)(r-14)$, then
\[
 \mu(X)\leq\rhoD(X)+4\beta-2+
                    2\left\lfloor\frac{\rhoD(X)}{\beta-1}\right\rfloor.
\]
For an integer lower bound $\mu(X)\geq u$, put
$v=\max(0,u-4\beta+2)$ and write $v=(\beta+1)q+s$, $0\leq s\leq\beta$.
Then
\begin{equation}
 \rhoD(X)\geq\min\bigl((\beta-1)(r-14),
                  (\beta-1)q+\min(s,\beta-1)\bigr).
 \label{u96:eq:flex96-inverse}
\end{equation}
\end{lemma}
\begin{proof}
If all increments $y_i$ for $15\leq i\leq r$ exceed $\beta$, the
recurrence $\rho_i\geq\rho_{i-1}+y_i-2$ contradicts the hypothesis.
Choose the last index $i\geq15$ with $y_i\leq\beta$.
Since $\beta\leq i-5$, Lemma~\ref{base:lem:small-increment} gives an interval
with at most $\beta$ deletions. The defect-counting argument in
Lemma~\ref{u96:lem:six-flexible98} uses only $i>\beta$ and gives at most
$2\beta$ defects; the zero-deletion case has one defect.
Every subsequent step adds at most one defect and raises $\rho$ by at
least $\beta-1$. Nonnegativity at index $i$ therefore gives
$r-i\leq\lfloor\rhoD(X)/(\beta-1)\rfloor$. Use
$\mu=\rhoD-2+2\delta_0$.
Finally, the least nonnegative integer $t$ satisfying
$t+2\lfloor t/(\beta-1)\rfloor\geq v$ is
$(\beta-1)q+\min(s,\beta-1)$: each complete block of $\beta-1$ in $t$
adds $\beta+1$ to the left side. This proves the stated inverse, including
the alternative where the small-defect hypothesis fails.
\end{proof}

\subsection{The lattice reduction at sizes 97 and 96}
\begin{proposition}\label{u96:prop:lattice96}
For $n\in\{97,96\}$, a feasible integer set of size $n$ containing zero
admits the common-core normalization in which $D$ is integral and
\[
 w\leq4z,\quad T\geq\kappa+G-3z\geq G-2z,
 \quad T\geq\lceil G/3\rceil,\quad T\geq z.
\]
The parameters are
\begin{center}
\begin{tabular}{rrrrrrrr}
\toprule
$n$&$d$&$e$&$m$&$\Lambda_0$&$\Omega_0$&$I$&$G$\\\midrule
97&89&81&80&1641&414&$k^2$&166\\
96&88&80&79&1601&432&$k(k+1)$&164\\\bottomrule
\end{tabular}
\end{center}
Here $(a,b)=(\lfloor m/2\rfloor-k,\lceil m/2\rceil+k)$, up to reflection,
$\Lambda=\Lambda_0+I$, and $\Omega=\Omega_0-I$.
\end{proposition}
\begin{proof}
The deficient-sign argument forces $|B|=8$: otherwise a one-sided core
of at least 80 nonzero elements would give
$s\geq1+T_{80}=3241>24\cdot97+77$. The unpaired-value bound is
$q^2\leq2288$ at $n=97$ and $q^2\leq2353$ at $n=96$.
Parity gives $q\leq46,47$, respectively, so both sides have at least
17,16 elements. The nonnegative budget applies and implies $k\leq20$.
For all these integer $k$, $\Omega<T_{\min(a,b)}-2$. Thus the simultaneous
inverse theorem supplies the common core lattice and
$h_P,h_Q\leq\kappa$, $z=h_P+h_Q-\kappa\leq\kappa$.

It remains to put the peeled elements on the same lattice. Apply
\eqref{u96:eq:span96} in the second-layer comparison, before asserting the
integrality of $D$. With $C=2m+4e+14-\Lambda$ this gives
\begin{align}
 |\FS(D_1)|-2s
 &\leq C+2\rho-\mu+\kappa+2z+\ell_2\notag\\
 &\leq C+3\Omega-\rho-\mu,\label{u96:eq:lattice96-three}\\
 |\FS(D_1)|-2s
 &\leq C+\Omega+\rho-\mu+2z.\label{u96:eq:lattice96-one}
\end{align}
These follow from $\rho+\kappa+\ell_2\leq\Omega$ and $z\leq\kappa$.
The maximum of each positive side is always defective; hence
$\rho\leq\mu$.

At $n=97$, when $z=0$, the first line above also gives the coarse bound
$C+2\Omega=-315-3I<0$ (use $\rho\leq\Omega$ and $\mu\geq0$).
When $z>0$, both side supports are non-progressions: a progression with
a larger integer step would have at least $|\FS|-1$ holes, contrary to
$h\leq\kappa<|\FS|-1$.
Consequently $\rho\geq m-28=52$, $\mu\geq m-6=74$, and
\eqref{u96:eq:lattice96-three} is at most $-27-4I<0$.

At $n=96$, equation~\eqref{u96:eq:lattice96-one} is at most
$C+\Omega+28=-649-2I<0$ for $z\leq14$.
For $z\geq15$, non-progression gives $\rho+\mu\geq2m-34=124$.
As $C+3\Omega=187-4I$, a remaining case must have $k\leq3$ and
$a,b\geq36$. The effective-hole test with $(r_0,C,\beta)=(36,109,7)$
gives $\mu(P)\geq a+110$, $\mu(Q)\geq b+110$.
Its strict hypotheses are verified by the row at $r=36$ in
Appendix~\ref{u96:app:upper96}; they continue to hold for larger $r$.
Thus $\rho+\mu\geq51+299=350>187$, again strictly.
The first-layer comparisons are
$1641>4\cdot89-10+414=760$ and
$1601>4\cdot88-10+432=774$.
The earlier overlap and hole-propagation arguments now apply, followed
by the four-extreme estimate. This proves all stated conclusions.
\end{proof}
The proposition assumes zero is already present. In its use below, zero
is adjoined only after the larger endpoint has been excluded. The result
at size 96 is a structural reduction, not an exclusion of that size.

\subsection{One model format for both exclusions}
The outer centers and eighteen covering units are those of
Appendix~\ref{u96:app:upper99}. Their exact labels and the new row formulas
are recorded in Appendix~\ref{u96:app:upper96}. We recall the validity issues
that are important for completeness.

There are two pairing branches, according as $\beta_p+\beta_q\leq0$
or $>0$. Reflect to put the contributing first-pair member on the positive
side when that pair is nonempty. If it is empty, retain both possibilities
for the second pair. With zero or one empty unit, the outer elements are
integral by the contributing-center argument in Lemma~\ref{u96:lem:one-empty}.
With two empty units, only weak outer order is imposed. Every retained
contributing center is still integral, since both $D$ and $\FS(D)$ are
integral.

In pairing branch A, if an unpaired center $p_5$ or $q_5$ is missing,
the second pair may have two contributors. Its individual centers are
dropped from the usual loss rows. If this paired unit is present, choose
one \emph{actually contributing} member, denoted by $s=p$ or $q$, and
retain the necessary inequality
\begin{equation}
 \beta_s+M_s\geq1.\label{u96:eq:chosen-member96}
\end{equation}
This row is not asserted to be sufficient for contribution. The choice
is made from the actual set, not from the solutions of the relaxation.
This distinction will justify the sole-contributor disjunction below.
The resulting numbers of modes with $j=0,1,2$ empty units are
$4,70,578$, including these member choices.

For every retained center let $\lambda_v=l_v+u_v$ be its number of hull
losses. The covering inequality gives
$\sum\lambda_v\leq\delta=c-jd$, also when a split pair has been dropped.
For a positive center the following relaxed quantile rows are available:
\begin{center}\small
\begin{tabular}{lll}
\toprule
Type&Lower row&Upper row\\\midrule
Ordinary&$-v-l_v\leq-b-5$&$v-u_v-S_p\leq-a-4$\\
Global&$-v-l_v-M_p\leq-d$&$v-u_v-S_p-M_q\leq-(d-1)$\\
Core&$-v-l_v-R_p\leq-(m+5)$&$v-u_v-S_p-R_q\leq-(m+4)$\\\bottomrule
\end{tabular}
\end{center}
Reflection gives the negative rows. Ordinary rows apply when the relevant
endpoint is inside the hull. Global rows apply to any nonempty center.
Core rows apply when its hull count exceeds four. Indeed, an out-of-hull
translate with more than four points must use a core point, since there
are only four removed positive and four removed negative elements.
The core-gap counting in the earlier quantile proof then gives these
rows; in-hull ordinary rows imply them using $R_p\geq a$, $R_q\geq b$.

Let $x_0=R_p<x_1<x_2<x_3<x_4=M_p$ be the positive core endpoint and the
four positive elements of $D\setminus E$; define $y_i$ by reflection.
All eight gaps are at least one. The mass and four nested-envelope
inequalities of Section~\ref{u96:sec:upper98} are imposed. In particular,
with $K=|2D|$ and $\varepsilon_p=\mu(P)+h_P$,
\begin{gather}
 S_p\geq T_a+\varepsilon_p+3R_p+6+M_p,\qquad
 S_p\geq T_{a-1}+4R_p+6+M_p,\label{u96:eq:mass96-short}\\
 S_p\geq T_a+\varepsilon_p+\sum_{i=1}^4x_i,\qquad
 S_p\geq T_{a-1}+R_p+\sum_{i=1}^4x_i,\label{u96:eq:mass96-full}\\
 S_p+S_q+\delta-2K\leq18d+84-jd+4z.
 \label{u96:eq:card96}
\end{gather}
The last inequality follows from the exact identity with $w$ in place
of $4z$. It does not rely on an estimate of the hull length of $2D$.

\subsection{Excluding size 98}
By the preceding upper bound, a hypothetical 98-element set contains
zero after adjoining it if necessary. Integralization and the corrected
Proposition~\ref{u96:prop:lattice98} give
\[
 d=90,\quad e=82,\quad m=81,\quad
 (a,b)=(40-k,41+k),\quad I=k(k+1),\quad
 \Omega=395-I,\quad G=168.
\]
For $z=0$, $c\leq227-I$. For $1\leq z\leq14$, both supports are
non-progressions, so $\rho\geq53$ and $c\leq174-I+2z\leq202-I$.
There are at most two empty units in this range. Set
$B_0=227-j\cdot90-I$. The zero-hole models use
\[
 \delta\leq B_0,\quad R_p+R_q+\delta\leq85+B_0,
 \quad \varepsilon_p,\varepsilon_q\geq0.
\]
The positive low-hole models retain a single variable $1\leq z\leq14$
and use
\begin{gather}
 \delta\leq B_0-53+2z,\qquad
 R_p+R_q+\delta\leq85+B_0+4z,\notag\\
 \varepsilon_p\geq a-3+z,\qquad
 \varepsilon_q\geq b-3+z.\label{u96:eq:low98}
\end{gather}
The mass and cardinality rows use that same $z$.

For one empty unit, $\delta\leq137$ or $112$, both less than
$2(d-4)=172$. Therefore at most one retained center can have
$\lambda_v\geq86$. If none does, all core quantiles are valid. If one
does, designate it and use both global endpoint rows for it, with
$\lambda_v\geq86$; all other centers use core rows.
For two empty units, $\delta\leq47$ or $22$, so every retained center
has hull count greater than four. In the positive-hole case the budget
is even below $\min(a,b)+5$, giving ordinary quantiles at all centers.

The finite low-hole root list is:
\begin{center}\small
\begin{tabular}{lrrlr}
\toprule
Regime&$j$&Range of $a$&Quantile rule&Roots\\\midrule
Z0&0&26--55&Global&120\\
P0&0&27--54&Global&112\\
Z1&1&29--52&Core&1680\\
P1&1&30--51&Core&1540\\
ZL&1&34--47&One designated large loss&16604\\
PL&1&36--45&One designated large loss&11860\\
Z2&2&34--47&Core&8092\\
P2&2&36--45&Ordinary&5780\\\bottomrule
\end{tabular}
\end{center}
Z means $z=0$, P means $1\leq z\leq14$, and L designates the large-loss
center. These are relaxations: a listed model need not impose the strict
small-loss bounds used to select it from an actual set.

For $z\geq15$, the initial penalty gives $k\leq16$ and $24\leq a,b\leq57$.
Use the hole rows in Appendix~\ref{u96:app:upper96} to define
\begin{align*}
 u(r,z)&=\max\bigl(\{r-3\}\cup
       \{r+C_\beta(r)+1:\beta\in\{7,9,10\},\ z>2\beta,
                                  \ C_\beta(r)\text{ listed}\}\bigr),\\
 v(r,u)&=\max\{r-14,\ u+2-2r,\ L_8(r,u),\ L_9(r,u)\},
\end{align*}
where $L_\beta$ is the right side of \eqref{u96:eq:flex96-inverse}.
Then $\mu(P)\geq u(a,z)$ and $\rho_P\geq v(a,u(a,z))$, with reflected
inequalities on $Q$. Consequently
\[
 c\leq C_k(z):=395-I-v(a,u(a,z))-v(b,u(b,z))
                         -\max(56,z,168-2z).
\]
Evaluate every integer $15\leq z\leq395-I$, $0\leq k\leq16$:
there are 4845 such states. Whenever $C_k(z)\geq0$, the four possible
all-active mass-bonus margins, using both orientations, have minimum
at least $63/2>0$. Thus no all-active state survives. A state with an
empty unit requires $C_k(z)\geq90$, and the only possibilities are
\[
 (k,z,C_k(z),u(a,z),u(b,z))
 =(0,56,91,184,193),\quad(0,57,90,184,193).
\]
Both are covered by a single family H1, in both orientations, using
\[
 \delta\leq1,\quad R_p+R_q+\delta\leq450,\quad
 (\varepsilon_p,\varepsilon_q)\geq(240,249),\quad w\leq228,
\]
or its reflection, and ordinary quantiles. Its cardinality row has
right side $1842$. There are $2\cdot70=140$ roots in H1.

\begin{proposition}\label{u96:prop:exclude98}
No feasible real set has 98 elements. Hence $N\leq97$.
\end{proposition}
\begin{proof}
The preceding classification assigns every actual set to one of 45928
roots. The model rows are necessary by the span, mass, envelope, loss,
and pairing arguments. The 179 integer certificates in
\texttt{upper97\_exact.json}, containing 5888 nonzero weights, give
inconsistency of every root. Explicitly, for its rows $Ax\leq b$, $x\geq0$,
the selected nonnegative integer vector $w$ satisfies
$w^{\mathsf T}A\geq0$ coordinatewise and $w^{\mathsf T}b<0$.
Thus $0\leq w^{\mathsf T}Ax\leq w^{\mathsf T}b<0$, a contradiction.
The independent checker regenerates the complete ordered root list and
checks each combination using arbitrary-precision integers.
\end{proof}

\subsection{Classifying all size-97 roots}
Zero can now be adjoined in a hypothetical 97-element set.
Proposition~\ref{u96:prop:lattice96} gives
$d=89$, $e=81$, $m=80$, $(a,b)=(40-k,40+k)$,
$I=k^2$, $\Omega=414-I$, and $G=166$.
For $z=0$, $c\leq248-I$; for $1\leq z\leq14$,
$c\leq196-I+2z\leq224-I$. Again there are at most two empty units.
Put $B_0=248-j\cdot89-I$. Zero-hole rows are
\[
 \delta\leq B_0,\quad R_p+R_q+\delta\leq84+B_0,
 \quad \varepsilon_p,\varepsilon_q\geq0.
\]
The positive-hole rows are
\begin{gather}
 \delta\leq B_0-52+2z,\quad R_p+R_q+\delta\leq84+B_0+4z,
 \quad1\leq z\leq14,\notag\\
 \varepsilon_p\geq a-3+z,\qquad
 \varepsilon_q\geq b-3+z.\label{u96:eq:low97}
\end{gather}
The cardinality row is
$S_p+S_q+\delta-2K\leq1686-j\cdot89+4z$.
With one empty unit, $\delta\leq159$ or $135$, both strictly less than
$2(d-4)=170$. At most one retained center has loss at least 85.
With two empty units, $\delta\leq70$ or $46$, both below 85, so every
retained center uses core quantiles. Unlike at size 98, ordinary
quantiles cannot be assumed throughout the positive two-empty family.

\begin{center}\small
\begin{tabular}{lrrlr}
\toprule
Regime&$j$&Range of $a$&Quantile rule&Roots\\\midrule
Z0&0&25--55&Global&124\\
P0&0&26--54&Global&116\\
Z1&1&28--52&Core&1750\\
P1&1&29--51&Core&1610\\
ZL&1&32--48&One designated large loss&20162\\
PL&1&33--47&One designated large loss&17790\\
Z2&2&32--48&Core&9826\\
P2&2&34--46&Core&7514\\\bottomrule
\end{tabular}
\end{center}
There are 58892 low-hole roots. Each large-loss family loops over every
retained center; the 70 one-empty modes together have 1186 such centers.

For $z\geq15$, the initial inequality $52+56\leq414-k^2$ gives $k\leq17$,
so $23\leq a,b\leq57$. Use the same functions $u,v$ as above and put
\[
 C_k(z)=414-k^2-v(a,u(a,z))-v(b,u(b,z))
                                  -\max(56,z,166-2z).
\]
The mass-test margin $M_k(z)$ is the minimum of the four expressions
\[
 \begin{split}
 F_A(a,b)+\tfrac12(u(a,z)+z)+u(b,z)+z,\\
 F_B(a,b)+\tfrac23(u(a,z)+z)+u(b,z)+z,
 \end{split}
\]
and their reflections, minus $1728+4z$.
The exact enumeration retains every state with $C_k(z)\geq0$ and finds
$C_k(z)\leq118<178$. Thus at most one unit is empty.
For no empty unit, states with $M_k(z)>0$ are excluded by
Lemma~\ref{u96:lem:mass-bonus}. There remain 128 $(k,z)$ states, or 251
oriented $(a,z)$ states. For one empty unit, $C_k(z)\geq89$ leaves
175 $(k,z)$ states, or 305 oriented states. In the latter cases
$\delta\leq C_k(z)-89\leq29<\min(a,b)+5$, so all retained centers obey
ordinary quantiles.

For each remaining high-hole parameter state, fix $z$ to that integer;
use $\delta\leq C_k(z)-j\cdot89$,
$\varepsilon_p\geq u(a,z)+z$, $\varepsilon_q\geq u(b,z)+z$, and retain
\eqref{u96:eq:span96-budget} and \eqref{u96:eq:card96} with the same fixed $z$.
Use global quantiles for H0, ordinary ones for H1. These give
$251\cdot4=1004$ H0 roots and $305\cdot70=21350$ H1 roots.
Together with the low-hole families there are 81246 roots.

\subsection{Three valid disjunctions and the 19 remaining roots}
Most size-97 roots are already inconsistent. The remaining 19 are covered
by finite binary trees. Each split has one of the following three proofs.

\paragraph{An ordinary endpoint or an outside endpoint.}
For a positive contributing center $v$, the lower endpoint satisfies
\begin{equation}
 -v-l_v\leq-b-5
 \quad\hbox{or}\quad (l_v\geq b+5\ \hbox{and}\ v\leq-1).
 \label{u96:eq:branch-normal-low96}
\end{equation}
If $v\geq0$, the ordinary counting proof applies. Otherwise the integral
center has $v\leq-1$, and all $b+4$ negative points and zero are lost
below 1. Similarly, the upper endpoint satisfies
\begin{equation}
 v-u_v-S_p\leq-a-4
 \quad\hbox{or}\quad (u_v\geq a+5\ \hbox{and}\ v\geq S_p+1).
 \label{u96:eq:branch-normal-high96}
\end{equation}
Reflection gives both negative-side versions. Only integrality of the
contributing center and $D$ is used, not integrality of every outer element.

\paragraph{Inner extreme or endpoint loss.}
For each $i=1,2,3,4$, the positive center satisfies
\begin{equation}
 l_v\geq i\ \hbox{or}\ v\geq y_{5-i}+1,
 \qquad
 u_v\geq i\ \hbox{or}\ v\leq S_p-x_{5-i}.
 \label{u96:eq:branch-inner96}
\end{equation}
For example, if the integer count $l_v<i$, at least one of the $i$
largest negative magnitudes remains above the lower endpoint; in
particular the smallest of those $i$ gives $v-y_{5-i}\geq1$.
The upper version is the reflected counting assertion. Again include the
versions with $p,q$ exchanged. These disjunctions apply to actual integer
loss counts. They are not asserted for arbitrary fractional loss variables.

\paragraph{Both paired members contribute, or the selected one is sole.}
Consider a two-empty mode with a split second pair and a member $s$
chosen from actual contributors as in \eqref{u96:eq:chosen-member96}.
Either the other member $t$ also contributes, in which case
$\beta_t+M_t\geq1$, or $s$ is the sole contributor. In the latter case
the paired unit's capacity loss is at least its individual hull loss.
Introduce $l_*,u_*\geq0$ for this hull loss and strengthen the total-loss
row to
\[
 \sum_{v\text{ already retained}}(l_v+u_v)+l_*+u_*\leq\delta.
\]
Its hull count is at least $d-\delta>4$, because the two-empty budgets
are at most 70 and 46. It therefore satisfies its own core quantiles:
\begin{equation}
 -\beta_s-l_*-R_s\leq-85,\qquad
 \beta_s-u_*-S_s-R_t\leq-84.\label{u96:eq:sole96}
\end{equation}
Only two unused nonnegative coordinates are needed for $l_*,u_*$.
All six trees using this split finish after its two children.
In particular, the argument does not replace actual contribution by the
merely necessary numerical row \eqref{u96:eq:chosen-member96}.

The trees and all assignments are explicit in
\texttt{upper96\_exact.json}. A summary of every branched root is printed
in Appendix~\ref{u96:app:upper96}. There are 19 such roots, 130 internal
nodes and 149 leaves. The other 81227 roots need no split. Thus there
are exactly 81376 certified leaf systems, counting the ordinary roots.
At every internal node the checker visits \emph{both} children and adds
the corresponding proved rows. No child is pruned on numerical evidence.

\begin{theorem}\label{u96:thm:upper96}
Every feasible finite real set has at most 96 elements. Consequently,
\[
 \boxed{N\leq96}.
\]
\end{theorem}
\begin{proof}
Proposition~\ref{u96:prop:exclude98} supplies $N\leq97$.
A hypothetical 97-element set can therefore be assumed to contain zero,
then integralized. Its common-lattice normalization satisfies one of the
81246 root models just classified. Follow any satisfied child at each
node of its finite tree. The disjunction proofs guarantee such a child
for the actual configuration, with the auxiliary hull losses chosen
from that configuration.

The selected leaf is inconsistent: the 303 certificates in
\texttt{upper96\_exact.json} have 10514 positive integer weights in total,
and certify $w^{\mathsf T}A\geq0$, $w^{\mathsf T}b<0$ at every leaf.
The standard-library verifier independently regenerates every root and
every allowed split and checks all these inequalities exactly. Hence
there is no feasible set of size 97. The general upper bound follows.

\end{proof}

The finite proof at sizes 98 and 97 consists of 127174 roots in the two exclusions and
127304 leaf contradictions. Numerical linear programming was used only
to discover multiplier vectors and useful disjunctions; none is called
by the delivered verification program. The universal reductions, the
assignment of actual sets to roots, and the exhaustiveness of each split
are separate mathematical arguments given above. 



\clearpage\part{The descent from 96 to 86}\label{part:upper86}

\section{The descending chain from 96 to 86}
\label{u86:sec:entry}
Theorem~\ref{u96:thm:upper96} establishes the global bound $N\le96$. We now exclude
sizes $96,95,\ldots,87$ in that order. The only cardinality premise at
stage $n$ is the already established bound $N\le n$. A feasible $n$-set
omitting zero would produce a feasible $(n+1)$-set by adjoining zero;
therefore a hypothetical size-$n$ set contains zero. Integralization
preserves its cardinality and this zero. No feasibility assertion for an
arbitrary subset is used.

\paragraph{Finite certificates.}
This part supplies the actual-set reductions, every added inequality
family, the exhaustive case and split rules, and the exact finite lemma
specifications. The accompanying certificate data and exact checkers are
part of the proof. Their complete verification is recorded in
Section~\ref{completion:verification}; its mathematical role is
Lemma~\ref{completion:coverage}.

\paragraph{Inputs already proved in the earlier chapters.}
We use the universal four-extreme covering and its 110 fixed labels with
26 disjoint nonpositive deletion groups; both peeling inequalities; the
increasing-chain, threshold-14, small-increment and flexible-defect inverse
lemmas; the parameterized effective-hole theorem with the corrected
endpoint bridge; the eighteen-unit pairing and contributor grammar; the
all-active mass certificate; and the size-96 core-lattice theorem.
These are mathematical lemmas of the preceding $N\le96$ proof, not a
second assumed upper bound. They are stated and proved in Parts~\ref{part:foundations} and~\ref{part:upper96}. The sharp four-extreme theorem
and all stronger ingredients first needed below are proved here.

\subsection{Core notation and the exact budget}\label{u86:sec:budget}
Let $A$ be feasible, integral, and contain zero. Remove its four largest
positive and four smallest negative elements, taking all elements of a
sign if fewer than four exist; call this removed set $B$, and put
$D=A\setminus B$, $d=|D|$. Peel the two largest positive and two smallest
negative elements twice more, obtaining
\[
 D\supset D_1\supset E=P\cup\{0\}\cup(-Q),\quad
 |D_1|=d-4,\quad e=|E|=d-8,
\]
where $P,Q$ are positive sets. The endpoint arguments below ensure both
signs have enough elements for these choices. Put
\[
 a=|P|,\quad b=|Q|,\quad m=a+b=e-1,\quad
 \Lambda=1+T_a+T_b,\quad s=|\FS(E)|.
\]
For any positive $r$-element set $X$, define
\begin{align*}
 \mu(X)&=|\FS(X)|-1-T_r,\\
 \delta_0(X)&=|\{x\in X:2x\notin R(X\cup\{0\})\}|,\\
 \rho(X)&=\mu(X)+2-2\delta_0(X).
\end{align*}
Write $\mu=\mu(P)+\mu(Q)$ and $\rho=\rho(P)+\rho(Q)$.
The increasing-chain bound gives $\mu(X)\ge0$. Theorems~\ref{u96:thm:phi14} and~\ref{u96:thm:nonap} give
\begin{equation}\label{u86:eq:input-rho}
 \rho(X)\ge0\ (r\ge14),\qquad
 \rho(X)\ge r-14\quad\hbox{if $r\ge6$ and $\FS(X)$ is not an AP}.
\end{equation}
The stability theorem of Carpenter--Defant--Kravitz \cite{CDK2026} additionally
gives $\mu(X)\ge r-3$ for a non-AP support in the range used here.
Since the maximum of $X$ is defective, $\rho(X)\le\mu(X)$.

The covering and peeling inequalities are
\begin{align}
 |\FS(D)|&\le2|2D|+18d+85,\label{u86:eq:cover}\\
 |\FS(F)|&\ge|\FS(F')|+|R(F)|-1,\label{u86:eq:peel}\\
 s&\le22d+2\Delta(D)+77\le24d+77.\label{u86:eq:corebound}
\end{align}
Here $F'$ is a four-extreme peel of $F$. The proof of \eqref{u86:eq:cover}
includes a fixed, explicitly checked 110-slot covering and 26 disjoint
deletion groups; it is not an asymptotic estimate. The last inequality
uses $\Delta(D)\le d$.

If $a,b\ge14$ and $E$ has an opposite pair, define
\begin{align*}
 \kappa&=s-\Lambda-\mu,&
 \vartheta&=\delta_0(P)+\delta_0(Q)+8-\Delta(D),\\
 c&=2|2D|+18d+85-|\FS(D)|,\\
 \ell_1&=|\FS(D)|-|\FS(D_1)|-|R(D)|+1,\\
 \ell_2&=|\FS(D_1)|-s-|R(D_1)|+1,\\
 g&=4d-10-(|R(D)|-|R(D_1)|),&
 T&=\kappa+\ell_1+\ell_2+g.
\end{align*}
All these quantities are nonnegative. For $\kappa$, use the elementary
sumset bound $|U+V|\ge|U|+|V|-1$. For $\vartheta$, zero is not a diagonal
exception, the eight members of $D\setminus E$ account for at most eight
exceptions, and any remaining exception is counted in $\delta_0(P)$ or
$\delta_0(Q)$. For $g$, count the $4(d-4)+6=4d-10$ distinct-element pairs
that meet the removed layer. The other signs follow from
\eqref{u86:eq:cover}--\eqref{u86:eq:peel}.
Direct cancellation gives the exact identity
\begin{equation}\label{u86:eq:budget}
 \boxed{\kappa+\rho+2\vartheta+c+\ell_1+\ell_2+g
       =22d+97-\Lambda=:\Omega.}
\end{equation}
For clarity, $c+\ell_1+\ell_2+g=22d+77+2\Delta(D)-s$ and
$\rho+2\vartheta=\mu+20-2\Delta(D)$, which proves the cancellation
without an inequality or a duplicated expenditure of slack.

\section{A sharp four-extreme defect inequality}\label{u86:sec:sharp}
\begin{theorem}\label{u86:thm:sharp}
Let $F=\{x_1<\cdots<x_m\}\subset\Z$, $m\ge7$, and let
\[
 a_1<a_2<x_1<\cdots<x_m<b_2<b_1,\qquad
 D=F\cup\{a_1,a_2,b_2,b_1\}.
\]
Let $\eta_F,\eta_D$ count the holes in the endpoint hulls of $R(F),R(D)$,
and put $g_H=4m+6-(|R(D)|-|R(F)|)$. Then
\begin{equation}\label{u86:eq:sharp}
 \boxed{\eta_D+\eta_F+g_H\ge2m-2.}
\end{equation}
The constant is sharp for every $m\ge7$.
\end{theorem}
\begin{proof}
Write $L=x_1+x_2$ and $U=x_{m-1}+x_m$. Let $W_+$ be
$((b_1+F)\cup(b_2+F))\cap(U,\infty)$, and let $W_-$ be the corresponding
union of the two lower translates below $L$. Put
\[
 I_+=(U,b_1+x_m]\cap\Z,\quad I_-=[a_1+x_1,L)\cap\Z,
 \quad H_\pm=|I_\pm|-|W_\pm|.
\]
Let $q$ count the distinct mixed sums $a_i+b_j$ that lie outside $[L,U]$
and are not already in $W_-\cup W_+$. These sums lie in $I_-\cup I_+$,
so $0\le q\le4$ and $q\le H_-+H_+$.

Translate $x_1$ to zero and write
\[
\begin{gathered}
 x_m=M,\quad x_2=s_0,\quad x_{m-1}=M-t_0,\\
 a_2=-w_0,\quad a_1=-w_0-v_-,\quad
 b_2=M+u,\quad b_1=M+u+v_+.
\end{gathered}
\]
All six gap parameters are positive integers. At most $u+1$ members of
$b_2+F$ lie above $U$: at most $u$ eligible points are at or below
$x_{m-1}$, followed by $x_m$. Thus the integers
\[
 S_+=u+1+m-|W_+|,\qquad S_-=w_0+1+m-|W_-|
\]
are nonnegative. Set $H=H_++H_-$ and $S=S_++S_-$. We claim
\begin{equation}\label{u86:eq:HS}
 H+S\ge2q-2.
\end{equation}
For $q\le2$ this follows from $H\ge q$. The same bound proves the case
$q=3,S\ge1$ and the case $q=4,S\ge2$.

Suppose first that $S=0$. Since $|W_\pm|\le2m$, we have
$u,w_0\le m-1$. The mixed value $B_0=a_2+b_2=M+u-w_0$ satisfies
\[
 M-m+2\le B_0\le M+m-2.
\]
Now $s_0\le M-m+2$ and $x_{m-1}\ge m-2$, so
$B_0\in[s_0,2M-t_0]=[L,U]$. In particular $q=4,S=0$ is impossible.

If $q=3,S=0$, the other mixed values
\[
 A_0=B_0-v_-,\qquad C_0=B_0+v_+-v_-,\qquad D_0=B_0+v_+
\]
are all outside $[L,U]$. Thus $A_0<L$ and $D_0>U$. Reflection allows us to
assume $C_0<L$. The definitions give
\[
 H_+=v_++t_0-m-1,\qquad H_-=v_-+s_0-m-1.
\]
Using integrality in $D_0>U$ and $C_0<L$ yields, respectively,
\[
 u-w_0\ge M-H_+-m,\qquad
 u-w_0\le t_0-M+H_--H_+-1.
\]
Consequently
\[
 H_-\ge2M-m-t_0+1=M+x_{m-1}-m+1\ge M-1\ge m-2\ge5.
\]
This proves \eqref{u86:eq:HS} in that case.

Finally, suppose $q=4,S=1$. Reflect if necessary so that
$S_+=1,S_-=0$. The cardinality bounds now give $u\le m$ and
$w_0\le m-1$, so $B_0\ge M-m+2\ge s_0$. Since $B_0$ lies outside the
old hull, $B_0>U$, whence
\[
 u-w_0>x_{m-1}\ge m-2,\qquad u-w_0\le m-1.
\]
Equality forces $u=m$, $w_0=1$, and $x_{m-1}=m-2$. The first $m-1$
members of $F$ are therefore $0,1,\ldots,m-2$, so $s_0=1$, and
$B_0=M+m-1=U+1$. Because $v_-\ge1$, the value $A_0=B_0-v_-$ is at most
$U$; it must consequently be less than 1. Thus $v_-\ge M+m-1$ and
\[
 H_-=v_-+s_0-m-1=v_--m\ge M-1\ge m-2\ge5.
\]
Together with $S=1$, this proves \eqref{u86:eq:HS} in the final case.

Put $h_-=L-(a_1+a_2)$, $h_+=(b_1+b_2)-U$, and $h=h_-+h_+$.
The definitions give the exact identity
\[
 2(|W_-|+|W_+|)=h+2m+2-H-S.
\]
Every new restricted sum outside $[L,U]$ lies in $W_-\cup W_+$, is one
of the $q$ new mixed values, or is one of $a_1+a_2,b_1+b_2$. The latter
values are distinct and outside $I_-\cup I_+$. Hence the number $N_{\rm out}$
of these sums satisfies
\[
 2N_{\rm out}=h+2m+6-H-S+2q\le h+2m+8.
\]
Write $N_0=|R(D)|-|R(F)|$. At most $\eta_F$ new sums are inside the old
hull, so $N_0\le N_{\rm out}+\eta_F$. Also
$\eta_D=h+\eta_F-N_0$. Therefore
\begin{align*}
 \eta_D+\eta_F+g_H
 &=h+2\eta_F+4m+6-2N_0\\
 &\ge h+4m+6-2N_{\rm out}\ge2m-2.
\end{align*}
For sharpness, take $F=[0,m-1]_{\Z}$ and
$\{a_1,a_2,b_2,b_1\}=\{-m-2,-1,m,2m+1\}$.
Here $R(F)=[1,2m-3]_{\Z}$ and $R(D)=[-m-3,3m+1]_{\Z}$.
For the latter assertion the four translates of $F$ fill the ranges
$[-m-2,-3]$, $[-1,m-2]$, $[m,2m-1]$, $[2m+1,3m]$;
$R(F)$ fills the central gap, and the sums of outer elements supply
$-m-3,-2,2m,3m+1$. Thus both hole counts vanish and $g_H=2m-2$.
\end{proof}

\section{Lattice reductions at sizes 95 and 94}\label{u86:sec:lattice}
This section extends the zero-containing common-core reduction from sizes
97 and 96 in Proposition~\ref{u96:prop:lattice96}. It does not exclude either new size.

\subsection{Explicit effective hole inputs}
For increasing prefixes $X_i$ of a positive set, write
$y_i=|\FS(X_i)|-|\FS(X_{i-1})|-i$; these increments are nonnegative
and sum to $\mu(X)$. For reproducibility, the exact finite parameter
test used below is as follows. For integers $C\ge0$ and $\beta\ge1$, define
\begin{align*}
 K_0&=\max\{k\ge3:k(k-3)/2\le C+1\},&D_0&=C+\max(3,K_0),\\
 E_0&=D_0+1,&
 M_0&=\max\left\{2\beta,1+\left\lfloor
 \frac{1+\sqrt{9+8(E_0+6\beta)}}2\right\rfloor\right\},\\
 J&=\max(K_0+1,2\beta+5,M_0).
\end{align*}
The parameterized hole test, Lemma~\ref{u96:lem:hole-parameter}, states
that if
\begin{equation}\label{u86:eq:holetest}
 (r-J+1)\beta>D_0,\qquad T_{J-1}-1>D_0,
\end{equation}
then a positive $r$-element set on a specified lattice with
$2\beta<h<|\FS(X)|-1$ satisfies $\mu(X)\ge r+C+1$.
No implicit ``sufficiently large'' threshold is used.

Here is the mechanism of that proved above test. If instead
$\mu\le r+C$, the last AP prefix and the nonnegative support increments
give a tail with $y_j\ge1$ and $\sum(y_j-1)\le D_0$, starting no later
than $K_0+1$. If there is no AP prefix of size at least three, it starts
at 4 with budget $C+3$. In the other case a last AP prefix of size $k$
followed by a non-AP support costs at least $r-1+k(k-3)/2$, so $k\le K_0$.
The first condition in \eqref{u86:eq:holetest} forces an index $i\ge J$ with
$1\le y_i\le\beta$. The small-increment inverse and Lev's
central-interval theorem \cite{Lev1998} identify a deleted integer interval
whose lower holes lie in $[1,2\beta]$ and number at most $\beta$.
The condition $h<|\FS|-1$ makes the final integer set primitive. A first
later element outside the prefix scale would make the two supports
disjoint and cost at least $T_{J-1}-1>D_0$; hence the scale is one.
Every later increment is at most $E_0$. The elementary endpoint bridge,
with threshold $M_0$, preserves these holes to the final support, leaving
at most $2\beta$ holes, a contradiction. The earlier chapters prove these inputs with the stated strict
thresholds.

Only $\beta=7$ and the following $C=C_7(r)$ rows are needed here:
\begin{center}\small
\begin{tabular}{rrrrrrrrr}
\toprule
$r$&31&32&33&34&35&36&37&38\\
$C_7(r)$&76&83&89&96&102&109&110&116\\\midrule
$r$&39&40&41&42&43&44&45&46\\
$C_7(r)$&122&128&129&135&142&147&149&155\\\bottomrule
\end{tabular}
\end{center}
Every row satisfies both strict inequalities in \eqref{u86:eq:holetest},
using exact integer square roots. A row remains valid at every larger
$r$ with its same $C,\beta,J$. The finite arithmetic checker must check
all sixteen rows, rather than using decimal square-root approximations.

\subsection{The lattice reduction}
We use the universal span estimate in Lemma~\ref{u96:lem:span96},
with $h(X)=\sum X+1-|\FS(X)|$. The notation $\rho(X)$ here denotes
the same defect $\rhoD(X)$ used in that lemma.

\begin{theorem}\label{u86:thm:lattice}
Let $A$ be a feasible integer set containing zero with $n\in\{95,94\}$.
With the peeling notation of Section~\ref{u86:sec:budget}, $|B|=8$ and there
is a common normalization of the core lattice in which $D$ is integral.
Let $h_P,h_Q,z$ be the hull-hole counts of $\FS(P),\FS(Q),\FS(E)$ on this
lattice, and let $w$ be the hull-hole count of $\FS(D)$. Then
\begin{equation}\label{u86:eq:newbudget}
 w\le4z,\qquad
 T\ge\kappa+G-3z\ge G-2z,\qquad
 T\ge\lceil G/3\rceil,\qquad T\ge z,
 \quad G=2d-10.
\end{equation}
The exact balanced parameters are
\begin{center}
\begin{tabular}{rrrrrrrrr}
\toprule
$n$&$d$&$e$&$m$&$(a_0,b_0)$&$\Lambda_0$&$\Omega_0$&$I$&$G$\\\midrule
95&87&79&78&$(39,39)$&1561&450&$k^2$&164\\
94&86&78&77&$(38,39)$&1522&467&$k(k+1)$&162\\\bottomrule
\end{tabular}
\end{center}
Here, after reflection, $(a,b)=(a_0-k,b_0+k)$,
$\Lambda=\Lambda_0+I$, and $\Omega=\Omega_0-I$.
\end{theorem}
\begin{proof}
\emph{Obtaining the core lattice.}
If $|B|\le7$, one sign is absent from $D$, and $E$ has at least 78
nonzero elements. The one-sided chain bound gives
$s\ge1+T_{78}=3082$, contradicting
$s\le24d+77\le24\cdot95+77=2357$. Thus $|B|=8$.

If $q$ counts the nonzero absolute values occurring only once in $E$,
the unpaired-value inequality is
$s\ge(e^2+3+q^2)/4$; the sharper version has an additional nonnegative
inversion count. Together with \eqref{u86:eq:corebound} and $d=e+8$, this gives
$q^2\le96e+1073-e^2$, namely 2416 at size 95 and 2477 at size 94.
Parity gives $q\le48$ and $q\le49$, respectively. Thus $E$ has at least
15 and 14 opposite pairs, respectively. In particular $a,b\ge14$, zero
is not a diagonal exception, and the nonnegative exact budget
\eqref{u86:eq:budget} applies.

Since $\Omega\ge0$, $k\le21$ in either case, improving the minimum side
size to 18 and 17, respectively. For every such integer $k$, the
Stanchescu sufficient margin $T_a-2-\Omega$ is
\[
 \frac{3k^2-79k+656}{2}\quad(n=95),\qquad
 \frac{3k^2-75k+544}{2}\quad(n=94).
\]
Their minima on $0\le k\le21$ are 68 and 38. Therefore the simultaneous
inverse theorem \cite[Corollary 2]{Stanchescu1996} applies to
$U=\FS(P)$ and $V=-\FS(Q)$: indeed
$|U+V|=|U|+|V|-1+\kappa$ and
$\kappa\le\Omega<T_a-2\le\min(|U|,|V|)-3$.
It places $U,V$ in progressions on one common lattice, with hull
cardinalities at most $|U|+\kappa,|V|+\kappa$. Since both contain zero,
normalize the common lattice to $\Z$. Then $E\subset\Z$ and
\begin{equation}\label{u86:eq:holes}
 h_P,h_Q\le\kappa,\quad z=h_P+h_Q-\kappa,\quad
 0\le z\le\min(h_P,h_Q)\le\kappa.
\end{equation}
No lattice membership of $D_1$ or $D$ has yet been assumed.

\emph{Forcing the second peeled layer onto this lattice.}
Put $R_p=\max P$, $R_q=\max Q$. Lemma~\ref{u96:lem:span96} and
\eqref{u86:eq:holes} give
\begin{equation}\label{u86:eq:Rsum}
 R_p+R_q\le m+\rho+\kappa+z+4.
\end{equation}
The integer set $R(E)$ has at most $2(R_p+R_q)+1$ values; adding four
new elements creates at most $4e+6$ further restricted sums, regardless
of their lattice membership. Thus, with $C=2m+4e+14-\Lambda$,
\begin{align}
 |\FS(D_1)|-2s
 &\le C+2\rho-\mu+\kappa+2z+\ell_2\notag\\
 &\le C+3\Omega-\rho-\mu,\label{u86:eq:latt3}\\
 |\FS(D_1)|-2s
 &\le C+\Omega+\rho-\mu+2z.\label{u86:eq:latt1}
\end{align}
The last two lines follow from $\rho+\kappa+\ell_2\le\Omega$ and
$z\le\kappa$. In particular neither uses $G$.

At size 95, $C=-1075-I$ and $C+3\Omega=275-4I$; at size 94 they are
$-1042-I$ and $359-4I$. If $z\le14$, use $\rho\le\mu$ in
\eqref{u86:eq:latt1}: its upper bound is $-597-2I$ or $-547-2I$, respectively,
and is strictly negative.

Suppose $z\ge15$. Both side supports have positive holes, and
$h\le\kappa<|\FS|-1$. They cannot be progressions with step one; a
progression with larger integer step would have at least $|\FS|-1$
holes. The non-progression inputs therefore give
\[
 \rho\ge m-28,\qquad\mu\ge m-6,\qquad\rho+\mu\ge2m-34.
\]
This already makes \eqref{u86:eq:latt3} strict except possibly $k=0,\ldots,6$
at size 95 and $k=0,\ldots,7$ at size 94.
At size 95 these remaining side sizes are at least 33. The effective-hole
row $(r,C,\beta)=(33,89,7)$ yields $\mu\ge78+180=258$ and
$\rho\ge50$. Hence $\rho+\mu\ge308>275-4I$.
At size 94 the exact comparisons are
\begin{center}\small
\begin{tabular}{rrrrrrrr}
\toprule
$k$&$a$&$b$&$C_7(a)$&$C_7(b)$&$\rho+\mu\ge$&$C+3\Omega$&Margin\\\midrule
0&38&39&116&122&366&359&7\\
1&37&40&110&128&366&351&15\\
2&36&41&109&129&366&335&31\\
3&35&42&102&135&365&311&54\\
4&34&43&96&142&366&279&87\\
5&33&44&89&147&364&239&125\\
6&32&45&83&149&360&191&169\\
7&31&46&76&155&359&135&224\\\bottomrule
\end{tabular}
\end{center}
Here the lower bound is $128+C_7(a)+C_7(b)$. Every margin is positive.
We have proved $|\FS(D_1)|<2s$ in all cases. An element
$x\in D_1\setminus E$ outside $\Z$ would supply the two disjoint
$s$-point sets $\FS(E)$ and $x+\FS(E)$ inside $\FS(D_1)$, a contradiction.
Thus $D_1\subset\Z$.

\emph{Forcing the first peeled layer.}
Writing $F_1=|\FS(D_1)|=s+|R(D_1)|-1+\ell_2$ gives
\[
 |\FS(D)|-F_1\le|R(D_1)|+4d-11+\Omega.
\]
The right side is less than $F_1$ because
\[
 1561>4\cdot87-10+450=788,\qquad
 1522>4\cdot86-10+467=801.
\]
Thus $|\FS(D)|<2|\FS(D_1)|$. The same disjoint-translate argument proves
$D\subset\Z$ in this normalization.

\emph{Propagating holes, only after integrality.}
Let $\alpha_2,\beta_2$ be the sums of the two negative and two positive
members of $D_1\setminus E$. The extreme translates
$\alpha_2+\FS(E)$ and $\beta_2+\FS(E)$ must overlap, by the strict
$2s$ bound. Their hulls cover the whole subset-sum hull of $D_1$, leaving
at most $2z$ holes. If $u_2=\beta_2-\alpha_2$, their union has at least
$s+u_2-z$ points. Since $R(D_1)$ has endpoints $\alpha_2,\beta_2$, its
hole count $\eta_2$ satisfies
\[
 |\FS(D_1)|=s+u_2-\eta_2+\ell_2,
 \qquad\eta_2\le z+\ell_2.
\]
Repeat at the first peeled layer, using at most $2z$ holes in
$\FS(D_1)$, to obtain $w\le4z$ and $\eta_1\le2z+\ell_1$ for $R(D)$.
Theorem~\ref{u86:thm:sharp}, now legitimately applied to $F=D_1$, gives
\[
 \eta_1+\eta_2+g\ge2(d-4)-2=2d-10=G.
\]
Thus $T\ge\kappa+G-3z\ge G-2z$. Also $T\ge\kappa\ge z$, so
$3T\ge G$ and all of \eqref{u86:eq:newbudget} follows.
\end{proof}

\begin{corollary}\label{u86:cor:budgets}
At size 96, the proved above common-lattice reduction similarly
allows $G=166$ in place of 164. At sizes 95 and 94 the new values are 164
and 162. The corresponding lower bounds $T\ge\lceil G/3\rceil$ are
56, 55 and 54. More generally, wherever the existing reduction has
established the same lattice and hole propagation, use $G=2d-10$.
\end{corollary}
If $j$ of the retained covering units are empty and $\delta=c-jd$ is
the residual covering slack, combining \eqref{u86:eq:Rsum},
\eqref{u86:eq:budget} and \eqref{u86:eq:newbudget} yields
\begin{equation}\label{u86:eq:joint}
 R_p+R_q+\delta\le m+4+\Omega-G-jd+4z.
\end{equation}
Indeed, $\rho+c\le\Omega-T$; substitute
$T\ge\kappa+G-3z$ into \eqref{u86:eq:Rsum}. Thus the new theorem lowers
this upper bound by exactly two. The same variable $z$ must be retained
in all simultaneous constraints.

\begin{remark}
Theorem~\ref{u86:thm:lattice} puts $D$ on the core lattice. It does not by
itself place the eight original outer elements $B$ on that lattice.
For real feasible sets, integralization transfers the existence problem
before the theorem is applied. To cover all real sets of an exact size by
a zero-containing endpoint argument, the larger sizes must first be
excluded. The cases 95 and 94 are structural reductions, not such
exclusions.
\end{remark}

\section{Further necessary inequalities}\label{u86:sec:cuts}
Throughout this section $D$ is already integral on the common lattice.
Write
\[
 E=P\cup\{0\}\cup(-Q),\quad |E|=e,\quad |D|=d=e+8,\quad K=|2D|,
\]
\[
 x_0=R_p<x_1<x_2<x_3<x_4=M_p,\qquad
 y_0=R_q<y_1<y_2<y_3<y_4=M_q,
\]
where the eight non-core elements are $x_1,\ldots,x_4$ and
$-y_1,\ldots,-y_4$. These statements allow holes in both core supports.

\begin{proposition}[Independent same-sign cutoffs]\label{u86:prop:tails}
For $0\le i,h\le4$,
\[
 K\le M_p+M_q+x_i+y_h+1+\binom{5-i}{2}+\binom{5-h}{2}.
\]
\end{proposition}
\begin{proof}
Every mixed-sign sum, and every same-sign sum with one summand at or
below the corresponding exposed endpoint, lies in
$[-M_q-y_h,M_p+x_i]_{\Z}$. Only pairs among the remaining positive
outliers, or among the remaining negative outliers, can be outside.
There are at most the displayed numbers of unordered pairs with
repetition. Counting these pairs and the integer interval proves the
claim, even if some pairs have equal sums.
\end{proof}

\begin{proposition}[Cardinality-aware envelopes]\label{u86:prop:225}
For $0\le j\le h\le4$ and $0\le k\le l\le4$, put
$t=8-j-k$, $r=d-t$, and $u=8-h-l$. Then
\begin{equation}\label{u86:eq:225}
 K\le x_j+x_h+y_k+y_l+1+ur+\frac{t(t+1)}2.
\end{equation}
These choices give 225 valid linear inequalities.
\end{proposition}
\begin{proof}
Let $C=E\cup\{x_1,\ldots,x_j,-y_1,\ldots,-y_k\}$. It has $r$ members
and lies in $[-y_k,x_j]$. Put $O=D\setminus C$, of size $t$.
The sumset $2C$ and the translates $C+x_i$ for $j<i\le h$ and
$C-y_i$ for $k<i\le l$ all lie in
$[-y_k-y_l,x_j+x_h]_{\Z}$. Each of the $u$ other translates of $C$ has
exactly $r$ points. Finally $2O$ has at most $t(t+1)/2$ points.
These sets cover $2D$, proving \eqref{u86:eq:225}. There are 15 choices of
$(j,h)$ and 15 of $(k,l)$.
\end{proof}

\subsection{An elementary odd-pivot defect bound}\label{u86:sec:odd-pivot}
\begin{lemma}\label{u86:lem:odd-pivot}
Let $X$ be a set of $r$ positive integers, let $R=\max X$, and assume
that $X$ contains an odd element. Then
\[
 \delta_0(X)\le R-r+2.
\]
If $R$ is odd, the stronger bound $\delta_0(X)\le R-r+1$ holds.
\end{lemma}
\begin{proof}
Let $t$ be the largest odd element, put $s=|X\cap(t,R]|$, and
$r_t=r-s$. The $s$ elements above $t$ are all even, so
$2s\le R-t+1$. For each defective $x<t$, define
\[
 f(x)=\begin{cases}2x,&x<t/2,\\2x-t,&x>t/2.\end{cases}
\]
Equality $x=t/2$ is impossible. Each image lies in $[1,t-1]$ and
is absent from $X$: otherwise $2x=0+f(x)$ or $2x=t+f(x)$ would be
a representation by two distinct members of $X\cup\{0\}$.
Both branches are injective; their images have opposite parity and
are therefore disjoint. There are only $t-r_t$ missing integers in
$[1,t]$, so at most that many defective elements lie below $t$.
There is at most one defect at $t$ and at most $s$ above it. Therefore
\[
 \delta_0(X)\le t-r_t+1+s=t-r+2s+1\le R-r+2.
\]
If $R$ is odd, $t=R$ and $s=0$, giving the stronger conclusion.
\end{proof}
In the strict common-core reductions, $h(X)<|\FS(X)|-1$ forces each
core to be primitive: a gcd of at least two would leave at least one
hole between every two successive support values. Thus each core has
an odd element, and the lemma gives the necessary linear bound
\begin{equation}\label{u86:eq:odd-rho}
 \rho(X)\ge\mu(X)+2r-2R-2.
\end{equation}
With a certified constant $\mu(X)\ge u$, replace $\mu(X)$ by $u$.
On each side independently, choose this bound or any bound proved above
constant lower bound for $\rho(X)$. The four choices in
$\rho(P)+\rho(Q)+c+T\le\Omega$ give four valid linear budget rows.
This lemma is proved without finite sampling or a solver. Its validity
by itself does not exclude the next cardinality.

\subsection{Leading twofold translate losses}
Let $S_p=\sum(D\cap\R_{>0})$, $S_q=-\sum(D\cap\R_{<0})$.
The leading positive covering slot is $c_p+2D$, where
$c_p=p_1-p_2-p_3-p_4$ for the four original positive outer elements
$p_1>p_2>p_3>p_4$. The reflected negative slot uses
$c_q=q_1-q_2-q_3-q_4$. Define the actual losses
\[
 L_p=K-|(c_p+2D)\cap\FS(D)\cap(0,\infty)|,
\]
and define $L_q$ by reflection. For a retained individual translate,
let $\lambda_v$ be the number of its points outside the appropriate
open half of the subset-sum hull. The covering keeps these
slots separate from the paired units. Thus, after $j$ empty units,
\begin{equation}\label{u86:eq:loss}
 L_p+L_q+\sum_v\lambda_v\le\delta=c-jd.
\end{equation}
This follows by subtracting the actual slot contributions from their
nominal capacities in the same union bound that proves
\eqref{u86:eq:cover}; discarded paired-unit deficits are nonnegative.

\begin{proposition}\label{u86:prop:leading}
If $\delta<K$, then $c_p,c_q$ are integers and
\begin{align*}
 c_p+2M_p+L_p&\ge K,&
 S_p+2M_q-c_p+L_p&\ge K-1,& S_p+L_p&\ge K,
\end{align*}
with the reflected inequalities on the negative side. If also
$\delta<2e+7$, then for $i=0,\ldots,4$,
\begin{align*}
 c_p+2x_i+L_p&\ge2e+2i+7,\\
 S_p+2y_i-c_p+L_p&\ge2e+2i+6,
\end{align*}
and likewise by reflection.
\end{proposition}
\begin{proof}
Since $L_p,L_q\le\delta<K$, each leading translate contributes at
least one point. All values in $2D$ and $\FS(D)$ are integers, so each
of $c_p,c_q$ is integral. The $K-L_p$ contributing positive values lie
in $[1,c_p+2M_p]$, in $[c_p-2M_q,S_p]$, and in $[1,S_p]$.
Counting integer points in these three intervals gives the first
three inequalities.

For the first nested inequality use
$F=E\cup\{-y_1,\ldots,-y_4,x_1,\ldots,x_i\}$. Then
$|F|=e+4+i$ and $|2F|\ge2|F|-1=2e+2i+7$ by the elementary ordered-chain
sumset bound. At most $L_p$ of these values fail to contribute through
the leading slot; all contributors lie in $[1,c_p+2x_i]$.
The hypothesis $L_p\le\delta<2e+7$ ensures that this interval has a
contributor, so its integer cardinality gives the assertion. For the
second use all four positive outliers and only the first $i$ negative
outliers. The contributing interval is $[c_p-2y_i,S_p]$, which has at
most $S_p+2y_i-c_p+1$ integer points. Reflection gives the remaining
claims.
\end{proof}
At size 96, $K\ge2d-1=175$ and $2e+7=167$. Thus bounds
$\delta<175$ and $\delta<167$ suffice for the respective hypotheses.
They must be established in each modeled regime; they are not automatic
from the cardinality alone.

\subsection{Overlap of retained translates}
The following generic form makes the overlap argument independent of
particular center labels or a numerical solver.

\begin{proposition}\label{u86:prop:overlap}
Let $v_1,\ldots,v_k$ be the integral centers of $k$ distinct retained
same-sign translate slots $v_i+D$, with minimum $v_{\min}$ and maximum
$v_{\max}$. Let $H$ be the target open half of the integer subset-sum
hull, and put $\lambda_i=|(v_i+D)\setminus H|$. Define
\[
 H_G=ke-(R_p+R_q+v_{\max}-v_{\min}+1).
\]
Their joint capacity deficit satisfies
\[
 kd-\left|H\cap\bigcup_{i=1}^k(v_i+D)\right|
 \ge H_G+\max_i\lambda_i.
\]
Disjoint groups of counted slots can be added in the total covering
budget. Their centers need not have distinct numerical values.
\end{proposition}
\begin{proof}
The union $U=\bigcup_i(v_i+D)$ is contained in the union of the core
hull $[v_{\min}-R_q,v_{\max}+R_p]_{\Z}$ and at most $k(d-e)$ non-core
points. Hence $kd-|U|\ge H_G$. Also $U\setminus H$ contains each
$(v_i+D)\setminus H$, so $|U\setminus H|\ge\max_i\lambda_i$.
Adding the two inequalities proves the claim. For disjoint groups of
slot labels, sum the corresponding union bounds along with all
ungrouped capacities and losses. No assumption that the core fills
its hull is made.
\end{proof}

Every proposition in this section is a necessary condition on an actual
configuration. The next section combines them with the original
rows and complete exact certificates to obtain an endpoint exclusion.

\section{An exact exclusion of size 96}\label{u86:sec:upper95}
\begin{theorem}\label{u86:thm:upper95}
No feasible real set has 96 elements. Consequently
\[
 \boxed{N\le95.}
\]
\end{theorem}
\begin{proof}
The preceding global bound $N\le96$ implies that a hypothetical size-96
example already contains zero: otherwise adjoining zero produces a
forbidden size-97 example. Integralize it, then apply Proposition~\ref{u96:prop:lattice96} and Corollary~\ref{u86:cor:budgets}. Thus
\[
 d=88,\quad e=80,\quad m=79,\quad a+b=79,\quad
 I=T_a+T_b-T_{39}-T_{40},\quad\Omega=432-I,\quad G=166.
\]
Every actual configuration falls in the exhaustive root families below.
Every root is contradicted by an exact positive-integer certificate,
either directly or at every leaf of a proved exhaustive binary tree.
The remaining subsections give the full finite reduction and verification
rule; the companion \texttt{exact95} directory supplies the complete
integer rows, root assignments, multipliers and trees. These finite data
are part of the proof.
\end{proof}

\subsection{The conservative low-hole root list}
Retain $19\le a\le60$, the side-size range given by the size-96
structural reduction. For $z=0$, even the older value $G=164$ gives
$c\le268-I$. For $1\le z\le14$, both supports are non-progressions and
$\rho\ge51$, so $c\le217-I+2z\le245-I$. If $j$ units are empty, put
$B_0=268-I-88j$. We enumerate
\[
 \begin{array}{lll}
 Z_j:&z=0,&0\le j\le3,\quad B_0\ge0,\\
 P_j:&1\le z\le14,&0\le j\le2,\quad B_0-23\ge0.
 \end{array}
\]
For every such $(a,j)$ use every pairing mode and every necessary
contributor choice from the eighteen-unit construction of
Appendix~\ref{u96:app:upper96}, with the center and pairing definitions in Appendix~\ref{u96:app:upper99}. The mode counts
for $j=0,1,2,3$ are exactly $4,70,578,2992$. Both pairing branches and
all choices of empty units are retained; split-pair contributor choices
are selected from actual contributors, as in that construction.
The complete counts are
\begin{center}
\begin{tabular}{lrrrrrrr}
\toprule
Regime&$Z_0$&$Z_1$&$Z_2$&$Z_3$&$P_0$&$P_1$&$P_2$\\\midrule
Roots&128&1820&11560&11968&128&1820&9248\\\bottomrule
\end{tabular}
\end{center}
Their sum is 36672. This is a conservative root list: using the stronger
$G$ later in a root's inequalities does not remove cases from the list.

The basic and extra model stages retain the valid older bounds
\[
 \delta\le B_0\ (Z),\qquad
 \delta\le B_0-51+2z\ (P),\qquad
 R_p+R_q+\delta\le83+B_0+4z.
\]
The overlap, group and leading stages reduce both relevant right sides
by two, using the proved $G=166$. The common cardinality row is
\[
 S_p+S_q+\delta-2K\le1668-88j+4z.
\]
Mass, inner-gap and quantile rows retain the same single $z$ coordinate.
For a retained positive center $v$ with lower and upper hull losses
$l_v,u_v$, the available quantile rows are
\begin{center}\small
\begin{tabular}{lll}
\toprule
Type&Lower row&Upper row\\\midrule
Ordinary&$v+l_v\ge b+5$&$v-u_v\le S_p-a-4$\\
Global&$v+l_v+M_p\ge88$&$v-u_v-M_q\le S_p-87$\\
Core&$v+l_v+R_p\ge84$&$v-u_v-R_q\le S_p-83$\\\bottomrule
\end{tabular}
\end{center}
Reflection supplies the negative rows. The global rows need only a
contributing center. Core rows require more than four points in the
hull contribution, ensured by a total loss cap below 84. Ordinary rows
hold at in-hull endpoints, and a cap below $\min(a,b)+5$ excludes every
out-of-hull exception. The exact generator selects only rows justified
by these bounds. With three empty units it keeps weak outer order;
with at most two it uses strict order only when the integrality argument
below applies.

\subsection{The high-hole root list}
For $z\ge15$, write $(a,b)=(39-k,40+k)$ up to reflection. For each side
size $r$ and integer $\beta$, choose a nonnegative $C_\beta(r)$ satisfying
\eqref{u86:eq:holetest}, if one is available, and set
\[
 u(r,z)=\max\left(\{r-3\}\cup
 \{r+C_\beta(r)+1:1\le\beta\le\lfloor(r-5)/2\rfloor,\ 2\beta<z\}\right).
\]
Only chosen constants satisfying the exact parameterized theorem are
used. The checker must verify all 646 nonnegative parameter rows that may
occur; it does not assume that a numerical optimization found the best
constant, or restrict to the 103-row table in Appendix~\ref{u96:app:upper96}.
For $2\le\beta\le10$, write
\[
 v_\beta=\max(0,u-4\beta+2)=(\beta+1)q_\beta+s_\beta,
 \qquad0\le s_\beta\le\beta,
\]
and define
\[
 L_\beta(r,u)=\min\bigl((\beta-1)(r-14),
       (\beta-1)q_\beta+\min(s_\beta,\beta-1)\bigr).
\]
Lemma~\ref{u96:lem:flex96} gives
\[
 \rho_*(r,u)=\max\{r-14,u+2-2r,L_2(r,u),\ldots,L_{10}(r,u)\}
 \le\rho(X).
\]
Consequently, with $u_a=u(a,z)$ and $u_b=u(b,z)$,
\begin{equation}\label{u86:eq:highcap}
 c\le C(a,z):=432-I-\rho_*(a,u_a)-\rho_*(b,u_b)
                  -\max(56,z,166-2z).
\end{equation}
The already known bound $\rho\ge51$ and $T\ge56$ excludes $k\ge19$.
The exact checker conservatively enumerates all 5662 pairs
\[
 0\le k\le18,\qquad15\le z\le432-k(k+1).
\]
After discarding $C(a,z)<0$, its largest $C(a,z)$ is 144. In particular
$j\le1$, and both leading-translate thresholds $175$ and $167$ are
satisfied. When $j=1$, $\delta\le144-88=56<84$, so the core quantiles
are valid.

For $j=0$, Lemmas~\ref{u96:lem:active-sharp} and~\ref{u96:lem:mass-bonus} give a further
explicit elimination. Put
\[
 F_A(a,b)=\tfrac12T_a+T_b+6a+8b+85,\qquad
 F_B(a,b)=\tfrac23T_a+T_b+6a+\tfrac{22b+236}{3}.
\]
The quantity $F=S_p+S_q+c-2(M_p+M_q+R_p+R_q)$ is at least one of
\[
 F_A(a,b)+\tfrac12(u_a+z)+u_b+z,\qquad
 F_B(a,b)+\tfrac23(u_a+z)+u_b+z,
\]
or one of their reflected versions. The basic envelope
$K\le M_p+M_q+R_p+R_q+21$ and the cardinality row give
$F\le1710+4z$. Thus a positive value of the minimum of those four
lower expressions minus $1710+4z$ excludes the state.

There remain exactly 795 unoriented $(a,j,z)$ states: 413 with $j=0$
and 382 with $j=1$. Group consecutive values of $z$ only when
$(a,j,u_a,u_b)$ is constant and the same affine part realizes
$\max(56,z,166-2z)$. The resulting 30 intervals, nine for $j=0$ and 21
for $j=1$, are listed in \texttt{high\_intervals.json}. The verifier must
expand every interval and checks equality with the complete set of
795 states. After both orientations and every pairing mode, they give
72 $H_0$ roots and 2940 $H_1$ roots, totaling 3012. Replacing integer
$z$ by a real variable in each interval only enlarges its model, so an
exact contradiction for that model excludes all covered integer states.

\subsection{Additional row validity and exhaustive splits}
The model stages use the inequalities of Section~\ref{u86:sec:cuts} and the
following elementary refinements, recorded to specify the complete new
mathematical content of the generators.

\paragraph{Tail sum cardinalities.}
Beyond cutoff $x_j$, the residual positive sums lie in
$[2x_{j+1},2M_p]$, so their number is at most
$2(M_p-x_{j+1})+1$, as well as the pair count in
Proposition~\ref{u86:prop:tails}. For four increasing outliers
$t_1<t_2<t_3<t_4$, their twofold sumset has cardinality at most each of
\[
 2(t_3-t_1)+5,\quad2(t_4-t_2)+5,\quad
 (t_2-t_1)+(t_4-t_3)+7.
\]
The first two bounds put the three-element twofold sumset in its integer
hull and count four pairs involving the remaining singleton. The third
counts at most three within-pair sums for each pair and bounds the
cross-pair sums by their hull, of length
$(t_2-t_1)+(t_4-t_3)+1$. Combine any positive and negative such bounds
with the central interval of Proposition~\ref{u86:prop:tails}.

\paragraph{Outer integrality with at most two empty units.}
Assume $c_p,c_q\in\Z$ and at most two units are empty. Every own center
$c_p+p_i$ and $c_q+q_i$ is an unpaired slot, so a nonintegral outer
element makes its own slot empty. There cannot be exactly one
nonintegral element on a side, because the integral identity
$c_p=p_1-p_2-p_3-p_4$ (or its reflection) would force it to be integral.
Thus any failure consists of exactly two elements on one side, and
exhausts all empty units. On that side, say positive, the always-unpaired
cross centers $c_q-p_3,c_q-p_4$ contribute, forcing $p_3,p_4\in\Z$.
In pairing branch A, $c_q-p_2$ is unpaired and contributes. In branch B,
$\alpha_p=p_2-p_3-p_4$ is unpaired and contributes. Either forces
$p_2\in\Z$, then $c_p$ forces $p_1\in\Z$, a contradiction.
The case of negative failures is its reflection.
Here branch A means $\beta_p+\beta_q\le0$, branch B means $>0$, where
$\beta_p=2p_2-p_1-p_3-p_4$ and $\beta_q$ is its reflection.
Leading-center integrality follows from the old contributing-center
argument with at most one empty unit; in the other applicable models it
follows from $\delta<K$ and Proposition~\ref{u86:prop:leading}. Therefore
strict outer integer gaps and the branch-B bound $\beta_p+\beta_q\ge1$
are imposed only in justified regimes.

\paragraph{Groups containing a leading slot.}
Group the leading slot $c_p+2D$ with $k$ retained own-center slots
$c_p+p_i+D$, and let $p_{\max}$ be their largest $p_i$.
The core pieces $c_p+2E$ and $c_p+p_i+E$ lie in the common integer hull
\[
 [c_p-2R_q,c_p+p_{\max}+R_p]_{\Z}.
\]
Since $|2E|\ge2e-1$, the remaining group points number at most
$(K-|2E|)+k(d-e)$. Subtracting this upper bound for the raw union from
the nominal capacity $K+kd$ proves a global deficit of at least
\[
 H_*=(k+2)e-2-p_{\max}-R_p-2R_q.
\]
The deficit in the actual positive subset-sum target is at least
$H_*+\max(L_p,\lambda_1,\ldots,\lambda_k)$: the part of the raw union
outside that target contains the loss set of every member, and hull
losses are a subset of actual-target losses. Disjoint groups and
ungrouped losses may be added in \eqref{u86:eq:loss}. This proves the
leading-group rows, including their constant $e(k+2)-2$.
For the grouped-translate recurrence, choose each auxiliary prefix
variable to be the actual deficit of that prefix union.
The deficit of a union of two disjoint groups of slot labels is at
least the sum of their separate deficits. Proposition~\ref{u86:prop:overlap}
therefore proves every recurrence row, and the final prefix deficits
fit in the total loss budget. The use of auxiliary variables is thus
existentially valid for every actual configuration.

\paragraph{Normal or outside endpoint.}
Only the following proved above disjunctions are used in the
trees. For a retained integral positive center $v$, either
\[
 v+l_v\ge b+5,
 \quad\hbox{or}\quad(l_v\ge b+5\ \hbox{and}\ v\le-1);
\]
and either
\[
 v-u_v\le S_p-a-4,
 \quad\hbox{or}\quad(u_v\ge a+5\ \hbox{and}\ v\ge S_p+1).
\]
For the first, if $v\ge0$, at most $v$ of the $b+5$ nonpositive
members of $D$ can translate to a positive value, proving the ordinary
row; if $v<0$, all $b+5$ are lost. For the second, if $v\le S_p$, at
most $S_p-v$ of the $a+4$ positive members fit below the upper endpoint;
if $v>S_p$, all these members and zero are lost. Integer centers give
the displayed strict outside bounds. Reflect for negative centers.
Thus each binary split covers every actual set, without requiring LP
loss variables themselves to be integral.


\subsection{Exact finite lemma for size 96}
\label{u86:finite:96}
The required finite assertion has 39684 canonical roots and 39809
leaf contradictions. The data are in \path{proofs/exact95}.
There are 36672 low and 3012 high roots. Of these, 39624 have direct
contradictions. The remaining 60 have complete binary trees with 125
internal splits and 185 leaves. The ordinary pool has 863 vectors and
the branch pool 81. Every root, current matrix, normal/outside split,
both children, and effective-hole witness is reconstructed.

Every leaf is checked by the integer criterion of
Lemma~\ref{completion:coverage}; every path row must be one of the
proved alternatives. The verifier \path{verify_exact95.py} regenerates
the complete ordered root universe and all required children, rather than
accepting stored counts or an optimizer status. These assertions are finite
mathematical proof obligations in the supplied data.



\section{A further exact exclusion of size 95}\label{u86:sec:upper94}
\begin{theorem}\label{u86:thm:upper94}
No feasible real set has 95 elements. Consequently
\[
 \boxed{N\le94.}
\]
\end{theorem}
\begin{proof}
Apply Theorem~\ref{u86:thm:upper95} first. A hypothetical size-95 example
contains zero by the adjoining-zero argument; integralization and
Theorem~\ref{u86:thm:lattice} give the common lattice. Its parameters are
\[
 d=87,\quad e=79,\quad m=78,\quad a+b=78,\quad
 I=(a-39)^2,\quad\Omega=450-I,\quad G=164.
\]
Use the same actual-set rows and finite verification rule as in
Section~\ref{u86:sec:upper95}, with the exact parameter changes detailed
below. The complete root list has 74106 members. Of these, 73951 have
direct exact integer contradictions and 155 have complete binary trees
with 488 internal splits and 643 contradictory leaves. Thus all 74594
leaf systems are exactly infeasible, excluding size 95.
The general upper bound follows.
\end{proof}

\subsection{All roots and their parameter changes}
For low holes, retain every $18\le a\le60$, all $j=0,1,2,3$, and every
pairing mode and contributor choice, subject to
\[
 B_0=286-I-87j\ge0\quad(z=0),\qquad
 B_0-22\ge0\quad(1\le z\le14).
\]
These follow from $c\le286-I$ for zero holes and
$c\le236-I+2z\le264-I$ for positive low holes, since $\rho\ge m-28=50$.
The exact root counts are
\begin{center}\small
\begin{tabular}{lrrrrrrrr}
\toprule
Regime&$Z_0$&$Z_1$&$Z_2$&$Z_3$&$P_0$&$P_1$&$P_2$&$P_3$\\\midrule
Roots&132&2030&12138&32912&132&1890&10982&8976\\\bottomrule
\end{tabular}
\end{center}
They sum to 69192. The loss and span rows are
\[
 \delta\le B_0\ (Z),\qquad\delta\le B_0-50+2z\ (P),\qquad
 R_p+R_q+\delta\le82+B_0+4z.
\]
The cardinality right side becomes $1650-87j+4z$.
For positive three-empty modes, the model explicitly retains the same
$1\le z\le14$ variable in its loss, span, mass and cardinality rows;
it does not import a zero-hole equality from the preceding template.

For high holes use the same $u(r,z)$ and $\rho_*(r,u)$ defined above,
now with the exact cap
\[
 C_{95}(a,z)=450-I-\rho_*(a,u_a)-\rho_*(b,u_b)
                      -\max(55,z,164-2z).
\]
The non-progression and $T$ lower bounds exclude $k\ge19$ for
$(a,b)=(39-k,39+k)$. The complete conservative enumeration
\[
 0\le k\le18,\qquad15\le z\le450-k^2
\]
contains 6175 states. The 620 effective-hole rows used in it are
rechecked against the two strict integer tests. Its maximum nonnegative
covering cap is 171, so at most one unit is empty. For $j=0$, apply the
same four all-active mass lower bounds, now against $1692+4z$.
Exactly 1552 unoriented $(a,j,z)$ states survive: 804 with $j=0$ and
748 with $j=1$. Consecutive states with the same side excess bounds
and affine $T$ piece form 48 intervals. They expand exactly to those
1552 states. Both orientations are kept when $a\ne b$; the balanced
case is kept once. Every pairing mode is included, giving 84 $H_0$
roots and 4830 $H_1$ roots, totaling 4914.

The common model construction in this endpoint uses the following
parameter substitutions in the already proved row families:
\begin{center}\small
\begin{tabular}{ll}
\toprule
Row or hypothesis&Size-95 value\\\midrule
Basic twofold lower bound&$K\ge2d-1=173$\\
Leading-translate threshold&$\delta<173$\\
Nested leading-translate threshold&$\delta<2e+7=165$\\
Core quantile threshold&Hull loss $<d-4=83$\\
Global quantile right sides&87 and 86\\
Core quantile right sides&83 and 82\\
Two-slot core overlap constant&$2e-1=157$\\
$k$-slot core overlap constant&$ke-1=79k-1$\\
Leading-core group constant&$(k+2)e-2=79(k+2)-2$\\\bottomrule
\end{tabular}
\end{center}
The 225 cardinality-aware envelopes use their actual $d=87$ value.
Each high interval has constant $u_a,u_b$ and an affine $T$ piece, so
its maximum residual cap is at an endpoint; the checker also tests it
at every integer $z$. Since 171 is below 173, all high leading slots
contribute. The stronger nested rows are included only for intervals
whose residual cap is below 165. At residual cap at least 83 use global
quantiles; below 83 use core quantiles; ordinary rows are used only
when their stronger loss-cap test applies. These separate threshold
checks prevent carrying an invalid size-96 row into size 95.

\subsection{One additional exhaustive disjunction}
In addition to the normal or outside splits already proved, the
size-95 trees permit this split for a retained center:
\[
 \text{both core quantile rows hold}\qquad\hbox{or}\qquad
 l_v+u_v\ge83.
\]
Indeed, if $l_v+u_v<83$, integrality of the actual loss count and $d=87$
give at least five points of $v+D$ in the target hull. Since only four
members on either side lie outside the core, the core-quantile
argument gives both rows. If the strict inequality fails, the second
child holds. The two children need not be disjoint; they cover every
actual configuration. This is a proved disjunction, not an inference
from an LP optimum.


\subsection{Exact finite lemma for size 95}
\label{u86:finite:95}
The required finite assertion has 74106 canonical roots and 74594
leaf contradictions. The data are in \path{proofs/exact94}.
There are 69192 low and 4914 high roots. The 155 tree roots consist of
147 low and eight high roots; the trees have 488 internal splits and
643 leaves. All other roots are direct contradictions. All 48 high-hole
intervals are expanded. The ordinary and branch pools have 1008 and 257
vectors, respectively.

Every leaf is checked by the integer criterion of
Lemma~\ref{completion:coverage}; every path row must be one of the
proved alternatives. The verifier \path{verify_exact94.py} regenerates
the complete ordered root universe and all required children, rather than
accepting stored counts or an optimizer status. These assertions are finite
mathematical proof obligations in the supplied data.



\section{The odd-pivot bound and an exact exclusion of size 94}\label{u86:sec:upper93}
\begin{theorem}\label{u86:thm:upper93}
No feasible real set has 94 elements. Consequently
\[
 \boxed{N\le93.}
\]
\end{theorem}
\begin{proof}
The preceding exclusion of 95 justifies zero at the new endpoint.
Integralization and Theorem~\ref{u86:thm:lattice} give
\[
 d=86,\quad e=78,\quad m=77,\quad a+b=77,\quad
 I=T_a+T_b-T_{38}-T_{39},\quad\Omega=467-I,\quad G=162.
\]
The same root-model construction is valid with these parameters,
retaining each threshold at its new value. Add the elementary
odd-pivot bound of Lemma~\ref{u86:lem:odd-pivot} in the one high-hole tree
where it is needed, and use the additional inner-extreme disjunction
proved below. The exhaustive exact verification supplies 118081 leaf
contradictions covering all 116832 root systems. Thus size 94 is
impossible; the general upper bound follows.
\end{proof}

\subsection{Complete low and high classifications}
The low-hole list contains every $17\le a\le60$, every $j=0,1,2,3$,
and every pairing mode and contributor choice, subject to
\[
 B_0=305-I-86j\ge0\quad(z=0),\qquad
 B_0-21\ge0\quad(1\le z\le14).
\]
Here $\rho\ge49$ when $z>0$, and the loss, span and cardinality rows are
\begin{gather*}
 \delta\le B_0\ (Z),\qquad\delta\le B_0-49+2z\ (P),\\
 R_p+R_q+\delta\le81+B_0+4z,\qquad
 S_p+S_q+\delta-2K\le1632-86j+4z.
\end{gather*}
The full low-hole counts are
\begin{center}\small
\begin{tabular}{lrrrrrrrr}
\toprule
Regime&$Z_0$&$Z_1$&$Z_2$&$Z_3$&$P_0$&$P_1$&$P_2$&$P_3$\\\midrule
Roots&136&2100&13872&41888&136&1960&12716&29920\\\bottomrule
\end{tabular}
\end{center}
Their total is 102728.

For high holes write $(a,b)=(38-k,39+k)$ up to reflection. Use the same
certified side bounds $u,\rho_*$ with
\[
 C_{94}(a,z)=467-I-\rho_*(a,u_a)-\rho_*(b,u_b)
                        -\max(54,z,162-2z).
\]
The inequalities $\rho\ge49$ and $T\ge54$ exclude $k\ge19$. Enumerate
all 6327 pairs $0\le k\le18$, $15\le z\le467-k(k+1)$. The checker
revalidates all 627 nonnegative effective-hole constants used.
The maximum nonnegative cap is 196, so $j\le2$. The all-active mass
comparison for $j=0$ now uses the upper bound $1674+4z$.
There remain 2242 unoriented states: 1157 with $j=0$, 980 with $j=1$,
and 105 with $j=2$. Their canonical grouping gives 54 intervals,
respectively 12, 34 and eight of these three types. Both orientations
and every mode give 96 $H_0$ roots, 4760 $H_1$ roots and 9248 $H_2$
roots, totaling 14104.

The thresholds are now $K\ge171$, nested threshold $2e+7=163$, and
core-quantile loss threshold $d-4=82$. The high cap 196 is \emph{not}
below the leading threshold. Therefore leading and nested capacity
rows are inserted only when the actual interval's residual cap is
below 171 or 163, respectively. High intervals use global quantiles
at cap at least 82, core quantiles below 82, and ordinary quantiles
only when their stronger test applies. The overlap constants become
$2e-1=155$, $ke-1=78k-1$, and $(k+2)e-2=78(k+2)-2$.
All 225 envelope rows use $d=86$.

\subsection{The parity row and inner-extreme splits}
The primitive-core hypothesis of Lemma~\ref{u86:lem:odd-pivot} holds here:
Stanchescu gave $h_P<|\FS(P)|-1$ and $h_Q<|\FS(Q)|-1$.
Thus the high-hole side bounds can also use
\[
 \rho(P)\ge u_a+2a-2R_p-2,\qquad
 \rho(Q)\ge u_b+2b-2R_q-2.
\]
Take either this expression or $\rho_*$ independently on the two sides
and substitute in $\rho+c+T\le\Omega$. Besides the original constant
row this supplies three additional valid affine rows. The finite proof
uses these rows in its parity-enhanced high-hole tree. Their proof is
the parity injection, not a numerical classification of dense cores.

The trees may also use, for $i=1,2,3,4$ and a retained positive center
$v$, the disjunctions
\[
 l_v\ge i\quad\hbox{or}\quad v\ge y_{5-i}+1,
 \qquad
 u_v\ge i\quad\hbox{or}\quad v\le S_p-x_{5-i}.
\]
If $l_v<i$, at least one of the $i$ largest negative magnitudes gives
a translated value above zero, forcing $v>y_{5-i}$, hence the first
integer inequality. If $u_v<i$, at least one of the $i$ largest positive
elements fits at or below $S_p$, giving the second. Reflection supplies
the negative versions. These cover actual integer losses; they do not
assert integrality of a solution to the relaxed LP. The earlier normal
or outside splits remain valid, as does the core or large-loss split
with 82 replacing 83.


\subsection{Exact finite lemma for size 94}
\label{u86:finite:94}
The required finite assertion has 116832 canonical roots and 118081
leaf contradictions. The data are in \path{proofs/exact93}.
There are 102728 low and 14104 high roots. Of these, 116566 are direct
contradictions, while 228 low and 38 high roots have trees with 1249
internal splits and 1515 leaves. The ordinary and branch pools contain
1363 and 671 vectors. Every selected parity and inner-extreme row must
be reconstructed on its stated actual-set domain.

Every leaf is checked by the integer criterion of
Lemma~\ref{completion:coverage}; every path row must be one of the
proved alternatives. The verifier \path{verify_exact93.py} regenerates
the complete ordered root universe and all required children, rather than
accepting stored counts or an optimizer status. These assertions are finite
mathematical proof obligations in the supplied data.




\section{Further common-core reductions at sizes 93 and 92}\label{u86:sec:lattice92}
\begin{theorem}\label{u86:thm:lattice92}
The zero-containing common-core conclusion of Theorem~\ref{u86:thm:lattice}
also holds at sizes 93 and 92, with the following parameters:
\begin{center}\small
\begin{tabular}{rrrrrrrrr}
\toprule
$n$&$d$&$e$&$m$&$(a_0,b_0)$&$\Lambda_0$&$\Omega_0$&$I$&$G$\\\midrule
93&85&77&76&$(38,38)$&1483&484&$k^2$&160\\
92&84&76&75&$(37,38)$&1445&500&$k(k+1)$&158\\\bottomrule
\end{tabular}
\end{center}
Again $(a,b)=(a_0-k,b_0+k)$ up to reflection,
$\Lambda=\Lambda_0+I$, and $\Omega=\Omega_0-I$.
All the inequalities in \eqref{u86:eq:newbudget} hold. This theorem alone
excludes neither cardinality and does not put the original outer set
$B$ on the core lattice.
\end{theorem}
\begin{proof}
\emph{Bootstrap.}
If $|B|<8$, the one-sided core has at least 76 nonzero members, giving
$s\ge1+T_{76}=2927>24\cdot93+77=2309\ge24d+77$, a contradiction.
Hence $|B|=8$. Before assuming any minimum side size, use the original
uniform defect bound, Lemma~\ref{base:lem:defect-excess}, which gives
$\Lambda\le22d+115$, or $I\le\Omega_0+18$.
At size 93 this implies $k\le22$ and $a\ge16$; at size 92 it gives
$k\le22$ and $a\ge15$. Thus threshold 14 is available on both sides.
The unpaired bound gives $q^2\le2536,2593$, respectively. Parity gives
$q\le50,49$ and at least 13 opposite pairs, so zero is not a diagonal
exception. The nonnegative exact budget \eqref{u86:eq:budget} now applies.
It refines the size-92 range to $k\le21$, $a\ge16$.
The bootstrap is necessary: 13 opposite pairs alone would not justify
threshold 14.

\emph{The size-93 core lattice.}
The sufficient Stanchescu margin is
\[
 T_a-2-\Omega=\frac{3k^2-77k+510}{2}\ge8\qquad(0\le k\le22).
\]
Thus the same inverse theorem gives \eqref{u86:eq:holes} and the strict
bounds $h_P<|\FS(P)|-1$, $h_Q<|\FS(Q)|-1$.

\emph{The size-92 core lattice when one support is not an AP.}
Here the crude margin
$f(k)=(3k^2-73k+402)/2$ can be $-21$. Retain the actual excess and defect
terms instead of incorrectly using that margin alone. If
$U=\FS(P)$ is not an AP, then $\mu(P)\ge a-3$ and $\rho(P)\ge a-14$;
using $\kappa\le\Omega-\rho$ gives
\[
 |U|-3-\kappa\ge f(k)+2a-17
 =\frac{3k^2-77k+516}{2}\ge11.
\]
If $U$ is an AP but $V=-\FS(Q)$ is not, use $\rho(Q)\ge b-14$ to get
\[
 |U|-3-\kappa\ge f(k)+b-14
 =\frac{3k^2-71k+450}{2}\ge15.
\]
The larger-side margin is positive since
$T_b-2\ge T_{38}-2=739>500\ge\Omega$. Thus Stanchescu gives
\eqref{u86:eq:holes} and both strict support-cardinality hole bounds.

\emph{The size-92 core lattice when both supports are APs.}
We use a short elementary fact. If two integer APs $U,V$, containing
zero, have sizes $u,v\ge2$ and
$|U+V|=u+v-1+\kappa<uv$, their hull-hole counts on their common gcd
lattice are at most $\kappa$. To prove this, reflect or translate to
write $U=p[0,u-1]_{\Z}$ and $V=q[0,v-1]_{\Z}$ with $\gcd(p,q)=1$.
A collision of two distinct pairs exists, so coprimality forces
$u-1\ge q$ and $v-1\ge p$. Counting the $V$ index modulo $p$, the
intervals in each residue class overlap or touch and give
\[
 |U+V|=pu+qv-pq.
\]
In detail, residue $r\in[0,p-1]$ contributes
$u+q\lfloor(v-1-r)/p\rfloor$ points; their sum is the displayed value.
Thus
\begin{align*}
 h_U&=(p-1)(u-1),&h_V&=(q-1)(v-1),\\
 \kappa-h_U&=(q-1)(v-p)\ge0,&
 \kappa-h_V&=(p-1)(u-q)\ge0,\\
 z&=h_U+h_V-\kappa=(p-1)(q-1).
\end{align*}
In our application,
$\kappa\le500<T_{16}T_{38}\le(u-1)(v-1)$, so the fact applies.
Also $(q-1)(v-1)=h_V\le500$ while $v-1\ge T_{38}=741$, forcing $q=1$
and $z=0$. All lattice and hole identities hold, although this special
case need not give $h_P<|U|-1$. It never enters a positive-hole case.
At this point only $E$, not yet $D_1$ or $D$, is asserted to be integral
on the normalized lattice.

\emph{Forcing $D_1$ onto the core lattice.}
Lemma~\ref{u96:lem:span96} applies because $a,b\ge16$. The already derived
comparisons \eqref{u86:eq:latt3}--\eqref{u86:eq:latt1} remain valid, with
\begin{center}\small
\begin{tabular}{rrrr}
\toprule
$n$&$C_0$&$C_0+\Omega_0$&$C_0+3\Omega_0$\\\midrule
93&$-1009$&$-525$&443\\
92&$-977$&$-477$&523\\\bottomrule
\end{tabular}
\end{center}
The bound $\rho\le\mu$ makes \eqref{u86:eq:latt1} strictly negative for
$z\le262$ at size 93 and $z\le238$ at size 92. This includes the
exceptional both-AP case, where $z=0$.

For larger holes, both side supports satisfy the strict hypotheses of
the effective-hole theorem and are non-progressions. Thus
$\rho+\mu\ge2m-34$, reducing possible failures of \eqref{u86:eq:latt3} to
$k=0,\ldots,9$ at either size. The following exact parameter rows in
\eqref{u86:eq:holetest}, followed by the flexible inverse $L_\gamma$ from
Section~\ref{u86:sec:upper95}, give
$\mu(X)\ge u_r=r+C_h+1$ and $\rho(X)\ge v_r$:
\begin{center}\small
\begin{tabular}{rrrrrr@{\qquad}rrrrrr}
\toprule
$r$&$C_h$&$\beta$&$\gamma$&$u_r$&$v_r$&
$r$&$C_h$&$\beta$&$\gamma$&$u_r$&$v_r$\\\midrule
28&59&6&5&88&48&38&126&9&7&165&105\\
29&64&6&5&94&52&39&134&9&7&174&112\\
30&70&7&6&101&57&40&143&9&7&184&120\\
31&76&7&6&108&62&41&151&9&9&193&128\\
32&83&7&6&116&68&42&160&10&9&203&136\\
33&89&7&6&123&73&43&169&10&9&213&144\\
34&96&8&6&131&79&44&179&10&10&224&153\\
35&103&8&6&139&85&45&188&10&10&234&162\\
36&111&8&7&148&92&46&198&11&10&245&171\\
37&118&8&7&156&98&47&208&11&10&256&180\\\bottomrule
\end{tabular}
\end{center}
The chosen $\beta$ are at most 11, while the holes are at least 239, so
$h>2\beta$ holds throughout. Each row satisfies the two strict integer
parameter tests. Substituting both sides gives these complete margins:
\begin{center}\small
\begin{tabular}{r|rrr|rrr}
\toprule
&\multicolumn{3}{c|}{$n=93$}&\multicolumn{3}{c}{$n=92$}\\
$k$&$C+3\Omega$&$u_a+v_a+u_b+v_b$&Margin&
$C+3\Omega$&$u_a+v_a+u_b+v_b$&Margin\\\midrule
0&443&540&97&523&524&1\\
1&439&540&101&515&526&11\\
2&427&544&117&499&528&29\\
3&407&545&138&475&531&56\\
4&379&549&170&443&535&92\\
5&343&553&210&403&541&138\\
6&299&561&262&355&547&192\\
7&247&566&319&299&554&255\\
8&187&574&387&235&562&327\\
9&119&582&463&163&572&409\\\bottomrule
\end{tabular}
\end{center}
Every margin is strict, including $524>523$ in the balanced size-92
case. Hence $|\FS(D_1)|<2s$, and the disjoint-translate argument forces
$D_1$ onto the same lattice.

\emph{Forcing $D$ and propagating holes.}
The first-layer comparisons are
\[
 1483>4\cdot85-10+484=814,\qquad
 1445>4\cdot84-10+500=826.
\]
Exactly as in Theorem~\ref{u86:thm:lattice}, they give
$|\FS(D)|<2|\FS(D_1)|$, putting $D$ on the lattice. The two strict
comparisons now give $w\le4z$, $\eta_2\le z+\ell_2$ and
$\eta_1\le2z+\ell_1$. Only at this stage apply
Theorem~\ref{u86:thm:sharp} to obtain $G=2d-10$, and hence all of
\eqref{u86:eq:newbudget}. No post-lattice estimate was used to force
integrality.

\emph{Primitivity in the size-92 both-AP branch.}
Here $z=0$ and, only after $D$ is integral, the post-lattice budget gives
$\kappa\le\Omega-158$. The already verified lower bound
$T_a-2-\Omega\ge-21$ implies $T_a-\kappa\ge139>0$.
Since $h_P,h_Q\le\kappa$ and $T_b\ge T_a$, both core hole counts are
strictly below their support cardinalities minus one. A coefficient gcd
at least two would leave at least one missing position between consecutive
support values, contradicting these strict bounds. Both cores are therefore
primitive on this lattice. This does not place the original outer eight
points on that lattice.
\end{proof}
All displayed finite margins, parameter tests and flexible inverses are
checked by \path{check_lattice93_92.py} in the companion
material. The elementary two-AP statement is proved above; tests on
examples are not substituted for its proof.

\section{Finite progression bounds and an exact exclusion of size 93}\label{u86:sec:upper92}
\begin{theorem}\label{u86:thm:upper92}
No feasible real set has 93 elements. Consequently
\[
 \boxed{N\le92.}
\]
\end{theorem}
\begin{proof}
Theorem~\ref{u86:thm:upper93} first justifies zero in a hypothetical size-93
example. Integralization and Theorem~\ref{u86:thm:lattice92} then give the
common lattice, with
\[
 d=85,\quad e=77,\quad m=76,\quad a+b=76,\quad
 I=(a-38)^2,\quad\Omega=484-I,\quad G=160.
\]
The exhaustive root reduction below has 157766 systems. Its complete
exact certificates give 160550 leaf contradictions. The new ingredients
needed by five strengthened high-hole roots are proved in this section:
finite three-term-progression bounds, affine partition bounds, exact-mass
defect inequalities and integer-width alternatives. Every other row is
covered by the preceding proofs. The actual-set-to-model reduction and
complete exact contradiction therefore exclude size 93.
The preceding exclusions give the stated general upper bound.
\end{proof}

\subsection{Complete endpoint classification}
For low holes, take every $16\le a\le60$, every $j=0,1,2,3$, and all
pairing modes and contributor choices, retaining precisely
\[
 B_0=324-I-85j\ge0\quad(z=0),\qquad
 B_0-20\ge0\quad(1\le z\le14).
\]
These follow from $G=160$ and, for positive holes, $\rho\ge m-28=48$.
The simultaneous loss, span and cardinality bounds are
\begin{gather*}
 \delta\le B_0\ (Z),\qquad\delta\le B_0-48+2z\ (P),\\
 R_p+R_q+\delta\le80+B_0+4z,\qquad
 S_p+S_q+\delta-2K\le1614-85j+4z.
\end{gather*}
The exact low-root counts are
\begin{center}\small
\begin{tabular}{lrrrrrrrr}
\toprule
Regime&$Z_0$&$Z_1$&$Z_2$&$Z_3$&$P_0$&$P_1$&$P_2$&$P_3$\\\midrule
Roots&148&2170&14450&50864&140&2030&13294&44880\\\bottomrule
\end{tabular}
\end{center}
They total 127976. Three-empty models retain weak outer order.

For high holes, write $(a,b)=(38-k,38+k)$ up to reflection and use
\[
 C_{93}(a,z)=484-I-\rho_*(a,u_a)-\rho_*(b,u_b)
                        -\max(54,z,160-2z).
\]
Since $\rho\ge48$ and $T\ge54$, $c\le382-k^2$, excluding $k\ge20$.
The complete conservative range is
$0\le k\le19$, $15\le z\le484-k^2$, containing 6930 states.
All 634 effective-hole constants are directly rechecked.
The maximum nonnegative cap is 220, so $j\le2$.
The all-active mass upper bound is now $1656+4z$.
Exactly 3372 unoriented states remain: 1646 with $j=0$, 1367 with $j=1$,
and 359 with $j=2$. Their canonical grouping produces 88 intervals,
respectively 14, 53 and 21. Use both orientations except in the balanced
case $a=b=38$, which is used once. All modes then give 108 $H_0$,
7140 $H_1$ and 22542 $H_2$ roots, totaling 29790.

At this endpoint $K\ge169$, nested leading rows require
$\delta<161$, and core quantiles require hull loss below 81.
Per-interval residual caps control these choices, exactly as in the
previous endpoints. In particular the high cap 220 does not imply
that every leading translate contributes. All added rows respect
their separate thresholds and the common $z$ variable.

\subsection{Finite three-term-progression bounds}\label{u86:sec:finiteAP}
Let $r_3(L)$ be the largest size of a subset of $[1,L]_{\Z}$ with no
nonconstant three-term arithmetic progression; set $r_3(0)=0$.
The finite inputs used by the certificates are
\begin{center}\small
\begin{tabular}{r|rrrrrrrrrrrrrrr}
\toprule
$L$&0&1&2&3&4&5&6&7&8&9&10&11&12&13&14\\\midrule
$r_3(L)$&0&1&2&2&3&4&4&4&4&5&5&6&6&7&8\\\midrule
$L$&15&16&17&18&19&20&21&22&23&24&25&26&27&28&29\\\midrule
$r_3(L)$&8&8&8&8&8&9&9&9&9&10&10&11&11&11&11\\\bottomrule
\end{tabular}
\end{center}
Together with $r_3(30)=12$, the additional bounds are
\begin{equation}\label{u86:eq:finiteAP}
 r_3(35)\le13,\qquad r_3(39)\le14,\qquad
 r_3(50)=16,\qquad r_3(70)\le20.
\end{equation}
These are finite, exactly verified statements, with the following
exhaustive algorithms and validity proofs; they are not empirical fits.

\paragraph{The table through 30.}
Start from the single empty set at length zero. At step $n$, retain
every old set and also adjoin $n$ whenever no pair $a<b$ in it satisfies
$a+n=2b$. Any new progression must have largest element $n$, so these
are exactly the forbidden pairs. Induction proves that the procedure
lists every AP-free set once, with or without its new largest point.
The integer-mask finite lemma reaches 880288 sets at length 30 and verifies
every maximum in the table. Its full finite loops are included as
\path{verify_apfree_blocks.py}.

\paragraph{The upper bounds at 35, 39 and 70.}
The bounded depth-first search records chosen points and forbidden
future points. It branches on exclusion or inclusion of the least
available future point $x$. Upon inclusion it forbids $2x-y$ for each
previously chosen $y<x$. Inductively these are exactly the future
completions of progressions whose first two points are chosen. It
therefore never omits an admissible set. A node is discarded only when
the number of chosen points plus either the number of available points
or a proved interval-partition bound cannot exceed the proposed bound.
The partition bound uses only the already verified table through 30:
a suffix of length $L$ can be partitioned into blocks of lengths at
most 30, with capacities equal to their table entries. This is an upper
bound on every possible completion, regardless of cross-block constraints.
Thus pruning is safe. The specified exhaustive calculation has respectively
1925, 6081 and 59770661 nodes and proves the three upper bounds in
\eqref{u86:eq:finiteAP}. Its complete implementation is
\path{verify_bounded_ap.py}.

\paragraph{The bound at 50.}
Any AP-free subset of $[1,50]$ of size at least 17 restricts to AP-free
sets in its two 25-point halves. Since $r_3(25)=10$, each half has size
between seven and ten. The exact length-25 enumeration has 140840
sets; its size-seven, eight, nine and ten classes have respectively
39332, 16061, 1940 and 33 members. For every pair of halves of combined
size at least 17, the checker examines cross-block progressions.
Two chosen points $a<b$ in the left half forbid $2b-a$ in the right;
two chosen points $a<b$ in the right half, in local coordinates, forbid
$25+2a-b$ in the left. Every progression crossing the division is of
one of these forms. All 69865347 candidate pairs fail at least one
cross-block test. Hence $r_3(50)\le16$. The explicitly checked set
\[
 \{1,2,4,5,10,11,13,14,28,29,31,32,37,38,40,41\}
\]
is AP-free, giving equality. This entire calculation is reproduced by
\path{verify_apfree50.py} using integer bit masks only.

\subsection{Affine and width-dependent partition bounds}\label{u86:sec:APaffine}
Restricting an AP-free set to consecutive blocks gives
$r_3(L_1+\cdots+L_t)\le\sum_i r_3(L_i)$. Thus any explicit partition
with proved block capacities gives a valid bound; the partition need
not be optimal. The companion files contain exact partitions proving
\begin{align}
 29r_3(L)&\le11L+61 &&(L\ge16),\label{u86:eq:AP29core}\\
 29r_3(L)&\le11L+50 &&(L\ge64),\label{u86:eq:AP29whole}\\
 50r_3(L)&\le16L+176 &&(L\ge1).\label{u86:eq:AP50}
\end{align}
Here is the finite-to-universal proof. For the first inequality, there
is one checked partition for each $L=16,\ldots,44$, using blocks of
length at most 30. Appending any number of length-29 blocks, each of
capacity 11, preserves the inequality and covers every $L\ge16$.
For the second, the same argument starts with $L=64,\ldots,92$.
For the third, checked partitions for $L=1,\ldots,50$ use capacities
derived from the table and the proved bounds at 35, 39 and 50. Append
length-50 blocks, each of capacity 16. The 58 witnesses for the first two affine bounds
and 50 witnesses for the third are included in
\texttt{apfree\_affine\_certificate.json} and
\texttt{apfree\_affine50\_16\_certificate.json}. The checker must verify
that every part is positive, the parts sum to the claimed length,
and their proved capacity sum satisfies the displayed integer inequality.
These checks establish \eqref{u86:eq:AP29core}--\eqref{u86:eq:AP50} in all their
stated ranges.

For the width splits also define a finite partition function. Let
$\mathcal C$ contain the proved capacities for lengths $1,\ldots,30$
and $(35,13),(39,14),(50,16),(70,20)$. For $1\le t\le70$, put
\[
 U(t)=\min_{(L,c)\in\mathcal C}\bigl(c+\max(0,t-L)\bigr),\qquad
 g(0)=0,\quad
 g(n)=\min_{1\le t\le\min(70,n)}\{g(n-t)+U(t)\}.
\]
Every $U(t)$ bounds $r_3(t)$, by restriction when $t\le L$, and by
adding at most $t-L$ further points when $t>L$. Induction using
consecutive block partitions proves $r_3(n)\le g(n)$. The width
certificates use these exactly computed values only through $n=201$.
No monotonicity of $g$ is assumed.

\subsection{Defect inequalities and exact core masses}\label{u86:sec:APdefects}
The diagonal-defective elements of an integer set are AP-free: if
$a<b<c$ were three such elements with $a+c=2b$, then $b$ would not be
defective. For a positive core $X$, its defective set together with
zero is also AP-free, since $(0,x,2x)$ would contradict the defectivity
of $x$ relative to $X\cup\{0\}$. With $r=|X|$ and $R=\max X$, this gives
\[
 \delta_0(X)+1\le r_3(R+1).
\]
For $r\ge15$, \eqref{u86:eq:AP29core} applies because $R+1\ge16$.
Using $\rho(X)=\mu(X)+2-2\delta_0(X)$ gives the core bounds
\begin{align}
 29\rho(X)+22R&\ge29\mu(X)-28,\label{u86:eq:APcore29}\\
 50\rho(X)+32R&\ge50\mu(X)-184.\label{u86:eq:APcore50}
\end{align}
The latter follows from \eqref{u86:eq:AP50} without a minimum-size condition.
Neither AP bound needs primitivity. The existing parity bound gives
$\rho(X)\ge\mu(X)+2r-2R-2$ when an odd element is present.
Each side may independently use a fixed $\rho_*$ bound, parity,
\eqref{u86:eq:APcore29}, or \eqref{u86:eq:APcore50}; substituting the chosen
pair in the exact budget and clearing positive denominators proves
all mixed core-defect rows.

For the full set $D$, write $W=M_p+M_q$ and retain the exact identity
\begin{equation}\label{u86:eq:jointtheta}
 \rho+2\vartheta=\mu+20-2\Delta(D),\qquad
 c+T+\rho+2\vartheta=\Omega.
\end{equation}
Its defect set lies in an integer interval of length $W+1$ and is
AP-free. At this endpoint $W+1\ge d=85\ge64$, so
\eqref{u86:eq:AP29whole} and \eqref{u86:eq:AP50} imply
\begin{align}
 29(c+T+\mu)-22W&\le29\Omega-458,\label{u86:eq:APwhole29}\\
 50(c+T+\mu)-32W&\le50\Omega-616.\label{u86:eq:APwhole50}
\end{align}
There is also a full-set parity bound. Translate $D$ by $M_q$ and
remove zero, obtaining a positive $(d-1)$-element set with maximum
$W$. It contains an odd element: otherwise the translated set,
including $M_q$ because $0\in D$, would be all even, forcing $D$ all
even, contrary to its primitive core. Translation preserves full
restricted-sum defects; zero adds precisely one defect. Therefore
Lemma~\ref{u86:lem:odd-pivot} yields $\Delta(D)\le W-d+4$, and
\begin{equation}\label{u86:eq:wholeparity}
 c+T+\mu-2W\le\Omega-2d-12.
\end{equation}
Equations \eqref{u86:eq:APwhole29}--\eqref{u86:eq:wholeparity} are joint bounds
on $\rho+2\vartheta$. One must not subtract fixed lower bounds for
$\rho$ again when using them.

The strengthened models also retain exact core mass. Put
\[
 C_p=\sum P=S_p-\sum_{i=1}^4x_i,\qquad
 C_q=\sum Q=S_q-\sum_{i=1}^4y_i,
\]
\[
 \mathcal B_p=\Omega+T_a-jd,\quad \mathcal B_q=\Omega+T_b-jd,\quad
 \mathcal B=\Omega+T_a+T_b-jd.
\]
Because $C_p-T_a=\mu(P)+h_P$, $h_P\le\kappa\le T$, and
$h_P+h_Q=\kappa+z$, we have
\begin{equation}\label{u86:eq:exactmass-input}
 T+\mu(P)\ge C_p-T_a,\qquad
 T+\mu\ge C_p+C_q-T_a-T_b-z.
\end{equation}
If a side inequality is written
$\rho(P)\ge\mu(P)-\lambda_pR_p+\nu_p$ and the other side has fixed
lower bound $v_q$, the original budget and \eqref{u86:eq:exactmass-input}
give
\begin{equation}\label{u86:eq:exactmass-one}
 \delta+C_p-\lambda_pR_p\le\mathcal B_p-v_q-\nu_p.
\end{equation}
Reflection gives its companion. For two side inequalities, they give
\begin{equation}\label{u86:eq:exactmass-two}
 \delta+C_p+C_q-z-\lambda_pR_p-\lambda_qR_q
 \le\mathcal B-\nu_p-\nu_q.
\end{equation}
For a joint full-set inequality
$\rho+2\vartheta\ge\mu-\lambda W+\nu$, they give
\begin{equation}\label{u86:eq:exactmass-whole}
 \delta+C_p+C_q-z-\lambda W\le\mathcal B-\nu.
\end{equation}
The parameter pairs $(\lambda,\nu)$ used are
\begin{center}\small
\begin{tabular}{lrrr}
\toprule
Bound&Parity&29-period AP bound&50-period AP bound\\\midrule
Individual core of size $r$&$(2,2r-2)$&$(22/29,-28/29)$&$(32/50,-184/50)$\\
Whole $D$&$(2,2d+12)$&$(22/29,458/29)$&$(32/50,616/50)$\\\bottomrule
\end{tabular}
\end{center}
These substitutions, after clearing denominators, are exactly the
additional mass-defect row families. The derivation uses the original
budget once and does not spend its defect or hole terms twice.

\subsection{Width alternatives and paired-unit alternatives}
In a strengthened high-hole root, let $u_p,u_q$ be its certified core
excess bounds and $v_p,v_q$ its fixed defect bounds. If $W\le t$, then
\[
 \Delta(D)\le r_3(W+1)\le r_3(t+1)\le g(t+1).
\]
Thus \eqref{u86:eq:jointtheta} and \eqref{u86:eq:exactmass-input} justify the
left-child rows
\begin{align*}
 \delta+T&\le\Omega-jd-u_p-u_q-20+2g(t+1),\\
 \delta+C_p+C_q-z&\le\mathcal B-20+2g(t+1).
\end{align*}
The implemented rows replace $T$ by its proved affine lower bound
on the current hole interval. The right child is simply $W\ge t+1$.
The integer width $W$ ensures that the two children are exhaustive.

If instead $R_p\le t$, the anchored defect set gives
$\delta_0(P)+1\le g(t+1)$ and hence
$\rho(P)\ge\mu(P)+4-2g(t+1)$. The corresponding left rows are
\begin{align*}
 \delta+T&\le\Omega-jd-u_p-v_q-4+2g(t+1),\\
 \delta+C_p&\le\mathcal B_p-v_q-4+2g(t+1).
\end{align*}
The right child has $R_p\ge t+1$, and reflection gives the other side.
The checked thresholds are $140\le t\le200$ for $W$ and
$40\le t\le120$ for a core width. Monotonicity is used only for the
true $r_3$ before applying its upper bound at $t+1$; the possibly
nonmonotone partition function $g$ is never assumed monotone.

The trees also extend the earlier paired-member disjunction to
three-empty modes. In a split second-pair case of pairing branch A,
choose a member $s$ from the \emph{actual contributors}. Either the
other member contributes too, implying $\beta_t+M_t\ge1$, or $s$ is
the sole contributor. In the latter case its lower and upper hull
losses may be included in the total-loss budget: the paired unit's
capacity deficit is at least this sole member's hull loss. Introduce
two fresh nonnegative coordinates for those losses. Use global quantiles,
or core quantiles only when the residual cap is below 81. These
alternatives are valid with two or three empty units and do not assert
integrality of every outer element in the three-empty case. The exact
checker verifies that the two new loss coordinates were previously
unused and adds the correct row to each child.


\subsection{Exact finite lemma for size 93}
\label{u86:finite:93}
The required finite assertion has 157766 canonical roots and 160550
leaf contradictions. The data are in \path{proofs/exact92}.
There are 127976 low and 29790 high roots. Exactly 157231 are direct
contradictions; 444 low and 91 high roots have trees. Their 2784 splits
comprise 2328 normal/outside, 231 inner-extreme, 146 core/large-loss,
28 paired-member, and 51 AP-width alternatives. The trees have 3319
leaves. Of the 535 tree roots, 524 use base rows, six add parity rows,
and five use the strengthened AP, exact-mass, and width families.
No ordinary root needs an AP row. The finite AP theorems and partition
witnesses in this chapter are nevertheless required by these five trees.

Every leaf is checked by the integer criterion of
Lemma~\ref{completion:coverage}; every path row must be one of the
proved alternatives. The verifier \path{verify_exact92.py} regenerates
the complete ordered root universe and all required children, rather than
accepting stored counts or an optimizer status. These assertions are finite
mathematical proof obligations in the supplied data.



\section{Full-set group deficits and an exact exclusion of size 92}\label{u86:sec:upper91}
\begin{theorem}\label{u86:thm:upper91}
No feasible real set has 92 elements. Consequently
\[
 \boxed{N\le91.}
\]
\end{theorem}
\begin{proof}
Theorem~\ref{u86:thm:upper92} first forces zero into a hypothetical size-92
example. Integralization and Theorem~\ref{u86:thm:lattice92} give the common
core lattice. The new group-deficit and retained-center integrality
lemmas below justify every additional necessary model row. The exhaustive
classification has 248060 original roots, and its complete finite
certificates give 248068 exact leaf contradictions. Therefore
size 92 is impossible. The earlier exclusions give the stated global
upper bound.
\end{proof}

\subsection{Using full translates in a common hull}
The previous leading-group estimate used the smaller core $E$ to measure
overlap. The following estimate uses the whole set $D$, including its
eight inner extreme elements. It introduces no new finite
progression-free bound.

\begin{proposition}\label{u86:prop:fullDgroups}
Let $D\subset\Z$, $|D|=d$, $\min D=-M_q$, $\max D=M_p$, and
$K=|2D|$. On the positive covering side write
\[
 F_p=\FS(D)\cap(0,\infty),\qquad
 U_0=c_p+2D,\quad U_i=c_p+p_i+D,\quad V_j=c_p-q_j+D,
\]
where $p_i>M_p$, $q_j>M_q$. Choose distinct components of the fixed
covering decomposition, with nominal capacities $K$ for $U_0$ and $d$
for each selected $U_i,V_j$. Their centers must be integral on the
normalized $D$ lattice, and each selected $d$-capacity slot must be
retained and genuinely contribute. Let $I,J$ be the selected own and
cross index sets, $k=|I|$, $l=|J|$, and let $p_*,q_*$ be the largest
selected outer elements when defined. The raw capacity deficit is at least
\begin{align*}
 H_{\rm own}&=K+kd-p_*-M_p-2M_q-1 &&(k>0,\ l=0),\\
 H_{\rm cross}&=K+ld-q_*-2M_p-M_q-1 &&(k=0,\ l>0),\\
 H_{\rm mixed}&=K+(k+l)d-p_*-q_*-M_p-M_q-1 &&(k,l>0).
\end{align*}
If $H$ denotes the applicable bound and $\lambda_i$ is any valid lower
bound for the actual loss of member $i$ outside $F_p$, then its actual
target capacity deficit $G_{I,J}$ satisfies
\[
 G_{I,J}\ge H+\lambda_i\quad\hbox{for every selected member }i,
 \qquad G_{I,J}\ge\sum_i\lambda_i.
\]
The leading component is included among these members. All statements
reflect to the negative target, interchanging $p$ and $q$.
\end{proposition}
\begin{proof}
The own-only raw union is contained in
$[c_p-2M_q,c_p+p_*+M_p]_{\Z}$: the leading component ends at
$c_p+2M_p$, which is smaller than its stated upper endpoint since
$p_*>M_p$, and each own component starts above $c_p-2M_q$.
Its hull has $p_*+M_p+2M_q+1$ points. The cross-only and mixed
hulls are respectively
\[
 [c_p-q_*-M_q,c_p+2M_p]_{\Z},\qquad
 [c_p-q_*-M_q,c_p+p_*+M_p]_{\Z}.
\]
Subtract their cardinalities from the total nominal capacities to obtain
the three bounds. They remain valid if their right-hand sides are negative.

Write $W$ for the raw union and $C$ for the total capacity. Its deficit
against the actual target is
\[
 C-|W\cap F_p|=(C-|W|)+|W\setminus F_p|.
\]
The outside union contains every member's outside set, proving
$G_{I,J}\ge H+\lambda_i$. Also
$|W\cap F_p|\le\sum_i|W_i\cap F_p|$, proving
$G_{I,J}\ge\sum_i\lambda_i$. Both parts use the same target $F_p$;
losses outside its hull are used only as lower bounds for actual losses.
\end{proof}

\subsection{One deficit variable for all subgroups and clipped hulls}
For one sign, let $\mathcal C_p$ contain the leading component, every
retained own slot, and the retained cross slots indexed 3 and 4 in
mathematical outer-element notation. Those two cross slots are always
unpaired in the fixed decomposition. Introduce an auxiliary $\Gamma_p$
whose realization in an actual configuration is
\[
 \Gamma_p=\sum_{W\in\mathcal C_p}|W|
                 -\left|\bigcup_{W\in\mathcal C_p}W\cap F_p\right|.
\]
For any selected subgroup $\mathcal S\subseteq\mathcal C_p$ containing
the leading component, let $G_{\mathcal S}$ be its actual deficit.
Then
\begin{equation}\label{u86:eq:fullDaux}
 \Gamma_p\ge G_{\mathcal S}+
       \sum_{W\in\mathcal C_p\setminus\mathcal S}\lambda_W.
\end{equation}
Indeed, the whole target union is contained in the subgroup target union
plus the separate target contributions of the outside components.
Thus every subgroup inequality simultaneously lower-bounds the same
auxiliary. Also $\Gamma_p\ge\sum_{W\in\mathcal C_p}\lambda_W$.
The corresponding negative family uses distinct covering resources, so
$\Gamma_p+\Gamma_q$ and all unmatched resource losses fit within the
residual covering slack $\delta$. Geometric overlap between groups
causes no problem; their component labels, hence their capacities,
must be disjoint.

\begin{proposition}\label{u86:prop:clippedfullD}
Suppose a selected subgroup of total capacity $C$ lies in an integral
hull $[L,U]_{\Z}$ and has a nonempty contribution to $F_p$.
Since $F_p\subseteq[1,S_p]_{\Z}$, its actual deficit also satisfies
\[
 G_{\mathcal S}\ge C-U,\qquad
 G_{\mathcal S}\ge C-S_p+L-1,\qquad
 G_{\mathcal S}\ge C-S_p.
\]
\end{proposition}
\begin{proof}
Nonemptiness implies $U\ge1$ and $L\le S_p$. The target union is contained
in $[L,U]_{\Z}\cap[1,S_p]_{\Z}$, whose cardinality is at most each of
$U-L+1$, $U$, $S_p-L+1$ and $S_p$. Subtract from $C$.
\end{proof}
These clipped bounds already measure actual target loss. Unlike a raw
hull bound in Proposition~\ref{u86:prop:fullDgroups}, they must not receive
an additional selected member's loss. Only the outside-component losses
from \eqref{u86:eq:fullDaux} are added. In the size-92 implementation every
clipped subgroup contains at least one retained own slot, establishing
the required nonemptiness. The applied hulls are the own-only or mixed
hulls above. Taking suffixes of the retained own and unpaired cross
slots supplies all the inequalities used by the certificate.

\subsection{Outer integrality from retained centers}
\begin{lemma}\label{u86:lem:twoemptyintegral}
If $D$ lies on the normalized integer lattice and at most two covering
units are empty, then all eight outer elements and both leading centers
$c_p,c_q$ are integral. No contribution from a leading $2D$ component
is required.
\end{lemma}
\begin{proof}
Order $p_1>p_2>p_3>p_4$ and $q_1>q_2>q_3>q_4$, and recall
\[
 c_p=p_1-p_2-p_3-p_4,\qquad
 \alpha_p=p_2-p_3-p_4,\qquad
 \beta_p=2p_2-p_1-p_3-p_4,
\]
with reflected definitions. Every retained unpaired $D$ slot has an
integral center, since it meets an integral subset-sum target.
The own slots $c_p+p_i,c_q+q_i$, the tail cross slots
$c_p-q_3,c_p-q_4,c_q-p_3,c_q-p_4$, and both $\alpha$ slots are
always unpaired. In branch A, $c_q-p_2$ is also unpaired; in branch B,
$\beta_p$ is unpaired. Reflection gives the opposite-side statements.

Writing $w_i=c_p+p_i$ gives
$c_p=-w_1+w_2+w_3+w_4$. Thus a complete own side makes its four
outer elements and leading center integral. Suppose first, after
reflection, that the $q$ own side is complete. If exactly one $p$ own
slot is missing, at least one of $c_p-q_3,c_p-q_4$ contributes because
only one other unit can be empty. It makes $c_p$ integral, then the
three retained own slots and the defining identity give every $p_i$.
If two $p$ own slots are missing, all cross and $\alpha$ slots that are
unpaired contribute. The centers $c_p-q_3,c_q-p_3,c_q-p_4$ give
$c_p,p_3,p_4\in\Z$. A retained own slot among $p_1,p_2$ then finishes
the argument. If both are missing, branch A gives $p_2$ from $c_q-p_2$;
branch B gives it from $\alpha_p$. The identity gives $p_1$.

It remains to have one missing own slot on each side. They consume both
empty units, so all unpaired cross, $\alpha$ and unpaired $\beta$ slots
contribute. At least one positive own slot with index $i\in\{3,4\}$
is retained. Adding its center $c_p+p_i$ to $c_q-p_i$ gives
$c_p+c_q\in\Z$. Let $k$ be the missing positive own index. For
$i\ne k$, one has $p_i\equiv-c_p\pmod\Z$.
If $k=1$, the integral center $\alpha_p\equiv c_p\pmod\Z$ gives
$c_p\in\Z$. If $k=3$ or 4, substituting the three known congruences
in the $c_p$ identity gives $p_k\in\Z$; then $c_q-p_k$ gives
$c_q\in\Z$ and hence $c_p\in\Z$.
If $k=2$, the same substitution gives $p_2\in\Z$.
Branch A uses $c_q-p_2$. In branch B,
$\alpha_p\equiv2c_p$ and $\beta_p\equiv3c_p\pmod\Z$; their
integral difference gives $c_p\in\Z$. In every case both leading
centers are integral, all retained own coordinates follow, and the
missing coordinate on each side follows from its defining identity.
\end{proof}
This strengthens the earlier conditional integrality argument. In
particular the zero-hole two-empty residual cap at size 92 can be 174,
which exceeds the general bound $K\ge167$; the preceding lemma justifies
all two-empty strict gaps even there. A redundant exact check of all
578 two-empty center patterns found an integer left inverse in every
case, but the proof above supplies the universal statement.

\paragraph{Individual certificates for the geometric roots.}
The new group rows are used only after an explicit integer-combination
certificate. Let
$\mathbf p=(p_1,p_2,p_3,p_4,q_1,q_2,q_3,q_4)^T$.
Each retained contributing center is an integer linear form in
$\mathbf p$: its translate of $D$ meets the integral target, and $D$ is
integral. Collect these coefficient rows in an integer matrix $C$.
Include a selected paired contributor's center where the covering mode
requires it. The stored certificate is an integer matrix $W$ with
\[
 WC=I_8.
\]
It follows that $\mathbf p=W(C\mathbf p)\in\Z^8$. The checker reconstructs
$C$ from the actual retained-center rows and checks this identity exactly;
it does not rely on a numerical rank computation or a discovered Smith
normal form. Integral outer coordinates give integral leading centers
even if a leading component happens to be empty.

Consequently every strict outer order relation can be strengthened to
an integer gap of at least 1. In pairing branch B the strict sign
$\beta_p+\beta_q>0$ can likewise be strengthened to
$\beta_p+\beta_q\ge1$. The six final low-hole trees use these strengthened
rows only after their individual $WC=I_8$ identities have been checked.
Thus four-empty roots and other cases lacking the earlier general
outer-integrality argument are not treated as integral without proof.

\subsection{The complete size 92 root classification}
The existing upper bound first forces a hypothetical size-92 example to
contain zero: otherwise adjoining zero gives the already excluded size
93. The real-to-integer reduction and Theorem~\ref{u86:thm:lattice92} give
\[
 d=84,\quad e=76,\quad m=75,\quad a+b=75,\qquad
 I=T_a+T_b-T_{37}-T_{38},\quad\Omega=500-I,\quad G=158.
\]
For low holes enumerate $16\le a\le59$. The structurally possible
side sizes $a=14,15$ and their reflections $a=60,61$ already have
negative low-hole caps and therefore contribute no roots. For $z=0$
allow $j=0,1,2,3,4$;
for $1\le z\le14$ allow $j=0,1,2,3$. Retain exactly
\[
 B_0=342-I-84j\ge0\quad(Z),\qquad B_0-19\ge0\quad(P).
\]
The positive-hole condition uses $\rho\ge m-28=47$. The corresponding
necessary simultaneous bounds are
\begin{gather*}
 \delta\le B_0\quad(Z),\qquad
 \delta\le B_0-47+2z\quad(P),\\
 R_p+R_q+\delta\le79+B_0+4z,\qquad
 S_p+S_q+\delta-2K\le1596-84j+4z.
\end{gather*}
\begin{samepage}
All pairing modes and selected contributors give these exact counts:
\begin{center}\small
\begin{tabular}{lrrrrr}
\toprule
Empty units $j$&0&1&2&3&4\\\midrule
$z=0$&152&2240&15028&59840&65280\\
$1\le z\le14$&144&2100&13872&47872&---\\\bottomrule
\end{tabular}
\end{center}
\end{samepage}
Their sum is 206528. Weak outer order remains in the base model whenever
the earlier argument does not prove outer integrality.

For high holes write $(a,b)=(37-k,38+k)$ up to reflection, so
$I=k(k+1)$. Use
\[
 C_{92}(a,z)=500-I-\rho_*(a,u_a)-\rho_*(b,u_b)
                  -\max(53,z,158-2z).
\]
The bounds $\rho\ge47$ and $T\ge53$ give
$c\le400-k(k+1)$, excluding $k\ge20$. Hence enumerate
$0\le k\le19$, $15\le z\le500-k(k+1)$, totaling 7060 states.
The effective-hole lemma's hypotheses are directly checked for all
640 constants used. The largest nonnegative cap is 244, so $j\le2$.
The all-active mass upper bound is $1638+4z$.
Exactly 4198 unoriented states survive: 2013 with $j=0$, 1586 with $j=1$,
and 599 with $j=2$. Their canonical consecutive-$z$ grouping gives
106 intervals, respectively 21, 56 and 29. Both orientations are distinct
because $a+b=75$ is odd. With every mode and contributor choice these
yield 168 $H_0$, 7840 $H_1$, and 33524 $H_2$ roots, totaling 41532.
Thus the complete endpoint has $206528+41532=248060$ original roots.

Here $K\ge167$, nested leading rows require $\delta<159$, and the
core-quantile threshold is hull loss below 80. Every quantile or leading-cardinality row is gated by its
own residual cap; strict outer gaps use Lemma~\ref{u86:lem:twoemptyintegral}
or an individual integer-center certificate. In particular the maximum cap 244 alone does not
establish that a leading component contributes. The integer-center
certificates and the nonempty own slots justify the full-set rows
independently of that stronger assertion.


\subsection{Exact finite lemma for size 92}
\label{u86:finite:92}
The required finite assertion has 248060 canonical roots and 248068
leaf contradictions. The data are in \path{proofs/exact91}.
There are 246487 ordinary direct roots, 1345 low and 222 high direct
geometric roots, and six low geometric tree roots. The latter have eight
splits and 14 leaves: four normal/outside and four inner-extreme/loss
alternatives. The only added row families are full-translate group
bounds, clipped hull bounds, and their simultaneous actual-deficit
realization. Every one of the 1573 geometric roots has a reconstructed
integer-center witness. No AP-width split is needed here; the earlier
size-93 exclusion still requires its own AP finite lemmas.

Every leaf is checked by the integer criterion of
Lemma~\ref{completion:coverage}; every path row must be one of the
proved alternatives. The verifier \path{verify_exact91.py} regenerates
the complete ordered root universe and all required children, rather than
accepting stored counts or an optimizer status. These assertions are finite
mathematical proof obligations in the supplied data.



\section{A common core lattice at size 91}\label{u86:sec:lattice91}
The next exclusion needs a new structural theorem. Its proof must handle
progression subset-sum supports directly: the previous sufficient
simultaneous-inverse inequality can fail at this size.

\subsection{Progression summands and a uniform repair}
We first record two elementary facts. If $W=\{w_1<\cdots<w_v\}\subset\Z$
and $J=[0,u-1]_{\Z}$, then
\begin{equation}\label{u86:eq:APgaps}
 |W+J|=v+u-1+\sum_{i=1}^{v-1}\min(w_{i+1}-w_i-1,u-1).
\end{equation}
This follows by adding the intervals $[w_i,w_i+u-1]$ in order: the
next interval adds $\min(w_{i+1}-w_i,u)$ new points. For a progression
with arbitrary integer step, apply this formula separately to each
occupied residue class. If there are $t$ classes, the sumset excess is
at least $(t-1)(u-1)$.

For positive integers $x_1<\cdots<x_b$, the largest consecutive gap in
their subset-sum support is
\begin{equation}\label{u86:eq:largestgap}
 \max_{1\le i\le b}\left(x_i-\sum_{j<i}x_j\right).
\end{equation}
Indeed, adjoining $[0,L-1]_{\Z}$ to the support produces one interval
exactly when $x_i\le L+\sum_{j<i}x_j$ for every $i$. Sufficiency is
induction on the two overlapping translated intervals. If a test fails,
the interval between the previous total and $x_i$ cannot be filled by
any later coefficient. This proves necessity and \eqref{u86:eq:largestgap}.
If the support has exactly one gap of length at least $u$, a prefix gap
attaining that size and its reflection about half the total must be the
same gap. Its endpoints are the previous total and $x_i$; equality with
the reflected gap forces $i=b$. Thus $x_b$ exceeds the preceding total,
the last two supports are disjoint, and
\begin{equation}\label{u86:eq:dominantmu}
 \mu(X)\ge 2(1+T_{b-1})-1-T_b=T_{b-2}.
\end{equation}

We now apply these facts at size 91. Write $a=37-k\le b=37+k$ and
\[
 d=83,\quad e=75,\quad m=74,\quad
 \Lambda=1407+k^2,\quad\Omega=516-k^2,\quad C=-945-k^2.
\]
Here $C=2m+4e+14-\Lambda$ is the coefficient in the pre-lattice
comparisons \eqref{u86:eq:latt1} and \eqref{u86:eq:latt3}.
The deficient-sign argument gives $|B|=8$: otherwise
$1+T_{75}=2851>24\cdot91+77=2261$.
The old uniform defect estimate first gives $k^2\le534$, hence
$a\ge14$ and the threshold-14 theorem applies. The unpaired-value bound
is $q^2\le2648$; parity gives $q\le50$, so there are at least 12
opposite pairs and zero is not a diagonal exception. The nonnegative
exact budget then gives $k\le22$, hence $a,b\ge15$.

Put $U=\FS(P)$, $V=\FS(Q)$, $u=|U|$, $v=|V|$. Reflection of $V$ is
a translate of $V$ by subset complementation. If both supports are
non-progressions, the sufficient small-side inverse margin is
\[
 T_a-2-\Omega+2a+b-31>0\qquad(0\le k\le22),
\]
whose minimum is 18. This uses $\mu_P\ge a-3$ and
$\rho_P+\rho_Q\ge a+b-28$. The larger-side comparison follows from
$T_b-2>\Omega$. Thus the strict simultaneous inverse theorem gives
one core lattice with $h_P,h_Q\le\kappa$ and strict support-cardinality
hole bounds.

Suppose instead that one support is a progression. Call it
$U=t[0,u-1]_{\Z}$, reflecting the names of the sides if necessary.
If this side has at least 37 coefficients, then
$u-1\ge T_{37}=703>\Omega$. Formula~\eqref{u86:eq:APgaps} forces the other
support into one residue class modulo $t$, with no gap as large as $u$.
On that lattice its hull holes equal $\kappa$; the progression has no
holes, so $z=0$.

It remains that the progression side has $15\le a\le37$, with
$37\le b\le59$. Throughout this range direct quadratic arithmetic gives
\begin{equation}\label{u86:eq:AP91range}
 \Omega<2T_a,\qquad \Omega-T_a\le47,
 \qquad u\ge T_a+1\ge121>2b.
\end{equation}
Since $\kappa\le\Omega$, at most two residue classes of $V$ occur.
If there is one class, normalize $t$. Formula~\eqref{u86:eq:APgaps} allows
at most one gap of $V$ as large as $u$. If such a gap existed,
\eqref{u86:eq:dominantmu} and $\delta_0(Q)\le b$ would give
\[
 \rho_Q\ge T_{b-2}+2-2b\ge558>516\ge\Omega,
\]
which is impossible. Hence $h_P=0$, $h_Q=\kappa$ and $z=0$.

If two residues occur, exactly one off-lattice coefficient would give
two disjoint translated supports and the same contradiction. At least
two coefficients therefore have the nonzero residue $r$. Their sum has
residue $2r\in\{0,r\}$, so $r$ has order two. Normalize the half-step:
$U=2[0,u-1]_{\Z}$ and $Q$ has $s\ge2$ odd coefficients. Let $w$ be
the least odd coefficient and let $h_0,h_1$ be the holes within the
even and odd fiber hulls of $V$, measured with step two. By
\eqref{u86:eq:APgaps} and $\kappa<2(u-1)$, no fiber gap is as large as $u$,
and therefore
\[
 \kappa=u-1+h_0+h_1,\qquad
 h_P=u-1,\qquad h_Q=\kappa-u+w,\qquad z=w-1.
\]
For the formula for $h_Q$, the least odd support element is $w$.
Complement symmetry gives the opposite fiber endpoints; the number of
integer positions outside their two hulls is $w-1$, whether the total
sum is even or odd.

We need $w\le u$. If $w\ge u+1>2b$, positivity and distinctness give
\[
 \sum Q\ge(b-s)(b-s+1)+sw+s(s-1).
\]
Since $|V|=\sum Q+1-h_Q$, the resulting lower bound is minimized at
$s=2$, and it yields
$|V|\ge b^2-3b+2u+6-\kappa$.
But the exact budget and $\rho_Q\ge\mu_Q+2-2b$ give
$\mu_Q+\kappa\le\Omega+2b-2$, hence
$|V|\le T_b+\Omega+2b-1-\kappa$.
The lower minus upper margin is
$(b^2-11b)/2+2u+7-\Omega\ge214$, a contradiction.
Thus $h_P,h_Q\le\kappa$ on this common half-step lattice.

In fact $w\le13$. If the odd number $w$ were at least 15, the same
mass estimate, minimized over real $s$, would imply
\[
 \sum Q-w-T_b
 \ge\min_s\{T_b+2s^2+s(15-2b-2)-15\}
 =7b-\frac{289}{8}.
\]
On the other hand the identity
$\mu_Q+\kappa=\sum Q+u-w-T_b$ and \eqref{u86:eq:AP91range} give
$\sum Q-w-T_b\le\Omega-u+2b-2\le2b+44$.
These inequalities contradict each other for $b\ge37$. Hence $z\le12$.

In this two-residue case a sharper span estimate forces the outer layers
onto the lattice without using the post-lattice constant. The progression
side has $R_P\le2(a+\mu_P)$, by subtracting the least possible sum
$T_{a-1}$ from its normalized full support total. The other side satisfies
$R_Q\le b+\mu_Q+\kappa-T_a-\mu_P+12$.
Substitution in the exact $D_1$ comparison gives
\[
 |\FS(D_1)|-2s
 \le2543+2a-4T_a-2T_b\le-1157.
\]
Also $\Lambda\ge1407>4\cdot83-10+516=838$, so the first-layer
comparison forces $D$ onto the same lattice. Only now is $G=156$
available. It gives
$T\ge\kappa+156-3z\ge T_a+120$, contradicting
$\Omega-T_a\le47$. Thus the two-residue case is excluded.
All remaining progression-support cases have $z=0$ and are forced
integral by the old comparison $C+\Omega=-429-2k^2<0$ and the same
first-layer inequality.

\subsection{A defect-sensitive span bound and the remaining forcing step}
\begin{lemma}\label{u86:lem:defectspan}
For a positive real set $Y$ of size at least 24,
$\rho(Y)\ge\delta_0(Y)-4$. Consequently, for a positive integer set
$X$ of size $r\ge25$ with hull-hole count $h$,
\[
 \max X\le r+\rho(X)+h+7-\delta_0(X).
\]
\end{lemma}
\begin{proof}
Write $s=|Y|$ and $\eta=\rho(Y)$. If $\eta<s-14$, the existing
flexible inverse bound with $\beta=2$ gives $\delta_0(Y)\le\eta+4$.
If $s-14\le\eta<2(s-14)$, use $\beta=3$ to get
$\delta_0(Y)\le\eta/2+6\le\eta+4$, since $\eta\ge10$.
In the remaining range, $\eta\ge2(s-14)\ge s-4\ge\delta_0(Y)-4$.
Now remove $R=\max X$, putting $Y=X\setminus\{R\}$. If $v$ old
diagonal defects are destroyed on adjoining $R$, the exact recurrence is
\[
 R=r+\rho(X)+h(X)+2-\rho(Y)-h(Y)-2v.
\]
Since $\delta_0(Y)\ge\delta_0(X)-1$ and $h(Y),v\ge0$, the first
part gives the stated bound. No primitivity assumption is used.
\end{proof}

Summing the new span bound for the two cores and using
$2(\delta_0(P)+\delta_0(Q))=\mu+4-\rho$ gives, before $D_1$ is
known to be integral,
\begin{equation}\label{u86:eq:defectforce}
 |\FS(D_1)|-2s
 \le C+16+3\rho-2\mu+\kappa+2z+\ell_2
 \le C+3\Omega+16-2\mu.
\end{equation}
The final step uses only $z\le\kappa$ and
$\rho+\kappa+\ell_2\le\Omega$, with each nonnegative term charged once.

For the remaining non-progression supports the old one-budget comparison
is at most $-429-2k^2+2z$, so it is strict for $z\le214+k^2$.
Otherwise $z\ge215+k^2$. Since $\rho+\mu\ge2m-34=114$, the old
three-budget comparison reduces the possible failures to $0\le k\le11$.
Both cores then have at least 26 elements. The effective-hole rows below
give $\mu(X)\ge r+C_h+1$; each $(C_h,\beta)$ satisfies the two strict
tests in \eqref{u86:eq:holetest}, and $2\beta<215$ throughout:
\begin{center}\small
\begin{tabular}{r|rrrrrrrrrrrr}
\toprule
$r$&26&27&28&29&30&31&32&33&34&35&36&37\\
$C_h$&48&53&59&64&70&76&83&89&96&103&111&118\\
$\beta$&5&6&6&6&6&7&7&7&7&8&8&8\\\midrule
$r$&38&39&40&41&42&43&44&45&46&47&48&\\
$C_h$&126&134&143&151&160&169&179&188&198&208&219&\\
$\beta$&8&9&9&9&9&10&10&10&10&11&11&\\\bottomrule
\end{tabular}
\end{center}
The resulting full list of forcing bounds in \eqref{u86:eq:defectforce} is
\begin{center}\small\setlength{\tabcolsep}{4pt}
\begin{tabular}{r|rrrrrrrrrrrr}
\toprule
$k$&0&1&2&3&4&5&6&7&8&9&10&11\\
$\mu_P+\mu_Q\ge$&312&313&313&315&316&319&321&325&328&333&337&343\\
$|\FS(D_1)|-2s\le$&$-5$&$-11$&$-23$&$-47$&$-77$&$-119$&$-167$&$-227$&$-293$&$-371$&$-455$&$-551$\\\bottomrule
\end{tabular}
\end{center}
All are strict. Two disjoint $s$-point translates would result from an
off-lattice element of $D_1\setminus E$, so $D_1$ is integral. The
first-layer inequality already displayed then forces $D$ to be integral.

\begin{theorem}\label{u86:thm:lattice91}
A zero-containing feasible integer set of size 91 has the common-core
lattice structure of Theorem~\ref{u86:thm:lattice}, with $d=83$, $e=75$,
$m=74$ and $G=156$. Both core supports satisfy the strict
support-cardinality hole bounds after this reduction.
\end{theorem}
\begin{proof}
The preceding arguments cover both non-progression supports and every
progression-support case. After $D$ is integral, hole propagation and
Theorem~\ref{u86:thm:sharp} give \eqref{u86:eq:newbudget} with $G=156$.
In a remaining $z=0$ case, $\kappa\le\Omega-156$. The minimum of
$T_a-2-\Omega$ is $-49$, so $T_a+\mu_P-\kappa\ge109>0$;
the larger-side comparison is stronger. Thus both hole counts are
strictly below the corresponding support cardinalities minus one.
Positive-hole cases already had these strict inverse bounds.
\end{proof}
The exact companion \path{verify_size91_lattice.py} checks the finite inverse
parameter tables and forcing margins. The uniform progression-support
repair above is an analytic proof, and no use of $G$ precedes
$D$-integrality. This theorem alone excludes no cardinality.

\section{Fixed-score losses and an exact exclusion of size 91}\label{u86:sec:upper90}
\begin{theorem}\label{u86:thm:upper90}
No feasible real set has 91 elements. Consequently
\[
 \boxed{N\le90.}
\]
\end{theorem}
\begin{proof}
The preceding upper bound forces zero in a hypothetical size-91 set.
Integralization and Theorem~\ref{u86:thm:lattice91} give the common core
lattice. The complete necessary-model classification below has 432560
roots. Its finite certificates give 432758 exact leaf
contradictions, using the new fixed-score inequalities proved below.
This excludes size 91 and proves the general upper bound.
\end{proof}

\subsection{Eight additional losses and exact fixed-score indicators}
The original covering has 110 labeled fixed scores, with at most
84 positive ones. Its 26 disjoint deletion groups have coordinatewise
nonpositive coefficient sums in the eight nonnegative outer gap variables;
each group therefore supplies at least one nonpositive score. These
original groups and all 110 coefficient vectors are regenerated exactly
by the accompanying fixed-score checker. We retain these labels even
when two labels happen to represent the same value: this weakens an upper
bound for a union and preserves the capacity accounting.

Suppose $z=0$. Hole propagation makes $\FS(D)=[-S_q,S_p]_{\Z}$.
Assume the largest positive own slot $u+D$, $u=c_p+p_1$, is empty,
its center is integral, and a lower own slot $v+D$, $v=c_p+p_i$ with
$i>1$, contributes. Then $u>S_p$. Indeed, $0<u\le S_p$ would give
the contribution $u+0$. If $u\le0$, choose $d'\in D$ with
$1\le v+d'\le S_p$. Since $u>v$, one has
$1\le u+d'\le d'\le M_p\le S_p$, again an integral contribution.
Thus the four fixed labels
\[
 C^+_{1i}=c_p+p_1+p_i\qquad(1\le i\le4)
\]
all exceed $S_p$. They are positive and hence distinct from the
nonpositive deletions. The reflected argument supplies four more losses
if the largest negative own slot is empty with a lower contributor.
These eight labels can therefore be charged as extra fixed losses.
The scalar definitions of $c$ and $\delta=c-jd$ still use the original
fixed capacity 84; the new loss is an additional term in its allocation.

For any other fixed label $f$, introduce $H_f\in[0,1]$, realized by an
actual configuration as the indicator that its score exceeds its own
target maximum $S_p$ or $S_q$. The high-loss auxiliary may equal
$\sum H_f$ and be added to every covering-deficit budget. For the seven
roots using the eight losses above, only the other 102 labels enter this
sum. Integral scores have the exhaustive alternatives
\[
 f\le S_{\rm side},\qquad
 \hbox{or}\quad f\ge S_{\rm side}+1\ \hbox{and}\ H_f\ge1.
\]
All score coefficients and their target sides are fixed by the original
covering, so each such split is reconstructed directly.

At the balanced root an additional nonpositive-loss auxiliary is used.
For all 110 labels introduce $N_f\in[0,1]$, realized as the indicator
of $f\le0$. Each of the 26 disjoint groups imposes
$\sum_{f\in\mathcal G}N_f\ge1$. Thus
\[
 L_{\le0}=\sum_fN_f-26\ge0
\]
is a valid realization of the extra nonpositive loss beyond the already
charged 26. It may be combined with high losses and the eight separately
charged positive labels, since these three sets of labels are disjoint.
The binary alternatives $f\le0$, $N_f\ge1$, or $f\ge1$ cover every
integral score. The LP may relax these indicators to real variables:
the stated actual realizations prove all necessary rows and both branches.

These rows are used only after individual retained-center identities
$WC=I_8$ establish outer integrality. The seven applicable zero-hole
roots have $a=34,\ldots,40$ and missing units
\texttt{pair1,pair2,p0,q0}; all required lower own slots are retained.
No statement about an arbitrary empty-largest slot is needed.

\subsection{Complete roots and the weak outer-order geometry}
At this endpoint $d=83$, $e=75$, $m=74$, $I=(a-37)^2$,
$\Omega=516-I$ and $G=156$. Enumerate low-hole sides $15\le a\le59$
and empty-unit counts $j=0,\ldots,4$, retaining exactly
\[
 B_0=360-I-83j\ge0\quad(Z),\qquad B_0-18\ge0\quad(P).
\]
The underlying bounds are
\begin{gather*}
 \delta\le B_0\quad(Z),\qquad
 \delta\le B_0-46+2z\quad(P),\qquad1\le z\le14\quad(P),\\
 R_p+R_q+\delta\le78+B_0+4z,\qquad
 S_p+S_q+\delta-2K\le1578-83j+4z.
\end{gather*}
The positive-hole four-empty wrapper starts with the more permissive
zero-hole cap, then installs the positive-hole mass penalties and exact
$z$ coefficients. Consequently its inherited optional quantile and
leading rows remain valid. The exact counts are
\begin{center}\small
\begin{tabular}{lrrrrr}
\toprule
$j$&0&1&2&3&4\\\midrule
$Z_j$&148&2310&15606&62832&119680\\
$P_j$&148&2310&15606&56848&76160\\\bottomrule
\end{tabular}
\end{center}
Their sum is 351648.

For high holes, $(a,b)=(37-k,37+k)$ up to reflection and
\[
 C_{91}(a,z)=516-k^2-\rho_*(a,u_a)-\rho_*(b,u_b)
                              -\max(52,z,156-2z).
\]
Since $\rho\ge46$, $T\ge52$ gives $c\le418-k^2$, so
$0\le k\le20$ suffices. Enumerate $15\le z\le516-k^2$: there are
7672 integer states, with 646 directly validated effective-hole constants.
The largest cover cap is 268, permitting $j\le3$. The all-active mass
upper bound is $1620+4z$.
The 5614 retained states have counts $(2568,2014,944,88)$ for
$j=0,1,2,3$. Their 110 canonical intervals have counts $(21,52,29,8)$.
Both orientations are used except at $a=b=37$. The resulting high roots
number $(164,7070,31790,41888)$, totaling 80912. All roots together number
$351648+80912=432560$.

The leading, nested and core-quantile thresholds are respectively
$165$, $157$ and $79$. With $j\le2$, Lemma~\ref{u86:lem:twoemptyintegral}
justifies strict outer gaps. For $j\ge3$ the ordinary and geometric
models keep weak real outer order. The actual residual caps are at most
111 in the low regime and 19 in the high regime; both are below
$K\ge165$. Therefore both leading components contribute and $c_p,c_q$
are integral. Every retained own or cross center is integral, aligning
the selected outer endpoints needed for Propositions~\ref{u86:prop:fullDgroups}
and \ref{u86:prop:clippedfullD}. Missing outer coordinates need not be
integral. The checker tests the caps at both ends of each affine interval.
Among $j\ge3$ roots, only the seven fixed-score roots restore all strict
outer gaps through
their separate $WC=I_8$ certificates.


\subsection{Exact finite lemma for size 91}
\label{u86:finite:91}
The required finite assertion has 432560 canonical roots and 432758
leaf contradictions. The data are in \path{proofs/exact90}.
There are 428966 ordinary roots, 3546 direct geometric roots, 41
geometric tree roots with 84 splits and 125 leaves, and seven fixed-score
tree roots with 114 splits and 121 leaves. The 198 splits comprise 74
normal/outside, four core/large-loss, 30 inner-extreme, 58 high-score,
and 32 nonpositive-score splits. No integer-coordinate split is selected.
The 110 fixed labels, all 26 disjoint groups, exact score rows, and all
retained-center identities must be reconstructed.

Every leaf is checked by the integer criterion of
Lemma~\ref{completion:coverage}; every path row must be one of the
proved alternatives. The verifier \path{verify_exact90.py} regenerates
the complete ordered root universe and all required children, rather than
accepting stored counts or an optimizer status. These assertions are finite
mathematical proof obligations in the supplied data.



\section{A sharper non-progression defect and the size 90 lattice}\label{u86:sec:lattice90}
The next structural reduction uses a sharper non-progression penalty and
applies the effective-hole theorem to the cores with their maxima removed.
The side of size 14 requires a separate argument.

\subsection{A sharp non-progression penalty}
\begin{theorem}\label{u86:thm:nonAPsharp}
If $X$ is a positive real set of size $r\ge21$ and $\FS(X)$ is not an
arithmetic progression, then $\rho(X)\ge r-5$. The constant is sharp.
\end{theorem}
\begin{proof}
First consider a deleted interval
$Y=[1,j+q]_{\Z}\setminus H$, where $j\ge10$, $|H|=q\le3$, and
the maximum $j+q$ is retained. The case $q=0$ has full support. For
$q\ge1$, Lev's central-interval theorem applies because
$j+q\le3j/2-2$. It gives
\[
 [2q+1,\sum Y-(2q+1)]_{\Z}\subseteq\FS(Y).
\]
Thus its total hole count is twice the number of missing subset sums in
$[1,2q]_{\Z}$, determined entirely by the deleted positions in that range.
The density of $Y$ forces gcd one, so a progression support is exactly a
full interval.

Every diagonal defect lies in $[1,q]$ or $[j,j+q]$. To see this, an
interior element $x$ with $q+1\le x\le j-1$ has more than $q$
disjoint pairs of distinct positions from $[0,j+q]$ centered at $x$;
deleting only $q$ positions cannot destroy all pairs. The low defects
depend only on positions through $2q$, and the high defects only on
depths through $2q$ below the retained maximum.

For $j\ge14$, these two bands are disjoint. Preserve the deleted low
positions and high depths, and let $l,u$ be their respective counts.
For the other $t=q-l-u$ deletions choose the depths
$2q+1,\ldots,2q+t$ in a new interval of 14 retained elements.
They fit between the two bands for every $q\le3$. The resulting set
$Y_0$ has exactly the same hole count and the same defect count.
Writing $h=j+q$ and measuring each deleted depth as $h-d$, one has
\[
 \mu(Y)=\sum_{d\in H}(h-d)-T_{q-1}-\operatorname{holes}(\FS(Y)).
\]
Each low deletion contributes an additional depth $j-14$ in $Y$;
the original middle-depth sum is at least the chosen one. If the
support is not full, at least one low deletion is necessary, so $l\ge1$.
Consequently
\[
 \rho(Y)-j\ge\rho(Y_0)-14+(l-1)(j-14)\ge\rho(Y_0)-14.
\]
This reduces all sizes to the finite cases $j=10,11,12,13,14$.
The finite checker must enumerate every $q=1,2,3$ and every
$H\subseteq[1,j+q-1]$ with $|H|=q$. It forms the subset-sum support
by repeated unions with its translates, and tests every diagonal by
all possible distinct-pair representations. The complete counts are
\begin{center}\small
\begin{tabular}{lrrrrr}
\toprule
$j$&10&11&12&13&14\\\midrule
All deletion sets&285&363&454&559&679\\
Non-full supports&130&154&180&208&238\\
Minimum $\rho-j$ for non-full support&$-5$&$-5$&$-5$&$-5$&$-5$\\\bottomrule
\end{tabular}
\end{center}
Thus all 2340 cases, including all 910 non-full ones, are covered.
This finite exhaustive check and the preceding signature reduction prove
$\rho(Y)\ge j-5$ for every non-progression deleted interval in this range.

For an arbitrary positive real set, use its increasing prefixes $X_i$
and increments $y_i$ from Section~\ref{u86:sec:lattice}. The exact
recurrence gives $\rho_i\ge\rho_{i-1}+y_i-2$, and globally
$\rho_i\ge-8$. If every $y_i\ge4$ for $10\le i\le r$, then
$\rho_r\ge-8+2(r-9)=2r-26\ge r-5$.
Otherwise take the last $j\ge10$ with $y_j\le3$. The
small-increment inverse theorem, valid since $3\le j-5$, makes $X_j$
a scaled deleted interval with at most three deletions and its maximum
retained. If its support is not a progression, the result just proved
and all later $y_i\ge4$ give $\rho_r\ge r-5$.
If its support is a progression, let $k>j$ be the first non-progression
support. A progression and its translate, if they intersect, have their
centers on the same progression step and their union is a progression.
Thus this first non-progression step has disjoint supports, and
\[
 y_k=|\FS(X_{k-1})|-k\ge1+T_{k-1}-k,
 \qquad\rho_k\ge T_{k-1}-k-9\ge k-5
\]
for $k\ge11$. Later increments again increase $\rho$ by at least two.
Finally $X=\{1,3,4,\ldots,r+1\}$ has two diagonal defects and
$\mu=r-3$, so $\rho=r-5$.
\end{proof}
The companion \path{verify_small_deleted.py} checks precisely this
finite base. It is separate from the earlier length-70 AP-free computation.

\subsection{The common lattice and the size 14 boundary}
At size 90 put
\begin{gather*}
 d=82,\quad e=74,\quad m=73,\quad a=36-k,\quad b=37+k,\\
 I=k(k+1),\quad\Lambda=1370+I,\quad\Omega=531-I,\quad C=-914-I.
\end{gather*}
The coarse unpaired bound gives $a\ge11$. For $a=11,12,13$, the
exact values of $\Omega$ are $-119,-69,-21$, each below the small-side
global bound $-8$ plus the other side's nonnegative defect. They are
impossible. Hence $a\ge14$, both defects are nonnegative, and $k\le22$.

Set $v_0(r)=r-5$ for $r\ge21$ and $v_0(r)=r-14$ for $14\le r<21$.
If both supports are non-progressions, the strict inverse test follows from
\[
 (a-3)+v_0(a)+v_0(b)>\Omega-T_a+2\qquad(0\le k\le22).
\]
The minimum margin is 7; the larger-side test is automatic.
Thus the common core lattice has $h_P,h_Q\le\kappa$, with strict
support-cardinality hole bounds.

If a progression side has at least 37 elements, the proof of
Section~\ref{u86:sec:lattice91} gives one residue and $z=0$, since
$T_{37}=703>531$. For a smaller progression side, the same repair
applies with the following endpoint-specific checks. One has
$\Omega<2T_a$ and $\Omega-T_a\le75$. A sole off-lattice generator
or a unique gap of length at least $u$ would again imply
$\rho_Q\ge558>531$, so neither is possible. In the two-residue case,
$\kappa\ge u-1\ge T_a$ forces $a\ge18$, whence
$u\ge172>2b$ and $37\le b\le55$.
The lower-minus-upper mass margin excluding $w>u$ is at least 301.
Also $\Omega-u\le74$. If $w\ge15$, the same quadratic estimate would
give $\sum Q-w-T_b\ge7b-289/8$, while the budget gives at most
$2b+72$, a contradiction for $b\ge37$. Hence $w\le13$ and $z\le12$.
The improved progression-side span now gives
\[
 |\FS(D_1)|-2s\le2513+2a-4T_a-2T_b\le-1089.
\]
The first-layer inequality is $1370>4\cdot82-10+531=849$.
Only after these establish $D$-integrality may we use $G=154$;
then $T\ge\kappa+154-3z\ge T_a+118$ contradicts
$\Omega-T_a\le75$. All remaining progression-support cases have $z=0$.

The boundary $a=14$, $b=59$ has $I=506$, $\Omega=25$.
It must be resolved before using Lemma~\ref{u96:lem:span96}: its prefix has
size 13 and need not have nonnegative $\rho$.
The exact maximum-removal identity gives instead
\[
 R\le r+\rho_X+h_X+2-\rho_{X\setminus\{R\}}
 \le r+\rho_X+h_X+10,
\]
using only the global lower bound $-8$ for that prefix. The other side
keeps the usual constant 2. Therefore the $D_1$ comparison gains 16 and
is still at most
\[
 -383-2\cdot506+2\cdot25+16=-1329<0.
\]
The first layer is strict as above, so $D$ is integral. Now
$T\ge\lceil154/3\rceil=52>\Omega=25$ excludes this boundary entirely.
All subsequent uses of the ordinary prefix span have $a\ge15$.

\subsection{Applying the effective-hole bound to the prefixes}
For $k\le21$, the one-budget $D_1$ comparison is at most
$-383-2I+2z$. It is strict for $z\le191+I$. Assume otherwise
$z\ge192+I$. The bound
$\kappa+\rho_P+\rho_Q\le\Omega$, together with $v_0$ above, excludes
$k\ge12$. Thus $0\le k\le11$ and both core sizes lie between 25 and 48.

For either core $X$ of size $r$, remove its maximum $R$ and write
$Y=X\setminus\{R\}$. The last two subset-sum translates cannot be
disjoint. If they were, \eqref{u86:eq:dominantmu} would give
$\rho_X\ge T_{r-2}+2-2r$. Adding the other core's effective-hole
defect bound and $\kappa\ge192+I$ exceeds $\Omega$ in every case;
the minimum surplus for the two choices is shown below.
The translates therefore intersect. Their hulls overlap, and every hole
of their union is a hole of one of the two constituent hull supports, so
$h_X\le2h_Y$. Hence $h_Y\ge\lceil z/2\rceil\ge96$.
The exact increment identity also gives
\[
 h_Y=h_X-R+r+(\mu_X-\mu_Y)
 \le\kappa+\mu_X\le\Omega-\rho_{\rm other}+2r-2.
\]
The latter upper bound is strictly below $T_{r-1}\le|\FS(Y)|-1$,
with the minimum strict margin recorded in the table. Thus the
effective-hole theorem applies to both prefixes on the same unit lattice.

For $26\le r\le48$, use the effective-hole rows in
Section~\ref{u86:sec:lattice91}, supplemented by $(r,C_h,\beta)=(24,39,5)$
and $(25,43,5)$. All have $2\beta<96$. Let $u_r=r+C_h+1$ and obtain
$v_r$ from the already proved flexible inverse bounds. Explicitly one may
take the maximum of $r-14$, $u_r+2-2r$ and, for $2\le\beta\le r-5$,
\[
 \min\!\left((\beta-1)(r-\max(14,\beta+4)),\,
 (\beta-1)q+\min(t,\beta-1)\right),
\]
where $\max(0,u_r-4\beta+2)=(\beta+1)q+t$, $0\le t\le\beta$.
The same inverse argument extends to the displayed range of $\beta$.
Start after $L=\max(14,\beta+4)$. If every later increment exceeds
$\beta$, then $\rho\ge(\beta-1)(r-L)$. Otherwise the last prefix
with increment at most $\beta$ has at most $2\beta$ diagonal defects
by the small-increment inverse theorem. If $s$ elements follow it,
then $s\le\lfloor\rho/(\beta-1)\rfloor$ and the final defect count
is at most $2\beta+s$. Thus
$\mu-4\beta+2\le\rho+2\lfloor\rho/(\beta-1)\rfloor$.
Exact integer inversion yields $(\beta-1)q+\min(t,\beta-1)$ above.
Taking the minimum with the no-small-increment alternative proves the
formula for every $2\le\beta\le r-5$. The size-90 arithmetic checker
repeats all its parameter tests.

The exact maximum-removal identity now strengthens the forcing bound to
\begin{equation}\label{u86:eq:prefixforce90}
 |\FS(D_1)|-2s
 \le C+3\Omega-\rho-\mu-2(\rho_{P'}+\rho_{Q'}),
\end{equation}
where $P',Q'$ have their maxima removed. Put
$U=u_a+u_b$, $V=v_a+v_b$, $V'=v_{a-1}+v_{b-1}$.
The complete finite checks are
\begin{center}\small
\begin{tabular}{r|rrrrrr}
\toprule
$k$&$U$&$V$&$V'$&Disjoint surplus&Prefix-hole margin&Forcing upper bound\\\midrule
0&304&190&177&284&127&$-169$\\
1&304&190&177&263&103&$-177$\\
2&305&191&178&247&82&$-197$\\
3&307&193&180&237&65&$-229$\\
4&309&196&182&232&51&$-270$\\
5&311&198&185&232&40&$-320$\\
6&314&201&188&237&32&$-380$\\
7&318&205&192&248&28&$-452$\\
8&322&210&196&264&27&$-533$\\
9&326&214&201&285&29&$-623$\\
10&331&219&206&311&34&$-723$\\
11&337&225&212&343&43&$-835$\\\bottomrule
\end{tabular}
\end{center}
The surplus and hole-margin columns are the minima over the two cores.
The last column is $679-4I-U-V-2V'$, the bound in
\eqref{u86:eq:prefixforce90}. Every value is strict. It forces $D_1$ onto
the core lattice, and the first-layer inequality forces $D$.

\begin{theorem}\label{u86:thm:lattice90}
A zero-containing feasible integer set of size 90 has the common-core
lattice structure of Theorem~\ref{u86:thm:lattice}, with $d=82$, $e=74$,
$m=73$ and $G=154$. The side of size 14 is impossible, and the two
remaining cores satisfy the strict support-cardinality hole bounds.
\end{theorem}
\begin{proof}
The preceding cases establish $D$-integrality before using $G$.
Hole propagation and Theorem~\ref{u86:thm:sharp} then give the full
post-lattice budget. In a remaining $z=0$ case,
$\kappa\le\Omega-154$ and $\Omega-T_a\le75$ imply
$T_a-\kappa\ge79>0$. This supplies both strict hole bounds; the
positive-hole cases already had them. The half-step progression case
and the size-14 boundary were excluded above.
\end{proof}
The \path{check90.py} verifies all displayed
finite arithmetic. The theorem is a structural reduction, not by itself
a cardinality exclusion.

\section{The prefix lattice reduction and an exact exclusion of size 90}\label{u86:sec:upper89}
\begin{theorem}\label{u86:thm:upper89}
No feasible real set has 90 elements. Consequently
\[
 \boxed{N\le89.}
\]
\end{theorem}
\begin{proof}
Theorem~\ref{u86:thm:upper90} justifies adjoining zero in a hypothetical
size-90 example. Integralization and Theorem~\ref{u86:thm:lattice90}, including
its separate size-14 argument, give the common lattice and every required
structural hypothesis. The complete necessary-model classification below
has 622492 roots, all excluded by 624102 exact leaf
contradictions. The preceding descending exclusions therefore yield
$N\le89$. This is an intermediate endpoint conclusion;
Section~\ref{u86:sec:upper88} excludes size 89 as well.
\end{proof}

\subsection{Complete endpoint classification}
Use the structural parameters $d=82$, $e=74$, $m=73$,
$I=T_a+T_b-T_{36}-T_{37}$, $\Omega=531-I$ and $G=154$.
The conservative low list considers $14\le a\le59$ and every
$j=0,\ldots,4$, retaining precisely
\[
 B_0=377-I-82j\ge0\quad(Z),\qquad B_0-17\ge0\quad(P).
\]
The size-14 structural boundary is already excluded, and its negative low
cap supplies no root. The necessary simultaneous bounds are
\begin{gather*}
 \delta\le B_0\quad(Z),\qquad
 \delta\le B_0-45+2z\quad(P),\qquad1\le z\le14\quad(P),\\
 R_p+R_q+\delta\le77+B_0+4z,\qquad
 S_p+S_q+\delta-2K\le1560-82j+4z.
\end{gather*}
The positive-hole four-empty wrapper uses the same conservative-template
argument as at size 91. All modes and contributor choices give
\begin{center}\small
\begin{tabular}{lrrrrr}
\toprule
$j$&0&1&2&3&4\\\midrule
$Z_j$&152&2380&17340&65824&152320\\
$P_j$&152&2380&16184&65824&130560\\\bottomrule
\end{tabular}
\end{center}
Their sum is 453116.

For high holes, write $(a,b)=(36-k,37+k)$ up to reflection and use
\[
 C_{90}(a,z)=531-k(k+1)-\rho_*(a,u_a)-\rho_*(b,u_b)
                              -\max(52,z,154-2z).
\]
The deliberately conservative domain is $0\le k\le22$,
$15\le z\le531-k(k+1)$, comprising 7843 integer states. All 713
effective-hole constants are rechecked directly. The largest cover cap
is 289, so $j\le3$. The all-active mass bound is $1602+4z$.
The 6564 retained unoriented states have counts $(2920,2223,1156,265)$;
their 163 intervals have counts $(33,67,45,18)$ for $j=0,1,2,3$.
Both side orientations are distinct. All modes then give high-root counts
$(264,9380,52020,107712)$, totaling 169376. Thus there are
$453116+169376=622492$ original roots.

The leading, nested and core-quantile thresholds are $163$, $155$ and
$78$. For $j\ge3$, the low residual cap is at most 131 and the high cap
at most 43, both below $K\ge163$. Hence leading centers are integral,
and the retained endpoints suffice for full-$D$ geometry while missing
outer coordinates keep weak order. For $j\le2$,
Lemma~\ref{u86:lem:twoemptyintegral} applies, with the redundant 578-mode
integer-left-inverse check matched to every ordinary two-empty model.

\subsection{Selected geometric and fixed-score proofs}
Of the original roots, 616746 have ordinary exact contradictions. Their
4626-vector pool has 103362 positive integer weights, maximum weight
1554 and minimum contradiction magnitude 1. The other 5746 original
roots divide into 5663 direct geometric proofs and 83 complete trees.
The direct models use only the full-$D$ raw and clipped group inequalities
already proved. Every used labeled row is matched against its freshly
regenerated complete model.

Every branch root has an explicit retained-center identity $WC=I_8$.
Only after checking it are strict outer gaps and fixed-score integer
dichotomies used. The 71 geometric trees have 236 splits and 307 leaves:
178 normal or outside alternatives, five core or large-loss alternatives,
and 53 inner-extreme alternatives. The 12 fixed-score trees have 1374
splits and 1386 leaves: 715 nonpositive-score, 635 high-score, and 24
inner-extreme alternatives. These are exactly the proved above
split types; no integer-coordinate or integer-form split is selected.

The fixed-score loss realization is the one in
Section~\ref{u86:sec:upper90}. Missing-largest charges are used only for
zero-hole roots with the required retained lower contributor. The two
positive-hole fixed roots receive no such charge. Their high and
nonpositive indicators remain valid, since a score above $S_p$ or $S_q$
or at most zero is outside its respective positive target even when that
target has holes. Every family excludes from its high-indicator sum any
label already charged separately, so the deficit allocation does not
double count. All 110 score formulas and the 26 disjoint deletion groups
are regenerated.

The direct proof records contain 123118 positive weights, with largest
weight 15258. The 307 geometric leaf vectors have 5410 weights and
largest weight 348; the 1386 fixed-score leaf vectors have 42822 weights
and largest weight 2082. All satisfy the exact nonnegative-column and
negative-right-side contradiction criterion.


\subsection{Exact finite lemma for size 90}
\label{u86:finite:90}
The required finite assertion has 622492 canonical roots and 624102
leaf contradictions. The data are in \path{proofs/exact89}.
The complete ordinary/direct/tree partition gives
\[616746+5663+307+1386=624102.\]
These are the leaf contradictions. There are 83 tree roots
and 1610 binary splits. The exact verification reconstructs every root,
every branch, all retained-center identities and high-hole intervals.
It also verifies the complete 2340-case deleted-interval base and the
size-90 structural scalar arithmetic, including the size-14 boundary.

Every leaf is checked by the integer criterion of
Lemma~\ref{completion:coverage}; every path row must be one of the
proved alternatives. The verifier \path{verify_exact89.py} regenerates
the complete ordered root universe and all required children, rather than
accepting stored counts or an optimizer status. These assertions are finite
mathematical proof obligations in the supplied data.



\section{The common core lattice at size 89}\label{u86:sec:lattice89}
The next endpoint requires a further structural reduction. The crude
simultaneous inverse criterion can fail even when both core supports are
non-progressions. A dense large summand repairs precisely those cases.
The proof below uses no finite AP70 table and no effective-hole table to
obtain the lattice. Effective-hole bounds enter only the later finite
high-hole classification.

\begin{theorem}\label{u86:thm:lattice89}
Suppose that a feasible integer set $A$ contains zero and has size 89.
Then $|B|=8$, and the two peeled layers lie on a common core lattice.
After normalization, $D$ is integral and
\[
 d=81,\quad e=73,\quad m=72,\quad
 (a,b)=(36-k,36+k),\quad I=k^2,
\]
\[
 \Lambda=1333+I,\qquad \Omega=546-I,\qquad C=-883-I.
\]
The core supports are individually primitive on that lattice. Their
hole counts satisfy
\[
 h_P,h_Q\le\kappa,\qquad
 z=h_P+h_Q-\kappa,\qquad 0\le z\le\min(h_P,h_Q),
\]
and the post-lattice conclusions are
\begin{equation}\label{u86:eq:budget89}
 w\le4z,\qquad T\ge\kappa+152-3z,\qquad
 T\ge z,\qquad T\ge51,\qquad T+\rho\le\Omega.
\end{equation}
Every surviving model has $a,b\ge15$. This theorem does not assert that
the original eight outer elements all belong to the core lattice.
\end{theorem}

\subsection{Bootstrap and the side of size 13}
If $|B|\le7$, one sign is absent from $D$ and the one-sided core has at
least 73 nonzero elements. Its support would satisfy
\[
 s\ge1+T_{73}=2702>24\cdot89+77=2213,
\]
contrary to \eqref{u86:eq:corebound}. Hence $|B|=8$. The unpaired-value
bound gives $q^2\le96\cdot73+1073-73^2=2752$.
Since $q$ is even, $q\le52$; thus the core has at least ten opposite
pairs and $a,b\ge10$. Zero is not a diagonal exception. The algebraic
identity \eqref{u86:eq:budget} is valid already here; only the small-side
$\rho$ term may be negative. Since $b\ge36$, the universal lower bound
$\rho_P\ge-8$ and $\rho_Q\ge0$ imply $\Omega\ge-8$.
Therefore $k\le23$ and $a\ge13$.

At $a=13,b=59,k=23$, one has $\Omega=17$. If the larger support were
not an AP, the non-progression theorem of Section~\ref{u86:sec:lattice90}
would give $\rho_Q\ge54$, incompatible with $\rho_P\ge-8$ and the
budget. Thus the larger support is an AP. Its size is at least 1771,
whereas $\kappa\le25$. The residue bound following \eqref{u86:eq:APgaps}
forces the smaller support into its lattice, and the same gap formula
gives $h_P=\kappa,h_Q=0,z=0$.

The maximum-deleted small prefix has size 12, so its $\rho$-value must
be bounded below by $-8$, not by zero. The exact maximum-removal recurrence from
Lemma~\ref{u96:lem:span96} gives $+10$ on the small side in place of $+2$.
The resulting $+16$ correction to \eqref{u86:eq:latt1} yields
\[
 |\FS(D_1)|-2s\le C+\Omega+\rho-\mu+16
 \le-337-2\cdot529+16=-1379<0.
\]
For the outer layer its available slack is at most
$\Omega-\rho\le25$; the strict comparison is
$\Lambda=1862>4\cdot81-10+25=339$.
Thus both layers are integral. Only now may one invoke $G=152$, giving
$T\ge51$, whereas the budget gives $T\le25$. This excludes $a=13$.
Henceforth $a,b\ge14$, $0\le k\le22$, and all pre-lattice budget
terms are nonnegative.

\subsection{A progression support and its half step alternative}
Put $U=\FS(P)$, $V=\FS(Q)$ and $u=|U|$, $v=|V|$.
If an AP side has at least 36 coefficients, its support has at least
667 points. Since $\Omega\le546$, the residue and gap arguments above
force the other support onto its lattice, with no long gap. Thus the
AP side has zero holes, the other has $\kappa$ holes, and $z=0$.

Otherwise name the smaller AP side $P$. Exact quadratic arithmetic for
$14\le a\le36$ gives
\begin{equation}\label{u86:eq:AP89range}
 \Omega<2T_a,\qquad \Omega-T_a\le102.
\end{equation}
The residue bound permits at most two classes for $V$ modulo the AP
step. With one residue, \eqref{u86:eq:APgaps} permits at most one gap of
length at least $u$. If it exists, the reflection argument in
\eqref{u86:eq:largestgap}--\eqref{u86:eq:dominantmu} gives
$\mu_Q\ge T_{b-2}$ and $\rho_Q\ge T_{b-2}+2-2b$.
A long gap also requires $\kappa\ge u-1\ge T_a$, so $\Omega\ge T_a$.
These rows have $16\le a\le32$ and $40\le b\le56$, and
\[
 T_{b-2}+2-2b-\Omega\ge133>0.
\]
This contradicts the budget. Therefore this case has $h_P=0$,
$h_Q=\kappa$ and $z=0$.

With two residues, exactly one off-lattice coefficient would likewise
give two disjoint copies of the remaining support and the same
contradiction. At least two have a common nonzero residue $r$, so
$2r=0$ modulo the AP step. Normalize the half-step. Then
$U=2[0,u-1]_{\Z}$ and $Q$ has $t\ge2$ odd coefficients. If $w_0$ is
the least odd coefficient and $\eta$ is the sum of the holes in the two
parity-fiber hulls, the fiber version of \eqref{u86:eq:APgaps} gives
\[
 \kappa=u-1+\eta,\qquad h_P=u-1,
 \qquad h_Q=\kappa-u+w_0,\qquad z=w_0-1.
\]
Here no fiber gap reaches $u$, since $\kappa<2(u-1)$; reflection leaves
exactly $w_0-1$ ordinary lattice positions outside the two fiber hulls.

We first show $w_0\le u$. In these rows $u\ge T_a+1>2b$.
If $w_0\ge u+1$, positivity and distinctness imply
\[
 \sum Q\ge(b-t)(b-t+1)+t w_0+t(t-1).
\]
After subtracting $w_0$, this quadratic is minimized over integers
$t\ge2$ at $t=2$. It follows that
\[
 |V|\ge b^2-3b+2u+6-\kappa.
\]
But $\mu_Q+\kappa\le\Omega+2b-2$ gives
$|V|\le T_b+\Omega+2b-1-\kappa$. Their difference is at least
\[
 \frac{b^2-11b}{2}+2u+7-\Omega\ge1065>0,
\]
a contradiction. Consequently $h_P,h_Q\le\kappa$.
In fact $w_0\le13$. If $w_0\ge15$, minimize the same mass expression
over real $t$ to obtain
\[
 \sum Q-w_0-T_b\ge7b-\frac{289}{8}.
\]
The identity $\mu_Q+\kappa=\sum Q+u-w_0-T_b$, the budget and
\eqref{u86:eq:AP89range} give the opposite bound
$\sum Q-w_0-T_b\le\Omega-u+2b-2\le2b+99$.
These are inconsistent for $b\ge40$, so $z\le12$.

The half-step case needs its sharper AP span estimate before lattice
forcing. Since $P$ is even, $R_P\le2(a+\mu_P)$, while
\[
 R_Q\le b+\mu_Q+\kappa-T_a-\mu_P+12.
\]
Substitute these into the exact pre-integrality comparison, using
$\mu\le\rho+2m-4$ and $\rho+\kappa+\ell_2\le\Omega$ once. This gives
\begin{align*}
 |\FS(D_1)|-2s
 &\le8m+2a-2T_a+30-\Lambda+\Omega\\
 &=2483+2a-4T_a-2T_b\le-1021<0.
\end{align*}
The outer layer is strict because
$\Lambda\ge1333>4\cdot81-10+546=860$. Thus $D$ is integral.
Now $G=152$ gives $T\ge\kappa+152-3z\ge T_a+116>\Omega$,
contradicting \eqref{u86:eq:AP89range}. All surviving AP cases therefore
have $z=0$; their ordinary lattice forcing is completed below.

\subsection{The dense large summand repair}
Suppose both supports are non-progressions. Write
\[
 \nu(r)=\begin{cases}r-5,&r\ge21,\\r-14,&14\le r<21.\end{cases}
\]
The large-side inverse condition is automatic since
$\kappa\le\Omega<T_b-1$. A deliberately conservative sufficient
small-side condition for Stanchescu's Corollary 2 is
\begin{equation}\label{u86:eq:inverse89}
 \mu_P+\rho_P+\rho_Q>\Omega-T_a+2.
\end{equation}
Using $\mu_P\ge a-3$ and $\rho_X\ge\nu(|X|)$ shows this condition
unless $9\le k\le17$. Whenever it holds, the corollary gives the
common-lattice hull bounds, including the strict inequalities
$h_P<u-1$ and $h_Q<v-1$.

When \eqref{u86:eq:inverse89} actually fails, subtracting the same lower
bounds gives the following exhaustive list. The last column bounds
$b-\lfloor\rho_Q/3\rfloor$ from below.
\begin{center}\small
\begin{tabular}{rrrrrr}
\toprule
$k$&$a$&$b$&$\Omega$&Upper bound on $\rho_Q$&Prefix size\\\midrule
9&27&45&465&43&31\\
10&26&46&446&53&29\\
11&25&47&425&60&27\\
12&24&48&402&64&27\\
13&23&49&377&65&28\\
14&22&50&350&63&29\\
15&21&51&321&58&32\\
16&20&52&290&59&33\\
17&19&53&257&48&37\\\bottomrule
\end{tabular}
\end{center}

\begin{lemma}\label{u86:lem:dense89}
Let $X$ be a positive real $r$-set with $45\le r\le53$,
$\eta=\rho(X)\le65$ and $r-\lfloor\eta/3\rfloor\ge27$.
There is a scale $t>0$ such that $X/t$ is integral and primitive, and
\[
 \FS(X/t)=[0,L]_{\Z}\setminus\bigl(H\cup(L-H)\bigr),
 \quad H\subseteq[1,8]_{\Z},\quad |H|\le4.
\]
In particular, the support has at most eight holes and contains
$[9,L-9]_{\Z}$.
\end{lemma}
\begin{proof}
Let $X_i$ be the increasing prefixes, with nonnegative excess increments
$y_i$. The recurrence is $\rho_i\ge\rho_{i-1}+y_i-2$;
threshold 14 gives $\rho_i\ge0$ for $i\ge14$.
If $y_i\ge5$ for all $15\le i\le r$, then
$\eta\ge3(r-14)>65$. Let $j\ge15$ be the last index with $y_j\le4$.
The recurrence gives $j\ge r-\lfloor\eta/3\rfloor\ge27$.
Set $f=y_j\le4$. The counted-hole lemma applies because $j\ge2f+5$:
\[
 X_j=t([1,j+f]_{\Z}\setminus F),\quad |F|=f,
 \quad \FS(X_j/t)=[0,L_j]_{\Z}\setminus(H\cup(L_j-H)),
\]
where $H\subseteq[1,2f]$ and $|H|\le f$.
The retained $j$ coefficients have gcd one, since a divisor at least two
would allow at most $\lfloor(j+f)/2\rfloor<j$ positions.

No later coefficient can leave $t\Z$. At the first such index $i$,
the two subset-sum translates are disjoint, and therefore
\[
 y_i=|\FS(X_{i-1})|-i\ge1+T_j-r,
 \qquad \rho_i\ge T_j-r-1\ge324.
\]
All later increments exceed four, so $\rho_r\ge\rho_i$, a contradiction.
The full coefficient set is integral and remains primitive on this scale.

For $i>j$, monotonicity gives $\rho_i\le\eta$ and hence $y_i\le67$.
We verify directly that the endpoint holes persist. Suppose there are
$m\ge27$ current coefficients, their support is
$[0,L]\setminus(H\cup(L-H))$, and the next integer coefficient is $x$.
Then $L\ge T_m$ and $x>\max X_m\ge m>8$.
If $x\ge L-16$, at least
$\min(x,L+1)-8\ge L-24\ge T_m-24>m+68$ support points of
$x+\FS(X_m)$ lie above $L$. This exceeds the allowed increment
$m+1+y_{m+1}\le m+68$. Thus $x\le L-17$.
The central intervals $[9,L-9]$ and $[x+9,x+L-9]$ overlap or are adjacent,
covering $[9,L+x-9]$. Because $x>8$, the lower holes remain exactly $H$;
reflection supplies the upper holes $L+x-H$. Induction proves the claim.
\end{proof}

Apply Lemma~\ref{u86:lem:dense89} to $Q$ and normalize its primitive scale.
If $U$ occupied $c$ residue classes modulo this scale, the ordered
sumset bound in each class would give
$|U+V|\ge u+c(v-1)$, hence $\kappa\ge(c-1)(v-1)$.
But $v-1\ge T_b>\Omega\ge\kappa$, so $c=1$ and all of $P$ lies
on the same lattice. We already have $h_Q\le8$.

The interval $J=[9,L_Q-9]\subseteq V$ has length
$\ell=L_Q-17=v+h_Q-18\ge v-18$.
If $g_i$ are consecutive gaps of $U$, \eqref{u86:eq:APgaps} and
$U+J\subseteq U+V$ imply
\[
 \sum_i\min(g_i-1,\ell-1)\le\kappa+18.
\]
Here $v\ge T_{45}+1=1036$, $\kappa\le465$, and
$\kappa+18<\ell-1$. Thus no term is truncated, proving
$h_P\le\kappa+18$. Consequently
\[
 \operatorname{diam}U\le u+\kappa+17\le u+v-3,
 \qquad \operatorname{diam}V\le v+7\le u+v-3.
\]
The first uses $\kappa+20\le v$, and the second $u\ge191$.
Stanchescu's Theorem 4 now applies directly: the stronger bound
$u+v-3$ covers its equal- and unequal-diameter alternatives.
This short-diameter theorem has no additional gcd premise.
It gives $|U+V|\ge u+v-1+\max(h_P,h_Q)$, hence
\[
 h_P,h_Q\le\kappa,\qquad
 z=h_P+h_Q-\kappa\le\min(h_P,h_Q)\le8.
\]
This repairs every failed inverse row without using a post-lattice
inequality to establish the lattice.

\subsection{Forcing the peeled layers and proving primitivity}
For a positive side $X$ of size $r$, maximum $R$, and maximum-deleted
prefix $Y$, let $v_0$ count old diagonal defects destroyed on reinsertion.
The exact recurrence is
\begin{equation}\label{u86:eq:removedmax89}
 R=r+\rho_X+h_X+2-\rho_Y-h_Y-2v_0.
\end{equation}
For $r\ge15$, $\rho_Y\ge0$. The ordinary one-budget comparison gives
\begin{equation}\label{u86:eq:D1ordinary89}
 |\FS(D_1)|-2s\le C+\Omega+\rho-\mu+2z
                 \le-337-2I+2z.
\end{equation}
It is strict if $z\le168+I$, including every surviving AP branch and
every dense-summand branch. At $a=14,b=58$, the prefix of size 13
requires the same $+16$ correction as before. Since $\Omega=62$,
\[
 |\FS(D_1)|-2s\le-337-2\cdot484+2\cdot62+16=-1165<0.
\]

It remains that $a\ge15$ and $z\ge169+I$. This cannot be an AP or a
repaired dense-summand case. Both supports are non-progressions.
The bounds $\kappa\ge z$ and $\rho_X\ge\nu(|X|)$ rule out $k\ge13$;
thus $0\le k\le12$ and both side sizes are at least 24.
For either side of size $r$, with other side of size $r'$, its final two
subset-sum translates cannot be disjoint. Otherwise
\eqref{u86:eq:dominantmu} gives
\[
 \rho_X\ge T_{r-2}+2-2r,\qquad \rho_{X'}\ge r'-5,
\]
and, for every such $k$ and either side,
\[
 T_{r-2}+2-2r+(r'-5)+169+I-\Omega\ge80>0,
\]
contradicting the budget. The final hulls therefore overlap and
$h_X\le2h_Y$. Retaining $-h_Y$ in \eqref{u86:eq:removedmax89} yields
$R_X\le r+\rho_X+h_X/2+2$.
Substitution in the original $D_1$ comparison now gives
\begin{align*}
 |\FS(D_1)|-2s
 &\le C+2\rho-\mu+z+\ell_2\\
 &\le C+\Omega+\rho-\mu\le-337-2I<0,
\end{align*}
using $z\le\kappa$ and $\rho+\kappa+\ell_2\le\Omega$ once.
Thus $D_1$ lies on the core lattice in every case.
Since $\Lambda\ge1333>4\cdot81-10+546$, the outer peeled layer is
strict as well. The standard overlap and hole propagation arguments
now give \eqref{u86:eq:budget89} with $G=152$.

The $a=14$ row is now excluded entirely. A non-AP large support costs
$\rho_Q\ge53$, so $T+\rho\ge51+53=104>62$.
An AP large support has $z=0$ and $T\ge152>62$.
No subsequent actual size-89 model uses a universal core-span row at
side size 14.

Finally, the strict Corollary 2 cases are already primitive. In the
surviving AP cases, $z=0$ and $\kappa\le\Omega-152<T_a$ by
\eqref{u86:eq:AP89range}. In the dense-summand cases $z\le8$ gives
$\kappa\le\Omega-128<T_a$. In either case also $\kappa<T_b$.
A coefficient gcd at least two would force at least $u-1$ or $v-1$
holes in its support, contradicting $h_X\le\kappa<T_{|X|}$.
This proves Theorem~\ref{u86:thm:lattice89}. The checks in
\path{check89.py} and \path{check_dense_repair.py} reproduce every
finite margin used above; the universal assertions are proved in the text.

\section{Translated core overlap and the size 89 endpoint}\label{u86:sec:upper88}
\begin{theorem}\label{u86:thm:upper88}
There is no feasible set of 89 real numbers. Consequently,
\[
 \boxed{N\le88.}
\]
\end{theorem}
\begin{proof}
Theorem~\ref{u86:thm:upper89} already excludes all larger cardinalities.
If a feasible size-89 set omitted zero, adjoining zero would give a
feasible size-90 set, a contradiction. Integralize without changing its
cardinality and apply Theorem~\ref{u86:thm:lattice89}.
The complete necessary-model partition and exact contradictions are
described below. Their finite certificate excludes every
model, so $N\le88$.

\end{proof}

\subsection{Complete low and high hole classification}
At this endpoint let $I=(a-36)^2$, $b=72-a$ and
$B_0=394-I-81j$. The low list uses $15\le a\le57$ and
$j=0,\ldots,4$, retaining $B_0\ge0$ for zero holes and
$B_0-16\ge0$ for $1\le z\le14$. Its simultaneous necessary bounds are
\begin{gather*}
 \delta\le B_0\quad(Z),\qquad
 \delta\le B_0-44+2z\quad(P),\\
 R_p+R_q+\delta\le76+B_0+4z,\qquad
 S_p+S_q+\delta-2K\le1542-81j+4z.
\end{gather*}
As at the preceding endpoint, the positive-hole four-empty wrapper
inherits optional rows from a more permissive zero-hole template before
installing its positive-hole mass penalties. This preserves their gates.
The mode counts, including selected contributor choices, are
\[
 (4,70,578,2992,10880)\qquad(j=0,1,2,3,4).
\]
The root counts are
\begin{center}\small
\begin{tabular}{lrrrrr}
\toprule
$j$&0&1&2&3&4\\\midrule
$Z_j$&156&2450&17918&74800&184960\\
$P_j$&156&2450&16762&68816&163200\\\bottomrule
\end{tabular}
\end{center}
Their sum is 531668.

For high holes, set $(a,b)=(36-k,36+k)$ up to reflection. With the
previously defined effective-hole excess bounds $u_a,u_b$ and flexible
defect bounds $\rho_*$, the cover cap is
\[
 C_{89}(a,z)=546-k^2-\rho_*(a,u_a)-\rho_*(b,u_b)
                         -\max(51,z,152-2z).
\]
Enumerate $0\le k\le21$ and $15\le z\le546-k^2$, comprising 8393
raw integer states. All 656 effective-hole constants used in this
classification are directly checked against their strict hypotheses.
The maximum cover cap is 311, so $j\le3$.
The all-active mass upper bound is $1584+4z$.
The 8319 retained unoriented states have counts
$(3527,2684,1534,574)$; the 195 intervals have counts
$(43,73,52,27)$, for $j=0,1,2,3$.
Use both side orientations unless $a=b=36$.
The resulting high-root counts are
$(340,9940,57800,149600)$, totaling 217680.
Thus the canonical original list contains exactly
\[
 531668+217680=749348\quad\hbox{roots}.
\]

The leading, nested and core-quantile thresholds are respectively
161, 153 and 77. For $j\ge3$, ordinary models keep weak real outer
order; their residual caps are at most 151 in the low regime and 68 in
the high regime, below $K\ge161$. The leading components therefore
contribute, and retained centers align the endpoints actually used in
full-$D$ geometry. They do not put missing outer coordinates on the
lattice. For $j\le2$, Lemma~\ref{u86:lem:twoemptyintegral} applies without
a leading-contribution premise, and all 578 two-empty center matrices
have independently reconstructed integer left inverses.
The checker validates all relevant gates on every ordinary model;
2660 fast/full representatives provide a separate parameter check.

\subsection{Actual group deficits and the final overlap inequality}
The direct selected models use the established full-$D$ group geometry,
raw overlap and clipped-hull inequalities. Their auxiliary variables
have one simultaneous realization against the actual positive target.
We record the small additional inequality needed for the final five
roots.

\begin{lemma}\label{u86:lem:translatedcore89}
Let $D$ be integral, $R_p,R_q\in\Z_{\ge0}$,
$E\subseteq D\cap[-R_q,R_p]$, $|D|=d$ and
$|E|=e$. Let $M_p=\max D>0$ and let an outer element $p>M_p$ have
$p-M_p\in\Z$. For any center $c$, put
\[
 A_0=c+2D,\qquad A_1=c+p+D,\qquad K=|2D|,
 \qquad H=2e-R_p-R_q-p+M_p-1.
\]
For a common target $F$, let $l_0=K-|A_0\cap F|$ and
$l_1=d-|A_1\cap F|$. Then the two-component capacity deficit satisfies
\[
 K+d-|(A_0\cup A_1)\cap F|\ge H+\max(l_0,l_1).
\]
The same conclusion holds with any lower bound $\lambda\le l_1$ in
place of $l_1$.
\end{lemma}
\begin{proof}
The sets $A_0,A_1$ contain $c+M_p+E$ and $c+p+E$, respectively.
Those two $e$-element sets lie in a translate of an integer interval of
cardinality $R_p+R_q+p-M_p+1$. Hence
$|A_0\cap A_1|\ge H$. This remains valid when $H<0$.
The exact identity
\[
 K+d-|(A_0\cup A_1)\cap F|
 =|A_0\cap A_1|+|(A_0\cup A_1)\setminus F|
\]
and $|(A_0\cup A_1)\setminus F|\ge\max(l_0,l_1)$ prove the claim.
No contribution by $A_0$ is required.
\end{proof}

For a larger labeled component family, partition off this pair and add
the individual lower losses of all remaining components. Capacity
deficits are superadditive under such a partition: the union has
cardinality at most the sum of the subunions' cardinalities.
The positive and negative families use disjoint covering labels and
combine through the established compact full-$D$ budget.
Reflection swaps $M_p,p$ with $M_q,q$.

In the implementation, start from the ordinary baseline
\[
 -G_s+L_s+\sum_{i\in J_s}\lambda_i\le0,
\]
where $G_s$ is the actual full-family deficit, $L_s$ is the leading loss,
and $\lambda_i$ is the sum of the $i$th $D$ component's two endpoint losses.
For each retained own component $i$, Lemma~\ref{u86:lem:translatedcore89}
gives the two alternative rows
\begin{align*}
 -G_s+L_s+\sum_{j\ne i}\lambda_j
                  -R_p-R_q-p_i+M_p&\le1-2e,\\
 -G_s+\lambda_i+\sum_{j\ne i}\lambda_j
                  -R_p-R_q-p_i+M_p&\le1-2e,
\end{align*}
and their negative-side reflections. These rows impose alternatives
on the same actual deficit; they do not sum overlapping pair groups.
They use the raw-overlap identity, so no member loss is added to an
already clipped hull bound. At size 89, $1-2e=-145$.
Three retained own slots on each sign produce exactly twelve added rows.
The required relative integrality follows from the selected roots'
checked integer left inverses.

Two other selected roots use whole-side deficits. On each sign, group
the leading component of capacity $K$ and all seven retained unpaired
$D$ components. If $G_s^{\rm all}$ is their actual target deficit and
the positive target has cardinality at most $S_s$, then
\[
 G_s^{\rm all}\ge K+7\cdot81-S_s,
 \qquad G_s^{\rm all}\ge L_s+\sum_{i\in J_s}\lambda_i.
\]
The two sign groups are disjoint in covering labels, giving
$G_p^{\rm all}+G_q^{\rm all}\le\delta$.
These use separate whole-side variables, replacing rather than adding
the old smaller-family losses. No maximum-member loss is added to the
clipped target-cardinality bound. The separately proved fixed-score losses
may further strengthen this whole-side covering budget: their fixed labels
are disjoint from the grouped components and from all precharged labels.

\subsection{Selected roots and exhaustive binary proofs}
The ordinary proof excludes 739722 roots. Its 5370 integer Farkas
vectors contain 119594 positive weights, with maximum weight 1440
and minimum contradiction magnitude 1. The other 9626 roots partition
into 9472 direct geometric proofs and 154 complete binary trees.
The direct roots comprise 6462 low and 3010 high cases.
Their 9472 proof records use 618 distinct vectors and 194302 positive
weights. Every used labeled row is matched to a regenerated model.

All 154 tree roots have checked retained-center identities $WC=I_8$.
Strict outer gaps and integer fixed-score splits are used only after
that root-specific check. Twelve high roots include selected $p/q$
variants; their selected beta centers are included in the matrices.
The two selected both/sole paired-component branches use the global
81/80 quantiles and reserve loss coordinates 52 and 53. Individual
paired-member losses are charged only in the applicable sole branch.
Every split keeps both children and is a proved catalog disjunction.
The complete census is 500 normal, 13 core, 682 inner-extreme,
two both-or-sole, 1300 fixed-high and 1515 fixed-nonpositive splits.
No arbitrary integer-coordinate or new midpoint split is selected.

The fixed-score rows retain the zero/positive-hole distinction.
Four of the thirteen low fixed-score roots have zero holes and the
required missing largest slots and retained lower contributors; they
receive eight missing-largest losses and 102 remaining high-label
candidates. Nine positive-hole roots, and the single high fixed-score
root, receive no such saving and use 110 high-label candidates.
Any additional nonpositive-label group losses are disjoint from the
high-label and separately precharged labels.
The two whole-side roots and five translated-overlap roots have zero
holes and the same valid eight-loss saving. No criterion sufficient
only for contribution is used to choose a contributor.

The complete selected tree partition is
\begin{center}\small
\begin{tabular}{lrrrr}
\toprule
Selected class&Roots&Nodes&Binary splits&Leaves\\\midrule
Low minimal geometry&56&360&152&208\\
High minimal geometry&77&1579&751&828\\
Low fixed score&13&4813&2400&2413\\
High fixed score&1&133&66&67\\
Low whole side&2&626&312&314\\
Low translated core overlap&5&667&331&336\\\midrule
Total&154&8178&4012&4166\\\bottomrule
\end{tabular}
\end{center}
The whole-side roots are original low indices 2427 and 2433.
The final five are indices 2428 through 2432, corresponding to
$a=34,35,36,37,38$, zero holes, branch A, missing units
\texttt{pair1}, \texttt{pair2}, \texttt{p0}, \texttt{q0}, contributor
choice zero, and no variant. Their added dependency consists exactly
of the twelve rows in Lemma~\ref{u86:lem:translatedcore89}, on top of the
already proved minimal geometry and fixed-score models.
Exploratory midpoint, dense-mass, expanded-side, restricted-pair and
nine-hull techniques are not dependencies of this proof.


\subsection{Exact finite lemma for size 89}
\label{u86:finite:89}
The required finite assertion has 749348 canonical roots and 753360
leaf contradictions. The data are in \path{proofs/exact88}.
The complete partition is
739722 ordinary roots, 9472 direct geometric roots, and 154 tree roots.
The trees have 4012 splits and 4166 leaves, so
749348=739722+9472+154 and
753360=739722+9472+4166. The source-to-model reconstruction includes the
dense-summand repair, all high-hole intervals, every translated-core
row, all 578 two-empty witnesses, and the 2340-case non-AP finite base.

Every leaf is checked by the integer criterion of
Lemma~\ref{completion:coverage}; every path row must be one of the
proved alternatives. The verifier \path{verify_exact88.py} regenerates
the complete ordered root universe and all required children, rather than
accepting stored counts or an optimizer status. These assertions are finite
mathematical proof obligations in the supplied data.


\section{The common core lattice at size 88}\label{u86:sec:lattice88}
The next structural reduction retains the order of implications used
above: obtain a common core lattice, force both peeled layers onto it,
and only then use the strengthened post-lattice budget. A primitive
large-side effective-hole argument repairs the new failed inverse rows.

\begin{theorem}\label{u86:thm:lattice88}
Let $A$ be a feasible integer set of size 88 containing zero. Then
$|B|=8$, and after a common normalization the peeled set $D$ is integral.
With $a\le b$, the parameters are
\[
 d=80,\quad e=72,\quad m=71,\quad
 (a,b)=(35-k,36+k),\quad I=k(k+1),
\]
\[
 \Lambda=1297+I,\qquad\Omega=560-I,\qquad C=-853-I.
\]
The core supports are individually primitive and satisfy
\[
 h_P,h_Q\le\kappa,\qquad z=h_P+h_Q-\kappa,
 \qquad0\le z\le\min(h_P,h_Q).
\]
Every surviving case has $a,b\ge15$, and
\begin{equation}\label{u86:eq:budget88}
 w\le4z,\qquad T\ge\kappa+150-3z,\qquad
 T\ge z,\qquad T\ge50,\qquad T+\rho\le\Omega.
\end{equation}
The original eight covering elements need not all be integral on this
lattice merely by this theorem.
\end{theorem}

\subsection{Bootstrap and the three small side boundaries}
The deficient-sign argument gives $|B|=8$, since otherwise
\[
 s\ge1+T_{72}=2629>24\cdot88+77=2189.
\]
The unpaired-value bound is $q^2\le2801$. Here $q$ is odd, so $q\le51$;
there are at least ten opposite pairs in $E$ and $a,b\ge10$.
The exact budget \eqref{u86:eq:budget} applies algebraically; its only
possibly negative terms are the two $\rho$ values. Since $b\ge36$,
$\rho_Q\ge0$ and $\rho_P\ge-8$ imply $k\le23$ and $a\ge12$.
As before, $\rho\le\mu$.

The boundary calculations are
\begin{center}\small
\begin{tabular}{rrrrrrr}
\toprule
$a$&$b$&$k$&$I$&$\Omega$&Lower $\rho_P$&Upper $\kappa$\\\midrule
12&59&23&552&8&$-8$&16\\
13&58&22&506&54&$-8$&62\\
14&57&21&462&98&0&98\\\bottomrule
\end{tabular}
\end{center}
Write $K_0$ for the last column. In all three rows,
$K_0<T_a-1$ and $K_0<T_b-1$. Stanchescu's strict Corollary 2 therefore
gives the common lattice, $h_P,h_Q\le\kappa$, and both strict
support-cardinality hole bounds.
The maximum-deleted small prefix has size 11, 12 or 13; use
$\rho_Y\ge-8$ in \eqref{u86:eq:removedmax89}, adding 16 to the ordinary
$D_1$ comparison. Hence
\[
 |\FS(D_1)|-2s\le-293-2I+2K_0+16,
\]
whose three values are $-1349,-1165,-1005$.
The outer-layer comparisons $\Lambda>4\cdot80-10+K_0$ are also strict.
Both peeled layers are consequently integral before $G=150$ is invoked.

Now $T\ge50$ excludes $a=12$, since its available budget is at most 16.
For $a=13$ or 14 with a non-AP larger support, respectively,
\[
 T+\rho\ge50-8+53=95>54,
 \qquad T+\rho\ge50+52=102>98.
\]
If that larger support is an AP, its strict hole bound forces step one
on the common lattice, so $h_Q=0$ and $z=0$. Then $T\ge150$ exceeds
the available budgets 62 and 98. All three boundaries are excluded.
Thus $0\le k\le20$, $a,b\ge15$, and every later maximum-deleted prefix
has size at least 14.

\subsection{Progression supports with the necessary extra defect term}
An AP side with at least 36 coefficients has support size at least 667,
whereas $\kappa\le\Omega\le560$. The residue and gap formula
\eqref{u86:eq:APgaps} forces the other support into the same lattice, with
no long gap: $h_{\rm AP}=0$, $h_{\rm other}=\kappa$, $z=0$.
This also handles two AP supports by choosing the larger side.
Otherwise $P$ is the smaller AP side and $Q$ is non-AP.
For every $0\le k\le20$,
\begin{equation}\label{u86:eq:AP88range}
 \Omega<2T_a,\qquad\Omega-T_a\le128.
\end{equation}
At most two residue classes of $V=\FS(Q)$ occur modulo the AP step.

In one residue, at most one gap can have length at least $u=|\FS(P)|$.
A unique long gap would give two disjoint final translates by
\eqref{u86:eq:largestgap}--\eqref{u86:eq:dominantmu}, hence
$\rho_Q\ge T_{b-2}+2-2b$. It also costs $\kappa\ge u-1\ge T_a$.
The rows with $\Omega\ge T_a$ have $15\le a\le32$, $39\le b\le56$,
and
\[
 T_{b-2}+2-2b-\Omega\ge79>0.
\]
Thus no long gap exists and $z=0$.

With two residues, one off-lattice coefficient would give the same
contradiction. At least two have the nonzero residue; their sum forces
that residue to have order two. Normalize the half-step. Exactly as in
Section~\ref{u86:sec:lattice89}, if $w_0$ is the least odd coefficient, then
\[
 h_P=u-1,\qquad h_Q=\kappa-u+w_0,\qquad z=w_0-1.
\]
The mass inequality for $t\ge2$ odd coefficients is
\[
 \sum Q\ge(b-t)(b-t+1)+tw_0+t(t-1).
\]
If $w_0>u$, use $u\ge T_a+1>2b$ to minimize at $t=2$.
The lower-minus-upper support-size margin becomes
\[
 (b^2-11b)/2+2u+7-\Omega\ge1003>0.
\]
Therefore $w_0\le u$ and $h_P,h_Q\le\kappa$.
If $w_0\ge15$, minimizing over real $t$ gives
$\sum Q-w_0-T_b\ge7b-289/8$; the budget and
\eqref{u86:eq:AP88range} give the upper bound $2b+125$.
These conflict for $b\ge39$. Hence $w_0\le13$ and $z\le12$.

The AP-specific span estimates and the same one-budget substitution give
\[
 |\FS(D_1)|-2s\le2453+2a-4T_a-2T_b\le-955<0.
\]
Also $\Lambda\ge1297>4\cdot80-10+560=870$, forcing $D$ integral.
Only now $G=150$ yields $T\ge T_a+114$.
At this endpoint that inequality alone is insufficient, since
$\Omega-T_a$ reaches 128. Retain the fact that $Q$ is non-AP:
\[
 T+\rho\ge T_a+114+(b-5),\qquad
 T_a+114+(b-5)-\Omega\ge28>0.
\]
This excludes the half-step case. Every surviving AP case has $z=0$.

\subsection{The primitive large summand repair}
Assume both supports are non-progressions, and retain the notation
$\nu(r)$ from Section~\ref{u86:sec:lattice89}. The large-side inverse
condition follows from $\Omega<T_b-1$. The sufficient small-side
condition is again
\[
 \mu_P+\rho_P+\rho_Q>\Omega-T_a+2.
\]
The lower bounds $\mu_P\ge a-3$ and $\rho_X\ge\nu(|X|)$ establish it
except possibly at $k=7,\ldots,18$. In the strict cases Stanchescu
supplies the common lattice and strict support-cardinality hole bounds.
Suppose instead that this sufficient condition actually fails. Put
$r_p=\nu(a)$, $r_q=b-5$, and
\begin{align*}
 R&=\Omega-T_a+2-(a-3)-r_p,&
 U_{\max}&=\Omega-r_p-r_q+3,\\
 K_{\max}&=\Omega-r_p-r_q,&
 V_{\min}&=T_b+b-2,\qquad U_{\min}=T_a+a-2.
\end{align*}
Then $\rho_Q\le R$, $u\le U_{\max}$, $\kappa\le K_{\max}$,
$v\ge V_{\min}$, and $u\ge U_{\min}$.

First scale by the coefficient gcd of $Q$ in the initial integer
realization. The set $Q$ is now primitive and integral; $P$ has not yet
been shown integral on this scale.
The effective-hole lemma of Section~\ref{u86:sec:lattice} has the following
primitive form: for admissible parameters $\beta,C_h$, if a positive
integer $r$-set is primitive and its support has $h>2\beta$ holes, then
$\mu\ge r+C_h+1$.
Indeed, in the proof of Lemma~\ref{u96:lem:hole-parameter}, the extra hypothesis
$h<|\FS|-1$ is used to obtain coefficient primitivity. When primitivity
is given directly, a positive-hole support cannot be an AP, and the
remaining tail, counted-hole, off-scale and endpoint-bridge arguments
apply with the same parameter tests. Thus this use on the intrinsic
primitive lattice of $Q$ does not assume an unproved upper bound on
its holes.

The required explicit witnesses are
\begin{center}\small
\begin{tabular}{rrrrrrrr}
\toprule
$k$&$a$&$b$&$R$&$\beta$&$C_h$&Forced lower $\rho_Q$&$H=2\beta$\\\midrule
7&28&43&52&3&51&53&6\\
8&27&44&66&3&70&67&6\\
9&26&45&77&4&82&78&8\\
10&25&46&85&4&93&86&8\\
11&24&47&90&4&97&91&8\\
12&23&48&92&4&100&93&8\\
13&22&49&91&4&98&92&8\\
14&21&50&87&4&91&88&8\\
15&20&51&89&4&92&90&8\\
16&19&52&79&3&79&80&6\\
17&18&53&66&2&61&67&4\\
18&17&54&50&2&38&51&4\\\bottomrule
\end{tabular}
\end{center}
Every $C_h$ is checked directly against both strict inequalities in
\eqref{u86:eq:holetest}, using its explicit $K_0,D_0,M_0,J_0$ definitions.
If $h_Q>2\beta$, insert $u_0=b+C_h+1$ into the old flexible defect
inverse: for each $\gamma=2,\ldots,10$, write
\[
 \max(0,u_0-4\gamma+2)=(\gamma+1)q+s,
 \qquad0\le s\le\gamma,
\]
and use
\[
 \rho_Q\ge\min\bigl((\gamma-1)(b-14),
                    (\gamma-1)q+\min(s,\gamma-1)\bigr),
\]
together with $b-14$ and $u_0+2-2b$.
The displayed maximum lower bounds all exceed $R$.
Consequently $h_Q\le H\le8$ on its primitive lattice.

Partition $U$ into $c$ residue classes modulo that lattice.
The ordered sumset bound gives $\kappa\ge(c-1)(v-1)$.
Since $K_{\max}<V_{\min}-1$, $c=1$ and $P$ is integral there too.
Let $A_0=\operatorname{diam}U$, $B_0=\operatorname{diam}V$, and
$h=h_Q\le H$. The checked inequality
\[
 V_{\min}>U_{\max}+K_{\max}-1
\]
has minimum margin 99. If $A_0>B_0$, the two disjoint sets
$V$ and $A_0+V$ would give $|U+V|\ge2v>u+v-1+\kappa$.
Hence $A_0\le B_0$. Their two endpoint hulls cover $[0,A_0+B_0]$,
so
\[
 |U+V|\ge|V\cup(A_0+V)|\ge A_0+B_0+1-2h=A_0+v-h.
\]
Thus $h_P\le\kappa+h_Q\le\kappa+H$.
The additional verified margins
\[
 V_{\min}-K_{\max}-H-2\ge536>0,
 \qquad U_{\min}-H-2\ge162>0
\]
put both support diameters at most $u+v-3$.
Stanchescu's short-diameter Theorem 4, which has no extra gcd premise,
now yields $h_P,h_Q\le\kappa$ and
$0\le z=h_P+h_Q-\kappa\le H\le8$.

\subsection{Layer forcing and final primitivity}
All remaining sides have size at least 15. The ordinary one-budget
comparison gives
\[
 |\FS(D_1)|-2s\le C+\Omega+\rho-\mu+2z
                    \le-293-2I+2z.
\]
This is strict for $z\le146+I$, including the AP and repaired cases.
If $z\ge147+I$, both supports are non-AP, and the defect lower bounds
with $\kappa\ge z$ exclude $k\ge13$.
For $0\le k\le12$, both side sizes are at least 23.
If either side's final translates were disjoint, its size $r$ and the
other side's size $r'$ would satisfy
\[
 T_{r-2}+2-2r+(r'-5)+147+I-\Omega\ge30>0,
\]
contradicting the exact budget. Their hulls therefore overlap.
For the maximum-deleted prefix $Y$ this gives $h_X\le2h_Y$;
\eqref{u86:eq:removedmax89} then yields
$R_X\le r+\rho_X+h_X/2+2$. Substitute once to obtain
\[
 |\FS(D_1)|-2s\le C+2\rho-\mu+z+\ell_2
       \le C+\Omega+\rho-\mu\le-293-2I<0.
\]
Hence $D_1$ is integral. The already stated strict outer-layer bound
forces $D$ onto that lattice. Hole propagation and the sharp
four-extreme theorem now give \eqref{u86:eq:budget88} with $G=150$.

The strict inverse cases were already primitive. In a surviving AP
case, $z=0$ and
$\kappa\le\Omega-150<T_a\le u-1$, and similarly for $Q$.
In a repaired case, $Q$ is intrinsically primitive and $\rho_Q\ge b-5\ge38$;
with $z\le8$, \eqref{u86:eq:AP88range} gives
\[
 \kappa\le\Omega-126-\rho_Q\le T_a+128-126-38=T_a-36<u-1.
\]
A nontrivial coefficient gcd for $P$ would contradict $h_P\le\kappa$.
This proves Theorem~\ref{u86:thm:lattice88}.
The standalone \path{check88.py} verifies every finite margin and
effective-hole witness used here; it does not itself exclude size 88.

\section{The size 88 endpoint and exhaustive loss covers}\label{u86:sec:upper87}
\begin{theorem}\label{u86:thm:upper87}
There is no feasible set of 88 real numbers. Consequently,
\[
 \boxed{N\le87.}
\]
\end{theorem}
\begin{proof}
Theorem~\ref{u86:thm:upper88} already excludes every larger cardinality.
A feasible size-88 set omitting zero would become a forbidden size-89
set after adjoining zero. Therefore a hypothetical size-88 example
contains zero; integralize it and apply Theorem~\ref{u86:thm:lattice88}.
The complete classification and necessary inequalities below assign it
to a root of the canonical finite universe.
The required exact contradiction at every root excludes
it, giving $N\le87$.

\end{proof}

\subsection{Three missing units give alignment without contribution}
At size 88 a low three-empty residual cap may reach 170, above
$|2D|\ge159$. Consequently the old leading-contribution shortcut is
not enough to obtain every outer coordinate's integrality.
An additional finite algebraic theorem resolves exactly this issue.

Recall the twenty center forms
\begin{align*}
 c_p&=p_1-p_2-p_3-p_4,&
 \mathtt{p0},\ldots,\mathtt{p3}&=c_p+p_1,\ldots,c_p+p_4,\\
 \mathtt{p4},\ldots,\mathtt{p7}&=c_p-q_1,\ldots,c_p-q_4,&
 \mathtt{p8}&=p_2-p_3-p_4,\\
 \mathtt{p9}&=2p_2-p_1-p_3-p_4,
\end{align*}
and their reflected $q$ forms. The paired indices are 4 and 9 in
branch A, and 4 and 5 in branch B.
The existing classification has eighteen covering units. Reflect to
choose an actual positive contributor of the first nonempty pair;
include each legitimate selected contributor of the second pair, with
the appropriate beta variant where required.
For exactly three missing units, the complete finite enumeration has 2992
ordered modes. For each, reconstruct the matrix $C$ of actually
contributing centers in the eight outer coordinates and check a supplied
integer matrix $W$ satisfying
\[
 WC=I_8.
\]
The checker enumerates the complete ordered mode list and verifies
every product entry with integers; the largest absolute witness
coefficient is six. Since $Cx$ is integral, $x=W(Cx)$ is integral.
The earlier analytic theorem covers at most two missing units.
Thus all eight outer coordinates are integral when $j\le3$.
This theorem supplies alignment, not leading contribution: the
leading loss inequalities still require their own strict cap gates.
For $j\ge4$, missing outer coordinates require a separate root-specific
witness before strict outer gaps or fixed-score integer branches are used.

\subsection{The canonical low and high families}
Put $b=71-a$, $I=T_a+T_b-T_{35}-T_{36}$, and
$B_0=410-I-80j$. The low zero-hole families retain $B_0\ge0$ for
$j=0,\ldots,5$. Positive low-hole families use $1\le z\le14$ and
$B_0-15\ge0$ for $j=0,\ldots,4$.
The side range is $15\le a\le56$ and the simultaneous rows include
\begin{gather*}
 \delta\le B_0\quad(Z),\qquad
 \delta\le B_0-43+2z\quad(P),\\
 R_p+R_q+\delta\le75+B_0+4z,\qquad
 S_p+S_q+\delta-2K\le1524-80j+4z.
\end{gather*}
The positive four-empty wrapper again starts from a more permissive
zero-hole template before installing its correct positive-hole rows.
The complete low counts are
\begin{center}\small
\begin{tabular}{lrrrrrr}
\toprule
$j$&0&1&2&3&4&5\\\midrule
$Z_j$&160&2520&18496&77792&217600&177072\\
$P_j$&160&2520&17340&71808&195840&0\\\bottomrule
\end{tabular}
\end{center}
They total 781308. The complete mode counts for $j=0,\ldots,5$ are
$4,70,578,2992,10880,29512$.

For high holes, $(a,b)=(35-k,36+k)$ up to reflection, and
\[
 C_{88}(a,z)=560-k(k+1)-\rho_*(a,u_a)-\rho_*(b,u_b)
                                  -\max(50,z,150-2z).
\]
The domain $0\le k\le20$, $15\le z\le560-k(k+1)$ has 8386 integer
states. All 630 effective-hole constants are checked against their
strict hypotheses. The maximum cover cap is 333, so $j\le4$.
The all-active mass upper bound is $1566+4z$.
The retained state-and-empty-count pairs number
$(3882,2881,1737,780,47)$, totaling 9327. Their 188 intervals have
counts $(51,60,45,26,6)$ for $j=0,\ldots,4$.
Both reflected orientations are distinct because $m=71$ is odd.
The high-root counts are $(408,8400,52020,155584,130560)$,
totaling 346972. Hence the canonical universe has
\[
 781308+346972=1128280\quad\hbox{original roots}.
\]

Leading, nested and core-quantile gates are respectively
$\delta<159$, $\delta<151$ and $\delta<76$, interpreted through the
relevant uniform residual cap. Ordinary $j\ge3$ models retain weak
outer order. For $j\ge4$, low caps are at most 90 and 10 for $j=4,5$;
the high four-empty cap is at most 13. Thus both leading components
contribute, and retained centers supply the relative alignment used
in full-$D$ geometry. They need not put missing outer coordinates on
the lattice. The finite gate check covers 2736 fast/full parameter
representatives; every certified current row is regenerated in the
complete arithmetic replay.

\subsection{Raw star overlap and a simultaneous denser core}
Let $A_0$ be a leading component and $A_1,\ldots,A_t$ distinct labeled
retained components in its group. Pointwise incidence counting gives
\begin{equation}\label{u86:eq:starraw88}
 \sum_{i=1}^t|A_0\cap A_i|
       \le\sum_{i=0}^t|A_i|-\left|\bigcup_{i=0}^t A_i\right|.
\end{equation}
Indeed, a point in the hub and $r$ leaves contributes $r$ on both
sides; a point outside the hub contributes zero on the left and a
nonnegative amount on the right.
For an actual common target $F$, define
$G_F=\sum_{i=0}^t|A_i|-|(\bigcup_{i=0}^tA_i)\cap F|$.
Its group deficit is therefore at least
\begin{equation}\label{u86:eq:startrim88}
 G_F\ge\sum_{i=1}^t|A_0\cap A_i|+|A_j\setminus F|
 \qquad(0\le j\le t).
\end{equation}
Outside $F$, adding one chosen member incidence to the hub-leaf count
never exceeds total component multiplicity. Endpoint-hull losses may
replace actual target losses as lower bounds.
Each choice of $j$ gives an alternative lower bound for the same
actual deficit; these alternatives are not added together.

Lemma~\ref{u86:lem:translatedcore89} and its cross-component reflection
supply
\[
 H_{\rm own}=2e-1-R_p-R_q-p_i+M_p,
 \qquad H_{\rm cross}=2e-1-R_p-R_q-q_j+M_q.
\]
The adopted size-88 full model sums at most one own and one cross raw
bound per row, retains one member loss, and adds unmatched component
losses. Its cheap scalar budget rows choose any nonempty subset of at
most four distinct overlap clusters across the two target signs and
drop other nonnegative losses. Negative $H$ values remain valid lower
bounds. The relevant alignment follows from the three-empty theorem,
or from leading and retained contribution when $j\ge4$.

Only the eight zero-hole loss-cover parents below need a denser nested
core and all retained leaves at once. Write the positive peeled points
as $x_1<x_2<x_3<x_4=M_p$, with $x_0=R_p$, and similarly
$y_0=R_q<y_1<\cdots<y_4=M_q$. Define
\[
 g_p=\max\bigl(0,\max_{1\le j\le4}(2j+R_p-x_j)\bigr),
 \quad
 g_q=\max\bigl(0,\max_{1\le k\le4}(2k+R_q-y_k)\bigr).
\]
Integral gaps give $0\le g_p,g_q\le4$.
Choose maximizing indices, allowing zero. Adjoining exactly those first
$j$ positive and first $k$ negative peeled points gives one simultaneous
set $E'\subseteq D$ with
\[
 2|E'|-\operatorname{diam}E'=2e-R_p-R_q+g_p+g_q.
\]
Using $E'$ in the two translated copies improves every own or cross
raw overlap bound to $H_i+g_p+g_q$ simultaneously.
Nonnegative overlap variables may be realized by the actual raw
intersections. Formula~\eqref{u86:eq:startrim88} permits their sum plus one
leading or leaf loss. No disjointness of the individual intersections
is assumed. The separately charged fixed-label losses remain in the
overall covering budget, not in these raw star rows.

\subsection{Retained-unit consequences and the fourteen-unit loss cover}
If a retained unpaired positive component $v+D$ contributes to the
actual target inside $[1,S_p]$, then $v$ is integral and
\[
 v+M_p\ge1,\qquad v-M_q\le S_p,\qquad
 \lambda_{\rm low}(v)+\lambda_{\rm high}(v)\le79.
\]
The two hull-loss sets are disjoint, and at least one of the 80 points
remains in the target hull. If a retained cross component contributes,
$c_p-q_j+M_p\ge1$ and $p_i>M_p$, so every retained own center is
positive and its low-normal quantile holds. Conversely, a contributing
own component gives $c_p+p_i-M_q\le S_p$; since $q_j>M_q$, every
retained cross center is below $S_p$, giving its high-normal quantile.
Reflection supplies the negative side. These statements apply only to
actual unpaired contributors, not arbitrarily chosen members of a pair.

The fixed-score argument of Section~\ref{u86:sec:upper90} supplies the same
mandatory saving $F=8$ for the eight zero-hole four-empty parents under
consideration. It uses their checked outer lattice and missing-largest
slots with retained lower contributors. In the common covering budget,
fixed labels are disjoint from component losses; hence the uniform
remaining cap $C$ bounds their sum. The usual quantiles may be recomputed
from that cap: global constants are 80/79, core constants 76/75 require
$C<76$, and normal quantiles require $C<\min(a,b)+5$.
This does not change any original lattice provenance.

Both paired units are missing in these eight roots, leaving fourteen
actual unpaired contributors. Let $L_i$ be each component's two-endpoint
hull-loss sum. Then $L_i$ are nonnegative integers and
\[
 \sum_{i=1}^{14}L_i\le C<152.
\]
If $C<76$, the all-small case is exhaustive. Otherwise use all fifteen
children: the all-small child $L_i\le75$ for every $i$, and one child
$L_i\ge76$ for each possible selected $i$.
At most one component can be large. The all-small child receives the
existing core quantiles for every component. In a selected-large child,
every other loss is at most $C-76<\min(a,b)+5$, so those thirteen
components receive normal quantiles. No unsupported new normal row is
imposed on the selected large component.

The two end parents have $C=70$ and one child each; the other six have
$C=76,80$ or 82 and fifteen children each. All 92 children are required.
They use the already proved dense-star, fixed-aware and retained-unit
rows, with fixed-score binary splits licensed by each root's original
$WC=I_8$ certificate. This finite cover is separate from the binary
branch count and cannot be replaced by one selected large-loss case.

\subsection{Positive holes and the three interval refinement}
For positive holes, $h_P,h_Q\ge z\ge1$. Individual primitivity implies
both supports are non-AP: a primitive AP support has step one and no
holes. If both side sizes are at least 21, the non-AP theorem
gives
\[
 \rho_P+\rho_Q\ge(a-5)+(b-5)=61.
\]
The older low-positive template used 43. The exact budget therefore
permits the additional row
\[
 \delta-2z\le\Omega-150-80j-61,
\]
improving its right side by exactly 18. The original labeled row remains
unchanged; new quantile and leading rows are appended only under their
new valid cap gates. A smaller cap does not retroactively alter the
center polynomials used in an existing lattice certificate.
This strengthening directly excludes 106 parents and supports 26
three-empty branch proofs. The latter first check the original-cap
$WC=I_8$ witness, then apply the positive-defect row. They use no
fixed-score layer.

Fourteen positive four-empty parents remain, with $29\le a\le42$,
branch A, missing units \texttt{pair1}, \texttt{pair2}, \texttt{p0},
\texttt{q0}, and contributor choice zero. Partition their integer holes
into exactly three closed intervals
\[
 [1,2]_{\Z},\qquad[3,4]_{\Z},\qquad[5,14]_{\Z}.
\]
Every parent requires all three children.
In a child with lower endpoint $l$, use each validated effective-hole
witness with $l>2\beta$ to obtain $\mu_X\ge r+C_h+1$.
At $l=1$ use only $\mu_X\ge r-3$; at $l=3$ use $\beta=1$;
at $l=5$ use $\beta=1,2$.
The usual strict hole hypothesis $h_X<|\FS(X)|-1$ is explicitly
available here: strict inverse cases already have it; in the repaired
cases $h_Q\le8$ and the post-lattice estimate gives
$h_P\le\kappa<T_a$; AP cases have $z=0$ and are absent.
Primitivity alone is not substituted for that hypothesis in this
small-hole application.

The 28 required witnesses are listed below. Here $C_1,C_2$ correspond
to $\beta=1,2$, and $u_r(l)$ is the resulting excess lower bound at
hole lower endpoint $l$.
\begin{center}\small
\begin{tabular}{rrrrrr}
\toprule
$r$&$C_1$&$C_2$&$u_r(1)$&$u_r(3)$&$u_r(5)$\\\midrule
29&13&28&26&43&58\\
30&14&30&27&45&61\\
31&15&31&28&47&63\\
32&16&32&29&49&65\\
33&17&33&30&51&67\\
34&18&35&31&53&70\\
35&18&37&32&54&73\\
36&19&39&33&56&76\\
37&19&41&34&57&79\\
38&20&41&35&59&80\\
39&21&42&36&61&82\\
40&22&44&37&63&85\\
41&23&46&38&65&88\\
42&24&48&39&67&91\\
\bottomrule
\end{tabular}
\end{center}
All strict parameter tests are explicit finite arithmetic obligations. Replacing the old per-side excess $r-3$ by $u_r$
changes the $S_s$ mass row by $-2(u_r-r+3)$ and the core mass row
by $-(u_r-r+3)$, preserving their hole coefficients. Recompute endpoint
quantiles on the same restricted child with the single 18-unit defect
improvement; then restore the old labeled cap and hole rows.
The $[3,4]$ and $[5,14]$ children are 28 direct raw-model contradictions.
The fourteen $[1,2]$ children use fixed-score trees with their original
outer-lattice witnesses. Their positive holes mean the mandatory
missing-largest saving is zero and all 110 high-label candidates remain.
No stronger non-AP inverse, seven-interval refinement, conditional-span
improvement or experimental core-mass table is used.

\subsection{The full exact partition and verification standard}
The ordinary scanner records an exact current-row certificate whenever
it certifies a root. Its terminal ledger contains 1127649 certified
roots and 631 explicit gaps, with no unattempted root. A numerical
proposal without a successful integer certificate remains a gap.
The complete gap list has 242 low and 389 high parents, partitioned as
follows:
\begin{center}\small
\begin{tabular}{lrrrrr}
\toprule
Supplemental class&Parents&Child models&Nodes&Binary splits&Leaves\\\midrule
Zero-hole three-empty&20&20&68&24&44\\
Zero-hole four/five-empty&68&68&548&240&308\\
Zero-hole loss cover&8&92&5112&2510&2602\\
Direct positive defect&106&106&106&0&106\\
Positive three-empty&26&26&144&59&85\\
Positive three-interval cover&14&42&876&417&459\\
High geometry first group&41&41&223&91&132\\
High geometry second group&170&170&922&376&546\\
High geometry terminal group&178&178&1448&635&813\\\midrule
Total&631&743&9447&4352&5095\\\bottomrule
\end{tabular}
\end{center}
The three high packages use only the adopted geometry and checked
outer-integrality wrapper. Their selected splits comprise 963 normal,
70 core-or-large-loss and 69 inner-quantile alternatives.
They use no fixed-score or additional inequality layer.
All selected contributor variants and original-cap gates remain linked
to their exact canonical descriptors. The last 178 IDs are checked to
be exactly the terminal 389 high gaps minus the previous 211.

The low union is checked against the 242 actual ordinary-low gaps,
not inferred from its count. Its 354 child models give 3604 leaves.
The high union has 389 child models and 1491 leaves.
The finite covers add 112 children beyond the original 631 parents;
therefore the exact accounting identities are
\begin{align*}
 743&=631+112,\\
 5095&=743+4352,\\
 1132744&=1127649+5095=1128280+4352+112.
\end{align*}
Counts alone do not prove infeasibility. Every ordinary and supplemental
leaf is checked on its regenerated current matrix by positive integer
weights $w$ with $w^{\mathsf T}A\ge0$ and $w^{\mathsf T}b<0$ on
$x\ge0$. Every required cover child and both children of every binary
split must be present and exactly linked to its original parent.




The portable ordinary certificate carries 10689 positive-integer labeled
Farkas vectors and exactly 1128280 unsigned 16-bit little-endian assignment
entries, totaling 2256560 bytes. Zero marks a supplemental parent. A
nonzero code is $2i+1+\varepsilon$, where $i$ is the zero-based vector index
and $\varepsilon=0,1$ selects the fast or full necessary model.
The decoder checks exact stream length, integer encoding, pool range,
vector identities and the 631 zero positions. All canonical descriptors
are regenerated in order and their complete manifest hash checked.
For every nonzero assignment, the verifier rebuilds the selected current
model and proves the exact Farkas contradiction. It does not need a stored
per-root matrix or the raw discovery journal to prove that contradiction.
The fast/full counts are 1086923 and 40726, summing to 1127649.




Supplemental verifiers reconstruct the full canonical manifest, actual
center polynomials, integer left inverses, strengthened rows and branch
formulas. Pure adapters are checked against frozen mathematical source
functions, and numerical optimizers are excluded from replay.
The integration verifies exact disjoint equality between the
supplemental parent IDs and the 631 compact zero positions, as well as
all source and proof identities. The earlier size-89 exclusion remains
a separate descending prerequisite.


The finite lemma requires exact equality between the supplemental
parent-ID set and all 631 zero positions in the ordinary assignment stream,
with disjoint supplemental parents and every required cover child.
The full entry point \path{proofs/exact87/verify_exact87.py} must check
structural and classification arithmetic, all nine supplemental families,
their exact complement, and every ordinary current-matrix contradiction.
A supplements-only check would leave the ordinary finite lemma unproved.


\section{The common core lattice at size 87}\label{u86:sec:lattice87}

This is a necessary-structure theorem, separate from the endpoint
infeasibility lemma. The last nonprimitive AP branch is discharged by the fully printed
analytic combinations in Appendix~\ref{u86ap:sec:certificate}.



\subsection{Bootstrap and the order of normalization}
The deficient-sign case has
$|\FS(E)|\ge1+T_{71}=2557>24\cdot87+77=2165$, so $d=79$, $e=71$ and
$m=70$. The inherited unpaired-element bound is $q^2\le2848$.
Its even parity gives $q\le52$, hence at least nine opposite pairs in
$E$ and at least nine elements on each sign. In particular zero is
not a diagonal defect. Reflect so that $a=35-k\le b=35+k$. Then
\[
 \Lambda=1261+k^2,\qquad \Omega=574-k^2,\qquad C=-823-k^2.
\]
The exact pre-lattice budget is
\begin{equation}\label{u86:eq:87prebudget}
 \rho_P+\rho_Q+\kappa+2\theta+c+\ell_1+\ell_2+g=\Omega.
\end{equation}
All terms other than a possible small-side $\rho$ are nonnegative.
Since $b\ge35$, $\rho_Q\ge0$; the global bound $\rho_P\ge-8$ gives
$k\le24$, so $a\ge11$.

Write $U=\FS(P)$, $V=\FS(Q)$, $u=1+T_a+\mu_P$ and
$v=1+T_b+\mu_Q$. Subset-sum reflection identifies $-V$ with a translate
of $V$, so $\kappa=|U+V|-u-v+1$. Once a common lattice is proved, put
$h_P=\sum P-T_a-\mu_P$, $h_Q=\sum Q-T_b-\mu_Q$, and
$z=h_P+h_Q-\kappa$. The inverse conclusion $h_P,h_Q\le\kappa$ then
gives $0\le z\le\min(h_P,h_Q)$.

For a positive side $X$ of size $r$, with maximum $R_X$, let
$Y=X\setminus\{R_X\}$ and let $\nu$ count its old diagonal defects
destroyed by adjoining $R_X$. The exact identity is
\begin{equation}\label{u86:eq:87delete}
 R_X=r+\rho_X+h_X+2-\rho_Y-h_Y-2\nu.
\end{equation}
Thus $R_X\le r+\rho_X+h_X+2$ for $r\ge15$. At the smaller boundaries
one must retain $\rho_Y\ge-8$, which adds eight to that side's bound
and sixteen to the first peeled-layer comparison.

Only after $D$ is integral and the requisite extreme translates
intersect may we use the sharp four-extreme theorem with $G=148$.
Writing $\mathcal T=\kappa+\ell_1+\ell_2+g$, its consequences are
\begin{equation}\label{u86:eq:87postbudget}
 \mathcal T\ge\kappa+148-3z,\quad
 \mathcal T\ge z,\quad \mathcal T\ge50,\quad
 c+\rho_P+\rho_Q+\mathcal T\le\Omega.
\end{equation}
The proof below respects this order; these post-lattice inequalities
are not used to obtain the initial common lattice.

\subsection{The four small-side boundaries}
The relevant $(a,b,\Omega,\kappa\text{ cap})$ are
\[
 (11,59,-2,6),\ (12,58,45,53),\ (13,57,90,98),\ (14,56,133,133).
\]
For $a=11,12$ the cap is less than both $T_a-1$ and $T_b-1$, so the
strict simultaneous inverse applies. For $a=13,14$, a non-AP larger
support has $\rho_Q\ge b-5$ and reduces the caps to $46,82$; the same
inverse again applies. If instead $V$ is an AP, its length exceeds the
raw cap plus one. Residue counting aligns $P$ with its step, and the
AP gap formula gives $h_P=\kappa$, $h_Q=z=0$.

Using the extra sixteen required by the deleted small prefix, the
first-layer comparison is bounded by
\[
 -249-2k^2+2\,\kappa_{\max}+16,
\]
whose four values are $-1373,-1185,-1005,-849$. The other layer is
strict from $\Lambda>4d-10+\kappa_{\max}$. These comparisons prove
$D$-integrality and overlap before \eqref{u86:eq:87postbudget} is used.
At $a=11$, $\mathcal T\ge50>6$ is impossible. At $a=12,13$, the
non-AP case gives $\mathcal T+\rho_P+\rho_Q\ge50-8+(b-5)>\Omega$;
the AP case has $\mathcal T\ge148$ and is also impossible.
For $a=14$ the AP case is excluded identically. Otherwise
$\rho_Q\ge51$ and $\rho_P\ge0$ imply $z\ge33$. The strict inverse
hole bound permits the effective-hole witness $(r,\beta,C_h)=(56,5,143)$.
It gives $\mu_Q\ge200$, and the old flexible inverse gives
$\rho_Q\ge134>133$. Hence all surviving sides have $a,b\ge15$.

\subsection{AP supports and the half-step case}
If the larger support is an AP, it has at least $T_{35}+1=631$ points,
whereas $\kappa\le574$. Residue counting and the exact AP gap formula
give a common lattice, $h_{\rm AP}=0$, $h_{\rm other}=\kappa$ and $z=0$.
Now suppose that only the smaller support $U$ is an AP. For
$0\le k\le20$, $\Omega<2T_a$, so the residue count permits at most
two residues of $V$ modulo its step.

In one residue, the AP gap formula is
$\kappa=\sum_i(\min(u,g_i)-1)$. A gap of length at least $u$ would be
unique and central by reflection. The largest-gap argument then forces
the final coefficient of $Q$ to exceed the sum of its predecessors.
The two final subset-sum translates are disjoint, giving
$\mu_Q\ge T_{b-2}$ and $\rho_Q\ge T_{b-2}+2-2b$.
The same disjoint-copy bound holds if exactly one coefficient is
off-lattice. In either case the essential extra cost is
$\kappa\ge u-1\ge T_a$. Over the relevant range $k=2,\ldots,20$,
\[
 T_a+T_{b-2}+2-2b-\Omega\ge549>0.
\]
Thus neither exception occurs. The one-residue case has only short
gaps and gives $h_P=0$, $h_Q=\kappa$, $z=0$.

If two or more coefficients occupy the nonzero residue, their sum
forces that residue to have order two. Normalize
$U=2[0,u-1]_{\Z}$ and let $w$ be the least odd coefficient of $Q$.
The two parity-fiber gap formulas, and reflection of their endpoints,
give exactly
\[
 h_P=u-1,\qquad h_Q=\kappa-u+w,\qquad z=w-1.
\]
These formulas alone do not yet give $h_Q\le\kappa$.
For $s\ge2$ odd coefficients, distinctness yields
\[
 \sum Q\ge(b-s)(b-s+1)+sw+s(s-1).
\]
If $w>u$, the quadratic is minimized at $s=2$ in this range;
comparison with the budget upper bound gives the positive margin
$(b^2-11b)/2+2u+7-\Omega\ge943$. Hence $w\le u$, establishing both
exact hole bounds. If $w\ge15$, minimize the same quadratic over real
$s$ to obtain
\[
 \sum Q-w-T_b\ge7b-289/8.
\]
But the exact mass identity and budget bound its left side by
$2b+\Omega-T_a-3$; the minimum contradiction margin is $343/8$.
Thus $w\le13$ and $z\le12$.

The AP-specific estimates
\[
 R_P\le2(a+\mu_P),\qquad
 R_Q\le b+\mu_Q+\kappa-T_a-\mu_P+12
\]
make the first peeled-layer comparison at most
$2423+2a-4T_a-2T_b\le-891$. The other layer is strict because
$1261>4\cdot79-10+574=880$. Therefore $D$ is integral before $G$ is
used. The post-lattice lower bound
$T_a+148-36+(b-5)-\Omega$ is positive except for equality at $k=11,12$.
If $z\le11$, it increases by three. At $z=12$, the effective-hole
witness $(\beta,C_h)=(2,33)$ gives $\rho_Q\ge42,43$ for $b=46,47$,
respectively, exceeding $b-5$ by one. This excludes the equalities too.
Every surviving AP branch therefore has $z=0$.

\subsection{Both supports non-AP and the inverse gap}
Put
\[
 n(r)=\begin{cases}r-5,&r\ge21,\\r-14,&15\le r\le20.\end{cases}
\]
The large-side strict inverse threshold is automatic. The sufficient
small-side test
\[
 \mu_P+\rho_P+\rho_Q>\Omega-T_a+2
\]
implies the strict simultaneous inverse bounds. Substitution of the
valid lower floors leaves possible failures only at $k=6,\ldots,19$.
If the test fails, define
\begin{align*}
 R_*&=\Omega-T_a+2-(a-3)-n(a),&\rho_Q&\le R_*,\\
 U_{\max}&=\Omega-n(a)-n(b)+3,&u&\le U_{\max},\\
 K_{\max}&=\Omega-n(a)-n(b),&\kappa&\le K_{\max},\\
 V_{\min}&=T_b+b-2,&v&\ge V_{\min},\qquad u\ge U_{\min}=T_a+a-2.
\end{align*}
For $k=7,\ldots,19$, work first on the intrinsic primitive lattice of
$Q$. The already proved primitive extension of the effective-hole
lemma, followed by the old flexible inverse, gives $h_Q\le H$ using
the following complete list. The ``forced'' column is the contradictory
old inverse floor if $h_Q>H$.
\begin{center}\small
\begin{tabular}{rrrrrr}
\toprule
$k$&$R_*$&$\beta$&$C_h$&Forced $\rho_Q$&$H$\\\midrule
7&73&4&88&79&8\\8&88&5&108&96&10\\9&100&6&132&115&12\\
10&109&6&133&117&12\\11&115&6&137&120&12\\12&118&6&143&127&12\\
13&118&6&149&132&12\\14&115&5&132&120&10\\15&118&5&136&123&10\\
16&109&5&139&127&10\\17&97&4&118&111&8\\18&82&3&95&93&6\\
19&64&2&66&72&4\\\bottomrule
\end{tabular}
\end{center}
Each parameter pair satisfies the strict effective-hole tests. For
these rows the four margins
\[
 V_{\min}-1-K_{\max},\quad V_{\min}-U_{\max}-K_{\max}+1,
\]
\[
 V_{\min}-K_{\max}-H-2,\qquad U_{\min}-H-2
\]
have minima $477,11,468,144$, respectively. The first aligns $P$ by
residue counting. The second forces $\operatorname{diam}U\le
\operatorname{diam}V$: otherwise two disjoint endpoint translates of
$V$ would exceed $u+v-1+\kappa$. The two endpoint hulls then give
$h_P\le\kappa+h_Q\le\kappa+H$. The last two margins put both
diameters at most $u+v-3$. Stanchescu's short-diameter theorem therefore
recovers the exact $h_P,h_Q\le\kappa$, with $z\le H\le12$.
No primitivity of $P$ has been assumed here.

At $k=6$ the diameter-order margin is negative, so a separate bridge
is required. Here $(a,b,\Omega,R_*)=(29,41,538,55)$. If every excess
increment $y_i$ for $15\le i\le41$ were at least five, the recurrence
would give $\rho_Q\ge81$. Thus the last $y_j\le4$ has
$j\ge41-\lfloor55/3\rfloor=23$. On its primitive scale, the
counted-hole lemma puts that prefix's support holes in
$H_0\cup(L-H_0)$, with $H_0\subset[1,8]_{\Z}$ and $|H_0|\le4$.
A first later off-scale coefficient would force
$\rho\ge T_{23}-41-1=234>55$. Later increments cannot undo this.
Every subsequent $y_i\le57$; for a current prefix of size $t\ge23$,
a new coefficient $x\ge L-16$ would add at least
$L-24\ge T_t-24>t+58$ support points, too many. Hence $x\le L-17$.
The central intervals join, and $x>8$ preserves the lower holes;
reflection preserves the upper holes. The final $V$ therefore contains
$[9,\sum Q-9]_{\Z}$ and has at most eight endpoint holes.
Residue counting aligns $P$, and this central interval gives
$h_P\le\kappa+18$. The checked inequalities
$\kappa+18<v-19$, $\kappa+20\le v$ and $u\ge10$ imply both
short-diameter hypotheses, so Stanchescu again gives the exact hole
bounds, now with $z\le8$.

\subsection{Forcing both layers and proving individual primitivity}
At this point every branch has a common core lattice and exact hole
bounds, and both maximum-deleted prefixes have nonnegative $\rho$.
The first-layer comparison is
\[
 |\FS(D_1)|-2s\le C+\Omega+\rho-\mu+2z
 \le-249-2k^2+2z.
\]
It is strict for $z\le124+k^2$, including all AP and dense-repair
branches. Otherwise $z\ge Z=125+k^2$. The pre-lattice budget excludes
$k\ge14$; this case retains the original strict inverse hole bounds.
If the last two translates for a side of size $r$ were disjoint, then
$\rho_X\ge T_{r-2}+2-2r$. Add $\kappa\ge Z$ and the other side's
non-AP floor. The elementary floor suffices except for a potentially
smaller disjoint side at $k=4,\ldots,8$. For these five cases the strict
hole hypotheses justify the effective-hole witnesses
\[
 (k,\beta,C_h)=(4,2,41),(5,3,53),(6,3,58),(7,3,54),(8,2,45),
\]
which give other-side floors $43,52,56,55,49$. Both side choices and
all $k=0,\ldots,13$ now have positive contradiction margin, with
minimum one at $k=4$ and disjoint-side size $31$. Thus the actual
supports intersect, not merely their hulls.

The maximum-deleted support hulls consequently cover the full side
hull, so $h_X\le2h_Y$. Equation~\eqref{u86:eq:87delete} improves to
$R_X\le r+\rho_X+h_X/2+2$. Substitution gives
\[
 |\FS(D_1)|-2s\le C+2\rho-\mu+z+\ell_2
 \le C+\Omega+\rho-\mu\le-249-2k^2<0.
\]
Here $z\le\kappa$ and $\rho+\kappa+\ell_2\le\Omega$ are used once.
An off-lattice peeled element would give two disjoint $s$-point
translates, contrary to this strict bound. The outer layer is likewise
strict from $1261>880$. The same comparisons force the extreme
translate intersections needed for propagation. Thus $D$ is integral
and $G=148$ is now available without circularity.

The original strict inverse branches are individually primitive:
a nonunit gcd would give at least $|\FS(X)|-1$ holes. In the dense
repairs $Q$ was primitive by construction; \eqref{u86:eq:87postbudget}
and both non-AP floors give
\[
 \kappa\le\Omega-(148-3H)-n(a)-n(b)<T_a\le u-1.
\]
The strict margins are at least eighteen in the table and eighty-one
at $k=6$, so $P$ is primitive too. In an AP branch whose smaller side
defines the step, the other side follows from $\kappa\le\Omega<T_b$.
When the larger AP side defines the step, a separate exception is
needed, as the next paragraph proves.

In the remaining branch the larger AP side defines the unit lattice.
All common-lattice, $D$-integrality, overlap and propagation premises have
already been proved, so Corollary~\ref{u86ap:cor:primitive} applies.
Its proof derives the integer cap $c\le6$ from any hypothetical
nonprimitive smaller side, then contradicts that cap with the four
analytic covering combinations of Lemma~\ref{u86ap:lem:small-slack}.
No primitive-core premise, finite side-size enumeration, or leading
contribution is used. This proves individual primitivity in the last branch.

We have proved the size-87 structural theorem: $a,b\ge15$, common
integral $D$, both primitive cores, exact core-hole bounds and the
required overlap and propagation. It asserts neither general outer-$B$
integrality nor contribution of a leading $2D$ component.

\section{The size 87 endpoint and selected actual-set refinements}\label{u86:sec:upper86}


\subsection{Canonical roots and strict contribution gates}
Use $I=(a-35)^2$, $\Omega=574-I$, $G=148$, and
$\delta=c-79j$, where $j$ counts missing ordinary cover units. The
canonical descriptors retain the previous covering branch, contributor
variants and missing-unit choices. Every matrix variable is relaxed to a nonnegative real. An integer
disjunction additionally requires a proof that its actual-set quantity
is integral; no linear-programming solution is assumed to be realized
by a set.

For $z=0$, put $B_0=426-I-79j$; then $\delta\le B_0$.
For $1\le z\le14$, both primitive core supports are non-AP because
$h_P,h_Q\ge z>0$. Thus
\[
 \rho_P+\rho_Q\ge n(a)+n(b),\qquad
 \delta-2z\le B_0-n(a)-n(b).
\]
The floor sum is $60$ when both sides have at least 21 elements and
$51$ otherwise. The uniform cap is $B_0-n(a)-n(b)+28$.
The scalar budget excludes zero-hole $j\ge6$ and positive-hole
$j\ge5$. In particular the balanced positive $j=5$ cap is $-1$;
the smaller-side range only reduces it further.

If $u_s$ is a proved lower bound on the side excess, the main positive
mass row is
\begin{equation}\label{u86:eq:87mainmass}
 S_s\ge T_r+6+u_s+z+M_s+3R_s.
\end{equation}
Here $S_s$ is the sum of the positive magnitudes on that side of $D$,
$M_s$ its outermost member, and $R_s$ its core maximum. The actual
core mass is
$C_s=S_s-M_s-x_{s,1}-x_{s,2}-x_{s,3}$, with the three inner peeled
magnitudes written in increasing order. Hence
$C_s\ge T_r+u_s+z$, and $x_{s,i}\ge R_s+i$ proves
\eqref{u86:eq:87mainmass}. This explains the constant six and the factor
three; doubling the core inequality alone would not give the main row.
The conservative span row is
$R_P+R_Q+\delta-4z\le74+B_0$. No conditional improved span is used.

For $z\ge15$, use the effective-hole floor
$u_s=u_{\rm old}(r_s,z)$ and
$v_s=\max(n(r_s),v_{\rm old}(r_s,u_s))$. The complete scalar cap is
\[
 c\le\Omega-v_P-v_Q-\max(50,z,148-2z).
\]
All $0\le k\le20$ and $15\le z\le\Omega$ are enumerated. The
established strict mass test omits only proved impossible $j=0$
states; it omits no $j\ge1$ state. Consecutive integers are merged
only when all active formulas and excess floors agree. The required finite
reconstruction has 8890 raw $(k,z)$ states, 11374 retained
$(k,z,j)$ states and 277 intervals, with maximum $c=354$. Thus
$j=0,\ldots,4$ is complete. The balanced high orientation is emitted
once. All 605 effective-hole constants are checked by strict integer
inequalities, without requiring optimality of the search that found them.

The full canonical counts are
\begin{center}\small
\begin{tabular}{lrrrrrrr}
\toprule
Family&$j=0$&$j=1$&$j=2$&$j=3$&$j=4$&$j=5$&Total\\\midrule
Zero holes&164&2590&19074&80784&228480&324632&655724\\
Positive low&164&2590&17918&74800&184960&0&280432\\
High holes&604&10780&69360&239360&359040&0&679144\\\midrule
All roots&&&&&&&1615300\\\bottomrule
\end{tabular}
\end{center}
Both low families together have 936156 roots. Their descriptors,
including complete high interval data, are compared individually;
agreement of aggregate counts alone is insufficient.

The final residual cap is recomputed after every budget change.
Leading endpoint rows require cap $<157$; nested leading rows require
cap $<149$; core quantiles require cap $<75$; normal quantiles require
cap $<\min(a,b)+5$. Equality does not license the stronger row.
Balanced positive $j=3$ has cap $157$ and zero-hole $j=3$ can have cap
$189$, so neither is granted a leading contributor by this test.
The retained-center theorem gives alignment instead.
For $j\ge4$ the maximal caps are $110,31,78,38$ for zero $j=4$,
zero $j=5$, positive $j=4$, high $j=4$, respectively. These permit
leading contribution but still do not make every missing outer
coordinate integral.

A supplemental geometry wrapper may restore all strict outer gaps
only after reconstructing genuine contributing centers and checking
an integer matrix identity $WC=I_8$. A center with an integral target
point and an integral $D$ point is integral; $W$ then proves the
outer-coordinate assertion. The original cap governs whether leading
centers may enter $C$. A later scalar strengthening cannot
retroactively add a leading center to this witness.

The exact cardinality identity, common to every hole regime, is
\begin{equation}\label{u86:eq:87converse}
 \delta+S_P+S_Q-2K=(18-j)79+84+w,\qquad0\le w\le4z,
\end{equation}
where $K=|D+D|$. Its upper and lower inequalities are both retained.
The lower one is an identity restriction, not a cover-loss budget:
fixed-label or overlap losses must never be injected into it.

\subsection{Zero-hole and small-positive-hole finite covers}
In the selected zero-hole loss cover, exactly
\texttt{pair1,pair2,p0,q0} are missing. The full outer integer witness
and a retained smaller own contributor on each side justify four extra
high fixed-label deletions per side. These eight losses are disjoint
from the 26 baseline nonpositive groups. Fourteen retained unpaired
ordinary units remain, seven on each sign. Their actual integral hull
losses satisfy
\[
 \sum_u L_u\le\delta-8\le C_0:=102-I<150.
\]
At most one has loss at least 75. Thus either all have loss at most 74,
licensing the $75/74$ core quantiles, or one specified unit has loss at
least 75 and every other has loss at most $C_0-75$, below the normal
quantile threshold. Every possible named-large alternative is retained;
when $C_0<75$ only the all-small alternative is needed. The eleven
selected parents require all 165 child models, not eleven selected
successful children. The converse row \eqref{u86:eq:87converse} remains
unchanged.

For the eleven remaining positive low parents, the full integer domain
$1\le z\le14$ is partitioned into $[1,2]$, $[3,4]$, $[5,14]$.
All three intervals are required for each parent. Here $j=4$,
$a=30,\ldots,40$ in the selected component, and the original budget is
$\delta-2z\le50-I$. No extra eighteen-unit saving is subtracted.
The broader applicability range $27\le a\le43$ gives
$\kappa\le92-I<T_{27}$, proving the strict effective-hole premise.
The old floor at each interval's lower endpoint raises $u_s$ in the
main and core mass rows, with multipliers two and one in their
respective inequality encodings. The three residual caps are
$54-I,58-I,78-I$, with the same strict quantile gates.

Positive-hole fixed saving is zero and all 110 high-label candidates
remain. Seventeen interval children are direct leaves. Sixteen use
checked outer witnesses and complete fixed-score trees; five of these,
only on $[1,2]$, use the mass-trimming refinement below. A replacement
of one interval child never deletes its two siblings. The 26 disjoint
nonpositive groups and high positive losses are charged once, and the
cardinality converse is protected throughout.

\subsection{Upper core mass and one-point trimming}
For $r$ distinct positive integers $x_1<\cdots<x_r=R$, unit gaps give
$x_i\le R-r+i$. Hence the elementary necessary inequality is
\begin{equation}\label{u86:eq:87uppermass}
 C=\sum_i x_i\le rR-T_{r-1}.
\end{equation}
It needs integer core points, not primitivity, zero holes or integral
initial outer coordinates. In the model it is
$S_s-M_s-x_{s,1}-x_{s,2}-x_{s,3}-rR_s\le-T_{r-1}$.

Let $U$ be the second-largest core element and $r\ge2$. The other
$r-2$ core elements have sum at least $T_{r-2}$, so
\[
 C\ge R+U+T_{r-2},\qquad
 R-U-2\ge2R-C+T_{r-2}-2.
\]
Deleting the maximum changes $2|F|-\operatorname{diam}F-1$ by exactly
$R-U-2$. On each sign choose independently either the untrimmed core
plus a prefix of zero to four peeled points, or the core with its
maximum deleted and no peeled point added. The two maximizing options
form one actual $F\subset D$ and give
\[
 2|F|-\operatorname{diam}F-1
 =2e-R_P-R_Q-1+g_P+g_Q,
\]
where $g_s\ge0$ dominates all original bounds $2i+R_s-x_{s,i}$ and
also $2R_s-C_s+T_{r_s-2}-2$. The new gains can exceed four, so exactly
the two old gain upper bounds are removed. Their domain remains
nonnegative real; no integer gain branch is authorized.

For each retained own or cross leaf, the leading component and the
leaf contain the corresponding translates of this same $F$. If their
relative displacement is the proved integer $t$, the intersection has
at least $2|F|-\operatorname{diam}F-|t|-1$ points. At a point in the
raw union, the number of hub-leaf incidences is at most total component
multiplicity minus one. Outside the target, one selected member-loss
indicator may additionally be charged. Summing pointwise proves that
the actual group deficit dominates all star intersection bounds plus
any one member loss. The same $F$ simultaneously realizes every star
auxiliary. Other fixed and unmatched losses keep their original
disjoint allocation. This argument is valid with positive holes and
does not assume that the trimmed set is non-AP.

\subsection{Dense spans and the stronger non-AP inverse}
Lev's theorem gives, for a positive integer $r$-set
of maximum $R\le3r/2-2$, writing $k=R-r$ and $C=\sum X$,
\[
 [2k+1,C-(2k+1)]_{\Z}\subseteq\FS(X).
\]
A lower hole lies in $[1,2k]_{\Z}$ and must be a missing generator,
so there are at most $k$; reflection gives at most $k$ upper holes.
Thus $h_X\le2(R-r)$. In a high interval with lower endpoint $z_0$,
set $q=\lfloor3r/2-2\rfloor$. If $z_0>2(q-r)$, the structural
inequality $h_X\ge z\ge z_0$ excludes $R\le q$, giving $R\ge q+1$.
The strict test is essential. This uses the same lattice throughout
and adds no gcd hypothesis to Lev's theorem.

The final ten high parents also use the following separately proved
non-AP inverse. Let $r\ge21$, $\FS(X)$ be non-AP and $\mu(X)\ge u$.
For $21\le L\le r$, put $M=T_L-r$. If $\mu\le M$, every prefix of
size at least $L$ has non-AP support: otherwise the first subsequent
non-AP prefix is a disjoint translate of its AP predecessor and adds
excess at least $1+T_L-r=M+1$. Fix $3\le\beta\le L-5$.
If there is no later small increment, the non-AP prefix floor and
recurrence imply
\[
 \rho\ge A:=r-5+(\beta-2)(r-L).
\]
If $i\ge L$ is the last increment of excess at most $\beta$, the
small-increment inverse gives at most $2\beta$ defects there. With
$t=r-i$, one has $d_0(X)\le2\beta+t$ and
$\rho\ge r-5+(\beta-2)t$. Eliminating $t$ using
$\mu=\rho-2+2d_0$ yields
\[
 \beta\rho\ge(\beta-2)(\mu-4\beta+2)+2(r-5).
\]
Thus this case has the floor
\[
 B:=\max\left(r-5,
 \left\lceil\frac{(\beta-2)(u-4\beta+2)+2(r-5)}{\beta}\right\rceil\right).
\]
The complementary case $\mu\ge M+1$ uses the old inverse, giving
$C:=\max(r-5,v_{\rm old}(r,M+1))$.
The minimum of $A,B,C$ is universally valid; maximize it over the
finite parameter choices and the old floor. This derivation needs no
AP-free enumeration and does not assume that the search is optimal.

The odd-pivot bound $d_0(X)\le R-r+2$ gives independently
$\rho\ge u+2r-2R-2$. In the selected compositions, each side of one
original budget may use its justified constant floor or this odd-pivot
floor. Each choice is a separate necessary inequality on the same
actual set; adding the budgets to each other would double-count.
The eight-parent component uses original-floor odd masks plus a
separate strengthened constant-floor budget. The final two also have
the derived joint forms. Every original row is retained.

\subsection{A full-D small-doubling alternative}
The structural theorem gives $D\subset\Z$, $|D|=79$, $0\in D$ and
primitive $P\subset D$. Thus $\gcd(D-D)=1$. Applying Stanchescu's
Corollary 2 to $D,D$ uses its equal-diameter parameter one and gives
\[
 K<3\cdot79-3=234\quad\Longrightarrow\quad
 \operatorname{diam}D=M_P+M_Q\le K-79.
\]
Every actual configuration therefore satisfies the exhaustive pair
\begin{equation}\label{u86:eq:87stanchsplit}
 \left\{K\le233,\ M_P+M_Q-K\le-79\right\}
 \quad\text{or}\quad K\ge234.
\end{equation}
The low child has both inequalities; the high child has no inverse
width restriction. This argument concerns the ordinary repeated-sum
set $D+D$, and needs no general outer-$B$ integrality or leading
contribution. In a wrapper those separate gates remain necessary for
any inherited stricter outer geometry.

\subsection{The final two parents use actual structural scalars}
The last two selected parents, canonical IDs 1570114 and 1570147, are
balanced high roots. Their refinement has one common actual-set
realization. For each core side of size $r$, assign
\[
 \mu_s=|\FS(X_s)|-1-T_r,\quad h_s=C_s-T_r-\mu_s,\quad
 d_s=d_0(X_s),\quad\rho_s=\mu_s+2-2d_s.
\]
Let $D_0=\Delta(D)$ be the actual full-$D$ diagonal-defect count and set
\[
 \kappa=h_P+h_Q-z,\qquad
 \mathcal T=\kappa+\ell_1+\ell_2+g,\qquad
 \theta=d_P+d_Q+8-D_0.
\]
These variables are nonnegative integers under the structural theorem.
Zero is nondefective because $E$ has at least nine opposite pairs.
The original $\delta=c-79j$ and its cardinality identity are preserved.
The prior common-$F$ construction simultaneously realizes the star
auxiliaries, which remain nonnegative real coordinates.

The scalar rows include the following actual identities and bounds:
\begin{gather*}
 \mu_s\ge u_s,\quad \rho_s\ge v_s,\quad
 \rho_s-\mu_s+2d_s=2,\quad 1\le d_s\le r,\\
 C_s-\mu_s\ge T_r+z,\qquad
 C_s-2\mu_s\le2T_r-1,\\
 C_s\le rR_s-T_{r-1},\qquad d_s\le R_s-r+2,\\
 \kappa+\mu_P+\mu_Q+z-C_P-C_Q=-T_a-T_b,\\
 \delta+\mathcal T+\rho_P+\rho_Q+2\theta
 \le574-I-79j,\\
 D_0-d_P-d_Q+\theta=8,\quad
 2\le D_0\le M_P+M_Q-75,\quad K\le79\cdot80/2.
\end{gather*}
The bounds on $\mathcal T$ are those in \eqref{u86:eq:87postbudget};
in addition $\mathcal T\ge\kappa$ follows from its definition and
nonnegative slack. The strict hole bound
$h_s<|\FS(X_s)|-1=T_r+\mu_s$ explains the integer minus one in the
second line. In these two balanced roots its hypothesis follows directly
from $h_s\le\kappa\le574<T_{35}+\mu_s$. To prove the full-$D$ parity bound, translate its minimum
to zero, apply the odd-pivot bound to its 78 positive members, and add
the defective minimum. Both parities occur since $D$ has gcd of
differences one. The two endpoints are always defective, explaining
$D_0\ge2$. The budget is the original exact budget with justified
nonnegative slack dropped, not a second copy of a fixed floor budget.

For each of the six distinct peeled inner points, choose its actual
nondefect indicator $\nu_{s,i}\in\{0,1\}$, $i=1,2,3$.
Adding points cannot create an old core defect. Of the eight newly
peeled points, at most $8-\sum_{s,i}\nu_{s,i}$ are defective, hence
\[
 \theta\ge\sum_{s,i}\nu_{s,i}.
\]
A point with several midpoint witnesses still uses this one indicator.
Opposite-sign inner labels refer to distinct positive/negative points,
so no point is charged twice.

For the five same-side endpoints $R<x_1<x_2<x_3<M$, write $w=M-R$.
Their integer gaps and midpoint witnesses imply
\begin{equation}\label{u86:eq:87nf}
 w+\sum_{i=1}^3\nu_i\ge7,\qquad
 w+2\sum_{i=1}^3\nu_i\ge8,\qquad
 \nu_i+x_{i+1}-x_{i-1}\ge3,
\end{equation}
where $x_0=R$, $x_4=M$. For the local inequality the neighboring
span is at least two, and equality forces a midpoint witness.
For $w\ge8$ the first two are immediate. For $w=4,5,6,7$, translating
$R$ to zero leaves respectively $\binom{w-1}{3}=1,4,10,20$ profiles.
Directly testing the distinct endpoint pairs gives minimum numbers of
witnessed inner points $3,2,1,1$, respectively. These 35 elementary
profiles prove both global constants. They use only witnesses among
these five points; such witnesses are also valid in $D$.

\subsection{The selected AP-free bounds and width alternatives}
The defective elements of any integer set contain no nontrivial
three-term AP: its middle element would not be defective. For a
positive core, adjoining zero to its defective set remains AP-free.
With $r_3(L)$ the maximal AP-free size in an $L$-point interval,
\[
 d_s+1\le r_3(R_s+1).
\]
Here is the boundary-localization argument used in the next inequalities.
Let $S\subset[0,L]_{\Z}$ contain $0,L$ and have $n$ elements. If a
diagonal-defective $x\in S$ satisfies $x\le L/2$, none of the $x$
disjoint pairs $\{x-i,x+i\}$, $1\le i\le x$, has both members in $S$.
Counting the points in $[0,2x]$ and then those above $2x$ gives
\[
 n\le x+1+(L-2x)=L-x+1,
\]
so $x\le L-n+1$. Reflection gives $x\ge n-1$ for a defective point
in the other half. All defects therefore lie in
\[
 [0,L-n+1]_{\Z}\ \cup\ [n-1,L]_{\Z},
\]
two intervals of $L-n+2$ points. The defective set is AP-free, so its
cardinality is at most $2r_3(L-n+2)$; overlap of the intervals only
weakens this upper bound. Apply this to $\{0\}\cup X_s$, where zero
itself is defective, and then to a translate of $D$. This gives
\[
 d_s\le2r_3(R_s-r+1)-1,\qquad
 D_0\le2r_3(M_P+M_Q-77).
\]
These facts require neither an odd element nor primitivity.

The finite dependencies are the exhaustive $r_3(0),\ldots,r_3(30)$
table and the following required finite upper bounds
\[
 r_3(31)\le12,\ r_3(35)\le13,\ r_3(39)\le14,\ r_3(40)\le15,
\]
\[
 r_3(50)\le16,\qquad r_3(58)\le19,\qquad r_3(70)\le20.
\]
The complete inclusion/exclusion verifier maintains an increasing
AP-free chosen set and forbids every future third term $2x-y$.
Branching on the least available point covers all possibilities;
available-position and contiguous-block capacity bounds justify every
prune. The finite upper bounds are mathematical proof dependencies,
not numerical optimizer conclusions.



Restricting an AP-free set to consecutive blocks and summing upper
capacities proves every partition bound $g(L)$ used by the model.
Checking $7r_3(L)\le2L+30$ through 70 and adjoining 70-blocks of
capacity 20 proves the same inequality for every positive $L$.
Together with boundary localization this gives
\begin{gather*}
 7d_s-2R_s\le25,\qquad 7d_s-4R_s\le57-4r,\\
 7D_0-2(M_P+M_Q)\le32,\qquad
 7D_0-4(M_P+M_Q)\le68-4\cdot79.
\end{gather*}
For an integer core span $R_s$, the pair $R_s\le t$ or $R_s\ge t+1$
is exhaustive. In the first child,
\[
 d_s\le\min\{r,g(t+1)-1,2g(t-r+1)-1,t-r+2\}.
\]
This uses monotonicity of the true $r_3$, not an assumed monotonicity
of a computed partition bound.

Only $r=35$, $t=77,79$ occur in the final two selected trees.
The explicit partitions $78=70+8$ and $80=70+10$, with capacities
$20+4$ and $20+5$, give $d_s\le23$ and $d_s\le24$ after subtracting
one for zero. The complementary children are $R_s\ge78$ and
$R_s\ge80$. All six selected AP-cap splits retain both children.
The trees also have two full-$D$ Stanchescu splits
\eqref{u86:eq:87stanchsplit}. Their remaining splits are the already
proved normal, inner and core-or-large alternatives.
No optional conditional-span split, full-$D$ finite-cap split,
peeled-$D$ tail cut or arbitrary integer-form split appears in these
selected trees. The final-two composition has no fixed-score layer.
Its 188 nodes, 93 binary splits and 95 exact leaves satisfy
$188=2+2\cdot93$ and $95=2+93$.

\subsection{Exact selected partition and dependency list}

The required ordinary classification has 1614434 exact contradictions
and 866 supplemental parents: 251 low and 615 high IDs. The following
is the complete supplemental partition. The actual canonical ID lists,
not the displayed totals, determine the parent-set union.


\begin{center}\small
\begin{longtable}{p{.43\linewidth}rrrr}
\toprule
Selected component&Parents&Children&Splits&Leaves\\\midrule
\endfirsthead
\toprule
Selected component&Parents&Children&Splits&Leaves\\\midrule
\endhead
Low geometry 27&27&27&44&71\\
Low fixed 77&77&77&228&305\\
Zero-hole five-missing 42&42&42&707&749\\
Zero-hole loss cover 11&11&165&2621&2786\\
Positive three-missing 24&24&24&69&93\\
Positive fixed 59&59&59&759&818\\
Positive three-interval 11&11&33&2959&2992\\
High geometry 9&9&9&10&19\\
High geometry 28&28&28&86&114\\
High geometry 85&85&85&143&228\\
High core mass 3&3&3&90&93\\
High geometry 107&107&107&274&381\\
Recovered high geometry 121&121&121&342&463\\
Recovered high core mass 1&1&1&15&16\\
Recovered high core mass 2&2&2&51&53\\
High geometry 86&86&86&312&398\\
High dense spans 9&9&9&251&260\\
New high core mass 2&2&2&49&51\\
High geometry 150&150&150&644&794\\
Final high core mass 2&2&2&28&30\\
High inverse and odd-pivot 8&8&8&388&396\\
High actual scalars and AP caps 2&2&2&93&95\\\midrule
Total&866&1042&10163&11205\\\bottomrule
\end{longtable}
\end{center}
The low totals are 427 child models, 7387 splits and 7814 leaves;
the high totals are 615 child models, 2776 splits and 3391 leaves.
The covers add 176 child models beyond their parent count. Thus
\[
 11205=1042+10163,\qquad
 1625639=1614434+11205=1615300+176+10163.
\]

The actual canonical parent IDs must form exactly the ordinary zero
positions, with pairwise disjointness. All required cover children and
both children of every binary split must be included.


The logical dependencies of the new endpoint are precisely:
\begin{enumerate}
\item The retained foundational real-to-integer reduction, covering,
peeling and budget identities; the descending exclusions proved in order
through size 88.
\item The complete size-87 lattice proof above, including the analytic four-combination
nonprimitive AP argument in Appendix~\ref{u86ap:sec:certificate}, the strengthened non-AP defect theorem,
old effective-hole and flexible-inverse lemmas, and Stanchescu's
common-lattice results.
\item The complete size-87 root classification,
605 effective-hole inputs, original-cap quantile gates, 2992-mode
retained-center identities, and full nonnegative-variable realization.
\item The already proved normal, inner and core-loss disjunctions,
checked outer integer witnesses, and the fixed-score 110-label/26-group
accounting used by the selected low forests.
\item The actual-set zero-hole loss cover and three-interval positive
cover; common-$F$ dense stars, one-point mass trimming and retained-unit
consequences, only at their selected levels.
\item The upper core-mass inequality, qualified dense-span exclusion
from Lev, strengthened non-AP inverse, original and joint odd-pivot
budget substitutions, and full-$D$ Stanchescu disjunction, where selected.
\item For the final two parents, the actual scalar and shared-indicator
identities, all 35 elementary five-point profiles, the AP-free finite
upper bounds and their block partitions, and exactly the selected
77/79 core-span alternatives described above.
\item Exact current-matrix Farkas arithmetic for every ordinary and
supplemental leaf, exhaustive covers and binary topology, immutable
source linkage, compact assignment equivalence, and final disjoint
integration. These are separate obligations, all required.
\end{enumerate}
The optional tail-eight and conditional-span research refinements are
not selected dependencies. Neither floating-point infeasibility,
source-code similarity alone, nor a partial successful run replaces
any item on this list.


\section{Composition of the ten exclusions}\label{u86:sec:composition}
\begin{theorem}[Global upper bound]
\label{u86:thm:upper86}
Every feasible finite real set has at most $86$ elements.
\end{theorem}
\begin{proof}
The earlier chapters prove $N\le96$. The size-$96$ actual-set reduction
and its finite lemma give $N\le95$. Repeating the proved descending
step for sizes $95,94,93,92,91,90,89,88$ gives $N\le87$.
A hypothetical feasible size-$87$ set therefore contains zero, since
otherwise adjoining zero would give a forbidden size-$88$ set.
Integralization supplies an integral size-$87$ set with zero.
The complete structural theorem maps it to one of the 1615300 canonical
necessary models. The ordinary assignment and all 22 supplemental
families cover that model and every required alternative. Their integer
row combinations yield the contradiction of
Lemma~\ref{completion:coverage}. Thus size $87$ is impossible, so
$N\le86$. None of these steps invokes SIZE86 or any subsequent exclusion
of size $86$.
\end{proof}


\section{Finite verification record}\label{completion:verification}

The finite premises used in the preceding proof were freshly recomputed
on 5 October 2026. The complete 19-stage endpoint replay exited with
status zero, and every added prerequisite and independent check passed.
Together with the actual-set reductions and exhaustive alternatives proved
above, these calculations discharge the finite premises of
Theorem~\ref{u86:thm:upper86} through Lemma~\ref{completion:coverage}.

\subsection{Companion files and reproduction}
The complete electronic companion is
\path{subset_sums_upper86_complete_verification_20261005.zip}.
Extract it completely and run the following command from its top-level
directory:
\begin{quote}\small\ttfamily
python3 -I -S -B verify\_complete\_upper86.py
\end{quote}
On Windows the corresponding command is
\begin{quote}\small\ttfamily
py -3 -I -S -B verify\_complete\_upper86.py
\end{quote}
Python 3.10 or later and its standard library suffice. No network access,
numerical solver, optimization package, or original research directory is
required. A complete run may take several hours. The supplied
\path{README_COMPLETE.txt} gives extraction, storage and launcher details.

The original generators, root descriptors, coefficient pools, compact
assignments, trees and manifests are retained under
\path{payload/subset_sum_verification/}. The paths beginning
\path{proofs/} in the mathematical chapters are relative to that directory.
The new complete exposition is included as
\path{proof/complete_upper86_proof.tex}, one independently compilable
\LaTeX{} file. Its compiled PDF is delivered separately alongside the archive.

The default entry point runs the earlier endpoint stages through size 88,
then the complete size-87 controller, then the added finite prerequisites
and independent checks. Each invocation creates its own new result
directory; it does not reuse endpoint receipts from a prior run.
The resulting \path{result.json}, aggregate log and individual step logs
record completion. A failed or interrupted run does not establish the
finite lemmas. The optional integrity-only and auxiliary-only modes do
not establish the complete chain either.

\subsection{The execution actually performed}
The validation of this edition was compositional:
\begin{enumerate}
\item One fresh invocation of the unchanged 19-stage endpoint replay
began at 16:53:13 UTC and finished at 20:21:50 UTC on 5 October 2026.
Every stage passed, the process exited zero, and the final source-integrity
check passed. Its elapsed time was 12\,517.34 seconds, approximately
3 hours 28 minutes 37 seconds.
\item All fourteen added auxiliary, independent and control subprocesses
were executed through the new public wrapper and exited zero. This
auxiliary run took 303.96 seconds. The public runner, all mathematical
checkers and all proof data in the delivered version are byte-identical
to those tested. The later changes to packaging and documentation do not
alter their mathematical computations.
\item Thirteen wrapper rejection controls and twelve mathematical
negative controls passed. Extraction and execution paths containing
spaces and non-ASCII characters were tested. The executed environment
was Linux/POSIX with Python 3.12.14. The Windows launcher and paths were
statically checked; no native-Windows execution is claimed.
\end{enumerate}
A second multi-hour invocation of the new wrapper in its default mode
was not performed. Its routing and failure behavior were tested
separately, and the completed component executions above cover all
required mathematical stages. The exact commands, elapsed times,
source-equivalence record and scope are recorded in
\path{REPRODUCTION_REPORT.json}. The companion's
\path{validation/historical_19_stage.log} and
\path{validation/auxiliary_result.json} preserve the corresponding
execution records. They document the calculations; the calculations,
not the existence of a record, supply the finite premises.

\subsection{Coverage of the final endpoint}
The size-87 replay regenerated and covered all 1\,615\,300 canonical
roots. Exactly 1\,614\,434 received ordinary contradictions; the
remaining 866 parents are the disjoint union of the 22 specified
supplemental families. These have 1\,042 child models,
10\,163 exhaustive binary splits and 11\,205 contradictory leaves.
Thus the final endpoint has
\[
 1\,614\,434+11\,205=1\,625\,639
\]
exact leaf contradictions. The ordinary replay used 65 consecutive
ranges whose union is exactly $[0,1615300)$; the final composition
checked the complete range union and every supplemental identity.
It was a complete fresh replay, rather than a sampled root test.

The additional checks also recomputed the finite structural and
strict-gate prerequisites, the additional AP-free bounds and block
partitions, fixed-label ownership, canonical root ordering, exact
supplemental-parent equality, selected split rows and independent
stored-row contradictions. The 605 effective-hole witnesses, all
2\,992 three-missing-unit integer-center identities, and all 35
five-point profiles were checked in their stated domains. The earlier
AP-free and non-AP finite bases are included in the 19-stage replay.
The four analytic combinations in Appendix~\ref{u86ap:sec:certificate}
were independently expanded; the former 32-case exception calculation
is retained only as an additional regression check.

At each finite leaf, its current necessary rows and the supplied
nonnegative integer multipliers satisfy the exact criterion in
Lemma~\ref{completion:farkas}. Exact child and parent identities, both
children of every selected split, and every leaf are checked; census
totals alone are not the test. The universal reductions, cited
mathematical inputs and applicability of the rows have their written
proofs in this supplement. No proof-assistant formalization is claimed.


\clearpage\appendix

\section{The fixed-slot deletion certificate}\label{base:app:slots}

Let
\[
 p_1>p_2>p_3>p_4>0,
 \qquad q_1>q_2>q_3>q_4>0,
\]
and define the eight positive gaps
\[
\begin{gathered}
 r_1=p_1-p_2,\ r_2=p_2-p_3,\ r_3=p_3-p_4,\ r_4=p_4,\\
 s_1=q_1-q_2,\ s_2=q_2-q_3,\ s_3=q_3-q_4,\ s_4=q_4.
\end{gathered}
\]
Put
\[
\begin{aligned}
 c_p&=p_1-p_2-p_3-p_4,&
 \alpha_p&=p_2-p_3-p_4,&
 \beta_p&=2p_2-p_1-p_3-p_4,\\
 c_q&=q_1-q_2-q_3-q_4,&
 \alpha_q&=q_2-q_3-q_4,&
 \beta_q&=2q_2-q_1-q_3-q_4.
\end{aligned}
\]
The two ordered eight-tuples are
\[
 B^+=(p_1,p_2,p_3,p_4,-q_1,-q_2,-q_3,-q_4),
\]
\[
 B^-=(q_1,q_2,q_3,q_4,-p_1,-p_2,-p_3,-p_4).
\]
For $\varepsilon\in\{+,-\}$, set
\[
 C^\varepsilon_{ij}=c_\varepsilon+B^\varepsilon_i+B^\varepsilon_j,
 \quad
 U^\varepsilon_i=\alpha_\varepsilon+B^\varepsilon_i,
 \quad
 V^\varepsilon_i=\beta_\varepsilon+B^\varepsilon_i,
\]
where $1\leq i\leq j\leq8$, and $U,V$ use $2\leq i\leq8$.  For a decreasing
quadruple $x=(x_1,x_2,x_3,x_4)$, define
\[
\begin{aligned}
 H_1(x)&=x_3-x_4,& H_2(x)&=2x_3-x_2-x_4,\\
 H_3(x)&=x_2+x_3-x_1-x_4,& H_4(x)&=2x_3-x_1-x_4,\\
 H_5(x)&=3x_3-x_1-x_2-x_4.
\end{aligned}
\]
Write $H_i^+=H_i(p)$ and $H_i^-=H_i(q)$.  The common slot universe has size
\[
 2\left(\binom{9}{2}+7+7+5\right)=110.
\]

The following table gives $15$ templates.  Every starred row is also used
after interchanging $+$ and $-$.  The four unstarred rows and the two
orientations of the eleven starred rows give $26$ groups.

\begingroup
\scriptsize
\setlength{\tabcolsep}{3.6pt}
\begin{longtable}{c>{\raggedright\arraybackslash}p{0.43\textwidth}p{0.39\textwidth}}
\toprule
&\textbf{sum of scores}&$(r_1,r_2,r_3,r_4;s_1,s_2,s_3,s_4)$\\
\midrule
\endfirsthead
\toprule
&\textbf{sum of scores}&$(r_1,r_2,r_3,r_4;s_1,s_2,s_3,s_4)$\\
\midrule
\endhead
1&$C^+_{25}+C^-_{25}$&$(0,0,-1,-2;0,0,-1,-2)$\\
$2^*$&$C^+_{48}+H^+_3$&$(0,0,0,-1;0,0,0,-1)$\\
3&$C^+_{58}+C^-_{58}$&$(0,-1,-2,-4;0,-1,-2,-4)$\\
$4^*$&$C^+_{77}+H^+_4$&$(0,-1,0,-2;0,0,-2,-2)$\\
5&$U^+_6+U^-_6$&$(0,0,-1,-2;0,0,-1,-2)$\\
6&$V^+_6+V^-_6$&$(-1,0,-1,-2;-1,0,-1,-2)$\\
$7^*$&$C^+_{15}+H^+_1+C^-_{55}$&$(0,-1,-1,-3;0,-1,-2,-3)$\\
$8^*$&$C^+_{26}+C^-_{57}+V^-_4$&$(0,0,-2,-3;0,0,-2,-3)$\\
$9^*$&$C^+_{35}+V^+_3+C^-_{46}$&$(0,0,0,-2;0,-1,-2,-2)$\\
$10^*$&$C^+_{36}+U^+_4+U^-_5$&$(0,0,-1,-2;0,0,-1,-2)$\\
$11^*$&$C^+_{44}+U^+_8+H^+_5$&$(0,-1,0,-1;0,0,0,-1)$\\
$12^*$&$C^+_{45}+V^+_8+C^-_{67}$&$(0,0,-3,-4;0,-1,-2,-4)$\\
$13^*$&$C^+_{47}+V^+_7+H^+_2$&$(0,0,0,-2;0,0,-2,-2)$\\
$14^*$&$C^+_{56}+V^+_2+C^-_{66}$&$(0,0,-2,-4;0,-2,-3,-4)$\\
$15^*$&$C^+_{68}+U^+_3+V^-_5$&$(0,0,-1,-3;-1,0,-1,-3)$\\
\bottomrule
\end{longtable}
\endgroup

Every gap is strictly positive, while each coefficient vector is
coordinatewise nonpositive and has a negative coordinate.  Hence all slots in
one group cannot have positive score.  The $26$ groups are mutually disjoint in the common universe, as is checked
directly by their labels. The accompanying exact-arithmetic script also
reconstructs every coefficient vector and checks this disjointness.

For the two translate pairings, first take
\[
 u_1=c_p-q_1,\qquad v_1=-c_q+p_1.
\]
Then
\[
 u_1-v_1=-(p_2+p_3+p_4+q_2+q_3+q_4)<0.
\]
If $\beta_p+\beta_q\leq0$, take
$u_2=\beta_p$ and $v_2=-\beta_q$.  Otherwise take
\[
 u_2=c_p-q_2,\qquad v_2=-c_q+p_2.
\]
The identity
\[
 u_2-v_2=-(\beta_p+\beta_q)-2(p_3+p_4+q_3+q_4)<0
\]
proves the needed order in the second branch.  The two pairs use distinct
table slots, so \eqref{base:eq:halfline} can be applied twice.


\section{Complete certificates for one empty covering unit}\label{u96:app:one-empty}
This appendix is part of the proof of Lemma~\ref{u96:lem:one-empty}. It gives
the models, a complete case table, and every multiplier; no external
certificate database or numerical solver is needed to specify the proof.

\subsection{Centers, modes, and constraints}
Index the ten positive-side and ten reflected negative-side centers from
zero, solely for this appendix:
\[
\begin{aligned}
v^p&=(c_p+p_1,c_p+p_2,c_p+p_3,c_p+p_4,
       c_p-q_1,c_p-q_2,c_p-q_3,c_p-q_4,\alpha_p,\beta_p),\\
v^q&=(c_q+q_1,c_q+q_2,c_q+q_3,c_q+q_4,
       c_q-p_1,c_q-p_2,c_q-p_3,c_q-p_4,\alpha_q,\beta_q).
\end{aligned}
\]
The first pair consists of indices 4 on the two sides. The second consists
of indices $j=9$ in branch A and $j=5$ in branch B. The other sixteen slots
are unpaired. A mode specifies the one missing unit and, when applicable,
the contributing side of the second pair. Labels $pi,qi$ refer to center
$v_i^p,v_i^q$; labels $1p,2p,2q$ refer to $v_4^p,v_j^p,v_j^q$ when used
as paired units. Such labels identify distinct covering slots even when
their numerical center values happen to agree.

Keep every unpaired center except the missing one. Keep $1p$ unless the
first pair is missing. Keep the selected member $2p$ or $2q$ unless the
second pair is missing. In branch A with missing $p5$ or $q5$, keep neither
member of the second pair: this is the deliberate relaxation allowing it
to split. These rules define the active list $\mathcal C$ for every mode.

All variables
\[
x=(p_1,\ldots,p_4,q_1,\ldots,q_4,M_p,M_q,R_p,R_q,S_p,S_q,
\delta,t,(l_v,u_v)_{v\in\mathcal C})
\]
are nonnegative real numbers. The following row names denote the indicated
inequalities in the orientation ``left side $\leq$ right side.''
\begin{center}\small
\begin{tabular}{ll}
\toprule
Row name&Constraint\\
\midrule
$D_i^p,D_i^q$ ($1\leq i\leq3$)&$p_{i+1}-p_i\leq-1$, $q_{i+1}-q_i\leq-1$\\
$D_4^p,D_4^q$&$M_p-p_4\leq-1$, $M_q-q_4\leq-1$\\
$E_p,E_q$&$R_p-M_p\leq-4$, $R_q-M_q\leq-4$\\
$G_p,G_q$&$-R_p+t\leq-39$, $-R_q-t\leq-47$\\
$H_p$&$2M_p+6R_p-2S_p+87t\leq-1556$\\
$H_q$&$2M_q+6R_q-2S_q-87t\leq-2252$\\
$Z,J$&$\delta\leq21$, $t\leq8$\\
$B$ (branch A)&$\beta_p+\beta_q\leq0$\\
$B$ (branch B)&$-\beta_p-\beta_q\leq-1$\\
$L_v$ (positive side)&$-v-l_v-t\leq-52$\\
$U_v$ (positive side)&$v-u_v-S_p+t\leq-43$\\
$L_v$ (negative side)&$-v-l_v+t\leq-44$\\
$U_v$ (negative side)&$v-u_v-S_q-t\leq-51$\\
$\Sigma$&$\sum_{v\in\mathcal C}(l_v+u_v)-\delta\leq0$\\
\bottomrule
\end{tabular}
\end{center}
In the multiplier tables, for example, $L_{p5}$ means the lower row for
$v^p_5$, and $U_{2q}$ means the upper row for the selected second negative
paired center. The objective is
\[
 f(x)=S_p+S_q+\delta-2(M_p+M_q)-2(R_p+R_q),\qquad F=95+f(x).
\]

\subsection{How to read and check a certificate}
Write the rows as $Ax\leq b$. Every displayed multiplier is a nonnegative
integer. A certificate with objective multiplier $D>0$ has row-weight
vector $w$ satisfying
\[
 Df+wA\geq0\quad\hbox{coefficientwise}.
\]
Since $x\geq0$, it proves $f(x)\geq-wb/D$, hence $F\geq95-wb/D$.
A contradiction certificate, marked $D=0$, satisfies
\[
 wA\geq0\quad\hbox{coefficientwise},\qquad wb=-K<0,
\]
and proves $0\leq wAx\leq-K$. Each certificate uses only rows present in
every mode to which it is assigned. The complete case table below contains
33 branch-A modes and 35 branch-B modes. A dash means no such mode is
required; for a missing second pair or a deliberately relaxed pair, the
single certificate is placed in the $p$ column without asserting that
the pair contributes on that side.
\begin{center}\small
\begin{tabular}{lrrrr}
\toprule
Missing unit&A: second $p$&A: second $q$&B: second $p$&B: second $q$\\
\midrule
First pair&1&2&10&11\\
Second pair&3&--&12&--\\
$p0$&4&5&13&14\\
$p1$&4&6&12&14\\
$p2$&4&6&12&14\\
$p3$&4&6&12&14\\
$p5$&3&--&--&--\\
$p6$&4&6&12&14\\
$p7$&4&6&12&14\\
$p8$&4&7&12&14\\
$p9$&--&--&15&16\\
$q0$&4&6&12&14\\
$q1$&4&6&12&14\\
$q2$&4&6&12&14\\
$q3$&4&6&12&14\\
$q5$&8&--&--&--\\
$q6$&4&6&17&14\\
$q7$&4&6&18&19\\
$q8$&9&6&12&14\\
$q9$&--&--&20&21\\
\bottomrule
\end{tabular}
\end{center}

\subsection{All twenty-one multiplier combinations}
The last column lists the nonzero weights as a formal sum of row names.
For $D>0$ the conclusion is the stated lower bound for $F$; for $D=0$ it
is the impossible inequality $0\leq-K$.
\begingroup\footnotesize
\setlength{\tabcolsep}{4pt}
\renewcommand{\arraystretch}{1.35}
\begin{longtable}{rrp{1.65cm}p{11.7cm}}
\toprule
ID&$D$&Conclusion&Integer row combination\\
\midrule
\endfirsthead
\toprule
ID&$D$&Conclusion&Integer row combination\\
\midrule
\endhead
1&13&$F\geq\frac{25847}{13}$&\( 10G_p+\allowbreak 10G_q+\allowbreak 6H_p+\allowbreak 6H_q+\allowbreak U_{p0}+\allowbreak 7D_{3}^{p}+\allowbreak 14D_{4}^{p}+\allowbreak 5L_{p5}+\allowbreak 3L_{2p}+\allowbreak U_{q0}+\allowbreak 7D_{3}^{q}+\allowbreak 14D_{4}^{q}+\allowbreak 2L_{q5}+\allowbreak 6L_{q8}+\allowbreak 13\Sigma \)\\
2&13&$F\geq\frac{25847}{13}$&\( 10G_p+\allowbreak 10G_q+\allowbreak 6H_p+\allowbreak 6H_q+\allowbreak U_{p0}+\allowbreak 7D_{3}^{p}+\allowbreak 14D_{4}^{p}+\allowbreak 2L_{p5}+\allowbreak 6L_{p8}+\allowbreak 3L_{2q}+\allowbreak U_{q0}+\allowbreak 7D_{3}^{q}+\allowbreak 14D_{4}^{q}+\allowbreak 5L_{q5}+\allowbreak 13\Sigma \)\\
3&14&$F\geq1982$&\( 2E_q+\allowbreak 8G_p+\allowbreak 16G_q+\allowbreak 6H_p+\allowbreak 7H_q+\allowbreak 2U_{p0}+\allowbreak 8D_{3}^{p}+\allowbreak 16D_{4}^{p}+\allowbreak 6L_{p8}+\allowbreak 4L_{1p}+\allowbreak 8D_{3}^{q}+\allowbreak 16D_{4}^{q}+\allowbreak 4L_{q5}+\allowbreak 4L_{q8}+\allowbreak 14\Sigma+\allowbreak 95J \)\\
4&0&$0\leq-171$&\( Z+\allowbreak L_{1p}+\allowbreak L_{2p}+\allowbreak L_{q5}+\allowbreak L_{q8}+\allowbreak \Sigma \)\\
5&11&$F\geq\frac{21575}{11}$&\( 8G_p+\allowbreak 8G_q+\allowbreak 5H_p+\allowbreak 5H_q+\allowbreak U_{p1}+\allowbreak 6D_{3}^{p}+\allowbreak 12D_{4}^{p}+\allowbreak 6L_{p8}+\allowbreak L_{1p}+\allowbreak 2L_{2q}+\allowbreak U_{q0}+\allowbreak 6D_{3}^{q}+\allowbreak 12D_{4}^{q}+\allowbreak 5L_{q5}+\allowbreak 11\Sigma \)\\
6&2&$F\geq2424$&\( 20E_p+\allowbreak 22E_q+\allowbreak 16G_p+\allowbreak 24G_q+\allowbreak H_q+\allowbreak 8Z+\allowbreak 2U_{p0}+\allowbreak 12D_{3}^{p}+\allowbreak 24D_{4}^{p}+\allowbreak 10L_{p8}+\allowbreak 4L_{1p}+\allowbreak 4L_{2q}+\allowbreak 12D_{3}^{q}+\allowbreak 24D_{4}^{q}+\allowbreak 8L_{q5}+\allowbreak 10\Sigma+\allowbreak 95J \)\\
7&4&$F\geq2011$&\( 8E_p+\allowbreak 20E_q+\allowbreak 6G_p+\allowbreak 24G_q+\allowbreak H_p+\allowbreak 2H_q+\allowbreak 2U_{p0}+\allowbreak 6D_{2}^{p}+\allowbreak 12D_{3}^{p}+\allowbreak 14D_{4}^{p}+\allowbreak 4L_{1p}+\allowbreak 4L_{2q}+\allowbreak 12D_{3}^{q}+\allowbreak 24D_{4}^{q}+\allowbreak 4L_{q5}+\allowbreak 4L_{q6}+\allowbreak 4\Sigma+\allowbreak 95J \)\\
8&10&$F\geq1940$&\( 6E_q+\allowbreak 4G_p+\allowbreak 16G_q+\allowbreak 4H_p+\allowbreak 5H_q+\allowbreak 2U_{p0}+\allowbreak 8D_{3}^{p}+\allowbreak 12D_{4}^{p}+\allowbreak 2L_{p8}+\allowbreak 4L_{1p}+\allowbreak 8D_{3}^{q}+\allowbreak 16D_{4}^{q}+\allowbreak 4L_{q6}+\allowbreak 4L_{q8}+\allowbreak 10\Sigma+\allowbreak 95J \)\\
9&0&$0\leq-171$&\( E_p+\allowbreak G_p+\allowbreak Z+\allowbreak D_{4}^{p}+\allowbreak L_{1p}+\allowbreak L_{2p}+\allowbreak L_{q5}+\allowbreak \Sigma \)\\
10&2&$F\geq2178$&\( 14E_p+\allowbreak 22E_q+\allowbreak 10G_p+\allowbreak 24G_q+\allowbreak H_q+\allowbreak 4Z+\allowbreak 2U_{p0}+\allowbreak 12D_{3}^{p}+\allowbreak 18D_{4}^{p}+\allowbreak 2L_{p9}+\allowbreak 6L_{2p}+\allowbreak 12D_{3}^{q}+\allowbreak 24D_{4}^{q}+\allowbreak 6L_{q6}+\allowbreak 6L_{q9}+\allowbreak 6\Sigma+\allowbreak 95J \)\\
11&2&$F\geq2178$&\( 22E_p+\allowbreak 14E_q+\allowbreak 24G_p+\allowbreak 10G_q+\allowbreak H_p+\allowbreak 4Z+\allowbreak 12D_{3}^{p}+\allowbreak 24D_{4}^{p}+\allowbreak 6L_{p6}+\allowbreak 6L_{p9}+\allowbreak 6L_{2q}+\allowbreak 2U_{q0}+\allowbreak 12D_{3}^{q}+\allowbreak 18D_{4}^{q}+\allowbreak 2L_{q9}+\allowbreak 6\Sigma \)\\
12&2&$F\geq2757$&\( 20E_p+\allowbreak 34E_q+\allowbreak 16G_p+\allowbreak 36G_q+\allowbreak H_q+\allowbreak 4Z+\allowbreak 2U_{p0}+\allowbreak 12D_{3}^{p}+\allowbreak 24D_{4}^{p}+\allowbreak 2L_{p9}+\allowbreak 6L_{1p}+\allowbreak 18D_{3}^{q}+\allowbreak 36D_{4}^{q}+\allowbreak 6L_{q6}+\allowbreak 6L_{q7}+\allowbreak 6L_{q9}+\allowbreak 6\Sigma+\allowbreak 95J \)\\
13&2&$F\geq2101$&\( 12E_p+\allowbreak 22E_q+\allowbreak 8G_p+\allowbreak 24G_q+\allowbreak H_q+\allowbreak 2Z+\allowbreak 2U_{p1}+\allowbreak 8D_{3}^{p}+\allowbreak 16D_{4}^{p}+\allowbreak 2L_{p9}+\allowbreak 4L_{1p}+\allowbreak 12D_{3}^{q}+\allowbreak 24D_{4}^{q}+\allowbreak 4L_{q6}+\allowbreak 4L_{q7}+\allowbreak 4L_{q9}+\allowbreak 4\Sigma+\allowbreak 95J \)\\
14&0&$0\leq-215$&\( Z+\allowbreak L_{p9}+\allowbreak L_{1p}+\allowbreak L_{2q}+\allowbreak L_{q7}+\allowbreak L_{q9}+\allowbreak \Sigma \)\\
15&2&$F\geq2662$&\( 20E_p+\allowbreak 34E_q+\allowbreak 16G_p+\allowbreak 36G_q+\allowbreak H_q+\allowbreak 2B+\allowbreak 4Z+\allowbreak 2U_{p0}+\allowbreak 12D_{3}^{p}+\allowbreak 24D_{4}^{p}+\allowbreak 6L_{1p}+\allowbreak 18D_{3}^{q}+\allowbreak 36D_{4}^{q}+\allowbreak 6L_{q6}+\allowbreak 6L_{q7}+\allowbreak 4L_{q9}+\allowbreak 6\Sigma+\allowbreak 95J \)\\
16&0&$0\leq-120$&\( B+\allowbreak Z+\allowbreak L_{1p}+\allowbreak L_{2q}+\allowbreak L_{q7}+\allowbreak \Sigma \)\\
17&2&$F\geq2691$&\( 20E_p+\allowbreak 34E_q+\allowbreak 16G_p+\allowbreak 36G_q+\allowbreak H_q+\allowbreak 10Z+\allowbreak 2U_{p0}+\allowbreak 6D_{3}^{p}+\allowbreak 24D_{4}^{p}+\allowbreak 2L_{p9}+\allowbreak 6L_{1p}+\allowbreak 18D_{3}^{q}+\allowbreak 36D_{4}^{q}+\allowbreak 12L_{q7}+\allowbreak 6L_{q9}+\allowbreak 12\Sigma+\allowbreak 95J \)\\
18&2&$F\geq2697$&\( 20E_p+\allowbreak 34E_q+\allowbreak 16G_p+\allowbreak 36G_q+\allowbreak H_q+\allowbreak 10Z+\allowbreak 2U_{p0}+\allowbreak 18D_{3}^{p}+\allowbreak 24D_{4}^{p}+\allowbreak 2L_{p9}+\allowbreak 6L_{1p}+\allowbreak 18D_{3}^{q}+\allowbreak 36D_{4}^{q}+\allowbreak 12L_{q6}+\allowbreak 6L_{q9}+\allowbreak 12\Sigma+\allowbreak 95J \)\\
19&0&$0\leq-215$&\( Z+\allowbreak L_{p9}+\allowbreak L_{1p}+\allowbreak L_{2q}+\allowbreak L_{q6}+\allowbreak L_{q9}+\allowbreak \Sigma \)\\
20&2&$F\geq2276$&\( 16E_p+\allowbreak 26E_q+\allowbreak 12G_p+\allowbreak 28G_q+\allowbreak H_q+\allowbreak 2B+\allowbreak 4Z+\allowbreak 2U_{p0}+\allowbreak 12D_{3}^{p}+\allowbreak 20D_{4}^{p}+\allowbreak 6L_{1p}+\allowbreak 14D_{3}^{q}+\allowbreak 28D_{4}^{q}+\allowbreak 6L_{q6}+\allowbreak 2L_{q7}+\allowbreak 4L_{q8}+\allowbreak 6\Sigma+\allowbreak 95J \)\\
21&0&$0\leq-171$&\( Z+\allowbreak L_{p9}+\allowbreak L_{1p}+\allowbreak L_{2q}+\allowbreak L_{q8}+\allowbreak \Sigma \)\\
\bottomrule
\end{longtable}
\endgroup

For example, certificate 8, the weakest lower bound, has $D=10$ and
$wb=-18450$. Its coefficients on the sixteen original variables in
$10f+wA$ are all zero, and its remaining coefficients on the loss variables
are nonnegative. It therefore gives $F\geq95+18450/10=1940>1836$.
Certificate 4 has $wb=-171$; its non-loss coefficients in $wA$ are
$2p_3+2p_4+2q_3+2q_4$, all other terms being nonnegative loss terms.
These examples illustrate both certificate types. Expanding every row in
the same manner verifies all entries. The companion verifier reconstructs
each model, checks that the mode list is exhaustive, and performs every
coefficient comparison using integers and rational numbers.


\clearpage
\section[Certificate specification for upper bound 100]{Complete certificate specification for the upper bound 100}\label{u96:app:upper100}
This appendix specifies all 2512 linear systems and all 42 integer
combinations used to exclude their feasible-set interpretations. It is
independent of floating-point optimization. The companion data file records
one successful combination per system; that assignment is a convenience,
not an additional premise. Trying the 42 printed combinations in order
provides a complete finite verification rule.

\subsection{Four regimes and all their modes}
Use exactly the centers, branches, paired units, missing-unit choices, and
active lists $\mathcal C$ defined in Appendix~\ref{u96:app:one-empty}. There
are 68 basic modes. In particular, the two possibly split branch-A second
pairs are still deliberately omitted from $\mathcal C$.
Put $a=a_0+t$, $b=b_0-t$, and let $f$ be the same linear form as in that
appendix. Now $F=d+f$, where $d=n-8$ depends on the regime.
\begin{center}\small
\begin{tabular}{lrrrrrrrrr}
\toprule
Regime&$n$&$a_0$&$b_0$&$t_{\max}$&$\delta_{\max}$&$s_0$&$A_p$&$A_q$&Upper $F$\\
\midrule
I&102&36&49&13&45&85&1308&2413&1874\\
Z&101&34&50&16&69&85&1138&2498&1800\\
P&101&36&48&12&41&87&1376&2420&1856\\
E&101&37&47&10&69&85&1393&2243&1800\\
\bottomrule
\end{tabular}
\end{center}
The letters designate size 102, zero holes with no retained exception,
positive holes, and zero holes with a retained exception, respectively.
Regimes I, Z and P use each of the 68 basic modes once. Regime E uses
every basic mode and every choice of one retained center and one of its
two endpoints. The two split-pair modes have 16 retained centers; the
other 66 have 17. Thus regime E has
\[
 2(2\cdot16+66\cdot17)=2308
\]
systems, and the total is $3\cdot68+2308=2512$.

All original variables and all losses are nonnegative. Keep the order
rows $D_i^p,D_i^q$, span rows $E_p,E_q$, branch row $B$, and total-loss
row $\Sigma$ from Appendix~\ref{u96:app:one-empty}. The rows that depend on
the regime are replaced as follows:
\begin{center}\small
\begin{tabular}{ll}
\toprule
Row&Constraint\\
\midrule
$G_p,G_q$&$-R_p+t\leq-a_0$, $-R_q-t\leq-b_0$\\
$H_p$&$2M_p+6R_p-2S_p+s_0t\leq-A_p$\\
$H_q$&$2M_q+6R_q-2S_q-s_0t\leq-A_q$\\
$Z,J$&$\delta\leq\delta_{\max}$, $t\leq t_{\max}$\\
$L_v$ (positive)&$-v-l_v-t\leq-(b_0+5)$\\
$U_v$ (positive)&$v-u_v-S_p+t\leq-(a_0+4)$\\
$L_v$ (negative)&$-v-l_v+t\leq-(a_0+5)$\\
$U_v$ (negative)&$v-u_v-S_q-t\leq-(b_0+4)$\\
\bottomrule
\end{tabular}
\end{center}
The mass rows use $2T_r\geq85r-1764$, the tangent at 42. In regime P
one also adds $\varepsilon_p\geq a-2$ and $\varepsilon_q\geq b-2$;
this accounts for its slope 87 and its larger constants.

In regime E, replace the \emph{single} quantile row at the designated
exceptional endpoint, and add a row $X$ for its unavoidable loss:
\begin{center}\small
\begin{tabular}{lll}
\toprule
Exceptional endpoint&Replacement quantile row&Additional row $X$\\
\midrule
Positive, lower&$-v-l_v-M_p\leq-d$&$-l_v-t\leq-(b_0+5)$\\
Positive, upper&$v-u_v-S_p-M_q\leq-(d-1)$&$-u_v+t\leq-(a_0+5)$\\
Negative, lower&$-v-l_v-M_q\leq-d$&$-l_v+t\leq-(a_0+5)$\\
Negative, upper&$v-u_v-S_q-M_p\leq-(d-1)$&$-u_v-t\leq-(b_0+5)$\\
\bottomrule
\end{tabular}
\end{center}
The designated quantile row keeps its original name $L_v$ or $U_v$.
Its opposite endpoint row remains unchanged. No inequality asserting the
actual location $v<0$ or $v>S$ is imposed: omitting it makes the model
weaker, so a contradiction still excludes the original configuration.

\subsection{Exact verification rule}
For each of the systems just specified, write its rows as $Ax\leq b$,
where $x\geq0$, and regard $f$ as a coefficient row. A candidate in the
table below specifies an objective multiplier $D\geq0$ and a nonnegative
integer combination $w$ of named rows. Discard that candidate if it calls
for a row absent from the current mode. Otherwise form, using integers,
\[
 q=Df+wA,\qquad h=wb.
\]
It certifies the current mode if $q\geq0$ coefficientwise and either
\[
 D=0,\ h<0,
 \qquad\hbox{or}\qquad
 D>0,\ d-h/D>F_{\rm upper}.
\]
The proof is the same elementary row-combination argument as in
Appendix~\ref{u96:app:one-empty}. Importantly, a candidate need not pass in
every mode; the assertion is that at least one passes in each mode.

The following pseudocode specifies the full check without relying on a
stored assignment of certificates:
\begin{quote}\small\ttfamily
for regime in (I, Z, P, E):\\
\hspace*{1em}for basic\_mode in the 68 generated modes:\\
\hspace*{2em}for endpoint in this regime's endpoint choices:\\
\hspace*{3em}construct A, b, f from the printed rows\\
\hspace*{3em}passed = false\\
\hspace*{3em}for (D, w) in the 42 printed candidates:\\
\hspace*{4em}if all named rows exist and the exact test passes:\\
\hspace*{5em}passed = true; break\\
\hspace*{3em}assert passed
\end{quote}
For I, Z and P there is a single choice ``no exception.'' For E the
choices are the $2|\mathcal C|$ retained endpoints. The delivered verifier
reconstructs the same exhaustive list and checks one recorded successful
integer combination for each system. Its checked outcomes are:
\begin{center}\small
\begin{tabular}{lrrrr}
\toprule
Regime&Systems&\shortstack{Contradiction\\combinations}&\shortstack{Lower-bound\\combinations}&\shortstack{Minimum\\certified $F$}\\
\midrule
I&68&8&60&$9378/5>1874$\\
Z&68&8&60&$9103/5>1800$\\
P&68&29&39&$1910>1856$\\
E&2308&351&1957&$10957/6>1800$\\
\bottomrule
\end{tabular}
\end{center}
These are counts of certificate types, not claims that a lower-bound
system is feasible. All 2512 interpretations as feasible sets are excluded.

\subsection{The forty-two integer combinations}
Row labels, including $L_{1p},L_{2p},L_{2q}$ and their upper counterparts,
have the same slot meanings as in Appendix~\ref{u96:app:one-empty}. The symbol
$X$ denotes the extra exception-loss row and is absent outside regime E.
All omitted row weights are zero.
\begingroup\footnotesize
\setlength{\tabcolsep}{4pt}
\renewcommand{\arraystretch}{1.35}
\begin{longtable}{rrp{13.8cm}}
\toprule
ID&$D$&Nonnegative integer row combination\\
\midrule
\endfirsthead
\toprule
ID&$D$&Nonnegative integer row combination\\
\midrule
\endhead
1&10&\( 8D_{3}^{p}+\allowbreak 12D_{4}^{p}+\allowbreak 8D_{3}^{q}+\allowbreak 16D_{4}^{q}+\allowbreak 6E_q+\allowbreak 4G_p+\allowbreak 16G_q+\allowbreak 4H_p+\allowbreak 5H_q+\allowbreak 93J+\allowbreak 2U_{p0}+\allowbreak 2L_{p8}+\allowbreak 4L_{q6}+\allowbreak 4L_{q8}+\allowbreak 4L_{1p}+\allowbreak 10\Sigma \)\\
2&0&\( D_{4}^{q}+\allowbreak Z+\allowbreak B+\allowbreak L_{q2}+\allowbreak L_{1p}+\allowbreak L_{2q}+\allowbreak \Sigma \)\\
3&0&\( D_{4}^{p}+\allowbreak E_p+\allowbreak G_p+\allowbreak Z+\allowbreak L_{q5}+\allowbreak L_{1p}+\allowbreak L_{2p}+\allowbreak \Sigma \)\\
4&6&\( 12D_{3}^{p}+\allowbreak 18D_{4}^{p}+\allowbreak 10E_p+\allowbreak 12D_{3}^{q}+\allowbreak 24D_{4}^{q}+\allowbreak 12E_q+\allowbreak 10G_p+\allowbreak 18G_q+\allowbreak 2H_p+\allowbreak 3H_q+\allowbreak 93J+\allowbreak 2U_{p0}+\allowbreak 2L_{p9}+\allowbreak 6L_{q6}+\allowbreak 6L_{q9}+\allowbreak 6L_{2p}+\allowbreak 6\Sigma \)\\
5&12&\( 10D_{2}^{p}+\allowbreak 12D_{3}^{p}+\allowbreak 14D_{4}^{p}+\allowbreak 12D_{3}^{q}+\allowbreak 24D_{4}^{q}+\allowbreak 12E_q+\allowbreak 6G_p+\allowbreak 24G_q+\allowbreak 5H_p+\allowbreak 6H_q+\allowbreak 95J+\allowbreak 2U_{p0}+\allowbreak 8L_{q5}+\allowbreak 4L_{1p}+\allowbreak 4L_{2q}+\allowbreak 12\Sigma \)\\
6&6&\( 12D_{3}^{p}+\allowbreak 24D_{4}^{p}+\allowbreak 12E_p+\allowbreak 12D_{3}^{q}+\allowbreak 18D_{4}^{q}+\allowbreak 10E_q+\allowbreak 18G_p+\allowbreak 10G_q+\allowbreak 3H_p+\allowbreak 2H_q+\allowbreak 6L_{p6}+\allowbreak 6L_{p9}+\allowbreak 2U_{q0}+\allowbreak 2L_{q9}+\allowbreak 6L_{2q}+\allowbreak 6\Sigma \)\\
7&13&\( 7D_{3}^{p}+\allowbreak 14D_{4}^{p}+\allowbreak 7D_{3}^{q}+\allowbreak 14D_{4}^{q}+\allowbreak 10G_p+\allowbreak 10G_q+\allowbreak 6H_p+\allowbreak 6H_q+\allowbreak U_{p0}+\allowbreak 5L_{p5}+\allowbreak U_{q0}+\allowbreak 2L_{q5}+\allowbreak 6L_{q8}+\allowbreak 3L_{2p}+\allowbreak 13\Sigma \)\\
8&13&\( 7D_{3}^{p}+\allowbreak 14D_{4}^{p}+\allowbreak 7D_{3}^{q}+\allowbreak 14D_{4}^{q}+\allowbreak 10G_p+\allowbreak 10G_q+\allowbreak 6H_p+\allowbreak 6H_q+\allowbreak U_{p0}+\allowbreak 2L_{p5}+\allowbreak 6L_{p8}+\allowbreak U_{q0}+\allowbreak 5L_{q5}+\allowbreak 3L_{2q}+\allowbreak 13\Sigma \)\\
9&11&\( 6D_{3}^{p}+\allowbreak 12D_{4}^{p}+\allowbreak 6D_{3}^{q}+\allowbreak 12D_{4}^{q}+\allowbreak 8G_p+\allowbreak 8G_q+\allowbreak 5H_p+\allowbreak 5H_q+\allowbreak U_{p1}+\allowbreak 6L_{p8}+\allowbreak U_{q0}+\allowbreak 5L_{q5}+\allowbreak L_{1p}+\allowbreak 2L_{2q}+\allowbreak 11\Sigma \)\\
10&2&\( 6D_{3}^{p}+\allowbreak 24D_{4}^{p}+\allowbreak 20E_p+\allowbreak 18D_{3}^{q}+\allowbreak 36D_{4}^{q}+\allowbreak 22E_q+\allowbreak 16G_p+\allowbreak 24G_q+\allowbreak H_q+\allowbreak 10Z+\allowbreak 93J+\allowbreak 2U_{p0}+\allowbreak 2L_{p9}+\allowbreak 12L_{q7}+\allowbreak 6L_{q9}+\allowbreak 6L_{1p}+\allowbreak 12\Sigma \)\\
11&14&\( 12D_{3}^{p}+\allowbreak 16D_{4}^{p}+\allowbreak 12D_{3}^{q}+\allowbreak 24D_{4}^{q}+\allowbreak 10E_q+\allowbreak 8G_p+\allowbreak 24G_q+\allowbreak 6H_p+\allowbreak 7H_q+\allowbreak 93J+\allowbreak 2U_{p1}+\allowbreak 2L_{p9}+\allowbreak 8L_{q6}+\allowbreak 4L_{q9}+\allowbreak 4L_{1p}+\allowbreak 14\Sigma \)\\
12&10&\( 8D_{3}^{p}+\allowbreak 12D_{4}^{p}+\allowbreak 8D_{3}^{q}+\allowbreak 16D_{4}^{q}+\allowbreak 2E_q+\allowbreak 4G_p+\allowbreak 12G_q+\allowbreak 4H_p+\allowbreak 5H_q+\allowbreak 87J+\allowbreak 2U_{p0}+\allowbreak 2L_{p8}+\allowbreak 4L_{q6}+\allowbreak 4L_{q8}+\allowbreak 4L_{1p}+\allowbreak 10\Sigma+\allowbreak 6X \)\\
13&0&\( D_{4}^{p}+\allowbreak Z+\allowbreak L_{q5}+\allowbreak L_{1p}+\allowbreak L_{2p}+\allowbreak \Sigma \)\\
14&6&\( 8D_{3}^{p}+\allowbreak 12D_{4}^{p}+\allowbreak 8D_{3}^{q}+\allowbreak 16D_{4}^{q}+\allowbreak 10E_q+\allowbreak 16G_q+\allowbreak 2H_p+\allowbreak 3H_q+\allowbreak 95J+\allowbreak 2U_{p0}+\allowbreak 2L_{p8}+\allowbreak 4L_{q6}+\allowbreak 4L_{q8}+\allowbreak 4L_{1p}+\allowbreak 6\Sigma+\allowbreak 2X \)\\
15&1432&\( 1238D_{3}^{p}+\allowbreak 2476D_{4}^{p}+\allowbreak 862E_p+\allowbreak 1238D_{3}^{q}+\allowbreak 2476D_{4}^{q}+\allowbreak 1748G_p+\allowbreak 922G_q+\allowbreak 625H_p+\allowbreak 631H_q+\allowbreak 182U_{p0}+\allowbreak 364L_{p5}+\allowbreak 1056L_{p8}+\allowbreak 170U_{q0}+\allowbreak 874L_{q5}+\allowbreak 534L_{2q}+\allowbreak 1432\Sigma+\allowbreak 558X \)\\
16&0&\( 2D_{4}^{p}+\allowbreak 2E_p+\allowbreak D_{4}^{q}+\allowbreak 2G_p+\allowbreak Z+\allowbreak L_{q5}+\allowbreak L_{1p}+\allowbreak L_{2p}+\allowbreak \Sigma \)\\
17&10&\( 8D_{3}^{p}+\allowbreak 12D_{4}^{p}+\allowbreak 8D_{3}^{q}+\allowbreak 16D_{4}^{q}+\allowbreak 6E_q+\allowbreak 4G_p+\allowbreak 16G_q+\allowbreak 4H_p+\allowbreak 5H_q+\allowbreak 95J+\allowbreak 2U_{p0}+\allowbreak 2L_{p8}+\allowbreak 4L_{q6}+\allowbreak 4L_{q8}+\allowbreak 4L_{1p}+\allowbreak 10\Sigma \)\\
18&8&\( 8D_{3}^{p}+\allowbreak 12D_{4}^{p}+\allowbreak 8D_{3}^{q}+\allowbreak 16D_{4}^{q}+\allowbreak 8E_q+\allowbreak 2G_p+\allowbreak 16G_q+\allowbreak 3H_p+\allowbreak 4H_q+\allowbreak 99J+\allowbreak 2U_{p0}+\allowbreak 2L_{p8}+\allowbreak 4L_{q6}+\allowbreak 4L_{q8}+\allowbreak 4L_{1p}+\allowbreak 8\Sigma+\allowbreak 6X \)\\
19&657&\( 657D_{3}^{p}+\allowbreak 1314D_{4}^{p}+\allowbreak 562E_p+\allowbreak 657D_{3}^{q}+\allowbreak 1314D_{4}^{q}+\allowbreak 934G_p+\allowbreak 378G_q+\allowbreak 281H_p+\allowbreak 282H_q+\allowbreak 95U_{p0}+\allowbreak 471L_{p5}+\allowbreak 93U_{q0}+\allowbreak 186L_{q5}+\allowbreak 564L_{q8}+\allowbreak 281L_{2p}+\allowbreak 657\Sigma+\allowbreak 93X \)\\
20&1058&\( 744D_{3}^{p}+\allowbreak 1190D_{4}^{p}+\allowbreak 744D_{3}^{q}+\allowbreak 1488D_{4}^{q}+\allowbreak 662G_p+\allowbreak 710G_q+\allowbreak 463H_p+\allowbreak 471H_q+\allowbreak 132U_{p1}+\allowbreak 446L_{p8}+\allowbreak 116U_{q0}+\allowbreak 314L_{q5}+\allowbreak 298L_{q6}+\allowbreak 132L_{1p}+\allowbreak 248L_{2q}+\allowbreak 1058\Sigma+\allowbreak 744X \)\\
21&2&\( 12D_{3}^{p}+\allowbreak 18D_{4}^{p}+\allowbreak 12E_p+\allowbreak 12D_{3}^{q}+\allowbreak 24D_{4}^{q}+\allowbreak 22E_q+\allowbreak 8G_p+\allowbreak 24G_q+\allowbreak H_q+\allowbreak 4Z+\allowbreak 97J+\allowbreak 2U_{p0}+\allowbreak 2L_{p9}+\allowbreak 6L_{q6}+\allowbreak 6L_{q9}+\allowbreak 6L_{2p}+\allowbreak 6\Sigma+\allowbreak 4X \)\\
22&6&\( 12D_{3}^{p}+\allowbreak 24D_{4}^{p}+\allowbreak 18E_p+\allowbreak 12D_{3}^{q}+\allowbreak 18D_{4}^{q}+\allowbreak 4E_q+\allowbreak 24G_p+\allowbreak 4G_q+\allowbreak 3H_p+\allowbreak 2H_q+\allowbreak 6L_{p6}+\allowbreak 6L_{p9}+\allowbreak 2U_{q0}+\allowbreak 2L_{q9}+\allowbreak 6L_{2q}+\allowbreak 6\Sigma \)\\
23&12&\( 10D_{2}^{p}+\allowbreak 12D_{3}^{p}+\allowbreak 14D_{4}^{p}+\allowbreak 12D_{3}^{q}+\allowbreak 24D_{4}^{q}+\allowbreak 8E_q+\allowbreak 6G_p+\allowbreak 20G_q+\allowbreak 5H_p+\allowbreak 6H_q+\allowbreak 85J+\allowbreak 2U_{p0}+\allowbreak 8L_{q5}+\allowbreak 4L_{1p}+\allowbreak 4L_{2q}+\allowbreak 12\Sigma+\allowbreak 8X \)\\
24&8&\( 12D_{3}^{p}+\allowbreak 16D_{4}^{p}+\allowbreak 6E_p+\allowbreak 12D_{3}^{q}+\allowbreak 24D_{4}^{q}+\allowbreak 12E_q+\allowbreak 8G_p+\allowbreak 20G_q+\allowbreak 3H_p+\allowbreak 4H_q+\allowbreak 89J+\allowbreak 2U_{p1}+\allowbreak 2L_{p9}+\allowbreak 8L_{q6}+\allowbreak 4L_{q9}+\allowbreak 4L_{1p}+\allowbreak 8\Sigma+\allowbreak 4X \)\\
25&1432&\( 1238D_{3}^{p}+\allowbreak 2476D_{4}^{p}+\allowbreak 1238D_{3}^{q}+\allowbreak 2476D_{4}^{q}+\allowbreak 862E_q+\allowbreak 922G_p+\allowbreak 1748G_q+\allowbreak 631H_p+\allowbreak 625H_q+\allowbreak 170U_{p0}+\allowbreak 874L_{p5}+\allowbreak 182U_{q0}+\allowbreak 364L_{q5}+\allowbreak 1056L_{q8}+\allowbreak 534L_{2p}+\allowbreak 1432\Sigma+\allowbreak 558X \)\\
26&2252&\( 1310D_{3}^{p}+\allowbreak 2620D_{4}^{p}+\allowbreak 1310D_{3}^{q}+\allowbreak 2620D_{4}^{q}+\allowbreak 162E_q+\allowbreak 1766G_p+\allowbreak 1796G_q+\allowbreak 1045H_p+\allowbreak 1023H_q+\allowbreak 162U_{p0}+\allowbreak 898L_{p5}+\allowbreak 206U_{q0}+\allowbreak 412L_{q5}+\allowbreak 1104L_{q8}+\allowbreak 574L_{2p}+\allowbreak 2252\Sigma+\allowbreak 2046X \)\\
27&1868&\( 1302D_{3}^{p}+\allowbreak 2604D_{4}^{p}+\allowbreak 1302D_{3}^{q}+\allowbreak 2604D_{4}^{q}+\allowbreak 538E_q+\allowbreak 1358G_p+\allowbreak 1812G_q+\allowbreak 849H_p+\allowbreak 835H_q+\allowbreak 170U_{p0}+\allowbreak 906L_{p5}+\allowbreak 198U_{q0}+\allowbreak 396L_{q5}+\allowbreak 1104L_{q8}+\allowbreak 566L_{2p}+\allowbreak 1868\Sigma+\allowbreak 1302X \)\\
28&111&\( 69D_{3}^{p}+\allowbreak 138D_{4}^{p}+\allowbreak 69D_{3}^{q}+\allowbreak 138D_{4}^{q}+\allowbreak 16E_q+\allowbreak 84G_p+\allowbreak 94G_q+\allowbreak 51H_p+\allowbreak 50H_q+\allowbreak 9U_{p0}+\allowbreak 18L_{p5}+\allowbreak 60L_{p8}+\allowbreak 11U_{q0}+\allowbreak 51L_{q5}+\allowbreak 29L_{2q}+\allowbreak 111\Sigma+\allowbreak 93X \)\\
29&657&\( 657D_{3}^{p}+\allowbreak 1314D_{4}^{p}+\allowbreak 657D_{3}^{q}+\allowbreak 1314D_{4}^{q}+\allowbreak 562E_q+\allowbreak 378G_p+\allowbreak 934G_q+\allowbreak 282H_p+\allowbreak 281H_q+\allowbreak 93U_{p0}+\allowbreak 186L_{p5}+\allowbreak 564L_{p8}+\allowbreak 95U_{q0}+\allowbreak 471L_{q5}+\allowbreak 281L_{2q}+\allowbreak 657\Sigma+\allowbreak 93X \)\\
30&2260&\( 1322D_{3}^{p}+\allowbreak 2644D_{4}^{p}+\allowbreak 1322D_{3}^{q}+\allowbreak 2644D_{4}^{q}+\allowbreak 170E_q+\allowbreak 1750G_p+\allowbreak 1788G_q+\allowbreak 1045H_p+\allowbreak 1023H_q+\allowbreak 170U_{p0}+\allowbreak 340L_{p5}+\allowbreak 1152L_{p8}+\allowbreak 214U_{q0}+\allowbreak 982L_{q5}+\allowbreak 554L_{2q}+\allowbreak 2260\Sigma+\allowbreak 2046X \)\\
31&5&\( 6D_{3}^{p}+\allowbreak 12D_{4}^{p}+\allowbreak 6D_{3}^{q}+\allowbreak 12D_{4}^{q}+\allowbreak 6E_q+\allowbreak 2G_p+\allowbreak 8G_q+\allowbreak 2H_p+\allowbreak 2H_q+\allowbreak Z+\allowbreak U_{p1}+\allowbreak 6L_{p8}+\allowbreak U_{q0}+\allowbreak 5L_{q5}+\allowbreak L_{1p}+\allowbreak 2L_{2q}+\allowbreak 6\Sigma \)\\
32&1880&\( 1128D_{3}^{p}+\allowbreak 2256D_{4}^{p}+\allowbreak 1128D_{3}^{q}+\allowbreak 2256D_{4}^{q}+\allowbreak 170E_q+\allowbreak 1370G_p+\allowbreak 1432G_q+\allowbreak 855H_p+\allowbreak 837H_q+\allowbreak 170U_{p1}+\allowbreak 1128L_{p8}+\allowbreak 206U_{q0}+\allowbreak 958L_{q5}+\allowbreak 170L_{1p}+\allowbreak 376L_{2q}+\allowbreak 1880\Sigma+\allowbreak 1674X \)\\
33&204&\( 120D_{3}^{p}+\allowbreak 240D_{4}^{p}+\allowbreak 120D_{3}^{q}+\allowbreak 240D_{4}^{q}+\allowbreak 14E_q+\allowbreak 150G_p+\allowbreak 152G_q+\allowbreak 93H_p+\allowbreak 91H_q+\allowbreak 18U_{p1}+\allowbreak 120L_{p8}+\allowbreak 22U_{q0}+\allowbreak 102L_{q5}+\allowbreak 18L_{1p}+\allowbreak 40L_{2q}+\allowbreak 204\Sigma+\allowbreak 186X \)\\
34&8&\( 10D_{2}^{p}+\allowbreak 12D_{3}^{p}+\allowbreak 14D_{4}^{p}+\allowbreak 12D_{3}^{q}+\allowbreak 24D_{4}^{q}+\allowbreak 16E_q+\allowbreak 2G_p+\allowbreak 24G_q+\allowbreak 3H_p+\allowbreak 4H_q+\allowbreak 97J+\allowbreak 2U_{p0}+\allowbreak 8L_{q5}+\allowbreak 4L_{1p}+\allowbreak 4L_{2q}+\allowbreak 8\Sigma+\allowbreak 4X \)\\
35&6&\( 12D_{3}^{p}+\allowbreak 18D_{4}^{p}+\allowbreak 4E_p+\allowbreak 12D_{3}^{q}+\allowbreak 24D_{4}^{q}+\allowbreak 18E_q+\allowbreak 4G_p+\allowbreak 24G_q+\allowbreak 2H_p+\allowbreak 3H_q+\allowbreak 93J+\allowbreak 2U_{p0}+\allowbreak 2L_{p9}+\allowbreak 6L_{q6}+\allowbreak 6L_{q9}+\allowbreak 6L_{2p}+\allowbreak 6\Sigma \)\\
36&4&\( 8D_{3}^{p}+\allowbreak 16D_{4}^{p}+\allowbreak 10E_p+\allowbreak 12D_{3}^{q}+\allowbreak 24D_{4}^{q}+\allowbreak 20E_q+\allowbreak 8G_p+\allowbreak 24G_q+\allowbreak H_p+\allowbreak 2H_q+\allowbreak 95J+\allowbreak 2U_{p1}+\allowbreak 2L_{p9}+\allowbreak 4L_{q6}+\allowbreak 4L_{q7}+\allowbreak 4L_{q9}+\allowbreak 4L_{1p}+\allowbreak 4\Sigma \)\\
37&2&\( 4D_{3}^{p}+\allowbreak 16D_{4}^{p}+\allowbreak 12E_p+\allowbreak 12D_{3}^{q}+\allowbreak 24D_{4}^{q}+\allowbreak 22E_q+\allowbreak 8G_p+\allowbreak 24G_q+\allowbreak H_q+\allowbreak 6Z+\allowbreak 85J+\allowbreak 2U_{p1}+\allowbreak 2L_{p9}+\allowbreak 8L_{q7}+\allowbreak 4L_{q9}+\allowbreak 4L_{1p}+\allowbreak 8\Sigma+\allowbreak 8X \)\\
38&4&\( 8D_{3}^{p}+\allowbreak 16D_{4}^{p}+\allowbreak 6E_p+\allowbreak 12D_{3}^{q}+\allowbreak 24D_{4}^{q}+\allowbreak 20E_q+\allowbreak 4G_p+\allowbreak 24G_q+\allowbreak H_p+\allowbreak 2H_q+\allowbreak 93J+\allowbreak 2U_{p1}+\allowbreak 2L_{p9}+\allowbreak 4L_{q6}+\allowbreak 4L_{q7}+\allowbreak 4L_{q9}+\allowbreak 4L_{1p}+\allowbreak 4\Sigma \)\\
39&2&\( 6D_{3}^{p}+\allowbreak 24D_{4}^{p}+\allowbreak 18E_p+\allowbreak 18D_{3}^{q}+\allowbreak 36D_{4}^{q}+\allowbreak 34E_q+\allowbreak 14G_p+\allowbreak 36G_q+\allowbreak H_q+\allowbreak 10Z+\allowbreak 103J+\allowbreak 2U_{p0}+\allowbreak 2L_{p9}+\allowbreak 12L_{q7}+\allowbreak 6L_{q9}+\allowbreak 6L_{1p}+\allowbreak 12\Sigma+\allowbreak 10X \)\\
40&2&\( 6D_{3}^{p}+\allowbreak 24D_{4}^{p}+\allowbreak 14E_p+\allowbreak 18D_{3}^{q}+\allowbreak 36D_{4}^{q}+\allowbreak 34E_q+\allowbreak 10G_p+\allowbreak 36G_q+\allowbreak H_q+\allowbreak 10Z+\allowbreak 99J+\allowbreak 2U_{p0}+\allowbreak 2L_{p9}+\allowbreak 12L_{q7}+\allowbreak 6L_{q9}+\allowbreak 6L_{1p}+\allowbreak 12\Sigma+\allowbreak 6X \)\\
41&4&\( 2D_{2}^{p}+\allowbreak 8D_{3}^{p}+\allowbreak 14D_{4}^{p}+\allowbreak 8E_p+\allowbreak 12D_{3}^{q}+\allowbreak 24D_{4}^{q}+\allowbreak 20E_q+\allowbreak 6G_p+\allowbreak 24G_q+\allowbreak H_p+\allowbreak 2H_q+\allowbreak 89J+\allowbreak 2U_{p0}+\allowbreak 4L_{q6}+\allowbreak 4L_{q7}+\allowbreak 4L_{1p}+\allowbreak 4L_{2q}+\allowbreak 4\Sigma+\allowbreak 4X \)\\
42&2&\( 8D_{3}^{p}+\allowbreak 16D_{4}^{p}+\allowbreak 10E_p+\allowbreak 12D_{3}^{q}+\allowbreak 24D_{4}^{q}+\allowbreak 22E_q+\allowbreak 6G_p+\allowbreak 24G_q+\allowbreak H_q+\allowbreak 2Z+\allowbreak 95J+\allowbreak 2U_{p1}+\allowbreak 2L_{p9}+\allowbreak 4L_{q6}+\allowbreak 4L_{q7}+\allowbreak 4L_{q9}+\allowbreak 4L_{1p}+\allowbreak 4\Sigma+\allowbreak 2X \)\\
\bottomrule
\end{longtable}
\endgroup
All multipliers are integers; the largest displayed one is 2644. The
finite certificate check uses no approximation, tolerance, or optimization
status flag. The exact source matrices and the checked assignment are
also supplied as readable Python and JSON in the verification archive.

\clearpage
\section[Certificate specification for upper bound 99]{Complete certificate specification for the upper bound 99}\label{u96:app:upper99}
This appendix gives the complete rule generating the 59206 systems and
all 52 integer combinations that certify them. The JSON file assigns a
successful combination to each system. That assignment is optional:
trying the printed combinations in order is another complete verification
procedure. The source model contains no optimization dependency.

\subsection{The 64 effective-hole rows}
For each row of the table use $(r,C_7,7)$ and $(r,C_9,9)$ in
Lemma~\ref{u96:lem:hole-parameter}. For either $\beta\in\{7,9\}$ define
\[
\begin{gathered}
 K=\left\lfloor\frac{3+\sqrt{9+8(C+1)}}2\right\rfloor,
 \qquad D=C+\max(3,K),\\
 M=\max\left(2\beta,1+\left\lfloor
       \frac{1+\sqrt{9+8(D+1+6\beta)}}2\right\rfloor\right),
 \qquad J=\max(K+1,2\beta+5,M).
\end{gathered}
\]
Every entry below satisfies $(r-J+1)\beta>D$ and $T_{J-1}-1>D$.
Thus $2\beta<h<|\FS|-1$ implies $\mu\geq r+C+1$.
\begin{center}\small
\begin{tabular}{rrr@{\qquad}rrr}
\toprule
$r$&$C_7$&$C_9$&$r$&$C_7$&$C_9$\\
\midrule
26&44&26&42&135&160\\
27&51&34&43&142&168\\
28&57&42&44&147&176\\
29&63&51&45&149&177\\
30&70&59&46&155&186\\
31&76&67&47&162&194\\
32&83&75&48&167&198\\
33&89&84&49&168&203\\
34&96&92&50&175&211\\
35&102&101&51&182&220\\
36&109&109&52&187&221\\
37&110&117&53&188&228\\
38&116&126&54&195&237\\
39&122&134&55&202&245\\
40&128&143&56&208&246\\
41&129&151&57&209&254\\
\bottomrule
\end{tabular}
\end{center}

\subsection{Centers, units, modes, and variables}
Use the twenty labeled centers and the two branch definitions in
Appendix~\ref{u96:app:one-empty}. Explicitly, with
$c_p=p_1-p_2-p_3-p_4$,
$\alpha_p=p_2-p_3-p_4$ and
$\beta_p=2p_2-p_1-p_3-p_4$, the positive list is
\[
 (c_p+p_1,c_p+p_2,c_p+p_3,c_p+p_4,
 c_p-q_1,c_p-q_2,c_p-q_3,c_p-q_4,\alpha_p,\beta_p).
\]
The negative list is obtained by interchanging $p,q$. Their labels are
$p_0,\ldots,p_9$ and $q_0,\ldots,q_9$; these labels are distinct from
the variables $p_1,\ldots,p_4$ and $q_1,\ldots,q_4$.
The first paired unit uses index 4. The second uses index 9 in branch A
($\beta_p+\beta_q\leq0$) and index 5 in branch B
($\beta_p+\beta_q>0$). There are sixteen other, unpaired units.

Choose one missing unit in Z, P, ZE, PE and ZZ, or any unordered pair of
missing units in EE. Retain every other unpaired center. If the first
pair is present, retain its positive member $\pi_p=c_p-q_1$.
If the second pair is present, retain either its positive member
$\sigma_p$ or its negative member $\sigma_q$, with both choices considered.
There is one exception: in branch A, when $p_5$ or $q_5$ is missing,
drop the second pair entirely from the retained list and do not duplicate
the mode. Likewise there is only one choice when the second pair is
itself missing. Thus the numbers of modes are 68 with one missing unit
and 547 with two. With one missing unit, 66 modes retain 17 centers and
two retain 16; choosing an exceptional center and an endpoint gives 2308
possibilities for each fixed $a$.

For each regime take every integer $a$ in the interval in
Section~\ref{u96:sec:upper99}, put $b=83-a$, $k=41-\min(a,b)$ and $I=k(k+1)$.
In ZE and PE additionally choose each retained center in turn and each
of its lower and upper endpoints. There are no exceptional-location
choices in ZZ or EE.

All variables $p_i,q_i,M_p,M_q,R_p,R_q,S_p,S_q,\delta,l_v,u_v$ are
nonnegative real variables. The implementation also includes a harmless
variable $t\geq0$ and a row $J:t\leq0$; all other coefficients of $t$
are zero. Unused loss variables can be set to zero.
Write
\[
 f=S_p+S_q+\delta-2(M_p+M_q)-2(R_p+R_q).
\]
Then $F=92+f$ except in EE, where $F=184+f$ and $\delta=0$.

\subsection{All coefficient rows}
In the first five regimes the order rows are
\[
 D_i^p:p_{i+1}-p_i\leq-1\ (1\leq i\leq3),\quad
 D_4^p:M_p-p_4\leq-1,
\]
with reflected $D_i^q$. In EE replace their right sides by zero.
In all regimes retain
\[
 E_p:R_p-M_p\leq-4,\quad E_q:R_q-M_q\leq-4,\quad
 G_p:-R_p\leq-a,\quad G_q:-R_q\leq-b.
\]
The exact mass rows are
\[
 H_p:2M_p+6R_p-2S_p\leq-2T_a-12-2\epsilon_p,
 \quad
 H_q:2M_q+6R_q-2S_q\leq-2T_b-12-2\epsilon_q,
\]
where $(\epsilon_p,\epsilon_q)=(a-2,b-2)$ in P and PE and $(0,0)$
otherwise. The budget row $Z$ is $\delta\leq92-I$ in Z, ZE and ZZ,
$\delta\leq65-I$ in P and PE, and $\delta\leq0$ in EE.
The branch row $B$ is $\beta_p+\beta_q\leq0$ in A,
$-\beta_p-\beta_q\leq-1$ in B outside EE, and
$-\beta_p-\beta_q\leq0$ in branch B of EE.
The total-loss row is
\[
 \Sigma:\sum_{v\text{ retained}}(l_v+u_v)-\delta\leq0.
\]
The ordinary quantile rows for positive centers are
\[
 L_v:-v-l_v\leq-b-5,\qquad U_v:v-u_v-S_p\leq-a-4,
\]
with $a,b$ and $p,q$ interchanged for negative centers.
In ZZ increase both right sides by 1 for every retained center, and add
$K_p:R_p\leq a+1$, $K_q:R_q\leq b+1$.

In ZE and PE replace only the designated exceptional endpoint row, and
add a row $X$ as follows:
\begin{center}\small
\begin{tabular}{lll}
\toprule
Endpoint&Replacement row&Row $X$\\
\midrule
Positive, lower&$-v-l_v-M_p\leq-92$&$-l_v\leq-b-5$\\
Positive, upper&$v-u_v-S_p-M_q\leq-91$&$-u_v\leq-a-5$\\
Negative, lower&$-v-l_v-M_q\leq-92$&$-l_v\leq-a-5$\\
Negative, upper&$v-u_v-S_q-M_p\leq-91$&$-u_v\leq-b-5$\\
\bottomrule
\end{tabular}
\end{center}
The replaced row keeps its name $L_v$ or $U_v$; the opposite row is
unchanged. The actual out-of-hull location inequality is omitted, which
is a relaxation. In EE replace all quantile rows by
\[
 L_v:-v-l_v+M_q\leq-1,\qquad
 U_v:v-u_v-S_p+M_p\leq0
\]
for positive centers, and their reflected versions for negative centers.
Since $\delta=0$, the total-loss row makes all retained losses zero.
These are exactly the full-translate constraints proved in the text.

\subsection{The integer verification rule and multiplier pool}
For a generated system $Ax\leq b$, $x\geq0$, a table entry gives $D\geq0$
and a nonnegative row combination $w$. A term naming a row absent from
the current mode makes that entry unavailable. Otherwise compute
\[
 q=Df+wA,\qquad h=wb.
\]
If $q$ is coordinatewise nonnegative, then for $D>0$ we have
$F\geq d_0-h/D$, where $d_0=92$ or 184. It is a certificate when this is
strictly greater than 1782, or 1838 in P and PE. For $D=0$, the condition
$h<0$ is an inconsistency certificate. Every entry below has integer
weights; there are 774 nonzero weights in all.
The verifier generates all modes independently of the assignment file,
checks the recorded case sequence, and evaluates these identities with
arbitrary-precision integers and exact fractions.

\begingroup\footnotesize
\setlength{\tabcolsep}{4pt}
\begin{longtable}{rrp{12.7cm}}
\toprule
ID&$D$&Nonzero row weights (omitted weights are zero)\\
\midrule\endfirsthead
\toprule ID&$D$&Nonzero row weights (continued)\\\midrule\endhead
\bottomrule\endfoot
1&10&\( 8D_3^p + 12D_4^p + 8D_3^q + 16D_4^q \)\newline\( + 6E_q + 4G_p + 16G_q + 4H_p \)\newline\( + 5H_q + 2U_{p_{0}} + 2L_{p_{8}} + 4L_{q_{6}} \)\newline\( + 4L_{q_{8}} + 4L_{\pi_p} + 10\Sigma \)\\[3pt]
2&0&\( Z + L_{q_{5}} + L_{\pi_p} + L_{\sigma_p} \)\newline\( + \Sigma \)\\[3pt]
3&2&\( 4D_3^p + 16D_4^p + 12E_p + 12D_3^q \)\newline\( + 24D_4^q + 22E_q + 8G_p + 24G_q \)\newline\( + H_q + 6Z + 2U_{p_{1}} + 2L_{p_{9}} \)\newline\( + 8L_{q_{7}} + 4L_{q_{9}} + 4L_{\pi_p} + 8\Sigma \)\newline\( + 8X \)\\[3pt]
4&6&\( 6D_3^p + 12D_4^p + 6E_p + 6D_3^q \)\newline\( + 12D_4^q + 4E_q + 12G_p + 4G_q \)\newline\( + 3H_p + 2H_q + 6L_{p_{8}} + 2U_{q_{0}} \)\newline\( + 6L_{q_{5}} + 2L_{\sigma_q} + 6\Sigma \)\\[3pt]
5&0&\( Z + \Sigma + X \)\\[3pt]
6&6&\( 12D_3^p + 24D_4^p + 12E_p + 12D_3^q \)\newline\( + 18D_4^q + 10E_q + 18G_p + 10G_q \)\newline\( + 3H_p + 2H_q + 6L_{p_{6}} + 6L_{p_{9}} \)\newline\( + 2U_{q_{0}} + 2L_{q_{9}} + 6L_{\sigma_q} + 6\Sigma \)\\[3pt]
7&12&\( 10D_2^p + 12D_3^p + 14D_4^p + 12D_3^q \)\newline\( + 24D_4^q + 12E_q + 6G_p + 24G_q \)\newline\( + 5H_p + 6H_q + 2U_{p_{0}} + 8L_{q_{5}} \)\newline\( + 4L_{\pi_p} + 4L_{\sigma_q} + 12\Sigma \)\\[3pt]
8&2&\( 12D_3^p + 18D_4^p + 14E_p + 12D_3^q \)\newline\( + 24D_4^q + 16E_q + 10G_p + 18G_q \)\newline\( + H_q + 4Z + 2U_{p_{0}} + 2L_{p_{9}} \)\newline\( + 6L_{q_{6}} + 6L_{q_{9}} + 6L_{\sigma_p} + 6\Sigma \)\\[3pt]
9&6&\( 6D_3^p + 12D_4^p + 4E_p + 6D_3^q \)\newline\( + 12D_4^q + 6E_q + 4G_p + 12G_q \)\newline\( + 2H_p + 3H_q + 2U_{p_{0}} + 6L_{p_{5}} \)\newline\( + 6L_{q_{8}} + 2L_{\sigma_p} + 6\Sigma \)\\[3pt]
10&10&\( 8D_3^p + 12D_4^p + 8D_3^q + 16D_4^q \)\newline\( + 2E_q + 4G_p + 12G_q + 4H_p \)\newline\( + 5H_q + 2U_{p_{0}} + 2L_{p_{8}} + 4L_{q_{6}} \)\newline\( + 4L_{q_{8}} + 4L_{\pi_p} + 10\Sigma + 6X \)\\[3pt]
11&6&\( 8D_3^p + 12D_4^p + 8D_3^q + 16D_4^q \)\newline\( + 10E_q + 16G_q + 2H_p + 3H_q \)\newline\( + 2U_{p_{0}} + 2L_{p_{8}} + 4L_{q_{6}} + 4L_{q_{8}} \)\newline\( + 4L_{\pi_p} + 6\Sigma + 2X \)\\[3pt]
12&12&\( 8D_3^p + 16D_4^p + 2E_p + 8D_3^q \)\newline\( + 16D_4^q + 8G_p + 12G_q + 5H_p \)\newline\( + 6H_q + 2U_{p_{0}} + 4L_{p_{5}} + 6L_{p_{8}} \)\newline\( + 4L_{q_{5}} + 4L_{\sigma_q} + 12\Sigma + 8X \)\\[3pt]
13&0&\( D_4^p + Z + L_{q_{5}} + L_{\pi_p} \)\newline\( + L_{\sigma_p} + \Sigma \)\\[3pt]
14&2&\( 6D_3^p + 24D_4^p + 20E_p + 18D_3^q \)\newline\( + 36D_4^q + 22E_q + 16G_p + 24G_q \)\newline\( + H_q + 10Z + 2U_{p_{0}} + 2L_{p_{9}} \)\newline\( + 12L_{q_{7}} + 6L_{q_{9}} + 6L_{\pi_p} + 12\Sigma \)\\[3pt]
15&8&\( 6D_3^p + 12D_4^p + 2E_p + 6D_3^q \)\newline\( + 12D_4^q + 4G_p + 8G_q + 3H_p \)\newline\( + 4H_q + 2U_{p_{1}} + 6L_{p_{8}} + 4L_{q_{5}} \)\newline\( + 2L_{\pi_p} + 2L_{\sigma_q} + 8\Sigma + 4X \)\\[3pt]
16&2&\( 8D_3^p + 16D_4^p + 14E_p + 8D_3^q \)\newline\( + 16D_4^q + 6E_q + 16G_p + 2G_q \)\newline\( + H_p + 4Z + 4L_{p_{5}} + 2U_{q_{0}} \)\newline\( + 4L_{q_{5}} + 6L_{q_{8}} + 4L_{\sigma_p} + 6\Sigma \)\\[3pt]
17&2&\( 8D_3^p + 16D_4^p + 12E_p + 12D_3^q \)\newline\( + 24D_4^q + 18E_q + 8G_p + 20G_q \)\newline\( + H_q + 2Z + 2U_{p_{1}} + 2L_{p_{9}} \)\newline\( + 4L_{q_{6}} + 4L_{q_{7}} + 4L_{q_{9}} + 4L_{\pi_p} \)\newline\( + 4\Sigma \)\\[3pt]
18&2&\( 12D_3^p + 18D_4^p + 8E_p + 12D_3^q \)\newline\( + 24D_4^q + 22E_q + 4G_p + 24G_q \)\newline\( + H_q + 4Z + 2U_{p_{0}} + 2L_{p_{9}} \)\newline\( + 6L_{q_{6}} + 6L_{q_{9}} + 6L_{\sigma_p} + 6\Sigma \)\\[3pt]
19&0&\( D_4^q + Z + L_{q_{5}} + L_{\pi_p} \)\newline\( + L_{\sigma_p} + \Sigma \)\\[3pt]
20&2&\( 6D_2^p + 12D_3^p + 14D_4^p + 10E_p \)\newline\( + 12D_3^q + 24D_4^q + 18E_q + 6G_p \)\newline\( + 20G_q + H_q + 2Z + 2U_{p_{0}} \)\newline\( + 4L_{q_{5}} + 4L_{q_{6}} + 4L_{\pi_p} + 4L_{\sigma_q} \)\newline\( + 4\Sigma \)\\[3pt]
21&10&\( 4D_3^p + 6D_4^p + 4D_3^q + 8D_4^q \)\newline\( + 2E_q + 4G_p + 12G_q + 4H_p \)\newline\( + 5H_q + 2U_{p_{1}} + 2L_{p_{8}} + 2L_{q_{6}} \)\newline\( + 2L_{q_{8}} + 2L_{\pi_p} + 2\Sigma \)\\[3pt]
22&2&\( 8D_3^p + 16D_4^p + 10E_p + 8D_3^q \)\newline\( + 16D_4^q + 12E_q + 12G_p + 8G_q \)\newline\( + H_p + 4Z + 4L_{p_{5}} + 2U_{q_{0}} \)\newline\( + 4L_{q_{5}} + 6L_{q_{8}} + 4L_{\sigma_p} + 6\Sigma \)\newline\( + 2X \)\\[3pt]
23&2&\( 12D_3^p + 24D_4^p + 22E_p + 12D_3^q \)\newline\( + 18D_4^q + 8E_q + 24G_p + 4G_q \)\newline\( + H_p + 4Z + 6L_{p_{6}} + 6L_{p_{9}} \)\newline\( + 2U_{q_{0}} + 2L_{q_{9}} + 6L_{\sigma_q} + 6\Sigma \)\\[3pt]
24&6&\( 6D_3^p + 12D_4^p + 6D_3^q + 12D_4^q \)\newline\( + 4E_q + 6G_p + 4G_q + 3H_p \)\newline\( + 2H_q + 6L_{p_{8}} + 2U_{q_{0}} + 6L_{q_{5}} \)\newline\( + 2L_{\sigma_q} + 6\Sigma \)\\[3pt]
25&2&\( 8D_3^p + 16D_4^p + 8E_p + 12D_3^q \)\newline\( + 24D_4^q + 22E_q + 4G_p + 24G_q \)\newline\( + H_q + 2Z + 2U_{p_{1}} + 2L_{p_{9}} \)\newline\( + 4L_{q_{6}} + 4L_{q_{7}} + 4L_{q_{9}} + 4L_{\pi_p} \)\newline\( + 4\Sigma \)\\[3pt]
26&12&\( 4D_3^p + 8D_4^p + 2E_p + 4D_3^q \)\newline\( + 8D_4^q + 14G_p + 6G_q + 6H_p \)\newline\( + 5H_q + 4L_{p_{8}} + 2U_{q_{0}} + 4L_{q_{5}} \)\newline\( + 2L_{q_{8}} + 4\Sigma \)\\[3pt]
27&20&\( 6D_3^p + 14D_4^p + 8D_3^q + 16D_4^q \)\newline\( + 2E_q + 14G_p + 22G_q + 9H_p \)\newline\( + 10H_q + 2B + 2U_{p_{0}} + 2L_{q_{7}} \)\newline\( + 4L_{q_{8}} + 6L_{\sigma_p} + 6\Sigma \)\\[3pt]
28&6&\( 6D_3^p + 12D_4^p + 4E_p + 6D_3^q \)\newline\( + 12D_4^q + 4G_p + 6G_q + 2H_p \)\newline\( + 3H_q + 2U_{p_{0}} + 6L_{p_{5}} + 6L_{q_{8}} \)\newline\( + 2L_{\sigma_p} + 6\Sigma \)\\[3pt]
29&10&\( 4D_3^p + 12D_4^p + 8D_3^q + 16D_4^q \)\newline\( + 6E_q + 4G_p + 16G_q + 4H_p \)\newline\( + 5H_q + 2U_{p_{0}} + 2L_{p_{8}} + 4L_{q_{7}} \)\newline\( + 4L_{q_{8}} + 4L_{\pi_p} + 10\Sigma + 10X \)\\[3pt]
30&16&\( 4D_3^p + 12D_4^p + 4D_2^q + 8D_3^q \)\newline\( + 12D_4^q + 2E_q + 10G_p + 18G_q \)\newline\( + 7H_p + 8H_q + 2U_{p_{0}} + 2L_{p_{8}} \)\newline\( + 4L_{q_{7}} + 4L_{\pi_p} + 4\Sigma \)\\[3pt]
31&2&\( 6D_2^p + 12D_3^p + 14D_4^p + 6E_p \)\newline\( + 12D_3^q + 24D_4^q + 22E_q + 2G_p \)\newline\( + 24G_q + H_q + 2Z + 2U_{p_{0}} \)\newline\( + 4L_{q_{5}} + 4L_{q_{6}} + 4L_{\pi_p} + 4L_{\sigma_q} \)\newline\( + 4\Sigma \)\\[3pt]
32&2&\( 6D_3^p + 18D_4^p + 14E_p + 12D_3^q \)\newline\( + 24D_4^q + 22E_q + 10G_p + 24G_q \)\newline\( + H_q + 4Z + 2U_{p_{0}} + 2L_{p_{9}} \)\newline\( + 6L_{q_{7}} + 6L_{q_{9}} + 6L_{\sigma_p} + 6\Sigma \)\newline\( + 6X \)\\[3pt]
33&2&\( 8D_3^p + 16D_4^p + 6E_p + 8D_3^q \)\newline\( + 16D_4^q + 14E_q + 2G_p + 16G_q \)\newline\( + H_q + 4Z + 2U_{p_{0}} + 4L_{p_{5}} \)\newline\( + 6L_{p_{8}} + 4L_{q_{5}} + 4L_{\sigma_q} + 6\Sigma \)\\[3pt]
34&9&\( 12D_3^p + 24D_4^p + 12D_3^q + 24D_4^q \)\newline\( + 14E_q + 20G_q + 3H_p + 4H_q \)\newline\( + 3Z + 3U_{p_{1}} + 12L_{p_{8}} + U_{q_{0}} \)\newline\( + 9L_{q_{5}} + 3L_{\pi_p} + 4L_{\sigma_q} + 12\Sigma \)\\[3pt]
35&18&\( 2D_2^p + 8D_3^p + 10D_4^p + 8D_3^q \)\newline\( + 16D_4^q + 2E_q + 12G_p + 20G_q \)\newline\( + 8H_p + 9H_q + 2U_{p_{0}} + 4L_{q_{6}} \)\newline\( + 4L_{q_{8}} + 4L_{\pi_p} + 4\Sigma \)\\[3pt]
36&20&\( 4D_3^p + 12D_4^p + 8D_3^q + 16D_4^q \)\newline\( + 2E_q + 14G_p + 22G_q + 9H_p \)\newline\( + 10H_q + 2U_{p_{1}} + 2L_{p_{9}} + 4L_{q_{7}} \)\newline\( + 4L_{q_{9}} + 4L_{\sigma_p} + 4\Sigma \)\\[3pt]
37&8&\( 8D_3^p + 12D_4^p + 8D_3^q + 16D_4^q \)\newline\( + 8E_q + 2G_p + 16G_q + 3H_p \)\newline\( + 4H_q + 2U_{p_{0}} + 2L_{p_{8}} + 4L_{q_{6}} \)\newline\( + 4L_{q_{8}} + 4L_{\pi_p} + 8\Sigma + 6X \)\\[3pt]
38&20&\( 8D_3^p + 16D_4^p + 2E_p + 6D_3^q \)\newline\( + 14D_4^q + 22G_p + 14G_q + 10H_p \)\newline\( + 9H_q + 2B + 2L_{p_{7}} + 4L_{p_{8}} \)\newline\( + 2U_{q_{0}} + 6L_{\sigma_q} + 6\Sigma \)\\[3pt]
39&27&\( 31D_3^p + 62D_4^p + 31D_3^q + 59D_4^q \)\newline\( + 31E_q + 55G_q + 9H_p + 13H_q \)\newline\( + 9U_{p_{0}} + 26L_{p_{5}} + 3L_{p_{6}} + U_{q_{0}} \)\newline\( + 2L_{q_{5}} + 27L_{q_{8}} + 11L_{\sigma_p} + 27\Sigma \)\newline\( + X \)\\[3pt]
40&12&\( 4D_3^p + 8D_4^p + 4D_3^q + 8D_4^q \)\newline\( + 2E_q + 6G_p + 14G_q + 5H_p \)\newline\( + 6H_q + 2U_{p_{0}} + 4L_{p_{5}} + 2L_{p_{8}} \)\newline\( + 4L_{q_{8}} + 4\Sigma \)\\[3pt]
41&2&\( 12D_3^p + 18D_4^p + 14E_p + 12D_3^q \)\newline\( + 24D_4^q + 22E_q + 10G_p + 24G_q \)\newline\( + H_q + 4Z + 2U_{p_{0}} + 2L_{p_{9}} \)\newline\( + 6L_{q_{6}} + 6L_{q_{9}} + 6L_{\sigma_p} + 6\Sigma \)\\[3pt]
42&2&\( 8D_3^p + 16D_4^p + 12E_p + 12D_3^q \)\newline\( + 24D_4^q + 22E_q + 8G_p + 24G_q \)\newline\( + H_q + 2Z + 2U_{p_{1}} + 2L_{p_{9}} \)\newline\( + 4L_{q_{6}} + 4L_{q_{7}} + 4L_{q_{9}} + 4L_{\pi_p} \)\newline\( + 4\Sigma \)\\[3pt]
43&18&\( 6D_2^p + 8D_3^p + 10D_4^p + 8D_3^q \)\newline\( + 16D_4^q + 2E_q + 12G_p + 20G_q \)\newline\( + 8H_p + 9H_q + 2U_{p_{0}} + 4L_{p_{5}} \)\newline\( + 4L_{q_{5}} + 4L_{\sigma_q} + 4\Sigma \)\\[3pt]
44&18&\( 8D_3^p + 16D_4^p + 2E_p + 2D_2^q \)\newline\( + 4D_3^q + 10D_4^q + 20G_p + 12G_q \)\newline\( + 9H_p + 8H_q + 4L_{p_{7}} + 2U_{q_{0}} \)\newline\( + 4L_{q_{5}} + 4L_{\sigma_p} + 4\Sigma \)\\[3pt]
45&13&\( D_2^p + D_3^p + 6D_4^p + 6D_3^q \)\newline\( + 12D_4^q + 10G_p + 10G_q + 6H_p \)\newline\( + 6H_q + U_{p_{1}} + U_{q_{0}} + 5L_{q_{7}} \)\newline\( + L_{\pi_p} + 2L_{\sigma_q} + 5\Sigma \)\\[3pt]
46&13&\( 3D_3^p + 9D_4^p + 6D_3^q + 9D_4^q \)\newline\( + 10G_p + 10G_q + 6H_p + 6H_q \)\newline\( + U_{p_{0}} + 3L_{p_{6}} + L_{p_{9}} + U_{q_{0}} \)\newline\( + 3L_{q_{7}} + L_{q_{9}} + 3\Sigma \)\\[3pt]
47&20&\( 8D_3^p + 16D_4^p + 2E_p + 4D_3^q \)\newline\( + 12D_4^q + 22G_p + 14G_q + 10H_p \)\newline\( + 9H_q + 4L_{p_{7}} + 4L_{p_{9}} + 2U_{q_{1}} \)\newline\( + 2L_{q_{9}} + 4L_{\sigma_q} + 4\Sigma \)\\[3pt]
48&14&\( 2D_3^p + 8D_4^p + 6D_3^q + 12D_4^q \)\newline\( + 2E_q + 8G_p + 16G_q + 6H_p \)\newline\( + 7H_q + 2U_{p_{8}} + 2L_{p_{9}} + 4L_{q_{7}} \)\newline\( + 2L_{q_{9}} + 2L_{\pi_p} + 4\Sigma \)\\[3pt]
49&2&\( 6D_3^p + 12D_4^p + 28E_p + 18D_3^q \)\newline\( + 24D_4^q + 30E_q + 24G_p + 32G_q \)\newline\( + H_q + 10Z + 2U_{p_{0}} + 2L_{p_{9}} \)\newline\( + 12L_{q_{3}} + 6L_{q_{9}} + 6L_{\pi_p} + 12\Sigma \)\\[3pt]
50&6&\( 6D_3^p + 12D_4^p + 4E_p + 6D_3^q \)\newline\( + 12D_4^q + 6E_q + 4G_p + 12G_q \)\newline\( + 2H_p + 3H_q + 2U_{p_{1}} + 6L_{p_{8}} \)\newline\( + 4L_{q_{5}} + 2L_{\pi_p} + 2L_{\sigma_q} + 6\Sigma \)\\[3pt]
51&8&\( 6D_3^p + 12D_4^p + 6D_3^q + 12D_4^q \)\newline\( + 4E_q + 2G_p + 12G_q + 3H_p \)\newline\( + 4H_q + 2U_{p_{1}} + 6L_{p_{8}} + 4L_{q_{5}} \)\newline\( + 2L_{\pi_p} + 2L_{\sigma_q} + 8\Sigma + 6X \)\\[3pt]
52&2&\( 6D_2^p + 12D_3^p + 14D_4^p + 10E_p \)\newline\( + 12D_3^q + 24D_4^q + 20E_q + 6G_p \)\newline\( + 22G_q + H_q + 2Z + 2U_{p_{0}} \)\newline\( + 4L_{q_{5}} + 4L_{q_{6}} + 4L_{\pi_p} + 4L_{\sigma_q} \)\newline\( + 4\Sigma + 2X \)\\[3pt]
\end{longtable}
\endgroup

\clearpage
\section[Certificate specification for upper bound 98]{Complete certificate specification for the upper bound 98}\label{u96:app:upper98}
This appendix specifies all 82257 systems and prints the 77 integer row
combinations certifying them. There are 1568 nonzero weights. The optional
JSON assignment selects a successful combination for each system; trying
every printed combination is an equivalent finite procedure.

\subsection{Effective-hole parameters}
Use the definitions of $K,D,M,J$ from Appendix~\ref{u96:app:upper99} with each
of the following $(r,C_7,7)$ and $(r,C_9,9)$ rows. All 66 rows satisfy
$(r-J+1)\beta>D$ and $T_{J-1}-1>D$, the strict hypotheses of
Lemma~\ref{u96:lem:hole-parameter}.
\begin{center}\small
\begin{tabular}{rrr@{\qquad}rrr}
\toprule
$r$&$C_7$&$C_9$&$r$&$C_7$&$C_9$\\\midrule
25&38&18&42&135&160\\
26&44&26&43&142&168\\
27&51&34&44&147&176\\
28&57&42&45&149&177\\
29&63&51&46&155&186\\
30&70&59&47&162&194\\
31&76&67&48&167&198\\
32&83&75&49&168&203\\
33&89&84&50&175&211\\
34&96&92&51&182&220\\
35&102&101&52&187&221\\
36&109&109&53&188&228\\
37&110&117&54&195&237\\
38&116&126&55&202&245\\
39&122&134&56&208&246\\
40&128&143&57&209&254\\
41&129&151&&&\\
\bottomrule
\end{tabular}
\end{center}

\subsection{All cases and variables}
The labeled centers, eighteen units, branch definitions, reflection,
retained-center ordering, and 68 one-empty or 547 two-empty modes are
exactly those defined in Appendix~\ref{u96:app:upper99}.
Here $b=82-a$ and $k=|41-a|$. Use each integer $a$ in the six regime
ranges of Section~\ref{u96:sec:upper98}. In ZE and PE, choose each retained
center and each of its two endpoints in turn as the exceptional location.
There are 2308 such mode/location choices for each fixed $a$.
There is no location enumeration in ZZ: its core-corrected rows are
imposed on every retained center.

There are precisely two refinements to this enumeration:
\begin{enumerate}
\item In ZE at $a=47$, branch A, missing unit $q_5$, and exceptional
lower endpoint of $\pi_p$ (the retained first-pair positive member),
replace that case by two cases. The first imposes $M_q-R_q\leq4$ and
uses upper bound $F\leq1758$; the second imposes $M_q-R_q\geq5$ and
uses $F\leq1764$. The ordinary inequality $M_q-R_q\geq4$ is always
present. Integrality makes this disjunction exhaustive.
\item In EE, branch A with missing units $p_0,q_5$, replace each of its
nine size choices by two cases, adding respectively
$\beta_p+M_p\geq1$ or $\beta_q+M_q\geq1$.
This is the exhaustive nonemptiness disjunction proved in the text.
\end{enumerate}
Thus the case counts are
$21\cdot68=1428$, $19\cdot68=1292$, $17\cdot2308+1=39237$,
$15\cdot2308=34620$, $11\cdot68=748$, and $9\cdot547+9=4932$.

All variables are nonnegative. The base variables are
$p_i,q_i,M_p,M_q,R_p,R_q,S_p,S_q,\delta$ and the endpoint losses
$l_v,u_v$. A harmless variable $t$ is forced to zero by $J:t\leq0$;
unused loss coordinates have zero coefficients and may be set to zero.
For every regime define
\[
 f=S_p+S_q+\delta-2M_p-2M_q-2R_p-2R_q.
\]
Then $F=91+f$ in the five one-empty regimes and $F=182+f$ in EE.

EE also has the nonnegative variables $x_1,x_2,x_3,y_1,y_2,y_3,K$
and $g^p_{ji},g^q_{ji}$ for $0\leq j\leq3$, $j<i\leq4$.
Set $x_0=R_p$, $x_4=M_p$, $y_0=R_q$, $y_4=M_q$ and $W_j=x_j+y_j$.
The model uses $56$ coordinates before these additions and $83$ after
(the surplus base coordinates are unused losses).

\subsection{All rows, with the labels used in the certificate table}
The order and core-size rows are
\[
\begin{gathered}
 D_i^p:p_{i+1}-p_i\leq-1\ (1\leq i\leq3),\quad
 D_4^p:M_p-p_4\leq-1,\\
 E_p:R_p-M_p\leq-4,\qquad G_p:-R_p\leq-a,
\end{gathered}
\]
and the reflected $q$ rows. In EE only, replace the right sides of all
$D_i^p,D_i^q$ by zero; all other displayed rows retain their constants.
The branch row $B$ is $\beta_p+\beta_q\leq0$ in A and
$-\beta_p-\beta_q\leq-1$ in B; in EE branch B its right side is zero.
The mass rows in every regime are
\[
\begin{split}
 H_p &: 2M_p+6R_p-2S_p\leq-2T_a-12-2\epsilon_p,\\
 \widehat H_p &: M_p+4R_p-S_p\leq-T_{a-1}-6,
\end{split}
\]
with reflected $q$ rows, where $(\epsilon_p,\epsilon_q)=(a-2,b-2)$ in
P and PE and $(0,0)$ otherwise.
The budget row $Z$ is $\delta\leq115-k^2$ in Z, ZE, ZZ,
$\delta\leq89-k^2$ in P, PE, and $\delta\leq24-k^2$ in EE.
The loss row is
\[
 \Sigma:\sum_{v\text{ retained}}(l_v+u_v)-\delta\leq0.
\]
The ordinary quantile rows for a positive center are
\[
 L_v:-v-l_v\leq-b-5,\qquad
 U_v:v-u_v-S_p\leq-a-4,
\]
with reflected versions for negative centers. In ZE and PE replace only
the selected exceptional row and add $X$ as follows:
\begin{center}\small
\begin{tabular}{lll}
\toprule
Endpoint&Replacement row&Row $X$\\\midrule
Positive, lower&$-v-l_v-M_p\leq-91$&$-l_v\leq-b-5$\\
Positive, upper&$v-u_v-S_p-M_q\leq-90$&$-u_v\leq-a-5$\\
Negative, lower&$-v-l_v-M_q\leq-91$&$-l_v\leq-a-5$\\
Negative, upper&$v-u_v-S_q-M_p\leq-90$&$-u_v\leq-b-5$\\
\bottomrule
\end{tabular}
\end{center}
The actual out-of-hull location is not imposed. In ZZ replace every
positive quantile row by
\[
 L_v:-v-l_v-R_p\leq-87,\qquad
 U_v:v-u_v-S_p-R_q\leq-86,
\]
with the reflected negative versions, and add
\[
 X_2:-\delta\leq-2(46-k),\qquad
 R:R_p+R_q+\delta\leq199-k^2.
\]
In EE all quantile rows remain ordinary, and add
$R:R_p+R_q+\delta\leq106-k^2$.
The special ZE split uses respectively
\[
 Y_4:M_q-R_q\leq4,\qquad Y_5:R_q-M_q\leq-5.
\]

The following additional rows occur only in EE:
\[
 I_i^p:x_{i-1}-x_i\leq-1\quad(1\leq i\leq4),\qquad
 \widetilde H_p:R_p+\sum_{i=1}^4x_i-S_p\leq-T_{a-1},
\]
and the reflected negative rows. For $0\leq j\leq3$ and $j<i\leq4$,
\[
 A^p_{ji}:g^p_{ji}-x_i+x_{i-1}\leq0,\qquad
 B^p_{ji}:g^p_{ji}-W_j\leq1,
\]
with reflected $q$ rows. Put $t_j=2(4-j)$. The envelope and cardinality
rows are
\[
\begin{split}
 C_j &: K-2W_j-\sum_{i=j+1}^4(g^p_{ji}+g^q_{ji})
              \leq1+\binom{t_j+1}{2},\qquad 0\leq j\leq3,\\
 M &: S_p+S_q+\delta-2K\leq1540.
\end{split}
\]
Finally the two special EE branches add respectively
\[
 P_p:p_1-2p_2+p_3+p_4-M_p\leq-1,\qquad
 P_q:q_1-2q_2+q_3+q_4-M_q\leq-1.
\]
These are all rows of the delivered generator; none requires a numerical
feasibility assumption.

\subsection{Exact verification rule and all multipliers}
For $Ax\leq b$, $x\geq0$, each entry gives an integer $D\geq0$ and
nonnegative integer row weights $w$. A row absent from the current case
makes that entry unavailable. Otherwise form
\[
 q=Df+wA,\qquad h=wb.
\]
If every coordinate of $q$ is nonnegative, then $D>0$ gives
$F\geq d_0-h/D$, where $d_0=91$ or $182$.
It certifies impossibility if that lower bound is strictly greater than
1764, or 1820 in P and PE, or 1758 in the special gap-four ZE case.
For $D=0$, the strict inequality $h<0$ certifies inconsistency.
Every assigned combination passes these tests over the integers.
The smallest strict lower-bound margin across all cases is $2/3$.
There is no use of floating-point tolerance in this verification.

\begingroup\footnotesize
\setlength{\tabcolsep}{4pt}
\begin{longtable}{rrp{12.7cm}}
\toprule
ID&$D$&Nonzero row weights (omitted weights are zero)\\
\midrule\endfirsthead
\toprule ID&$D$&Nonzero row weights (continued)\\\midrule\endhead
\bottomrule\endfoot
1&10&\( 8D_3^p + 12D_4^p + 8D_3^q + 16D_4^q \)\newline\( + 6E_q + 4G_p + 16G_q + 4H_p \)\newline\( + 5H_q + 2U_{p_{0}} + 2L_{p_{8}} + 4L_{q_{6}} \)\newline\( + 4L_{q_{8}} + 4L_{\pi_p} + 10\Sigma \)\\[3pt]
2&0&\( Z + L_{q_{5}} + L_{\pi_p} + L_{\sigma_p} \)\newline\( + \Sigma \)\\[3pt]
3&1&\( 3D_3^p + 12D_4^p + 10E_p + 9D_3^q \)\newline\( + 18D_4^q + 11E_q + 8G_p + 13G_q \)\newline\( + 5Z + U_{p_{0}} + L_{p_{9}} + 6L_{q_{7}} \)\newline\( + 3L_{q_{9}} + 3L_{\pi_p} + 6\Sigma + \widehat H_q \)\\[3pt]
4&6&\( 5D_2^p + 6D_3^p + 7D_4^p + 6D_3^q \)\newline\( + 12D_4^q + 6E_q + 8G_p + 18G_q \)\newline\( + U_{p_{0}} + 4L_{q_{5}} + 2L_{\pi_p} + 2L_{\sigma_q} \)\newline\( + 6\Sigma + 5\widehat H_p + 6\widehat H_q \)\\[3pt]
5&3&\( 6D_3^p + 12D_4^p + 9E_p + 6D_3^q \)\newline\( + 9D_4^q + 2E_q + 15G_p + 4G_q \)\newline\( + 3L_{p_{6}} + 3L_{p_{9}} + U_{q_{0}} + L_{q_{9}} \)\newline\( + 3L_{\sigma_q} + 3\Sigma + 3\widehat H_p + 2\widehat H_q \)\\[3pt]
6&5&\( 3D_3^p + 6D_4^p + E_p + 3D_3^q \)\newline\( + 6D_4^q + 11G_p + 6G_q + 3L_{p_{8}} \)\newline\( + U_{q_{0}} + 3L_{q_{5}} + L_{\sigma_q} + 5\Sigma \)\newline\( + 5\widehat H_p + 4\widehat H_q \)\\[3pt]
7&1&\( 6D_3^p + 9D_4^p + 7E_p + 6D_3^q \)\newline\( + 12D_4^q + 11E_q + 5G_p + 13G_q \)\newline\( + 2Z + U_{p_{0}} + L_{p_{9}} + 3L_{q_{6}} \)\newline\( + 3L_{q_{9}} + 3L_{\sigma_p} + 3\Sigma + \widehat H_q \)\\[3pt]
8&6&\( 4D_3^p + 8D_4^p + 4D_3^q + 8D_4^q \)\newline\( + E_q + 12G_p + 9G_q + 2L_{p_{5}} \)\newline\( + U_{q_{0}} + 2L_{q_{5}} + 3L_{q_{8}} + 2L_{\sigma_p} \)\newline\( + 6\Sigma + 4X + 6\widehat H_p + 5\widehat H_q \)\\[3pt]
9&1&\( 4D_3^p + 8D_4^p + 4E_p + 6D_3^q \)\newline\( + 12D_4^q + 11E_q + 2G_p + 13G_q \)\newline\( + Z + U_{p_{1}} + L_{p_{9}} + 2L_{q_{6}} \)\newline\( + 2L_{q_{7}} + 2L_{q_{9}} + 2L_{\pi_p} + 2\Sigma \)\newline\( + \widehat H_q \)\\[3pt]
10&5&\( 4D_3^p + 6D_4^p + 4D_3^q + 8D_4^q \)\newline\( + E_q + 6G_p + 11G_q + U_{p_{0}} \)\newline\( + L_{p_{8}} + 2L_{q_{6}} + 2L_{q_{8}} + 2L_{\pi_p} \)\newline\( + 5\Sigma + 3X + 4\widehat H_p + 5\widehat H_q \)\\[3pt]
11&3&\( 4D_3^p + 8D_4^p + 4E_p + 4D_3^q \)\newline\( + 8D_4^q + 3E_q + 6G_p + 9G_q \)\newline\( + U_{p_{0}} + 2L_{p_{5}} + 3L_{p_{8}} + 2L_{q_{5}} \)\newline\( + 2L_{\sigma_q} + 3\Sigma + X + 2\widehat H_p \)\newline\( + 3\widehat H_q \)\\[3pt]
12&5&\( 4D_3^p + 6D_4^p + 4D_3^q + 8D_4^q \)\newline\( + 3E_q + 3G_p + 8G_q + U_{p_{0}} \)\newline\( + L_{p_{8}} + 2L_{q_{6}} + 2L_{q_{8}} + 2L_{\pi_p} \)\newline\( + 2\Sigma + 4\widehat H_p + 5\widehat H_q + 3X_2 \)\\[3pt]
13&6&\( 8D_3^p + 12D_4^p + 8D_3^q + 16D_4^q \)\newline\( + 10E_q + 16G_q + 2H_p + 3H_q \)\newline\( + 2U_{p_{0}} + 2L_{p_{8}} + 4L_{q_{6}} + 4L_{q_{8}} \)\newline\( + 4L_{\pi_p} + 6\Sigma + 2X \)\\[3pt]
14&18&\( 12D_3^p + 22D_4^p + 12D_3^q + 24D_4^q \)\newline\( + 6G_p + 18G_q + 7H_p + 9H_q \)\newline\( + 4U_{p_{1}} + 10L_{p_{8}} + 6L_{q_{5}} + 2L_{q_{6}} \)\newline\( + 4L_{\pi_p} + 4L_{\sigma_q} + 18\Sigma + 12X \)\\[3pt]
15&3&\( 4D_3^p + 8D_4^p + 5E_p + 4D_3^q \)\newline\( + 8D_4^q + E_q + 11G_p + 3G_q \)\newline\( + 2L_{p_{5}} + U_{q_{0}} + 2L_{q_{5}} + 3L_{q_{8}} \)\newline\( + 2L_{\sigma_p} + 3\Sigma + 3\widehat H_p + 2\widehat H_q \)\\[3pt]
16&0&\( Z + \Sigma + X \)\\[3pt]
17&1&\( 6D_3^p + 8D_4^p + 6E_p + 6D_3^q \)\newline\( + 12D_4^q + 11E_q + 4G_p + 13G_q \)\newline\( + 3Z + U_{p_{1}} + L_{p_{9}} + 4L_{q_{6}} \)\newline\( + 2L_{q_{9}} + 2L_{\pi_p} + 4\Sigma + 4X \)\newline\( + \widehat H_q \)\\[3pt]
18&8&\( 6D_3^p + 12D_4^p + E_p + 6D_3^q \)\newline\( + 9D_4^q + 17G_p + 12G_q + 3L_{p_{6}} \)\newline\( + 3L_{p_{9}} + U_{q_{0}} + L_{q_{9}} + 3L_{\sigma_q} \)\newline\( + 8\Sigma + 5X + 8\widehat H_p + 7\widehat H_q \)\\[3pt]
19&0&\( D_4^p + Z + L_{q_{5}} + L_{\pi_p} \)\newline\( + L_{\sigma_p} + \Sigma \)\\[3pt]
20&0&\( D_4^q + Z + L_{q_{5}} + L_{\pi_p} \)\newline\( + L_{\sigma_p} + \Sigma \)\\[3pt]
21&1&\( 6D_3^p + 9D_4^p + 7E_p + 6D_3^q \)\newline\( + 12D_4^q + 8E_q + 5G_p + 10G_q \)\newline\( + 2Z + U_{p_{0}} + L_{p_{9}} + 3L_{q_{6}} \)\newline\( + 3L_{q_{9}} + 3L_{\sigma_p} + 3\Sigma + \widehat H_q \)\\[3pt]
22&1&\( 4D_3^p + 8D_4^p + 6E_p + 6D_3^q \)\newline\( + 12D_4^q + 9E_q + 4G_p + 11G_q \)\newline\( + Z + U_{p_{1}} + L_{p_{9}} + 2L_{q_{6}} \)\newline\( + 2L_{q_{7}} + 2L_{q_{9}} + 2L_{\pi_p} + 2\Sigma \)\newline\( + \widehat H_q \)\\[3pt]
23&1&\( 6D_3^p + 10D_4^p + 5E_p + 7D_3^q \)\newline\( + 14D_4^q + 13E_q + 3G_p + 15G_q \)\newline\( + 2Z + B + U_{p_{0}} + 3L_{q_{6}} \)\newline\( + L_{q_{7}} + 2L_{q_{8}} + 3L_{\pi_p} + 3\Sigma \)\newline\( + \widehat H_q \)\\[3pt]
24&1&\( 3D_2^p + 6D_3^p + 7D_4^p + 5E_p \)\newline\( + 6D_3^q + 12D_4^q + 9E_q + 3G_p \)\newline\( + 11G_q + Z + U_{p_{0}} + 2L_{q_{5}} \)\newline\( + 2L_{q_{6}} + 2L_{\pi_p} + 2L_{\sigma_q} + 2\Sigma \)\newline\( + \widehat H_q \)\\[3pt]
25&0&\( 2D_4^p + 2E_p + D_4^q + E_q \)\newline\( + Z + L_{q_{5}} + L_{\pi_p} + L_{\sigma_p} \)\newline\( + \Sigma \)\\[3pt]
26&1&\( 3D_3^p + 12D_4^p + 10E_p + 9D_3^q \)\newline\( + 18D_4^q + 17E_q + 4G_p + 9G_q \)\newline\( + 5Z + U_{p_{0}} + L_{p_{9}} + 6L_{q_{7}} \)\newline\( + 3L_{q_{9}} + 3L_{\pi_p} + 6\Sigma + \widehat H_q \)\\[3pt]
27&0&\( 8D_3^p + 24D_4^p + 8D_2^q + 16D_3^q \)\newline\( + 24D_4^q + 24G_p + 38G_q + 3H_q \)\newline\( + 4U_{p_{0}} + 4L_{p_{8}} + 8L_{q_{7}} + 8L_{\pi_p} \)\newline\( + 8\Sigma + 8I_1^p + 16\widetilde H_p + 6I_1^q \)\newline\( + 2I_3^q + 4I_4^q + 14\widetilde H_q + 8A_{12}^p \)\newline\( + 8A_{13}^p + 8A_{14}^p + 8A_{12}^q + 8A_{13}^q \)\newline\( + 8A_{14}^q + 8C_1 + 16A_{23}^p + 16A_{24}^p \)\newline\( + 16A_{23}^q + 16A_{24}^q + 16C_2 + 16A_{34}^p \)\newline\( + 16A_{34}^q + 16C_3 + 20M \)\\[3pt]
28&3&\( 4D_3^p + 8D_4^p + E_p + 4D_3^q \)\newline\( + 8D_4^q + 5E_q + 3G_p + 11G_q \)\newline\( + U_{p_{0}} + 2L_{p_{5}} + 3L_{p_{8}} + 2L_{q_{5}} \)\newline\( + 2L_{\sigma_q} + 3\Sigma + 2\widehat H_p + 3\widehat H_q \)\\[3pt]
29&1&\( 6D_3^p + 9D_4^p + 4E_p + 6D_3^q \)\newline\( + 12D_4^q + 11E_q + 2G_p + 13G_q \)\newline\( + 2Z + U_{p_{0}} + L_{p_{9}} + 3L_{q_{6}} \)\newline\( + 3L_{q_{9}} + 3L_{\sigma_p} + 3\Sigma + \widehat H_q \)\\[3pt]
30&5&\( 4D_3^p + 6D_4^p + 4D_3^q + 8D_4^q \)\newline\( + 2E_q + 6G_p + 12G_q + U_{p_{0}} \)\newline\( + L_{p_{8}} + 2L_{q_{6}} + 2L_{q_{8}} + 2L_{\pi_p} \)\newline\( + 5\Sigma + 4X + 4\widehat H_p + 5\widehat H_q \)\\[3pt]
31&7&\( 4D_3^p + 8D_4^p + 4D_3^q + 8D_4^q \)\newline\( + E_q + 5G_p + 10G_q + U_{p_{0}} \)\newline\( + 2L_{p_{5}} + 3L_{p_{8}} + 2L_{q_{5}} + 2L_{\sigma_q} \)\newline\( + 3\Sigma + 6\widehat H_p + 7\widehat H_q + 4X_2 \)\\[3pt]
32&6&\( 6D_3^p + 12D_4^p + 6D_3^q + 12D_4^q \)\newline\( + 4E_q + 6G_p + 4G_q + 3H_p \)\newline\( + 2H_q + 6L_{p_{8}} + 2U_{q_{0}} + 6L_{q_{5}} \)\newline\( + 2L_{\sigma_q} + 6\Sigma \)\\[3pt]
33&0&\( 6D_3^p + 12D_4^p + 3D_3^q + 6D_4^q \)\newline\( + 8G_p + 12G_q + 2U_{p_{0}} + 3L_{p_{6}} \)\newline\( + 3L_{p_{7}} + 2L_{p_{9}} + 3\Sigma + 2I_3^p \)\newline\( + 8\widetilde H_p + 2I_1^q + 10\widetilde H_q + 4A_{12}^p \)\newline\( + 4A_{13}^p + 4A_{14}^p + 4A_{12}^q + 4A_{13}^q \)\newline\( + 4B_{14}^q + 4C_1 + 10A_{23}^p + 10A_{24}^p \)\newline\( + 10A_{23}^q + 10A_{24}^q + 10C_2 + 6A_{34}^p \)\newline\( + 6A_{34}^q + 6C_3 + 10M \)\\[3pt]
34&4&\( 5D_2^p + 6D_3^p + 7D_4^p + 6D_3^q \)\newline\( + 12D_4^q + 8E_q + 4G_p + 16G_q \)\newline\( + U_{p_{0}} + 4L_{q_{5}} + 2L_{\pi_p} + 2L_{\sigma_q} \)\newline\( + 4\Sigma + 2X + 3\widehat H_p + 4\widehat H_q \)\\[3pt]
35&5&\( 2D_3^p + 6D_4^p + 4D_3^q + 8D_4^q \)\newline\( + 3E_q + 6G_p + 13G_q + U_{p_{0}} \)\newline\( + L_{p_{8}} + 2L_{q_{7}} + 2L_{q_{8}} + 2L_{\pi_p} \)\newline\( + 5\Sigma + 5X + 4\widehat H_p + 5\widehat H_q \)\\[3pt]
36&1&\( 3D_3^p + 9D_4^p + 7E_p + 6D_3^q \)\newline\( + 12D_4^q + 11E_q + 5G_p + 13G_q \)\newline\( + 2Z + U_{p_{0}} + L_{p_{9}} + 3L_{q_{7}} \)\newline\( + 3L_{q_{9}} + 3L_{\sigma_p} + 3\Sigma + 3X \)\newline\( + \widehat H_q \)\\[3pt]
37&3&\( 3D_3^p + 6D_4^p + 2E_p + 3D_3^q \)\newline\( + 6D_4^q + 4G_p + 6G_q + U_{p_{0}} \)\newline\( + 3L_{p_{5}} + 3L_{q_{8}} + L_{\sigma_p} + 3\Sigma \)\newline\( + 2\widehat H_p + 3\widehat H_q \)\\[3pt]
38&8&\( 8D_3^p + 12D_4^p + 8D_3^q + 16D_4^q \)\newline\( + 8E_q + 2G_p + 16G_q + 3H_p \)\newline\( + 4H_q + 2U_{p_{0}} + 2L_{p_{8}} + 4L_{q_{6}} \)\newline\( + 4L_{q_{8}} + 4L_{\pi_p} + 8\Sigma + 6X \)\\[3pt]
39&3&\( 3D_3^p + 6D_4^p + 3E_p + 6D_3^q \)\newline\( + 12D_4^q + 7E_q + 9G_p + 5G_q \)\newline\( + 2U_{q_{0}} + 3L_{q_{6}} + 3L_{q_{7}} + 2L_{q_{9}} \)\newline\( + 3\Sigma + 3\widehat H_p + \widehat H_q \)\\[3pt]
40&8&\( 6D_3^p + 12D_4^p + 4E_p + 6D_3^q \)\newline\( + 9D_4^q + 13G_p + 8G_q + 3L_{p_{6}} \)\newline\( + 3L_{p_{9}} + U_{q_{0}} + L_{q_{9}} + 3L_{\sigma_q} \)\newline\( + 3\Sigma + 8\widehat H_p + 7\widehat H_q + 5X_2 \)\\[3pt]
41&10&\( 6D_3^p + 12D_4^p + 6D_3^q + 12D_4^q \)\newline\( + 2E_q + 4G_p + 12G_q + 4H_p \)\newline\( + 5H_q + 2U_{p_{1}} + 6L_{p_{8}} + 4L_{q_{5}} \)\newline\( + 2L_{\pi_p} + 2L_{\sigma_q} + 10\Sigma \)\\[3pt]
42&0&\( 4D_3^p + 8D_4^p + 3D_3^q + 7D_4^q \)\newline\( + 12G_p + 8G_q + H_p + B \)\newline\( + L_{p_{7}} + 2L_{p_{8}} + U_{q_{0}} + 3L_{\sigma_q} \)\newline\( + 3\Sigma + I_3^p + 6\widetilde H_p + I_1^q \)\newline\( + 7\widetilde H_q + 4A_{12}^p + 4A_{13}^p + 4A_{14}^p \)\newline\( + 4A_{12}^q + 4A_{13}^q + 2A_{14}^q + 2B_{14}^q \)\newline\( + 4C_1 + 7A_{23}^p + 7A_{24}^p + 7A_{23}^q \)\newline\( + 7A_{24}^q + 7C_2 + 5A_{34}^p + 5A_{34}^q \)\newline\( + 5C_3 + 8M \)\\[3pt]
43&0&\( 12D_3^p + 24D_4^p + 12D_3^q + 24D_4^q \)\newline\( + 34G_p + 12G_q + 4H_p + 12L_{p_{8}} \)\newline\( + 6U_{q_{0}} + 12L_{q_{5}} + 6L_{q_{8}} + 12\Sigma \)\newline\( + 2I_2^p + 4I_3^p + 16\widetilde H_p + 18\widetilde H_q \)\newline\( + 3A_{01}^p + 3A_{02}^p + 3A_{03}^p + 3A_{04}^p \)\newline\( + 3A_{01}^q + 3A_{02}^q + 3A_{03}^q + 3B_{04}^q \)\newline\( + 3C_0 + 15A_{12}^p + 15A_{13}^p + 15A_{14}^p \)\newline\( + 15A_{12}^q + 15A_{13}^q + 12A_{14}^q + 3B_{14}^q \)\newline\( + 15C_1 + 18A_{23}^p + 18A_{24}^p + 18A_{23}^q \)\newline\( + 18A_{24}^q + 18C_2 + 12A_{34}^p + 12A_{34}^q \)\newline\( + 12C_3 + 24M \)\\[3pt]
44&0&\( D_2^p + 2D_3^p + 5D_4^p + 4D_3^q \)\newline\( + 8D_4^q + 2E_q + 4G_p + 10G_q \)\newline\( + H_q + U_{p_{0}} + 2L_{q_{7}} + 2L_{q_{8}} \)\newline\( + 2L_{\pi_p} + 2\Sigma + I_1^p + 3\widetilde H_p \)\newline\( + I_3^q + 2I_4^q + 2\widetilde H_q + 2A_{12}^p \)\newline\( + 2A_{13}^p + 2A_{14}^p + 2A_{12}^q + 2A_{13}^q \)\newline\( + 2A_{14}^q + 2C_1 + 3A_{23}^p + 3A_{24}^p \)\newline\( + 3A_{23}^q + 3A_{24}^q + 3C_2 + 3A_{34}^p \)\newline\( + 3A_{34}^q + 3C_3 + 4M \)\\[3pt]
45&9&\( 12D_3^p + 24D_4^p + 12D_3^q + 24D_4^q \)\newline\( + 14E_q + 20G_q + 3H_p + 4H_q \)\newline\( + 3Z + 3U_{p_{1}} + 12L_{p_{8}} + U_{q_{0}} \)\newline\( + 9L_{q_{5}} + 3L_{\pi_p} + 4L_{\sigma_q} + 12\Sigma \)\\[3pt]
46&5&\( 3D_3^p + 6D_4^p + 3D_3^q + 6D_4^q \)\newline\( + E_q + 2G_p + 7G_q + U_{p_{0}} \)\newline\( + 3L_{p_{5}} + 3L_{q_{8}} + L_{\sigma_p} + 3\Sigma \)\newline\( + 4\widehat H_p + 5\widehat H_q + 2X_2 \)\\[3pt]
47&5&\( 3D_3^p + 6D_4^p + 3D_3^q + 6D_4^q \)\newline\( + E_q + 2G_p + 7G_q + U_{p_{1}} \)\newline\( + 3L_{p_{8}} + 2L_{q_{5}} + L_{\pi_p} + L_{\sigma_q} \)\newline\( + 3\Sigma + 4\widehat H_p + 5\widehat H_q + 2X_2 \)\\[3pt]
48&1&\( 6D_3^p + 9D_4^p + 7E_p + 6D_3^q \)\newline\( + 12D_4^q + 11E_q + G_p + 6G_q \)\newline\( + 2Z + U_{p_{0}} + L_{p_{9}} + 3L_{q_{6}} \)\newline\( + 3L_{q_{9}} + 3L_{\sigma_p} + 3\Sigma + \widehat H_q \)\\[3pt]
49&1&\( 3D_2^p + 6D_3^p + 7D_4^p + 5E_p \)\newline\( + 6D_3^q + 12D_4^q + 11E_q + G_p \)\newline\( + 6G_q + Z + U_{p_{0}} + 2L_{q_{5}} \)\newline\( + 2L_{q_{6}} + 2L_{\pi_p} + 2L_{\sigma_q} + 2\Sigma \)\newline\( + \widehat H_q \)\\[3pt]
50&1&\( 4D_3^p + 8D_4^p + 6E_p + 6D_3^q \)\newline\( + 12D_4^q + 11E_q + G_p + 6G_q \)\newline\( + Z + U_{p_{1}} + L_{p_{9}} + 2L_{q_{6}} \)\newline\( + 2L_{q_{7}} + 2L_{q_{9}} + 2L_{\pi_p} + 2\Sigma \)\newline\( + \widehat H_q \)\\[3pt]
51&6&\( 5D_2^p + 6D_3^p + 7D_4^p + 6D_3^q \)\newline\( + 12D_4^q + 4E_q + 8G_p + 16G_q \)\newline\( + U_{p_{0}} + 4L_{q_{5}} + 2L_{\pi_p} + 2L_{\sigma_q} \)\newline\( + 6\Sigma + 4X + 5\widehat H_p + 6\widehat H_q \)\\[3pt]
52&0&\( 24D_3^p + 72D_4^p + 48D_3^q + 96D_4^q \)\newline\( + 70G_p + 124G_q + H_p + 16H_q \)\newline\( + 12U_{p_{1}} + 12L_{p_{9}} + 24L_{q_{7}} + 24L_{q_{9}} \)\newline\( + 24L_{\sigma_p} + 24\Sigma + 18I_1^p + 46\widetilde H_p \)\newline\( + 18I_3^q + 36I_4^q + 28\widetilde H_q + 28A_{12}^p \)\newline\( + 28A_{13}^p + 28A_{14}^p + 28A_{12}^q + 28A_{13}^q \)\newline\( + 28A_{14}^q + 28C_1 + 46A_{23}^p + 46A_{24}^p \)\newline\( + 46A_{23}^q + 46A_{24}^q + 46C_2 + 46A_{34}^p \)\newline\( + 46A_{34}^q + 46C_3 + 60M \)\\[3pt]
53&0&\( 18D_3^p + 42D_4^p + 30D_3^q + 60D_4^q \)\newline\( + 46G_p + 76G_q + H_p + 10H_q \)\newline\( + 6B + 6U_{p_{1}} + 6L_{q_{6}} + 12L_{q_{7}} \)\newline\( + 6L_{q_{8}} + 12L_{\pi_p} + 12\Sigma + 12I_1^p \)\newline\( + 28\widetilde H_p + 12I_3^q + 24I_4^q + 16\widetilde H_q \)\newline\( + 16A_{12}^p + 16A_{13}^p + 16A_{14}^p + 16A_{12}^q \)\newline\( + 16A_{13}^q + 16A_{14}^q + 16C_1 + 28A_{23}^p \)\newline\( + 28A_{24}^p + 28A_{23}^q + 28A_{24}^q + 28C_2 \)\newline\( + 28A_{34}^p + 28A_{34}^q + 28C_3 + 36M \)\\[3pt]
54&2&\( 2D_3^p + 2D_4^p + 2D_3^q + 4D_4^q \)\newline\( + E_q + 4G_p + G_q + U_{q_{0}} \)\newline\( + 2L_{q_{6}} + L_{q_{8}} + 2\Sigma + 2X \)\newline\( + 2\widehat H_p + \widehat H_q \)\\[3pt]
55&8&\( 6D_3^p + 12D_4^p + 6D_3^q + 12D_4^q \)\newline\( + 4E_q + 2G_p + 12G_q + 3H_p \)\newline\( + 4H_q + 2U_{p_{1}} + 6L_{p_{8}} + 4L_{q_{5}} \)\newline\( + 2L_{\pi_p} + 2L_{\sigma_q} + 8\Sigma + 6X \)\\[3pt]
56&0&\( 6D_3^p + 12D_4^p + 3D_3^q + 6D_4^q \)\newline\( + 8G_p + 12G_q + 2U_{p_{0}} + 3L_{p_{6}} \)\newline\( + 3L_{p_{7}} + 2L_{\sigma_p} + 3\Sigma + 2I_3^p \)\newline\( + 8\widetilde H_p + 2I_1^q + 10\widetilde H_q + 4A_{12}^p \)\newline\( + 4A_{13}^p + 4A_{14}^p + 4A_{12}^q + 4A_{13}^q \)\newline\( + 4B_{14}^q + 4C_1 + 10A_{23}^p + 10A_{24}^p \)\newline\( + 10A_{23}^q + 10A_{24}^q + 10C_2 + 6A_{34}^p \)\newline\( + 6A_{34}^q + 6C_3 + 10M \)\\[3pt]
57&0&\( D_2^p + 2D_3^p + 5D_4^p + 4D_3^q \)\newline\( + 8D_4^q + 2E_q + 4G_p + 10G_q \)\newline\( + H_q + U_{p_{0}} + 2L_{p_{5}} + 2L_{q_{7}} \)\newline\( + 2L_{\sigma_q} + 2\Sigma + I_1^p + 3\widetilde H_p \)\newline\( + I_3^q + 2I_4^q + 2\widetilde H_q + 2A_{12}^p \)\newline\( + 2A_{13}^p + 2A_{14}^p + 2A_{12}^q + 2A_{13}^q \)\newline\( + 2A_{14}^q + 2C_1 + 3A_{23}^p + 3A_{24}^p \)\newline\( + 3A_{23}^q + 3A_{24}^q + 3C_2 + 3A_{34}^p \)\newline\( + 3A_{34}^q + 3C_3 + 4M \)\\[3pt]
58&0&\( 2D_3^p + 6D_4^p + 4D_3^q + 8D_4^q \)\newline\( + 2E_q + 4G_p + 10G_q + H_q \)\newline\( + U_{p_{1}} + 2L_{q_{7}} + 2L_{q_{8}} + 2L_{\pi_p} \)\newline\( + 2\Sigma + I_1^p + 3\widetilde H_p + I_3^q \)\newline\( + 2I_4^q + 2\widetilde H_q + 2A_{12}^p + 2A_{13}^p \)\newline\( + 2A_{14}^p + 2A_{12}^q + 2A_{13}^q + 2A_{14}^q \)\newline\( + 2C_1 + 3A_{23}^p + 3A_{24}^p + 3A_{23}^q \)\newline\( + 3A_{24}^q + 3C_2 + 3A_{34}^p + 3A_{34}^q \)\newline\( + 3C_3 + 4M + P_p \)\\[3pt]
59&0&\( 24D_3^p + 32D_4^p + 24D_3^q + 48D_4^q \)\newline\( + 38G_p + 60G_q + H_p + 8H_q \)\newline\( + 4U_{p_{1}} + 4L_{p_{9}} + 16L_{q_{6}} + 8L_{q_{9}} \)\newline\( + 8L_{\pi_p} + 16\Sigma + 10I_1^p + 22\widetilde H_p \)\newline\( + 10I_3^q + 20I_4^q + 12\widetilde H_q + 12A_{12}^p \)\newline\( + 12A_{13}^p + 12A_{14}^p + 12A_{12}^q + 12A_{13}^q \)\newline\( + 12A_{14}^q + 12C_1 + 22A_{23}^p + 22A_{24}^p \)\newline\( + 22A_{23}^q + 22A_{24}^q + 22C_2 + 22A_{34}^p \)\newline\( + 22A_{34}^q + 22C_3 + 28M \)\\[3pt]
60&0&\( 4D_3^p + 8D_4^p + 6D_3^q + 12D_4^q \)\newline\( + 6E_q + 4G_p + 16G_q + H_q \)\newline\( + 2U_{p_{8}} + 2L_{p_{9}} + 2L_{q_{6}} + 2L_{q_{7}} \)\newline\( + 2L_{q_{9}} + 2L_{\pi_p} + 2\Sigma + 4\widetilde H_p \)\newline\( + 4\widetilde H_q + 4A_{12}^p + 4A_{13}^p + 4A_{14}^p \)\newline\( + 4A_{12}^q + 4A_{13}^q + 4A_{14}^q + 4C_1 \)\newline\( + 4A_{23}^p + 4A_{24}^p + 4A_{23}^q + 4A_{24}^q \)\newline\( + 4C_2 + 4A_{34}^p + 4A_{34}^q + 4C_3 \)\newline\( + 6M \)\\[3pt]
61&3&\( 3D_3^p + 6D_4^p + 2E_p + 3D_3^q \)\newline\( + 6D_4^q + 3E_q + 4G_p + 9G_q \)\newline\( + U_{p_{0}} + 3L_{p_{5}} + 3L_{q_{8}} + L_{\sigma_p} \)\newline\( + 3\Sigma + 2\widehat H_p + 3\widehat H_q \)\\[3pt]
62&0&\( 6D_3^p + 18D_4^p + 12D_3^q + 24D_4^q \)\newline\( + 22G_p + 32G_q + H_p + 4H_q \)\newline\( + 2B + 2U_{p_{0}} + 6L_{q_{7}} + 4L_{q_{9}} \)\newline\( + 6L_{\sigma_p} + 6\Sigma + 4I_1^p + 12\widetilde H_p \)\newline\( + 4I_3^q + 8I_4^q + 8\widetilde H_q + 8A_{12}^p \)\newline\( + 8A_{13}^p + 8A_{14}^p + 8A_{12}^q + 8A_{13}^q \)\newline\( + 8A_{14}^q + 8C_1 + 12A_{23}^p + 12A_{24}^p \)\newline\( + 12A_{23}^q + 12A_{24}^q + 12C_2 + 12A_{34}^p \)\newline\( + 12A_{34}^q + 12C_3 + 16M \)\\[3pt]
63&0&\( 12D_3^p + 24D_4^p + 12D_3^q + 24D_4^q \)\newline\( + 12G_p + 34G_q + 4H_q + 6U_{p_{0}} \)\newline\( + 12L_{p_{5}} + 6L_{p_{8}} + 12L_{q_{8}} + 12\Sigma \)\newline\( + 18\widetilde H_p + 2I_2^q + 4I_3^q + 16\widetilde H_q \)\newline\( + 3A_{01}^p + 3A_{02}^p + 3A_{03}^p + 3B_{04}^p \)\newline\( + 3A_{01}^q + 3A_{02}^q + 3A_{03}^q + 3A_{04}^q \)\newline\( + 3C_0 + 15A_{12}^p + 15A_{13}^p + 12A_{14}^p \)\newline\( + 3B_{14}^p + 15A_{12}^q + 15A_{13}^q + 15A_{14}^q \)\newline\( + 15C_1 + 18A_{23}^p + 18A_{24}^p + 18A_{23}^q \)\newline\( + 18A_{24}^q + 18C_2 + 12A_{34}^p + 12A_{34}^q \)\newline\( + 12C_3 + 24M \)\\[3pt]
64&0&\( 6D_3^p + 18D_4^p + 18D_3^q + 36D_4^q \)\newline\( + 34G_q + 7H_q + 6U_{p_{3}} + 6L_{q_{6}} \)\newline\( + 6L_{q_{7}} + 6L_{\pi_p} + 6L_{\sigma_q} + 6\Sigma \)\newline\( + 30\widetilde H_p + 8I_2^q + 16I_3^q + 22\widetilde H_q \)\newline\( + 15A_{01}^p + 15A_{02}^p + 15A_{03}^p + 15B_{04}^p \)\newline\( + 15A_{01}^q + 15A_{02}^q + 15A_{03}^q + 15A_{04}^q \)\newline\( + 15C_0 + 21A_{12}^p + 21A_{13}^p + 12A_{14}^p \)\newline\( + 9B_{14}^p + 21A_{12}^q + 21A_{13}^q + 21A_{14}^q \)\newline\( + 21C_1 + 30A_{23}^p + 30A_{24}^p + 30A_{23}^q \)\newline\( + 30A_{24}^q + 30C_2 + 6A_{34}^p + 6A_{34}^q \)\newline\( + 6C_3 + 36M \)\\[3pt]
65&1&\( 6D_3^p + 9D_4^p + 6E_p + 6D_3^q \)\newline\( + 12D_4^q + 11E_q + 4G_p + 13G_q \)\newline\( + 2Z + U_{p_{0}} + L_{p_{9}} + 3L_{q_{6}} \)\newline\( + 3L_{q_{9}} + 3L_{\sigma_p} + 3\Sigma + 2X \)\newline\( + \widehat H_q \)\\[3pt]
66&0&\( 12D_3^p + 24D_4^p + 18D_3^q + 36D_4^q \)\newline\( + 12G_p + 38G_q + 5H_q + 6U_{p_{1}} \)\newline\( + 6L_{p_{8}} + 6L_{q_{6}} + 6L_{q_{7}} + 6L_{\pi_p} \)\newline\( + 6\Sigma + 24\widetilde H_p + 4I_2^q + 8I_3^q \)\newline\( + 20\widetilde H_q + 6A_{01}^p + 6A_{02}^p + 6A_{03}^p \)\newline\( + 6B_{04}^p + 6A_{01}^q + 6A_{02}^q + 6A_{03}^q \)\newline\( + 6A_{04}^q + 6C_0 + 18A_{12}^p + 18A_{13}^p \)\newline\( + 12A_{14}^p + 6B_{14}^p + 18A_{12}^q + 18A_{13}^q \)\newline\( + 18A_{14}^q + 18C_1 + 24A_{23}^p + 24A_{24}^p \)\newline\( + 24A_{23}^q + 24A_{24}^q + 24C_2 + 12A_{34}^p \)\newline\( + 12A_{34}^q + 12C_3 + 30M + 6P_q \)\\[3pt]
67&1&\( 4D_3^p + 8D_4^p + 6E_p + 6D_3^q \)\newline\( + 12D_4^q + 11E_q + 4G_p + 13G_q \)\newline\( + Z + U_{p_{1}} + L_{p_{9}} + 2L_{q_{6}} \)\newline\( + 2L_{q_{7}} + 2L_{q_{9}} + 2L_{\pi_p} + 2\Sigma \)\newline\( + \widehat H_q \)\\[3pt]
68&1&\( D_2^p + 4D_3^p + 7D_4^p + 5E_p \)\newline\( + 6D_3^q + 12D_4^q + 11E_q + 3G_p \)\newline\( + 13G_q + Z + U_{p_{0}} + 2L_{q_{6}} \)\newline\( + 2L_{q_{7}} + 2L_{\pi_p} + 2L_{\sigma_q} + 2\Sigma \)\newline\( + 2X + \widehat H_q \)\\[3pt]
69&3&\( 4D_3^p + 8D_4^p + 5E_p + 4D_3^q \)\newline\( + 8D_4^q + 4E_q + 11G_p + 6G_q \)\newline\( + 2L_{p_{5}} + U_{q_{0}} + 2L_{q_{5}} + 3L_{q_{8}} \)\newline\( + 2L_{\sigma_p} + 3\Sigma + 3\widehat H_p + 2\widehat H_q \)\\[3pt]
70&3&\( 4D_3^p + 8D_4^p + E_p + 4D_3^q \)\newline\( + 8D_4^q + 5E_q + 3G_p + 11G_q \)\newline\( + U_{p_{0}} + 3L_{p_{8}} + 2L_{q_{5}} + 2L_{q_{8}} \)\newline\( + 2L_{\pi_p} + 3\Sigma + 2\widehat H_p + 3\widehat H_q \)\\[3pt]
71&1&\( 4D_3^p + 8D_4^p + 5E_p + 6D_3^q \)\newline\( + 12D_4^q + 11E_q + 3G_p + 13G_q \)\newline\( + Z + U_{p_{1}} + L_{p_{9}} + 2L_{q_{6}} \)\newline\( + 2L_{q_{7}} + 2L_{q_{9}} + 2L_{\pi_p} + 2\Sigma \)\newline\( + X + \widehat H_q \)\\[3pt]
72&2&\( D_2^p + 3D_3^p + 5D_4^p + 3D_3^q \)\newline\( + 6D_4^q + 4E_q + 8G_q + U_{p_{1}} \)\newline\( + 2L_{p_{8}} + 2L_{q_{5}} + L_{\pi_p} + L_{\sigma_q} \)\newline\( + 2\Sigma + \widehat H_p + 2\widehat H_q \)\\[3pt]
73&3&\( 4D_3^p + 8D_4^p + 4E_p + 4D_3^q \)\newline\( + 8D_4^q + 5E_q + 6G_p + 11G_q \)\newline\( + U_{p_{0}} + 2L_{p_{5}} + 3L_{p_{8}} + 2L_{q_{5}} \)\newline\( + 2L_{\sigma_q} + 3\Sigma + 2\widehat H_p + 3\widehat H_q \)\\[3pt]
74&3&\( 4D_3^p + 6D_4^p + 4D_3^q + 8D_4^q \)\newline\( + 2G_p + 11G_q + U_{p_{0}} + L_{p_{8}} \)\newline\( + 2L_{q_{6}} + 2L_{q_{8}} + 2L_{\pi_p} + 3\Sigma \)\newline\( + X + 2\widehat H_p + 3\widehat H_q + 5Y_5 \)\\[3pt]
75&0&\( 12D_3^p + 24D_4^p + 12D_3^q + 24D_4^q \)\newline\( + 12G_p + 34G_q + 4H_q + 6U_{p_{1}} \)\newline\( + 12L_{p_{8}} + 6L_{q_{5}} + 6L_{q_{8}} + 6L_{\pi_p} \)\newline\( + 12\Sigma + 18\widetilde H_p + 2I_2^q + 4I_3^q \)\newline\( + 16\widetilde H_q + 3A_{01}^p + 3A_{02}^p + 3A_{03}^p \)\newline\( + 3B_{04}^p + 3A_{01}^q + 3A_{02}^q + 3A_{03}^q \)\newline\( + 3A_{04}^q + 3C_0 + 15A_{12}^p + 15A_{13}^p \)\newline\( + 12A_{14}^p + 3B_{14}^p + 15A_{12}^q + 15A_{13}^q \)\newline\( + 15A_{14}^q + 15C_1 + 18A_{23}^p + 18A_{24}^p \)\newline\( + 18A_{23}^q + 18A_{24}^q + 18C_2 + 12A_{34}^p \)\newline\( + 12A_{34}^q + 12C_3 + 24M \)\\[3pt]
76&0&\( 2D_3^p + 10D_4^p + 8D_3^q + 16D_4^q \)\newline\( + 14G_p + 10G_q + U_{p_{0}} + 2L_{p_{5}} \)\newline\( + L_{p_{8}} + 3U_{q_{0}} + 6L_{q_{7}} + 5L_{q_{8}} \)\newline\( + 6\Sigma + 2I_1^p + 12\widetilde H_p + 2I_3^q \)\newline\( + 10\widetilde H_q + 6A_{12}^p + 6A_{13}^p + 2A_{14}^p \)\newline\( + 4B_{14}^p + 6A_{12}^q + 6A_{13}^q + 6A_{14}^q \)\newline\( + 6C_1 + 12A_{23}^p + 12A_{24}^p + 12A_{23}^q \)\newline\( + 12A_{24}^q + 12C_2 + 8A_{34}^p + 8A_{34}^q \)\newline\( + 8C_3 + 13M \)\\[3pt]
77&0&\( 8D_3^p + 16D_4^p + 3D_3^q + 10D_4^q \)\newline\( + 10G_p + 14G_q + 3U_{p_{0}} + L_{p_{6}} \)\newline\( + 5L_{p_{7}} + 5L_{p_{8}} + U_{q_{0}} + 2L_{q_{5}} \)\newline\( + L_{q_{8}} + 5\Sigma + 2I_3^p + 10\widetilde H_p \)\newline\( + 2I_1^q + 12\widetilde H_q + 6A_{12}^p + 6A_{13}^p \)\newline\( + 6A_{14}^p + 6A_{12}^q + 6A_{13}^q + 2A_{14}^q \)\newline\( + 4B_{14}^q + 6C_1 + 12A_{23}^p + 12A_{24}^p \)\newline\( + 12A_{23}^q + 12A_{24}^q + 12C_2 + 8A_{34}^p \)\newline\( + 8A_{34}^q + 8C_3 + 13M \)\\[3pt]
\end{longtable}
\endgroup

\section[Certificate specification for upper bound 96]{Certificate specification for the upper bound 96}\label{u96:app:upper96}
This appendix specifies the additional finite proof in
Section~\ref{u96:sec:upper96}. The names \texttt{upper97} and
\texttt{upper96} refer to the resulting upper bound, so they exclude
sets of sizes 98 and 97 respectively. The two JSON files contain the
complete integer multiplier pools, ordered root assignments, and all
binary trees. They are required companions to the new finite proof;
unlike the earlier multiplier tables, the new pools are not duplicated
in print. The LaTeX source itself compiles without the JSON files.

\subsection{The 103 effective-hole rows}
The following table gives $C_\beta(r)$ in
Lemma~\ref{u96:lem:hole-parameter}, with its parameter $b=\beta$.
A dash means that the row is not used. Each entry asserts
\[
 h>2\beta,\quad h<|\FS(X)|-1
 \quad\Longrightarrow\quad \mu(X)\geq r+C_\beta(r)+1.
\]
For independent verification, compute
\begin{align*}
 K&=\left\lfloor\frac{3+\sqrt{9+8(C+1)}}2\right\rfloor,
 &D&=C+\max(3,K),\\
 M&=\max\left(2\beta,
       1+\left\lfloor\frac{1+\sqrt{9+8(D+1+6\beta)}}2\right\rfloor\right),
 &J&=\max(K+1,2\beta+5,M).
\end{align*}
Every listed entry satisfies $(r-J+1)\beta>D$ and $T_{J-1}-1>D$.
The verifier uses integer square roots. At size 98 only rows $r\geq24$
are needed (101 entries); size 97 uses all 103.

\begin{longtable}{rrrr}
\toprule
$r$&$C_7(r)$&$C_9(r)$&$C_{10}(r)$\\\midrule
\endfirsthead
\toprule
$r$&$C_7(r)$&$C_9(r)$&$C_{10}(r)$\\\midrule
\endhead
23&25&3&--\\
24&32&11&--\\
25&38&18&4\\
26&44&26&12\\
27&51&34&21\\
28&57&42&30\\
29&63&51&39\\
30&70&59&48\\
31&76&67&57\\
32&83&75&66\\
33&89&84&75\\
34&96&92&85\\
35&102&101&94\\
36&109&109&103\\
37&110&117&113\\
38&116&126&122\\
39&122&134&132\\
40&128&143&141\\
41&129&151&150\\
42&135&160&160\\
43&142&168&169\\
44&147&176&179\\
45&149&177&188\\
46&155&186&198\\
47&162&194&207\\
48&167&198&215\\
49&168&203&217\\
50&175&211&227\\
51&182&220&236\\
52&187&221&239\\
53&188&228&246\\
54&195&237&255\\
55&202&245&264\\
56&208&246&265\\
57&209&254&274\\
\bottomrule
\end{longtable}

\subsection{Root variables, center labels, and modes}
All variables are nonnegative. The 83 coordinates consist of
\[
 (p_1,\ldots,p_4,q_1,\ldots,q_4,M_p,M_q,R_p,R_q,S_p,S_q,\delta,t),
\]
40 reserved endpoint-loss coordinates, six inner elements
$x_1,x_2,x_3,y_1,y_2,y_3$, the variable $K$, and twenty envelope variables
$g_{jis}$, where $0\leq j\leq3$, $j<i\leq4$, and $s=p,q$.
The coordinate $t$ equals the common hole parameter $z$ in positive-hole
low models and all size-97 high models. It is fixed to zero in zero-hole
models and in the coarse size-98 H1 family, whose bounds were stated
separately. Unused loss coordinates are harmless nonnegative variables.

With the earlier definitions $c_p=p_1-p_2-p_3-p_4$,
$\alpha_p=p_2-p_3-p_4$, and
$\beta_p=2p_2-p_1-p_3-p_4$, the ordered positive center labels are
\[
 \begin{array}{c|cccccccccc}
 \text{label}&p_0&p_1&p_2&p_3&p_4&p_5&p_6&p_7&p_8&p_9\\\hline
 \text{center}&c_p+p_1&c_p+p_2&c_p+p_3&c_p+p_4&
 c_p-q_1&c_p-q_2&c_p-q_3&c_p-q_4&\alpha_p&\beta_p
 \end{array}
\]
Here a center label $p_i$ is a string used in the programs; in this table
it must not be confused with the outer element $p_i$ in the formulas.
Negative center labels $q_0,\ldots,q_9$ are obtained by swapping $p,q$.
The first pair has index 4 on both sides. The second pair has index 9 in
branch A and index 5 in branch B. Remove those indices from each side
to obtain the sixteen unpaired units. Enumerate every $j$-element set
of missing units among these sixteen and the two pairs.

If pair 1 is present, its retained member is positive after reflection.
If pair 2 is present and cannot split, enumerate its positive and negative
member choices, denoted by 0 and 1. If it is missing, keep the dummy choice
0. In branch A, when a missing unit is $p_5$ or $q_5$, drop pair 2 from
the ordinary retained-center list; if pair 2 is nevertheless present,
enumerate the two choices $s=p,q$ of an actually contributing member and
add \eqref{u96:eq:chosen-member96}. This gives exactly $4,70,578$ modes for
$j=0,1,2$.

The code first lists the surviving unpaired centers, then pair 1, then
the chosen pair-2 member if retained. Assign the loss coordinates to that
list in order. Labels \texttt{pair1p}, \texttt{pair2p}, and
\texttt{pair2q} remove ambiguity. The separate member-choice variable
for a split pair is called \texttt{variant}.

\subsection{A complete description of the root rows}
The outer orders are $p_i\geq p_{i+1}+1$, $p_4\geq M_p+1$ and their
reflections for $j\leq1$; replace these gaps by zero for $j=2$.
Always impose $M_p\geq R_p+4$, $M_q\geq R_q+4$,
$R_p\geq a$, $R_q\geq b$.
Branch A has $\beta_p+\beta_q\leq0$; branch B has
$\beta_p+\beta_q\geq1$ for $j\leq1$ and its relaxation $\geq0$ for $j=2$.
Add the total-loss bound $\sum(l_v+u_v)\leq\delta$ and the regime's
quantile, mass, loss-budget, span-budget, and hole-parameter rows from
Section~\ref{u96:sec:upper96}. The original mass rows are multiplied by two
in the code; this is just a positive rescaling. The quantities
$\epsilon_p,\epsilon_q$ in the mass formulas are replaced by the affine
lower bounds specified for that regime; they are not additional coordinates.

For $x_0=R_p$, $x_4=M_p$, $y_0=R_q$, $y_4=M_q$, impose
$x_i-x_{i-1}\geq1$, $y_i-y_{i-1}\geq1$ for $1\leq i\leq4$ and both
mass rows \eqref{u96:eq:mass96-full} on both sides. Write $W_j=x_j+y_j$,
$t_j=2(4-j)$. For each $j=0,1,2,3$, the envelope rows are
\begin{gather*}
 0\leq g_{jip}\leq x_i-x_{i-1},\qquad g_{jip}\leq W_j+1,\\
 0\leq g_{jiq}\leq y_i-y_{i-1},\qquad g_{jiq}\leq W_j+1,
 \qquad j<i\leq4,\\
 K\leq2W_j+1+\frac{t_j(t_j+1)}2+
                \sum_{i=j+1}^4(g_{jip}+g_{jiq}).
\end{gather*}
These rows hold for an actual set by choosing each $g$ equal to the
indicated minimum of a gap and $W_j+1$. No equality with that minimum is
needed in the relaxed model. Add \eqref{u96:eq:card96}.
The only additional root row in an L family is
$l_v+u_v\geq d-4$ for its designated center.

The high-state selection is itself explicit: loop over all
$0\leq k\leq16$, $15\leq z\leq395-k(k+1)$ at size 98, and all
$0\leq k\leq17$, $15\leq z\leq414-k^2$ at size 97. Compute $u,v,C,M$
as in the proof, with exact integer and rational arithmetic. The size-97
loop has 5415 states, including those immediately rejected by $C<0$.
No sampled subrange or numerical tolerance is used.

\subsection{The complete list of roots that require a tree}
Every branched root is in pairing branch A. In the following table,
``choice'' is the ordinary second-pair choice, ``member'' is the
actually contributing split-pair member, and a dash means absent.
The large center is $p_8$ in every L row. Each split has both children.

\begin{center}\small
\begin{tabular}{lrllrrr}
\toprule
Regime&$a$&Missing units&Member&Choice&Splits&Depth\\\midrule
ZL&38&$p_0$&--&1&4&4\\
ZL&39&$p_0$&--&1&39&8\\
PL&37&$p_0$&--&1&1&1\\
PL&38&$p_0$&--&1&1&1\\
PL&39&$p_0$&--&1&1&1\\
PL&40&$p_0$&--&1&1&1\\
Z2&38&$p_0,p_8$&--&1&4&3\\
Z2&39&$p_0,p_8$&--&1&64&10\\
Z2&40&$p_0,p_8$&--&1&3&3\\
Z2&40&$p_0,q_5$&q&0&1&1\\
Z2&41&$p_0,q_5$&q&0&1&1\\
Z2&42&$p_0,q_5$&q&0&1&1\\
P2&37&$p_0,p_8$&--&1&1&1\\
P2&38&$p_0,p_8$&--&1&2&2\\
P2&39&$p_0,p_8$&--&1&1&1\\
P2&40&$p_0,p_8$&--&1&2&2\\
P2&39&$p_0,q_5$&q&0&1&1\\
P2&40&$p_0,q_5$&q&0&1&1\\
P2&41&$p_0,q_5$&q&0&1&1\\
\bottomrule
\end{tabular}
\end{center}

The 130 internal nodes consist of 13 ordinary/outside splits, 111
inner-extreme/loss splits, and six both/sole-pair splits. The maximum
depth is ten. The two additional sole-pair loss variables use reserved
coordinates 52 and 53 (zero-based); all ordinary retained losses in
these two-empty modes have indices at most 47.

For complete reproducibility, \texttt{branches()} enumerates, in retained
endpoint-row order, first the ordinary/outside split and then the four
inner splits with $i=1,2,3,4$. It finally adds the both/sole split exactly
when $j=2$ and a split-pair member was selected. A tree node stores the
zero-based index in that list. Its left and right children contain the
rows in the same order as equations
\eqref{u96:eq:branch-normal-low96}--\eqref{u96:eq:sole96}:
ordinary versus outside; loss versus inner extreme; both versus sole.
The JSON is therefore an explicit finite tree specification, not a
record of a solver-dependent search procedure.

\subsection{Exact contradiction check and data integrity}
Each certificate stores a list of distinct row labels with strictly
positive integer weights. All root and branch rows have integer
coefficients. At a leaf, add its path rows to its root matrix $A$ and
right side $b$, then compute
\[
 c_i=\sum_r w_rA_{ri}\quad(1\leq i\leq83),\qquad
 v=\sum_r w_rb_r.
\]
The verifier requires every $c_i\geq0$ and $v<0$, and checks the stored
positive value $-v$. This suffices because every coordinate is nonnegative.
There is no objective-bound denominator in these two new pools: their
historical field \texttt{D} is always zero.

The checker independently generates the full ordered root sequence,
compares every regime, side size, high-hole parameter, missing-unit mode,
large center, and member choice with the assignment, and recursively
checks both children of every tree. Repeated split indices on a path,
duplicate row labels, absent children, and noninteger weights are rejected.
The complete verification checks 45928 size-98 roots and 81376 size-97 leaves.
The least size-97 contradiction magnitude is 1. All arithmetic uses
Python's arbitrary-precision integers or exact fractions.

The proof-model and certificate hashes below identify the exact companion
companion data used in the proof. The archive's \texttt{SHA256SUMS.txt} also
covers the complete source and historical verification dependencies.

\noindent\texttt{upper97\_model.py}\par\noindent{\footnotesize\ttfamily 445525a759e423c2705aa00b1dc5f0eb\allowbreak{}f5eba5963c6320cc17cf6d87c6f4079a}\par\smallskip
\noindent\texttt{upper97\_exact.json}\par\noindent{\footnotesize\ttfamily b42b641a9154d8607a8b2570d5686628\allowbreak{}482b6ff0df4fbfb8b10c85aadddcefcd}\par\smallskip
\noindent\texttt{verify\_upper97.py}\par\noindent{\footnotesize\ttfamily cdec064f37a0a65aad991692aa90ffee\allowbreak{}62f7b28cb86f2eb8e0acad02fb2ebbea}\par\smallskip
\noindent\texttt{upper96\_model.py}\par\noindent{\footnotesize\ttfamily 455d71883018823e8eceb4cff557fb69\allowbreak{}bf108c3130003539f885d3e3c44a573c}\par\smallskip
\noindent\texttt{upper96\_exact.json}\par\noindent{\footnotesize\ttfamily f30e9448567cb33fd08aa9ea10800401\allowbreak{}9122b7f185e8642670d22912607ed411}\par\smallskip
\noindent\texttt{verify\_upper96.py}\par\noindent{\footnotesize\ttfamily 2bfd1f4a05557dc1f33371a98ce65311\allowbreak{}cd6aca74b15d4ea78fe57f5a7bdc2704}\par\smallskip



\clearpage

\section[Four analytic covering combinations at size 87]{An analytic certificate for the size-87 AP primitivity exception}
\label{u86ap:sec:certificate}

The following argument proves individual primitivity in the remaining
AP branch of the size-87 common-lattice proof. It uses four explicitly displayed row
combinations, valid uniformly in the side sizes. No enumeration or
primitive-core hypothesis enters the covering argument.

\begin{lemma}[Small covering slack is impossible at zero holes]
\label{u86ap:lem:small-slack}
In the size-87 covering setup, suppose $D$ is integral on the common
core lattice, $a+b=70$, $a,b\ge15$, and the extreme-translate propagation
conclusion $w\le4z$ holds with $z=0$. Then the covering slack $c$ cannot satisfy
$c\le6$.
\end{lemma}
\begin{proof}
Here $|D|=79$. Every empty ordinary covering unit costs 79 in $c$,
so $c\le6$ forces all eighteen ordinary units to contribute. Write
$\delta=c$, and use the outer coordinates
\[
 p_1>p_2>p_3>p_4>M_p,\qquad
 q_1>q_2>q_3>q_4>M_q,
\]
where $M_p,M_q$ are the positive and negative endpoint magnitudes of
$D$. Put
\[
 c_p=p_1-p_2-p_3-p_4,\qquad c_q=q_1-q_2-q_3-q_4.
\]
The four own centers $b_i=c_p+p_i$ on the positive side are all
contributing. Since the actual target and $D$ are integral, each $b_i$
is integral. The identity
\[
 c_p=-b_1+b_2+b_3+b_4
\]
then makes $c_p$ and every $p_i$ integral. Reflection proves the same
for $c_q,q_i$. No leading component is assumed to contribute.

Let $R_p,R_q$ be the core maxima, $S_p,S_q$ the sums of the positive
and negative magnitudes in $D$, and $K=|D+D|$. Distinct integral
points give
\[
 M_p\ge R_p+4,\quad M_q\ge R_q+4,\quad R_p\ge a,\quad R_q\ge b,
\]
\[
 S_p\ge T_a+M_p+3R_p+6,\qquad
 S_q\ge T_b+M_q+3R_q+6.
\]
The last inequalities count the core mass and the three inner peeled
points individually; they do not require a primitive core. The integer
hull and the exact zero-hole cardinality identity give
\[
 K\le2M_p+2M_q+1,\qquad S_p+S_q+\delta-2K=1506.
\]
Indeed propagation gives $w=0$, and $18\cdot79+84=1506$.

\paragraph{The centers and the four actual alternatives.}
Define
\[
 \alpha_p=p_2-p_3-p_4,\quad
 \beta_p=2p_2-p_1-p_3-p_4,
\]
with their reflected $q$ versions. For $i=0,\ldots,9$, the centers
on the positive side are
\[
 (\gamma_{p,i})_{i=0}^9=
 (c_p+p_1,c_p+p_2,c_p+p_3,c_p+p_4,
 c_p-q_1,c_p-q_2,c_p-q_3,c_p-q_4,\alpha_p,\beta_p).
\]
Reflection interchanges $p,q$ to define $\gamma_{q,i}$. Branch A means
$\beta_p+\beta_q\le0$ and has second-pair index $j=9$;
branch B means $\beta_p+\beta_q>0$ and has $j=5$.
The sixteen unpaired centers are exactly
\[
 \gamma_{s,i},\qquad s\in\{p,q\},\quad i\notin\{4,j\}.
\]
All these unpaired components contribute. Both paired units in fact have
exactly one contributing member. Put
$V=p_3+p_4+q_3+q_4$. The first paired centers have sum
$-(p_2+p_3+p_4+q_2+q_3+q_4)<-(M_p+M_q)$.
If both met their respective open half-lines, their sum would exceed
$-(M_p+M_q)$, a contradiction.
In branch A, the unpaired centers $c_p-q_2,c_q-p_2$ both contribute, so
\[
 c_p-q_2+c_q-p_2>-(M_p+M_q).
\]
It follows that
\[
 \beta_p+\beta_q
 =-(c_p-q_2+c_q-p_2)-2V
 <M_p+M_q-2V<-(M_p+M_q).
\]
In branch B, the second paired centers have sum
$-(\beta_p+\beta_q)-2V<-(M_p+M_q)$ directly.
Thus at most one member of each pair meets its open half-line, and the
nonempty-unit condition makes that member unique. Reflection lets us
choose the actual contributing member of the first pair to be
$\gamma_{p,4}$. Denote this selected member by 1. Choose an actual
contributing member of the second paired unit, denoted by 2.
Choice 0 selects $\gamma_{p,j}$ and choice 1 selects $\gamma_{q,j}$.
Thus A0, A1, B0, B1 cover every actual configuration. Reflection can
interchange $a,b$, so no order between these two sizes is imposed.
The selection is made from actual contribution, not from a necessary
inequality satisfied by an arbitrary relaxation point.

For each of these eighteen selected components $v$, take its actual
lower and upper hull losses $l_v,u_v\ge0$. For a center $v$ on side $s$,
with the opposite side denoted by $t$, its actual nonempty integral
translate satisfies the global endpoint-loss inequalities
\[
 v+l_v+M_s\ge79,\qquad S_s-v+u_v+M_t\ge78.
\]
These are the global quantile rows proved in
Appendix~\ref{u96:app:upper96}; they require a nonempty integral
translate, not an in-hull center or a primitive core. Because each pair
has a unique contributing member, its other member contributes nothing.
Consequently the covering count uses precisely these eighteen actual
components. Their hull losses are conservative lower bounds for the
corresponding actual-target deficits, so their sum, together with any
chosen conservative leading losses, fits the same single covering budget.

Write $L_p,L_q\ge0$ for the two leading losses and put
\[
 \mathcal L=L_p+L_q+\sum_v(l_v+u_v).
\]
The covering inequality gives $\mathcal L\le\delta$. We need no other
inequality for $L_p,L_q$.

\paragraph{All elementary rows.}
The names below denote inequalities whose left and right sides are added
with the displayed positive weights. Let $r_p=a$, $r_q=b$. For $s=p,q$,
let $s_i$ mean $p_i$ or $q_i$ accordingly:
\begin{align*}
 \mathsf C &: \delta\le E,\qquad E=6,\\
 \mathsf O_s &: M_s-s_4\le-1,\\
 \mathsf D_s &: s_4-s_3\le-1,\\
 \mathsf H_s &: R_s-M_s\le-4,\\
 \mathsf J_s &: -R_s\le-r_s,\\
 \mathsf F_s &:2M_s+6R_s-2S_s\le-2T_{r_s}-12,\\
 \mathsf Q &:S_p+S_q+\delta-2K\le1506,\\
 \mathsf K &:K-2M_p-2M_q\le1,\\
 \mathsf B &: \mathcal L-\delta\le0,\\
 \mathsf L_v &: -v-l_v-M_s\le-79,\\
 \mathsf U_v &: v-u_v-S_s-M_t\le-78.
\end{align*}
The independent coordinates $p_i,q_i,M_s,R_s,S_s,\delta,K,l_v,u_v,L_s$
are nonnegative. Centers $v$ are linear forms and need not be nonnegative.
Subscripts
$p,i$ and $q,i$ on loss rows refer to $\gamma_{p,i},\gamma_{q,i}$;
subscripts 1 and 2 refer to the selected pair members.
Every row used below is among the retained centers in its stated mode.

\paragraph{Mode A0.}
Add the following rows:
\[
 \mathsf C+2\mathsf O_p+\mathsf L_1+\mathsf L_2
 +2\mathsf O_q+\mathsf L_{q,5}+\mathsf L_{q,8}+\mathsf B.
\]
Exact cancellation gives
\[
 2p_3+2q_3+\mathcal L-l_1-l_2-l_{q,5}-l_{q,8}
 \le E-320.
\]
Each subtracted loss is a distinct term of $\mathcal L$, so the left
side is nonnegative; the right side is at most $-314$.

\paragraph{Mode A1.}
Add
\begin{align*}
 &56\mathsf C+6\mathsf H_p+24\mathsf J_p+12\mathsf J_q
 +3\mathsf F_p+2\mathsf F_q+14\mathsf Q+28\mathsf K\\
 &\quad+8\mathsf U_{p,0}+74\mathsf D_p+152\mathsf O_p
 +70\mathsf L_{p,8}+16\mathsf L_1+26\mathsf L_2\\
 &\quad+10\mathsf U_{q,0}+70\mathsf D_q+148\mathsf O_q
 +62\mathsf L_{q,5}+70\mathsf B.
\end{align*}
The result is
\begin{align*}
 &4p_3+8q_3+70\mathcal L
 -8u_{p,0}-70l_{p,8}-16l_1-26l_2-10u_{q,0}-62l_{q,5}\\
 &\hspace{12mm}\le56E+5434-24a-12b-6T_a-4T_b.
\end{align*}
The subtracted loss variables are distinct, each with coefficient at
most 70. Hence every coefficient on the left is nonnegative. Since
$b=70-a$ and $E\le6$, the right side is at most
\[
 -5a^2+267a-5010
 =-5\left(a-\frac{267}{10}\right)^2-\frac{28911}{20}<0.
\]

\paragraph{Mode B0.}
Add
\begin{align*}
 &11\mathsf C+12\mathsf H_p+12\mathsf H_q+12\mathsf J_p
 +18\mathsf J_q+\mathsf F_q+4\mathsf Q+8\mathsf K\\
 &\quad+4\mathsf U_{p,0}+19\mathsf D_p+46\mathsf O_p
 +4\mathsf L_{p,9}+12\mathsf L_1+2\mathsf U_{q,0}\\
 &\quad+6\mathsf L_{q,2}+32\mathsf D_q+74\mathsf O_q
 +9\mathsf L_{q,6}+15\mathsf L_{q,7}+14\mathsf L_{q,9}
 +15\mathsf B.
\end{align*}
This yields
\begin{align*}
 &2p_3+4q_3+15\mathcal L
 -4u_{p,0}-4l_{p,9}-12l_1-2u_{q,0}\phantom{{}+11E}\\
 &\qquad-6l_{q,2}-9l_{q,6}-15l_{q,7}-14l_{q,9}\\
 &\hspace{12mm}\le11E+545-12a-18b-2T_b.
\end{align*}
All subtracted losses are distinct and have coefficients at most 15.
The right side, using $a=70-b$ and $E\le6$, is at most
\[
 -229-b^2-7b<0.
\]

\paragraph{Mode B1.}
Add
\[
 \mathsf C+2\mathsf O_p+\mathsf L_{p,9}+\mathsf L_1
 +\mathsf L_2+\mathsf L_{q,2}+3\mathsf O_q
 +\mathsf L_{q,9}+\mathsf B.
\]
The sum is
\[
 2p_3+2q_3+\mathcal L-l_{p,9}-l_1-l_2-l_{q,2}-l_{q,9}
 \le E-400\le-394.
\]
Again all left-side coefficients are nonnegative. The four modes exhaust
the actual covering alternatives, proving the lemma.
\end{proof}

\begin{corollary}[The larger-AP branch is individually primitive]
\label{u86ap:cor:primitive}
In the larger-side AP branch of the size-87 common-lattice proof, the
smaller core side is primitive on the common lattice.
\end{corollary}
\begin{proof}
All common-lattice, $D$-integrality and translate-overlap conclusions have
already been proved before this step; thus $G=148$ is available without
circularity. Let $r=\min(a,b)$, $k=35-r$, and
$\Omega=574-k^2$. The larger AP support defines the unit lattice, so
$h_{\rm large}=z=0$ and $h_{\rm small}=\kappa$.
If the smaller side had gcd at least 2, its subset-sum support, of size
$1+T_r+\mu_{\rm small}$, would have at least
$T_r+\mu_{\rm small}$ holes. Therefore $\kappa\ge T_r$.
The post-lattice budget, $\rho_p,\rho_q\ge0$, and
$T\ge\kappa+148$ give
\[
 c\le\Omega-148-T_r=:E_*(k)
 =\frac{-408+71k-3k^2}{2}
 =\frac{145}{24}-\frac32\left(k-\frac{71}{6}\right)^2.
\]
The integer $E_*(k)$ is at most 6. If it is negative, $c\ge0$ is already
impossible. Otherwise $c\le6$ contradicts
Lemma~\ref{u86ap:lem:small-slack}. This proves primitivity without
enumerating side sizes, proving the smaller support AP, or singling out
its possible gcd values.
\end{proof}



\clearpage
\begin{thebibliography}{9}
\addcontentsline{toc}{section}{References}
\bibitem{CDK2026}
R.~Carpenter, C.~Defant and N.~Kravitz,
\emph{Sets with few subset sums}, arXiv:2605.05498 (2026),
Theorem~1.1 and Lemma~2.1.
\url{https://arxiv.org/abs/2605.05498}.
\bibitem{BGGH2013}
E.~Balandraud, B.~Girard, S.~Griffiths and Y.~Ould Hamidoune,
\emph{Subset sums in abelian groups},
European Journal of Combinatorics \textbf{34} (2013), 1269--1286,
Lemma~3.6 and Theorem~1.5.
\url{https://arxiv.org/abs/1112.1929}.
\bibitem{Stanchescu1996}
Y.~Stanchescu, \emph{On addition of two distinct sets of integers},
Acta Arithmetica \textbf{75} (1996), 191--194, Theorem~4 and Corollary~2.
\url{https://doi.org/10.4064/aa-75-2-191-194}.
\bibitem{Lev1998}
V.~F.~Lev, \emph{On consecutive subset sums},
Discrete Mathematics \textbf{187} (1998), 151--160, Theorem~1.
\url{https://doi.org/10.1016/S0012-365X(98)80006-X}.
\end{thebibliography}
\end{document}
